% Generated mechanically from the frozen spaces-more-morphisms.tex target.
% Do not edit this generated file; edit the bound locale source instead.

% OK, start here.
%

\setcounter{chapter}{75}
\chapter{空间态射进阶}



\phantomsection
\label{spaces-more-morphisms-section-phantom}


\section{引言}
\label{spaces-more-morphisms-section-introduction}

\noindent
本章继续研究代数空间态射的性质。一个基本参考文献是 \cite{Kn}。





\section{约定}
\label{spaces-more-morphisms-section-conventions}

\noindent
我们始终假设所有概形均包含在一个大 fppf 位点 $\Sch_{fppf}$ 中，
并且所考虑的所有环 $A$ 都满足：$\Spec(A)$（同构于）该大位点的一个对象。

\medskip\noindent
设 $S$ 为概形，$X$ 为 $S$ 上的代数空间。在本章及下一章中，
我们将以 $X \times_S X$ 表示 $X$ 与自身的乘积（在 $S$ 上代数空间的范畴中），
而不用 $X \times X$。







\section{根态射}
\label{spaces-more-morphisms-section-radicial}

\noindent
与概形态射的情形不同，根态射并不等同于普遍单射态射。
事实上，前者略强一些。

\begin{definition}
\label{spaces-more-morphisms-definition-radicial}
设 $S$ 为概形，$f : X \to Y$ 为 $S$ 上代数空间的态射。
若 $f$ 满足下述条件，则称其为{\it 根态射}：对任意态射
$\Spec(K) \to Y$（其中 $K$ 为域），约化
$(\Spec(K) \times_Y X)_{red}$ 或为空，或可由 $K$ 的某个
纯不可分域扩张之谱表示。
\end{definition}

\begin{lemma}
\label{spaces-more-morphisms-lemma-radicial-implies-universally-injective}
代数空间的根态射是普遍单射的。
\end{lemma}

\begin{proof}
设 $S$ 为概形，$f : X \to Y$ 为 $S$ 上代数空间的根态射。
由定义显然可知，给定态射 $\Spec(K) \to Y$，它至多有一个
到 $X$ 的提升。因此，由《空间的态射》引理
\ref{spaces-morphisms-lemma-universally-injective}，
$f$ 是普遍单射的。
\end{proof}

\begin{example}
\label{spaces-more-morphisms-example-universally-injective-not-radicial}
普遍单射与根态不再等价。例如，《空间》例
\ref{spaces-example-Qbar} 中的态射
$$
X = [\Spec(\overline{\mathbf{Q}})/
\text{Gal}(\overline{\mathbf{Q}}/\mathbf{Q})]
\longrightarrow
S = \Spec(\mathbf{Q})
$$
是普遍单射的，但不是上述意义下的根态射。
\end{example}

\noindent
不过，逆向蕴含在许多情形下仍然成立。

\begin{lemma}
\label{spaces-more-morphisms-lemma-when-universally-injective-radicial}
设 $S$ 为概形，$f : X \to Y$ 为 $S$ 上代数空间的普遍单射态射。
\begin{enumerate}
\item 若 $f$ 良态，则 $f$ 为根态射。
\item 若 $f$ 拟分离，则 $f$ 为根态射。
\item 若 $f$ 局部分离，则 $f$ 为根态射。
\end{enumerate}
\end{lemma}

\begin{proof}
设 $\mathcal{P}$ 为代数空间态射的一个性质，它在基变换与复合下稳定，
且对闭浸入成立。假设 $f : X \to Y$ 具有 $\mathcal{P}$ 且普遍单射。
那么，在定义 \ref{spaces-more-morphisms-definition-radicial} 的情形中，态射
$(\Spec(K) \times_Y X)_{red} \to \Spec(K)$
普遍单射且具有 $\mathcal{P}$。因此，为证明
$$
\mathcal{P} + \text{普遍单射}
\Rightarrow
\text{根态}
$$
只需证明：任意非空约化代数空间 $X$，若其在域上的结构态射
$X \to \Spec(K)$ 普遍单射且具有 $\mathcal{P}$，则该空间可由某个域的谱表示。
% P09-ERR-0506
确实，这时 $X \to \Spec(K)$ 将是概形态射，而概形态射的根态与普遍单射等价；
见《态射》引理 \ref{morphisms-lemma-universally-injective}。

\medskip\noindent
先证明 (1)。假设 $f$ 良态且普遍单射。由《良态空间》诸引理
\ref{decent-spaces-lemma-base-change-relative-conditions}、
\ref{decent-spaces-lemma-composition-relative-conditions} 及
\ref{decent-spaces-lemma-properties-trivial-implications}
（最后一个引理说明浸入是良态的），第一段的讨论适用。
设 $X$ 为域 $K$ 上非空、良态、约化且普遍单射的代数空间。
特别地，$|X|$ 是单点集。由《良态空间》引理
\ref{decent-spaces-lemma-when-field}，可知
$X \cong \Spec(L)$，其中 $K \subset L$ 是某个扩张，如所需。

\medskip\noindent
拟分离态射是良态的，见《良态空间》引理
\ref{decent-spaces-lemma-properties-trivial-implications}。
故 (1) 蕴含 (2)。

\medskip\noindent
下面证明 (3)。回顾：分离性公理在基变换与复合下稳定，并且闭浸入是分离的；
见《空间的态射》诸引理
\ref{spaces-morphisms-lemma-base-change-separated}、
\ref{spaces-morphisms-lemma-composition-separated} 及
\ref{spaces-morphisms-lemma-immersions-monomorphisms}。
因此，证明第一段的讨论适用。
设 $X$ 为域 $K$ 上约化、普遍单射且局部分离的代数空间。
特别地，$|X|$ 是单点集，因而 $X$ 拟紧；见《空间的性质》引理
\ref{spaces-properties-lemma-quasi-compact-space}。
由《空间的性质》引理
\ref{spaces-properties-lemma-quasi-compact-affine-cover}，
可找到满平展态射 $U \to X$，其中 $U$ 仿射。考虑概形态射
$$
j :
U \times_X U
\longrightarrow
U \times_{\Spec(K)} U
$$
由于 $X \to \Spec(K)$ 普遍单射，$j$ 是满射；
由于 $X \to \Spec(K)$ 局部分离，$j$ 是浸入。
满浸入是闭浸入，见《概形》引理
\ref{schemes-lemma-immersion-when-closed}。
因此 $R = U \times_X U$ 作为仿射概形的闭子概形而仿射，
特别地，$R$ 拟紧。于是 $X = U/R$ 拟分离，结论由 (2) 得出。
\end{proof}

\begin{remark}
\label{spaces-more-morphisms-remark-weakly-radicial}
设 $X \to Y$ 为代数空间的态射。对根态射的某些应用，
只需要求对每个态射 $\Spec(K) \to Y$（其中 $K$ 为域）：
\begin{enumerate}
\item 空间 $|\Spec(K) \times_Y X|$ 为单点集；
\item 存在单态射 $\Spec(L) \to \Spec(K) \times_Y X$；以及
\item $K \subset L$ 纯不可分。
\end{enumerate}
若以后有需要，我们可以称这样的态射为{\it 弱根态射}。% P09-ERR-0507
例如，若 $X \to Y$ 是满的弱根态射，则映射
$X(k) \to Y(k)$ 对每个代数闭域 $k$ 都是满射。注意，例
\ref{spaces-more-morphisms-example-universally-injective-not-radicial} 中态射的基变换
$X_{\overline{\mathbf{Q}}} \to \Spec(\overline{\mathbf{Q}})$
是弱根态射，但不是根态射。引理
\ref{spaces-more-morphisms-lemma-when-universally-injective-radicial} 的对应结论是：
若 $X \to Y$ 具有性质 ($\beta$) 且普遍单射，则它是弱根态射
（证明略）。
\end{remark}

\begin{lemma}
\label{spaces-more-morphisms-lemma-check-universally-injective}
设 $S$ 为概形，$f : X \to Y$ 为 $S$ 上代数空间的态射。假设
\begin{enumerate}
\item $f$ 局部有限型；
\item 对每个平展态射 $V \to Y$，映射
$|X \times_Y V| \to |V|$ 是单射。
\end{enumerate}
则 $f$ 普遍单射。
\end{lemma}

\begin{proof}
由《空间的态射》引理
\ref{spaces-morphisms-lemma-universally-injective-local}，
该问题在 $Y$ 上是平展局部的。因此可假设 $Y$ 是概形。
于是 $Y$ 特别是良态的；由《良态空间》引理
\ref{decent-spaces-lemma-conditions-on-point-in-fibre-and-qf}，
$f$ 局部拟有限。设 $y \in Y$ 为一点，并令 $X_y$ 为概形论纤维。
假设 $X_y$ 非空。由《域上的空间》引理
\ref{spaces-over-fields-lemma-locally-quasi-finite-over-field}，
$X_y$ 是在 $\kappa(y)$ 上局部拟有限的概形。
由于 $|X_y| \subset |X|$ 是 $|X| \to |Y|$ 在 $y$ 上的纤维，
$X_y$ 有唯一一点 $x$。空间
$X_y \times_{\Spec(\kappa(y))} \Spec(k)$
对任意有限可分扩张 $k/\kappa(y)$ 亦如此，
因为可将 $k$ 实现为位于 $y$ 上方的、某个 $Y$ 上平展概形中的一点的剩余域；
见《态射进阶》引理
\ref{more-morphisms-lemma-realize-prescribed-residue-field-extension-etale}。
于是由《簇》引理
\ref{varieties-lemma-characterize-geometrically-disconnected}，
$X_y$ 几何连通。这蕴含有限扩张 $\kappa(x)/\kappa(y)$
是纯不可分的。

\medskip\noindent
因此（在 $Y$ 为概形的情形），对每个 $y \in Y$，
纤维 $X_y$ 或为空，或
$(X_y)_{red} = \Spec(\kappa(x))$，其中
$\kappa(y) \subset \kappa(x)$ 纯不可分。
故 $f$ 为根态射（略去若干细节），从而由引理
\ref{spaces-more-morphisms-lemma-radicial-implies-universally-injective} 也普遍单射。
\end{proof}




\section{单态射}
\label{spaces-more-morphisms-section-monomorphisms}

\noindent
本节续接《空间的态射》第
\ref{spaces-morphisms-section-monomorphisms} 节。
我们想知道代数空间的每个单态射是否都可表。如果你能证明这一点，
或能给出反例，请发邮件至
\href{mailto:stacks.project@gmail.com}{stacks.project@gmail.com}。
目前已知下列情形成立：
\begin{enumerate}
\item 对局部有限型的单态射成立
（更一般地，由《空间的态射》引理
\ref{spaces-morphisms-lemma-locally-quasi-finite-separated-representable}，
任意分离且局部拟有限的态射都可表；又由《空间的态射》引理
\ref{spaces-morphisms-lemma-monomorphism-loc-finite-type-loc-quasi-finite}，
局部有限型的单态射局部拟有限）；
\item 若靶是若干零维局部环之谱的不交并，则成立
（《良态空间》引理
\ref{decent-spaces-lemma-monomorphism-toward-disjoint-union-dim-0-rings}）；
\item 对平坦单态射成立（见下文）。
\end{enumerate}

\begin{lemma}[David Rydh]
\label{spaces-more-morphisms-lemma-flat-case}
代数空间的平坦单态射可由概形表示。
\end{lemma}

\begin{proof}
设 $f : X \to Y$ 为代数空间的平坦单态射。
要证明 $f$ 可表，必须证明 $X \times_Y V$ 对每个概形
$V$（它映到 $Y$）都是概形。由于成为概形是局部性质（《空间的性质》引理
\ref{spaces-properties-lemma-subscheme}），可假设 $V$ 仿射。
因此可假设 $Y = \Spec(B)$ 是仿射概形。
接着，以一个拟紧开子空间替换 $X$，可假设 $X$ 拟紧。
空间 $X$ 是分离的，因为 $X \to X \times_{\Spec(B)} X$ 是同构。
应用《空间的极限》引理 \ref{spaces-limits-lemma-enough-local}，
可归约到下述情形：$B$ 是局部环，$X \to \Spec(B)$ 是平坦单态射，
并且存在点 $x \in X$ 映到 $\Spec(B)$ 的闭点。
由于分离代数空间的平坦态射可提升一般化（见《良态空间》引理
\ref{decent-spaces-lemma-generalizations-lift-flat}），
$X \to \Spec(B)$ 是满射。因此
$\{X \to \Spec(B)\}$ 是 fpqc 覆盖。
态射 $X \to \Spec(B)$ 沿 $X \to \Spec(B)$ 基变换后成为同构，
故由 fpqc 下降，它本身是同构；见《空间上的下降》引理
\ref{spaces-descent-lemma-descending-property-isomorphism}。
\end{proof}

\noindent
下面是上述引理的一个（某种意义下的）变体。

\begin{lemma}
\label{spaces-more-morphisms-lemma-ui-case}
设 $S$ 为概形，$f : X \to Y$ 为代数空间的拟紧单态射，
并且对每个 $T \to Y$，映射
$$
\mathcal{O}_T \to f_{T,*}\mathcal{O}_{X \times_Y T}
$$
都是单射。则 $f$ 是同构（因而可由概形表示）。
\end{lemma}

\begin{proof}
该问题在 $Y$ 上是平展局部的，故可假设
$Y = \Spec(A)$ 仿射。于是 $X$ 拟紧；可选取仿射概形
$U = \Spec(B)$ 及满平展态射 $U \to X$
（《空间的性质》引理
\ref{spaces-properties-lemma-quasi-compact-affine-cover}）。
注意 $U \times_X U = \Spec(B \otimes_A B)$。
因此，拟凝聚 $\mathcal{O}_X$-模的范畴等价于
模下降数据范畴 $DD_{B/A}$，其中下降数据相对于环同态 $A \to B$。
见《空间的性质》命题
\ref{spaces-properties-proposition-quasi-coherent}、
《下降》定义 \ref{descent-definition-descent-datum-modules} 及
《下降》小节 \ref{descent-subsection-descent-modules-morphisms}。
另一方面，
$$
A \to B
$$
是普遍单射的环同态。事实上，给定一个 $A$-模 $M$，
由引理的假设可知
$A \oplus M \to B \otimes_A (A \oplus M)$ 是单射。
因此，由《下降》定理 \ref{descent-theorem-descent}，
$DD_{B/A}$ 等价于 $A$-模范畴。于是沿
$f : X \to \Spec(A)$ 的拉回给出拟凝聚模范畴的等价。
特别地，$f^*$ 在拟凝聚模上正合，因而 $f$ 平坦
（略去一个小细节）。此外，$f$ 显然满射
（例如，因为 $\Spec(B) \to \Spec(A)$ 满射）。
故 $\{X \to \Spec(A)\}$ 是 fpqc 覆盖。
态射 $X \to \Spec(A)$ 沿 $X \to \Spec(A)$ 基变换后成为同构，
所以由 fpqc 下降，它本身是同构；见《空间上的下降》引理
\ref{spaces-descent-lemma-descending-property-isomorphism}。
\end{proof}

\begin{lemma}
\label{spaces-more-morphisms-lemma-flat-surjective-monomorphism}
代数空间的拟紧、平坦、满的单态射是同构。
\end{lemma}

\begin{proof}
这样的态射满足引理 \ref{spaces-more-morphisms-lemma-ui-case} 的假设。
\end{proof}







\section{浸入的余法层}
\label{spaces-more-morphisms-section-conormal-sheaf}

\noindent
设 $S$ 为概形，$i : Z \to X$ 为 $S$ 上代数空间的闭浸入。
令 $\mathcal{I} \subset \mathcal{O}_X$ 为对应的拟凝聚理想层；
见《空间的态射》引理
\ref{spaces-morphisms-lemma-closed-immersion-ideals}。
考虑短正合列
$$
0 \to \mathcal{I}^2 \to \mathcal{I} \to \mathcal{I}/\mathcal{I}^2 \to 0
$$
，它由 $X$ 上的拟凝聚层组成。
由于层 $\mathcal{I}/\mathcal{I}^2$ 被 $\mathcal{I}$ 零化，
由《空间的态射》引理
\ref{spaces-morphisms-lemma-i-star-equivalence}，它对应于 $Z$ 上的一个层。
这个拟凝聚 $\mathcal{O}_Z$-模称为{\it $Z$ 在 $X$ 中的余法层}。
沿用《空间的态射》第
\ref{spaces-morphisms-section-closed-immersions-quasi-coherent} 节中所述的
记号滥用，通常仍将它记作 $\mathcal{I}/\mathcal{I}^2$。

\medskip\noindent
若 $i : Z \to X$ 是（局部闭）浸入，则把 $i$ 的余法层定义为闭浸入
$i : Z \to X \setminus \partial Z$ 的余法层；见《空间的态射》注
\ref{spaces-morphisms-remark-immersion}。
它通常记作 $\mathcal{I}/\mathcal{I}^2$，其中 $\mathcal{I}$
是闭浸入 $i : Z \to X \setminus \partial Z$ 的理想层。

\begin{definition}
\label{spaces-more-morphisms-definition-conormal-sheaf}
设 $i : Z \to X$ 为浸入。称 $\mathcal{C}_{Z/X}$ 为
{\it $Z$ 在 $X$ 中的余法层}，亦称{\it $i$ 的余法层}；它是上述拟凝聚
$\mathcal{O}_Z$-模 $\mathcal{I}/\mathcal{I}^2$。
\end{definition}

\noindent
在 \cite[IV Definition 16.1.2]{EGA} 中，该层记作
$\mathcal{N}_{Z/X}$。我们不采用这一约定，因为希望把记号
$\mathcal{N}_{Z/X}$ 留给{\it 浸入的法层}。若余法层有限表示，
则法层定义为
$$
\mathcal{N}_{Z/X} =
\SheafHom_{\mathcal{O}_Z}(\mathcal{C}_{Z/X}, \mathcal{O}_Z) =
\SheafHom_{\mathcal{O}_Z}(\mathcal{I}/\mathcal{I}^2, \mathcal{O}_Z)
$$
（否则法层甚至可能不是拟凝聚的）。稍后将回到法层
（此处插入未来引用）。% P09-ERR-0508

\begin{lemma}
\label{spaces-more-morphisms-lemma-etale-conormal}
设 $S$ 为概形，$i : Z \to X$ 为浸入。设
$\varphi : U \to X$ 为平展态射，其中 $U$ 是概形。
令 $Z_U = U \times_X Z$；它是 $U$ 的局部闭子概形。则典范地有
$$
\mathcal{C}_{Z/X}|_{Z_U} = \mathcal{C}_{Z_U/U}
$$
并且该等式关于 $U$ 具有函子性。
\end{lemma}

\begin{proof}
取闭子空间 $T \subset X$，使 $i$ 定义到 $X \setminus T$ 中的闭浸入。
令 $\mathcal{I}$ 为 $X \setminus T$ 上定义 $Z$ 的拟凝聚理想层。
则引理只是说，$\mathcal{I}|_{U \setminus \varphi^{-1}(T)}$
是浸入 $Z_U \to U \setminus \varphi^{-1}(T)$ 的理想层。
这由《空间的态射》引理
\ref{spaces-morphisms-lemma-closed-immersion-ideals} 中
$\mathcal{I}$ 的构造立即可见。
\end{proof}

\begin{lemma}
\label{spaces-more-morphisms-lemma-conormal-functorial}
设 $S$ 为概形，并设
$$
\xymatrix{
Z \ar[r]_i \ar[d]_f & X \ar[d]^g \\
Z' \ar[r]^{i'} & X'
}
$$
为 $S$ 上代数空间的交换图。假设 $i$、$i'$ 为浸入。% P09-ERR-0509
则存在典范的
$\mathcal{O}_Z$-模映射
$$
f^*\mathcal{C}_{Z'/X'}
\longrightarrow
\mathcal{C}_{Z/X}
$$
\end{lemma}

\begin{proof}
先找开子空间 $U' \subset X'$ 和 $U \subset X$，使
$g(U) \subset U'$，且 $i(Z) \subset U$ 和
$i(Z') \subset U'$ 都是闭的（略去存在性的证明）。% P09-ERR-0510
将 $X$ 替换为 $U$，将 $X'$ 替换为 $U'$，可假设 $i$ 与 $i'$
都是闭浸入。令 $\mathcal{I}' \subset \mathcal{O}_{X'}$ 和
$\mathcal{I} \subset \mathcal{O}_X$ 为分别对应于 $i'$ 与 $i$
的拟凝聚理想层；见《空间的态射》引理
\ref{spaces-morphisms-lemma-closed-immersion-ideals}。
考虑复合
$$
g^{-1}\mathcal{I}' \to g^{-1}\mathcal{O}_{X'}
\xrightarrow{g^\sharp} \mathcal{O}_X \to
\mathcal{O}_X/\mathcal{I} = i_*\mathcal{O}_Z
$$
由于 $g(i(Z)) \subset Z'$，可知该复合为零（见《空间的态射》引理
\ref{spaces-morphisms-lemma-closed-immersion-ideals} 中关于分解的陈述）。
因此得到交换图
$$
\xymatrix{
0 \ar[r] &
\mathcal{I} \ar[r] &
\mathcal{O}_X \ar[r] &
i_*\mathcal{O}_Z \ar[r] &
0 \\
0 \ar[r] &
g^{-1}\mathcal{I}' \ar[r] \ar[u] &
g^{-1}\mathcal{O}_{X'} \ar[r] \ar[u] &
g^{-1}i'_*\mathcal{O}_{Z'} \ar[r] \ar[u] &
0
}
$$
由于 $g^{-1}$ 是正合函子，下行正合。由正合性还可看出
$(g^{-1}\mathcal{I}')^2 = g^{-1}((\mathcal{I}')^2)$。
故该图诱导映射
$g^{-1}(\mathcal{I}'/(\mathcal{I}')^2) \to \mathcal{I}/\mathcal{I}^2$。
将其拉回（例如使用 $i^{-1}$）到 $Z$，得到
$i^{-1}g^{-1}(\mathcal{I}'/(\mathcal{I}')^2) \to \mathcal{C}_{Z/X}$。
由于 $i^{-1}g^{-1} = f^{-1}(i')^{-1}$，这给出映射
$f^{-1}\mathcal{C}_{Z'/X'} \to \mathcal{C}_{Z/X}$，
进而诱导所需映射。
\end{proof}

\begin{lemma}
\label{spaces-more-morphisms-lemma-conormal-functorial-more}
设 $S$ 为概形。定义 \ref{spaces-more-morphisms-definition-conormal-sheaf} 的余法层及
引理 \ref{spaces-more-morphisms-lemma-conormal-functorial} 的函子性满足下列性质：% P09-ERR-0511
\begin{enumerate}
\item 若 $Z \to X$ 是 $S$ 上概形的浸入，则该余法层与
《态射》定义 \ref{morphisms-definition-conormal-sheaf} 中的余法层相同。
\item 若引理 \ref{spaces-more-morphisms-lemma-conormal-functorial} 中所有空间都是概形，则映射
$f^*\mathcal{C}_{Z'/X'} \to \mathcal{C}_{Z/X}$
与《态射》引理
\ref{morphisms-lemma-conormal-functorial} 中构造的映射相同。
\item 给定交换图
$$
\xymatrix{
Z \ar[r]_i \ar[d]_f & X \ar[d]^g \\
Z' \ar[r]^{i'} \ar[d]_{f'} & X' \ar[d]^{g'} \\
Z'' \ar[r]^{i''} & X''
}
$$
则映射 $(f' \circ f)^*\mathcal{C}_{Z''/X''} \to \mathcal{C}_{Z/X}$
等于映射 $f^*\mathcal{C}_{Z'/X'} \to \mathcal{C}_{Z/X}$
与下述映射沿 $f$ 的拉回之复合：
$(f')^*\mathcal{C}_{Z''/X''} \to \mathcal{C}_{Z'/X'}$。
\end{enumerate}
\end{lemma}

\begin{proof}
略。注意，(1) 是引理 \ref{spaces-more-morphisms-lemma-etale-conormal} 的特殊情形。
\end{proof}

\begin{lemma}
\label{spaces-more-morphisms-lemma-conormal-functorial-flat}
设 $S$ 为概形，并设
$$
\xymatrix{
Z \ar[r]_i \ar[d]_f & X \ar[d]^g \\
Z' \ar[r]^{i'} & X'
}
$$
为 $S$ 上代数空间的纤维积图。假设 $i$、$i'$ 为浸入。% P09-ERR-0509
则引理 \ref{spaces-more-morphisms-lemma-conormal-functorial} 的典范映射
$f^*\mathcal{C}_{Z'/X'} \to \mathcal{C}_{Z/X}$ 是满射。
若 $g$ 平坦，则它是同构。
\end{lemma}

\begin{proof}
选择交换图
$$
\xymatrix{
U \ar[r] \ar[d] & X \ar[d] \\
U' \ar[r] & X'
}
$$
其中 $U$、$U'$ 是概形，且水平箭头满且平展；见《空间》引理
\ref{spaces-lemma-lift-morphism-presentations}。
利用引理 \ref{spaces-more-morphisms-lemma-etale-conormal} 和
\ref{spaces-more-morphisms-lemma-conormal-functorial-more}，问题归约为概形态射的情形。
对概形而言，这就是《态射》引理
\ref{morphisms-lemma-conormal-functorial-flat}。
\end{proof}

\begin{lemma}
\label{spaces-more-morphisms-lemma-transitivity-conormal}
设 $S$ 为概形，$Z \to Y \to X$ 为代数空间的浸入。
则存在典范正合列
$$
i^*\mathcal{C}_{Y/X} \to
\mathcal{C}_{Z/X} \to
\mathcal{C}_{Z/Y} \to 0
$$
其中各映射来自引理 \ref{spaces-more-morphisms-lemma-conormal-functorial}，
而 $i : Z \to Y$ 是第一个态射。
\end{lemma}

\begin{proof}
设 $U$ 为概形，$U \to X$ 为满平展态射。
借助引理 \ref{spaces-more-morphisms-lemma-etale-conormal} 和
\ref{spaces-more-morphisms-lemma-conormal-functorial-more}，上述序列的正合性立即转化为
概形浸入
$Z \times_X U \to Y \times_X U \to U$
对应序列的正合性。因此，引理由《态射》引理
\ref{morphisms-lemma-transitivity-conormal} 得出。
\end{proof}








\section{浸入的法锥}
\label{spaces-more-morphisms-section-normal-cone}

\noindent
设 $S$ 为概形，$i : Z \to X$ 为 $S$ 上代数空间的闭浸入。
令 $\mathcal{I} \subset \mathcal{O}_X$ 为对应的拟凝聚理想层；
见《空间的态射》引理
\ref{spaces-morphisms-lemma-closed-immersion-ideals}。
考虑拟凝聚分次 $\mathcal{O}_X$-代数层
$\bigoplus_{n \geq 0} \mathcal{I}^n/\mathcal{I}^{n + 1}$。
由于每个层 $\mathcal{I}^n/\mathcal{I}^{n + 1}$ 都被
$\mathcal{I}$ 零化，由《空间的态射》引理
\ref{spaces-morphisms-lemma-i-star-equivalence}，
该分次代数对应于一个拟凝聚分次 $\mathcal{O}_Z$-代数层。
这个拟凝聚分次 $\mathcal{O}_Z$-代数称为{\it $Z$ 在 $X$ 中的余法代数}。
沿用《空间的态射》第
\ref{spaces-morphisms-section-closed-immersions-quasi-coherent} 节中所述的
记号滥用，通常简称为
$\bigoplus_{n \geq 0} \mathcal{I}^n/\mathcal{I}^{n + 1}$。

\medskip\noindent
若 $i : Z \to X$ 是（局部闭）浸入，则把 $i$ 的余法代数定义为闭浸入
$i : Z \to X \setminus \partial Z$ 的余法代数；见《空间的态射》注
\ref{spaces-morphisms-remark-immersion}。
它通常记作
$\bigoplus_{n \geq 0} \mathcal{I}^n/\mathcal{I}^{n + 1}$，
其中 $\mathcal{I}$ 是闭浸入
$i : Z \to X \setminus \partial Z$ 的理想层。

\begin{definition}
\label{spaces-more-morphisms-definition-conormal-algebra}
设 $i : Z \to X$ 为浸入。称 $\mathcal{C}_{Z/X, *}$ 为
{\it $Z$ 在 $X$ 中的余法代数}，亦称{\it $i$ 的余法代数}；
它是上述拟凝聚分次 $\mathcal{O}_Z$-代数层
$\bigoplus_{n \geq 0} \mathcal{I}^n/\mathcal{I}^{n + 1}$。
\end{definition}

\noindent
因此 $\mathcal{C}_{Z/X, 1} = \mathcal{C}_{Z/X}$ 是该浸入的余法层。
此外，$\mathcal{C}_{Z/X, 0} = \mathcal{O}_Z$，而
$\mathcal{C}_{Z/X, n}$ 是由下述性质刻画的拟凝聚
$\mathcal{O}_Z$-模：
\begin{equation}
\label{spaces-more-morphisms-equation-conormal-in-degree-n}
i_*\mathcal{C}_{Z/X, n} = \mathcal{I}^n/\mathcal{I}^{n + 1}
\end{equation}
其中 $i : Z \to X \setminus \partial Z$，而 $\mathcal{I}$ 是上述
$i$ 的理想层。最后，注意存在典范满射
\begin{equation}
\label{spaces-more-morphisms-equation-conormal-algebra-quotient}
\text{Sym}^*(\mathcal{C}_{Z/X}) \longrightarrow \mathcal{C}_{Z/X, *}
\end{equation}
，它是拟凝聚分次 $\mathcal{O}_Z$-代数的映射，并且在次数
$0$ 和 $1$ 上是同构。

\begin{lemma}
\label{spaces-more-morphisms-lemma-etale-conormal-algebra}
设 $S$ 为概形，$i : Z \to X$ 为 $S$ 上代数空间的浸入。
设 $\varphi : U \to X$ 为平展态射，其中 $U$ 是概形。
令 $Z_U = U \times_X Z$；它是 $U$ 的局部闭子概形。则典范地有
$$
\mathcal{C}_{Z/X, *}|_{Z_U} = \mathcal{C}_{Z_U/U, *}
$$
并且该等式关于 $U$ 具有函子性。
\end{lemma}

\begin{proof}
取闭子空间 $T \subset X$，使 $i$ 定义到 $X \setminus T$ 中的闭浸入。
令 $\mathcal{I}$ 为 $X \setminus T$ 上定义 $Z$ 的拟凝聚理想层。
该引理由如下事实得出：
$\mathcal{I}|_{U \setminus \varphi^{-1}(T)}$ 是浸入
$Z_U \to U \setminus \varphi^{-1}(T)$ 的理想层。
这由《空间的态射》引理
\ref{spaces-morphisms-lemma-closed-immersion-ideals} 中
$\mathcal{I}$ 的构造立即可见。
\end{proof}

\begin{lemma}
\label{spaces-more-morphisms-lemma-conormal-algebra-functorial}
设 $S$ 为概形，并设
$$
\xymatrix{
Z \ar[r]_i \ar[d]_f & X \ar[d]^g \\
Z' \ar[r]^{i'} & X'
}
$$
为 $S$ 上代数空间的交换图。假设 $i$、$i'$ 为浸入。% P09-ERR-0512
则存在典范的分次 $\mathcal{O}_Z$-代数映射
$$
f^*\mathcal{C}_{Z'/X', *}
\longrightarrow
\mathcal{C}_{Z/X, *}
$$
\end{lemma}

\begin{proof}
先找开子空间 $U' \subset X'$ 和 $U \subset X$，使
$g(U) \subset U'$，且 $i(Z) \subset U$ 和
$i(Z') \subset U'$ 都是闭的（略去存在性的证明）。% P09-ERR-0513
将 $X$ 替换为 $U$，将 $X'$ 替换为 $U'$，可假设 $i$ 与 $i'$
都是闭浸入。令 $\mathcal{I}' \subset \mathcal{O}_{X'}$ 和
$\mathcal{I} \subset \mathcal{O}_X$ 为分别对应于 $i'$ 与 $i$
的拟凝聚理想层；见《空间的态射》引理
\ref{spaces-morphisms-lemma-closed-immersion-ideals}。
考虑复合
$$
g^{-1}\mathcal{I}' \to g^{-1}\mathcal{O}_{X'}
\xrightarrow{g^\sharp} \mathcal{O}_X \to
\mathcal{O}_X/\mathcal{I} = i_*\mathcal{O}_Z
$$
由于 $g(i(Z)) \subset Z'$，可知该复合为零（见《空间的态射》引理
\ref{spaces-morphisms-lemma-closed-immersion-ideals} 中关于分解的陈述）。
因此得到交换图
$$
\xymatrix{
0 \ar[r] &
\mathcal{I} \ar[r] &
\mathcal{O}_X \ar[r] &
i_*\mathcal{O}_Z \ar[r] &
0 \\
0 \ar[r] &
g^{-1}\mathcal{I}' \ar[r] \ar[u] &
g^{-1}\mathcal{O}_{X'} \ar[r] \ar[u] &
g^{-1}i'_*\mathcal{O}_{Z'} \ar[r] \ar[u] &
0
}
$$
由于 $g^{-1}$ 是正合函子，下行正合。由正合性还可看出，
$(g^{-1}\mathcal{I}')^n = g^{-1}((\mathcal{I}')^n)$
对所有 $n \geq 1$ 成立。
故该图诱导映射
$g^{-1}((\mathcal{I}')^n/(\mathcal{I}')^{n + 1}) \to
\mathcal{I}^n/\mathcal{I}^{n + 1}$。
将其拉回（例如使用 $i^{-1}$）到 $Z$，得到
$i^{-1}g^{-1}((\mathcal{I}')^n/(\mathcal{I}')^{n + 1}) \to
\mathcal{C}_{Z/X, n}$。
由于 $i^{-1}g^{-1} = f^{-1}(i')^{-1}$，这给出各映射
$f^{-1}\mathcal{C}_{Z'/X', n} \to \mathcal{C}_{Z/X, n}$，
它们诱导所需映射。
\end{proof}

\begin{lemma}
\label{spaces-more-morphisms-lemma-conormal-algebra-functorial-flat}
设 $S$ 为概形，并设
$$
\xymatrix{
Z \ar[r]_i \ar[d]_f & X \ar[d]^g \\
Z' \ar[r]^{i'} & X'
}
$$
为 $S$ 上代数空间的笛卡尔方块，其中
$i$、$i'$ 为浸入。% P09-ERR-0514
则引理 \ref{spaces-more-morphisms-lemma-conormal-algebra-functorial} 的典范映射
$f^*\mathcal{C}_{Z'/X', *} \to \mathcal{C}_{Z/X, *}$ 是满射。
若 $g$ 平坦，则它是同构。
\end{lemma}

\begin{proof}
该陈述可在对 $X'$ 作平展局部化后检验。
在此情形可假设 $X' \to X$ 是概形的态射；% P09-ERR-0515
于是 $Z$ 和 $Z'$ 是概形，结论由概形的情形得出；
见《除子》引理
\ref{divisors-lemma-conormal-algebra-functorial-flat}。
\end{proof}

\noindent
对于代数空间上的锥与向量丛，我们采用与概形情形相同的约定
（概形情形采用 EGA 的约定）；见《构造》第
\ref{constructions-section-cone} 节和第
\ref{constructions-section-vector-bundle} 节。
特别地，向量丛是一个非常一般的对象
（并不要求局部同构于仿射空间丛）。

\begin{definition}
\label{spaces-more-morphisms-definition-normal-cone}
设 $S$ 为概形，$i : Z \to X$ 为 $S$ 上代数空间的浸入。
称{\it 法锥 $C_ZX$} 为 $Z$ 在 $X$ 中的如下锥：
$$
C_ZX = \underline{\Spec}_Z(\mathcal{C}_{Z/X, *})
$$
，见《空间的态射》定义
\ref{spaces-morphisms-definition-relative-spec}。
{\it $Z$ 在 $X$ 中的法丛}是向量丛
$$
N_ZX = \underline{\Spec}_Z(\text{Sym}(\mathcal{C}_{Z/X}))
$$
\end{definition}

\noindent
因此 $C_ZX \to Z$ 是 $Z$ 上的锥，而
$N_ZX \to Z$ 是 $Z$ 上的向量丛。此外，分次代数的典范满射
(\ref{spaces-more-morphisms-equation-conormal-algebra-quotient}) 定义了典范闭浸入
\begin{equation}
\label{spaces-more-morphisms-equation-normal-cone-in-normal-bundle}
C_ZX \longrightarrow N_ZX
\end{equation}
，它是 $Z$ 上锥的态射。










\section{态射的微分层}
\label{spaces-more-morphisms-section-sheaf-differentials}

\noindent
建议读者参阅交换代数章中的相应一节
（《代数》第 \ref{algebra-section-differentials} 节）、
概形态射章中的相应一节
（《态射》第 \ref{morphisms-section-sheaf-differentials} 节），
以及《位点上的模》第
\ref{sites-modules-section-differentials} 节。
我们先证明：概形态射的微分层概念与相应的小平展（带环）位点态射
给出的概念一致。

\medskip\noindent
为清楚陈述下一个引理，我们暂时重新以 $\mathcal{F}^a$
表示一个 $\mathcal{O}_{X_\etale}$-模层，它伴随于拟凝聚
$\mathcal{O}_X$-模 $\mathcal{F}$，而后者定义在概形 $X$ 上；见
《下降》定义 \ref{descent-definition-structure-sheaf}。

\begin{lemma}
\label{spaces-more-morphisms-lemma-match-modules-differentials}
设 $f : X \to Y$ 为概形的态射。设
$f_{small} : X_\etale \to Y_\etale$ 为相伴的小平展位点态射，见
《下降》注 \ref{descent-remark-change-topologies-ringed}。
则存在与泛导子相容的典范同构
$$
(\Omega_{X/Y})^a = \Omega_{X_\etale/Y_\etale}
$$
。这里第一个模是 $X_\etale$ 上的层，它伴随于拟凝聚
$\mathcal{O}_X$-模 $\Omega_{X/Y}$，见
《态射》定义 \ref{morphisms-definition-sheaf-differentials}；
第二个模则来自《位点上的模》定义
\ref{sites-modules-definition-module-differentials}。
\end{lemma}

\begin{proof}
设 $h : U \to X$ 为平展态射。此时自然映射
$h^*\Omega_{X/Y} \to \Omega_{U/Y}$ 是同构，见
《关于态射的更多内容》引理
\ref{more-morphisms-lemma-sheaf-differentials-etale-localization}。
这意味着存在自然的 $\mathcal{O}_{Y_\etale}$-导子
$$
\text{d}^a : \mathcal{O}_{X_\etale} \longrightarrow (\Omega_{X/Y})^a
$$
，因为我们刚刚看到，$(\Omega_{X/Y})^a$ 在任意对象
$U$（它属于 $X_\etale$）上的值典范地等同于
$\Gamma(U, \Omega_{U/Y})$。由
$\text{d}_{X/Y} :
\mathcal{O}_{X_\etale}
\to
\Omega_{X_\etale/Y_\etale}$
的泛性质，存在唯一的 $\mathcal{O}_{X_\etale}$-线性映射
$c : \Omega_{X_\etale/Y_\etale} \to (\Omega_{X/Y})^a$，使得
$\text{d}^a = c \circ \text{d}_{X/Y}$。

\medskip\noindent
反过来，设 $\mathcal{F}$ 为
$\mathcal{O}_{X_\etale}$-模，
而 $D : \mathcal{O}_{X_\etale} \to \mathcal{F}$ 为
$\mathcal{O}_{Y_\etale}$-导子。只需把
$D$ 限制到小 Zariski 位点 $X_{Zar}$，即 $X$ 的相应位点。
由于 $X_{Zar}$ 上的层与 $X$ 上的层一致，见
《下降》注 \ref{descent-remark-Zariski-site-space}，
可见 $D|_{X_{Zar}} : \mathcal{O}_X \to \mathcal{F}|_{X_{Zar}}$
正是一个“通常的” $Y$-导子。因此得到映射
$\psi : \Omega_{X/Y} \longrightarrow \mathcal{F}|_{X_{Zar}}$，
使得 $D|_{X_{Zar}} = \psi \circ \text{d}$。特别地，对
$\mathcal{F} = \Omega_{X_\etale/Y_\etale}$ 应用此结论，得到映射
$$
c' :
\Omega_{X/Y}
\longrightarrow
\Omega_{X_\etale/Y_\etale}|_{X_{Zar}}
$$
考虑《下降》注
\ref{descent-remark-change-topologies-ringed} 和引理
\ref{descent-lemma-compare-sites} 中讨论的带环位点态射
$\text{id}_{small, \etale, Zar} : X_\etale \to X_{Zar}$。
由于限制函子 $\mathcal{F} \mapsto \mathcal{F}|_{X_{Zar}}$
等于 $\text{id}_{small, \etale, Zar, *}$，又由于
$\text{id}_{small, \etale, Zar}^*$ 是
$\text{id}_{small, \etale, Zar, *}$ 的左伴随，且
$(\Omega_{X/Y})^a = \text{id}_{small, \etale, Zar}^*\Omega_{X/Y}$，
可见 $c'$ 伴随于映射
$$
c'' :
(\Omega_{X/Y})^a
\longrightarrow
\Omega_{X_\etale/Y_\etale}.
$$
我们断言 $c''$ 与 $c'$ 互逆。% P09-ERR-0516
该断言即完成引理的证明。
为验证这一点，只需说明：
$c''(\text{d}(f)) = \text{d}_{X/Y}(f)$ 且
$c(\text{d}_{X/Y}(f)) = \text{d}(f)$，其中 $f$ 是
$\mathcal{O}_X$ 在 $X$ 的一个开集上的局部截面。
验证略去。
\end{proof}

\noindent
这样便可作如下定义。另一种做法见注 \ref{spaces-more-morphisms-remark-alternative}。

\begin{definition}
\label{spaces-more-morphisms-definition-sheaf-differentials}
设 $S$ 为概形。设 $f : X \to Y$ 为 $S$ 上代数空间的态射。
{\it 微分层 $\Omega_{X/Y}$}，即 $X$ 在 $Y$ 上的微分层，% P09-ERR-0517
是带环拓扑斯态射
$$
(f_{small}, f^\sharp) :
(X_\etale, \mathcal{O}_X)
\to
(Y_\etale, \mathcal{O}_Y)
$$
的微分层（《位点上的模》定义
\ref{sites-modules-definition-sheaf-differentials}）；该态射见
《空间的性质》引理
\ref{spaces-properties-lemma-morphism-ringed-topoi}。
{\it 泛 $Y$-导子}记为
$\text{d}_{X/Y} : \mathcal{O}_X \to \Omega_{X/Y}$。
\end{definition}

\noindent
由引理 \ref{spaces-more-morphisms-lemma-match-modules-differentials}，
当 $X$ 与 $Y$ 可表时，这一定义不与已有概念冲突。
此后若 $X$ 与 $Y$ 可表，我们不再区分上述微分层与
《态射》定义 \ref{morphisms-definition-sheaf-differentials}
中的微分层。我们希望把它同概形态射通常的微分模联系起来。
关键引理如下。

\begin{lemma}
\label{spaces-more-morphisms-lemma-localize-differentials}
设 $S$ 为概形。设 $f : X \to Y$ 为 $S$ 上代数空间的态射。
考虑任意交换图
$$
\xymatrix{
U \ar[d]_a \ar[r]_\psi & V \ar[d]^b \\
X \ar[r]^f & Y
}
$$
，其中竖直箭头是代数空间的平展态射。则
$$
\Omega_{X/Y}|_{U_\etale} = \Omega_{U/V}
$$
特别地，若 $U$、$V$ 为概形，则它等于概形态射
$U \to V$ 通常的微分层。
\end{lemma}

\begin{proof}
由《空间的性质》引理
\ref{spaces-properties-lemma-etale-morphism-topoi}
和公式 (\ref{spaces-properties-equation-restrict})，
可把 $X_\etale$ 上层到 $U_\etale$ 的限制视为由
$a_{small}$ 所作的拉回。对 $b$ 亦然。由
《位点上的模》引理
\ref{sites-modules-lemma-localize-differentials}，有
$$
\Omega_{X/Y}|_{U_\etale} =
\Omega_{\mathcal{O}_{U_\etale}/
a_{small}^{-1}f_{small}^{-1}\mathcal{O}_{Y_\etale}}
$$
由于
$a_{small}^{-1}f_{small}^{-1}\mathcal{O}_{Y_\etale}
= \psi_{small}^{-1}b_{small}^{-1}\mathcal{O}_{Y_\etale}
= \psi_{small}^{-1}\mathcal{O}_{V_\etale}$，
引理得证。
\end{proof}

\begin{lemma}
\label{spaces-more-morphisms-lemma-module-differentials-quasi-coherent}
设 $S$ 为概形。设 $f : X \to Y$ 为 $S$ 上代数空间的态射。
则 $\Omega_{X/Y}$ 是拟凝聚 $\mathcal{O}_X$-模。
\end{lemma}

\begin{proof}
选择引理 \ref{spaces-more-morphisms-lemma-localize-differentials} 中的一个图，
使 $a$ 与 $b$ 都满，且 $U$ 与 $V$ 都是概形。
于是 $\Omega_{X/Y}|_U = \Omega_{U/V}$，它是拟凝聚的
（例如见《态射》引理
\ref{morphisms-lemma-differentials-diagonal}）。
故由《空间的性质》引理
\ref{spaces-properties-lemma-characterize-quasi-coherent}，
$\Omega_{X/Y}$ 是拟凝聚的。
\end{proof}

\begin{remark}
\label{spaces-more-morphisms-remark-alternative}
既然已知 $\Omega_{X/Y}$ 拟凝聚，就可以尝试用另一种方式构造它。
例如可利用《空间的性质》第
\ref{spaces-properties-section-quasi-coherent-presentation} 节的结果，
通过粘合来构造微分层。
例如，若 $Y$ 是概形，而 $U \to X$ 是从概形到 $X$ 的满平展态射，
则 $\Omega_{U/Y}$ 是拟凝聚 $\mathcal{O}_U$-模；又因
$s, t : R \to U$ 都平展，得到同构
$$
\alpha : s^*\Omega_{U/Y} \to \Omega_{R/Y} \to t^*\Omega_{U/Y}
$$
，这里使用了《态射》引理
\ref{morphisms-lemma-triangle-differentials-smooth}。
可以检验它满足余圈条件，于是构造完成。
若 $Y$ 不是概形，则把 $\Omega_{U/Y}$ 定义为映射
$(U \to Y)^*\Omega_{Y/S} \to \Omega_{U/S}$ 的余核，
再按上述方式进行。这种两步过程略显笨拙。% P09-ERR-0518
另一种可能是对引理 \ref{spaces-more-morphisms-lemma-localize-differentials}
中任意一个图粘合各层 $\Omega_{U/V}$，但这也不很优雅。
不过两种方法都可行，并会给出微分层略为初等的构造。
\end{remark}

\begin{lemma}
\label{spaces-more-morphisms-lemma-functoriality-differentials}
设 $S$ 为概形。设
$$
\xymatrix{
X' \ar[d] \ar[r]_f & X \ar[d] \\
Y' \ar[r] & Y
}
$$
为代数空间的交换图。映射
$f^\sharp : \mathcal{O}_X \to f_*\mathcal{O}_{X'}$
同映射
$f_*\text{d}_{X'/Y'} : f_*\mathcal{O}_{X'} \to f_*\Omega_{X'/Y'}$
的复合是一个 $Y$-导子。因此得到典范的
$\mathcal{O}_X$-模映射
$\Omega_{X/Y} \to f_*\Omega_{X'/Y'}$；再由
$f_*$ 与 $f^*$ 的伴随性，得到典范的
$\mathcal{O}_{X'}$-模同态
$$
c_f : f^*\Omega_{X/Y} \longrightarrow \Omega_{X'/Y'}.
$$
它由下述性质唯一刻画：
$f^*\text{d}_{X/Y}(t)$ 映到
$\text{d}_{X'/Y'}(f^* t)$，
其中 $t$ 是 $\mathcal{O}_X$ 的任意局部截面。
\end{lemma}

\begin{proof}
这是《位点上的模》引理
\ref{sites-modules-lemma-functoriality-differentials-ringed-topoi}
的特例。
\end{proof}

\begin{lemma}
\label{spaces-more-morphisms-lemma-check-functoriality-differentials}
设 $S$ 为概形。设
$$
\xymatrix{
X'' \ar[d] \ar[r]_g & X' \ar[d] \ar[r]_f & X \ar[d] \\
Y'' \ar[r] & Y' \ar[r] & Y
}
$$
为 $S$ 上代数空间的交换图。则有
$$
c_{f \circ g} = c_g \circ g^* c_f
$$
，作为映射 $(f \circ g)^*\Omega_{X/Y} \to \Omega_{X''/Y''}$。
\end{lemma}

\begin{proof}
略。提示：使用 $c_f, c_g, c_{f \circ g}$
对局部截面的作用来刻画这些映射。
\end{proof}

\begin{lemma}
\label{spaces-more-morphisms-lemma-triangle-differentials}
设 $S$ 为概形。
设 $f : X \to Y$、$g : Y \to B$ 为 $S$ 上代数空间的态射。
则存在典范正合列
$$
f^*\Omega_{Y/B} \to \Omega_{X/B} \to \Omega_{X/Y} \to 0
$$
，其中各映射由应用引理
\ref{spaces-more-morphisms-lemma-functoriality-differentials} 得到。
\end{lemma}

\begin{proof}
这由概形情形的相应结果经平展局部化得出；见
《态射》引理 \ref{morphisms-lemma-triangle-differentials}
以及引理 \ref{spaces-more-morphisms-lemma-localize-differentials}。
\end{proof}

\begin{lemma}
\label{spaces-more-morphisms-lemma-immersion-differentials}
设 $S$ 为概形。若 $X \to Y$ 是 $S$ 上代数空间的浸入，
则 $\Omega_{X/S}$ 为零。% P09-ERR-0519
\end{lemma}

\begin{proof}
这由概形情形的相应结果经平展局部化得出；见
《态射》引理 \ref{morphisms-lemma-immersion-differentials}
以及引理 \ref{spaces-more-morphisms-lemma-localize-differentials}。
\end{proof}

\begin{lemma}
\label{spaces-more-morphisms-lemma-differentials-relative-immersion}
设 $S$ 为概形。设 $B$ 为 $S$ 上的代数空间。
设 $i : Z \to X$ 为 $B$ 上代数空间的浸入。
存在典范正合列
$$
\mathcal{C}_{Z/X} \to i^*\Omega_{X/B} \to \Omega_{Z/B} \to 0
$$
，其中第一箭头由 $\text{d}_{X/B}$ 诱导，
第二箭头来自引理 \ref{spaces-more-morphisms-lemma-functoriality-differentials}。
\end{lemma}

\begin{proof}
这是《态射》引理
\ref{morphisms-lemma-differentials-relative-immersion}
的代数空间版本；由该引理经平展局部化即可推出，见引理
\ref{spaces-more-morphisms-lemma-localize-differentials} 和
\ref{spaces-more-morphisms-lemma-etale-conormal}。
不过，我们应当确认第一箭头可在整体上定义。
因此在此说明“由 $\text{d}_{X/B}$ 诱导”的含义。
也就是说，可假设 $i$ 是闭浸入，方法是把 $X$ 替换为一个开子空间。
设 $\mathcal{I} \subset \mathcal{O}_X$
为对应于 $Z \subset X$ 的拟凝聚理想层。
则 $\text{d}_{X/S} : \mathcal{I} \to \Omega_{X/S}$% P09-ERR-0520
把子层 $\mathcal{I}^2 \subset \mathcal{I}$ 映到
$\mathcal{I}\Omega_{X/S}$。因此它诱导映射
$\mathcal{I}/\mathcal{I}^2 \to
\Omega_{X/S}/\mathcal{I}\Omega_{X/S}$，
该映射是 $\mathcal{O}_X/\mathcal{I}$-线性的。
由《空间的态射》引理
\ref{spaces-morphisms-lemma-i-star-equivalence}，
这对应于所需的映射
$\mathcal{C}_{Z/X} \to i^*\Omega_{X/S}$。
\end{proof}

\begin{lemma}
\label{spaces-more-morphisms-lemma-differentials-relative-immersion-section}
设 $S$ 为概形。设 $B$ 为 $S$ 上的代数空间。
设 $i : Z \to X$ 为 $B$ 上代数空间的浸入，并假设
$i$（平展局部地）有左逆。则典范列
$$
0 \to \mathcal{C}_{Z/X} \to i^*\Omega_{X/B} \to \Omega_{Z/B} \to 0
$$
（平展局部地）分裂正合；该列来自引理
\ref{spaces-more-morphisms-lemma-differentials-relative-immersion}。
\end{lemma}

\begin{proof}
说明如下：我们断言，若 $g : X \to Z$ 是 $i$
在 $B$ 上的左逆，则 $i^*c_g$ 是映射
$i^*\Omega_{X/B} \to \Omega_{Z/B}$ 的右逆。
说明这一点后，结论由概形态射的相应结果经平展局部化得出，见引理
\ref{spaces-more-morphisms-lemma-localize-differentials} 和
\ref{spaces-more-morphisms-lemma-etale-conormal}。
\end{proof}

\begin{lemma}
\label{spaces-more-morphisms-lemma-base-change-differentials}
设 $S$ 为概形。
设 $X \to Y$ 为 $S$ 上代数空间的态射。
设 $g : Y' \to Y$ 为 $S$ 上代数空间的态射。
设 $X' = X_{Y'}$ 为 $X$ 的基变换。
以 $g' : X' \to X$ 表示投影。
则映射
$$
(g')^*\Omega_{X/Y} \to \Omega_{X'/Y'}
$$
是同构；该映射来自引理
\ref{spaces-more-morphisms-lemma-functoriality-differentials}。
\end{lemma}

\begin{proof}
这由概形情形的相应结果与平展局部化得出；见
《态射》引理 \ref{morphisms-lemma-base-change-differentials}
以及引理 \ref{spaces-more-morphisms-lemma-localize-differentials}。
\end{proof}

\begin{lemma}
\label{spaces-more-morphisms-lemma-differential-product}
设 $S$ 为概形。
设 $f : X \to B$ 和 $g : Y \to B$ 为 $S$ 上具有相同靶的
代数空间态射。
设 $p : X \times_B Y \to X$ 和 $q : X \times_B Y \to Y$
为投影态射。引理
\ref{spaces-more-morphisms-lemma-functoriality-differentials} 中的映射
$$
p^*\Omega_{X/B} \oplus q^*\Omega_{Y/B}
\longrightarrow
\Omega_{X \times_B Y/B}
$$
给出同构。
\end{lemma}

\begin{proof}
这由概形情形的相应结果与平展局部化得出；见
《态射》引理 \ref{morphisms-lemma-differential-product}
以及引理 \ref{spaces-more-morphisms-lemma-localize-differentials}。
\end{proof}

\begin{lemma}
\label{spaces-more-morphisms-lemma-finite-type-differentials}
设 $S$ 为概形。
设 $f : X \to Y$ 为 $S$ 上代数空间的态射。
若 $f$ 局部有限型，则 $\Omega_{X/Y}$
是有限型 $\mathcal{O}_X$-模。
\end{lemma}

\begin{proof}
这由概形情形的相应结果与平展局部化得出；见
《态射》引理 \ref{morphisms-lemma-finite-type-differentials}
以及引理 \ref{spaces-more-morphisms-lemma-localize-differentials}。
\end{proof}

\begin{lemma}
\label{spaces-more-morphisms-lemma-finite-presentation-differentials}
设 $S$ 为概形。
设 $f : X \to Y$ 为 $S$ 上代数空间的态射。
若 $f$ 局部有限表示，则 $\Omega_{X/Y}$
是有限表示的 $\mathcal{O}_X$-模。
\end{lemma}

\begin{proof}
这由概形情形的相应结果与平展局部化得出；见
《态射》引理 \ref{morphisms-lemma-finite-presentation-differentials}
以及引理 \ref{spaces-more-morphisms-lemma-localize-differentials}。
\end{proof}

\begin{lemma}
\label{spaces-more-morphisms-lemma-smooth-omega-finite-locally-free}
设 $S$ 为概形。
设 $f : X \to Y$ 为 $S$ 上代数空间的光滑态射。
则微分模 $\Omega_{X/Y}$ 是有限局部自由的。
\end{lemma}

\begin{proof}
由引理 \ref{spaces-more-morphisms-lemma-localize-differentials}，该陈述在
$X$ 与 $Y$ 上都是平展局部的。
故它由概形情形得出，见《态射》引理
\ref{morphisms-lemma-smooth-omega-finite-locally-free}。
\end{proof}













\section{平展位点的拓扑不变性}
\label{spaces-more-morphisms-section-topological-invariance}

\noindent
我们证明位点 $X_{spaces, \etale}$ 是一个“拓扑不变量”。
于是 $X_\etale$ 也是拓扑不变量；它由
$X_{spaces, \etale}$ 中的可表对象组成。见引理
\ref{spaces-more-morphisms-lemma-topological-invariance}。

\begin{theorem}
\label{spaces-more-morphisms-theorem-topological-invariance}
设 $S$ 为概形。
设 $f : X \to Y$ 为 $S$ 上代数空间的态射。
假设 $f$ 整、普遍单射且满。函子
$$
V \longmapsto V_X = X \times_Y V
$$
定义范畴等价
$Y_{spaces, \etale} \to X_{spaces, \etale}$。
\end{theorem}

\begin{proof}
态射 $f$ 可表且为普遍同胚，见《空间的态射》第
\ref{spaces-morphisms-section-universal-homeomorphisms} 节。

\medskip\noindent
先证明该函子忠实。
设 $V', V$ 为 $Y_{spaces, \etale}$ 的对象，而
$a, b : V' \to V$ 为 $Y$ 上两个不同的态射。
由于 $V', V$ 在 $Y$ 上平展，如下等化子
$$
E =  V' \times_{(a, b), V \times_Y V, \Delta_{V/Y}} V
$$
（即 $a, b$ 的等化子）也在 $Y$ 上平展。因此 $E \to V'$ 是平展单态射
（即开浸入）；它为同构当且仅当它满。
由于 $X \to Y$ 是普遍同胚，可见此条件成立当且仅当
$E_X = V'_X$，即当且仅当 $a_X = b_X$。

\medskip\noindent
其次证明该函子全忠实。
设 $V', V$ 为 $Y_{spaces, \etale}$ 的对象，而
$c : V'_X \to V_X$ 为 $X$ 上的态射。我们要构造
态射 $a : V' \to V$，它定义在 $Y$ 上，并使得 $a_X = c$。
取满平展态射 $a' : V'' \to V'$，使 $V''$ 为分离代数空间。
若能构造态射 $a'' : V'' \to V$，使得
$a''_X = c \circ a'_X$，则两个复合
$$
V'' \times_{V'} V'' \xrightarrow{\text{pr}_i} V'' \xrightarrow{a''} V
$$
由第一段已经证明的函子忠实性而相等。
因此 $a''$ 唯一地经态射 $a : V' \to V$ 分解，
因为作为层的 $V'$ 是 $V''$ 关于等价关系
$V'' \times_{V'} V''$ 的商。因此可以假设 $V'$ 分离。
此时图
$$
\Gamma_c \subset (V' \times_Y V)_X
$$
既开又闭（细节略）。由于 $X \to Y$ 是普遍同胚，存在既开又闭的子空间
$\Gamma \subset V' \times_Y V$，使得 $\Gamma_X = \Gamma_c$。
投影 $\Gamma \to V'$ 是平展态射，其到 $X$ 的基变换是同构。
因此 $\Gamma \to V'$ 平展、普遍单射且满，故由
《空间的态射》引理
\ref{spaces-morphisms-lemma-etale-universally-injective-open}
可知它是同构。于是 $\Gamma$ 正是所需态射
$a : V' \to V$ 的图。

\medskip\noindent
最后证明该函子本质满。
设 $U$ 为 $X_{spaces, \etale}$ 的对象。
需要找到对象 $V$，它属于 $Y_{spaces, \etale}$，
使得 $V_X \cong U$。取满平展态射 $U' \to U$，
使得 $U' \cong V'_X$ 且
$U' \times_U U' \cong V''_X$，其中 $V'', V'$ 是
$Y_{spaces, \etale}$ 的某些对象。
于是由函子的全忠实性，得到态射% P09-ERR-0521
$s, t : V'' \to V'$，其中
$t_X = \text{pr}_0$ 且 $s_X = \text{pr}_1$；这里它们作为态射
$U' \times_U U' \to U'$ 来看。
利用
$(\text{pr}_0, \text{pr}_1) : U' \times_U U' \to U' \times_S U'$
是平展等价关系，以及 $U' \to V'$ 和
$U' \times_U U' \to V''$ 都普遍单射且满，可推出
$(t, s) : V'' \to V' \times_S V'$ 是平展等价关系。
于是商 $V = V'/V''$（见《空间》定理
\ref{spaces-theorem-presentation}）
是代数空间 $V$，它位于 $Y$ 上。存在态射
$V' \to V$，使得 $V'' = V' \times_V V'$。
因而得到态射 $V \to Y$（见《空间上的下降》引理
\ref{spaces-descent-lemma-fpqc-universal-effective-epimorphisms}）。
基变换到 $X$ 后，可见存在态射 $U' \to V_X$
以及相容同构
$U' \times_{V_X} U' = U' \times_U U'$；
这蕴含 $V_X \cong U$（再次使用刚才引用的引理）。

\medskip\noindent
选择概形 $W$ 和满平展态射 $W \to Y$。
选择概形 $U'$ 和满平展态射 $U' \to U \times_X W_X$。
注意，$U'$ 与 $U' \times_U U'$ 都是在 $X$ 上平展的概形，
且它们到 $X$ 的结构态射经概形 $W_X$ 分解。
因此由《平展上同调》定理
\ref{etale-cohomology-theorem-topological-invariance}，
存在概形 $V', V''$，它们在 $W$ 上平展，且到
$W_X$ 的基变换分别同构于 $U'$ 和 $U' \times_U U'$。
这就完成了证明。
\end{proof}

\begin{lemma}
\label{spaces-more-morphisms-lemma-topological-invariance}
在定理 \ref{spaces-more-morphisms-theorem-topological-invariance}
的假设和记号下，范畴等价
$Y_{spaces, \etale} \to X_{spaces, \etale}$
限制为范畴等价
$Y_\etale \to X_\etale$ 和
$Y_{affine, \etale} \to X_{affine, \etale}$。
\end{lemma}

\begin{proof}
这不过是下述陈述：给定对象
$V \in Y_{spaces, \etale}$，$V$ 为（仿射）概形，
当且仅当 $V \times_Y X$ 为（仿射）概形。
由于 $V \times_Y X \to V$ 整、普遍单射且满
（它是 $X \to Y$ 的基变换），这由《空间的极限》引理
\ref{spaces-limits-lemma-integral-universally-bijective-scheme}
和命题 \ref{spaces-limits-proposition-affine} 得出。
\end{proof}

\begin{remark}
\label{spaces-more-morphisms-remark-topological-invariance-etale-site}
\begin{reference}
Lenny Taelman 于 2016 年 5 月 1 日所发电子邮件。
\end{reference}
代数空间的普遍同胚未必可表，见《空间的态射》例
\ref{spaces-morphisms-example-universal-homeomorphism}。
事实上，定理 \ref{spaces-more-morphisms-theorem-topological-invariance}
对于一般普遍同胚并不成立。为说明这一点，设 $k$ 为特征 $0$
的代数闭域，并令
$$
\mathbf{A}^1 \to X \to \mathbf{A}^1
$$
如《空间的态射》例
\ref{spaces-morphisms-example-universal-homeomorphism}
所述。回忆第一个态射平展，并把
$t$ 与 $-t$（其中 $t \in \mathbf{A}^1_k \setminus \{0\}$）
等同；第二个态射则是我们的普遍同胚。
由于 $\mathbf{A}^1_k$ 没有非平凡的连通有限平展覆盖
（因为 $k$ 是特征零的代数闭域；细节略），
只需构造一个非平凡连通有限平展覆盖
$Y \to X$。为此，令 $Y$ 为原点加倍的仿射直线
（《概形》例 \ref{schemes-example-affine-space-zero-doubled}）。
于是 $Y = Y_1 \cup Y_2$，其中
$Y_i = \mathbf{A}^1_k$，并沿
$\mathbf{A}^1_k \setminus \{0\}$ 粘合。
为定义态射 $Y \to X$，使用态射
$$
Y_1 \xrightarrow{1} \mathbf{A}^1_k \to X
\quad\text{且}\quad
Y_2 \xrightarrow{-1} \mathbf{A}^1_k \to X.
$$
它们在 $Y_1 \cap Y_2$ 上粘合，这是由 $X$ 的构造保证的，
因而定义态射 $Y \to X$。事实上，我们断言
$$
\xymatrix{
Y \ar[d] & Y_1 \amalg Y_2 \ar[l] \ar[d] \\
X & \mathbf{A}^1_k \ar[l]
}
$$
是笛卡尔方块。细节略去；例如可使用《群胚》引理
\ref{groupoids-lemma-criterion-fibre-product}。
由于 $\mathbf{A}^1_k \to X$ 平展且满，
这证明 $Y \to X$ 是次数为 $2$ 的有限平展态射，
从而给出所需例子。

\medskip\noindent
还可用如下更简单的论证。概形 $Y$ 上有群
$G = \{+1, -1\}$ 的自由作用，其中 $-1$
通过交换 $Y_1$ 与 $Y_2$ 并改变坐标符号而作用。
于是 $X = Y/G$（见《空间》定义
\ref{spaces-definition-quotient}），故 $Y \to X$ 有限平展。
还可直接证明存在普遍同胚 $X \to \mathbf{A}^1_k$，
方法是在仿射空间上使用 $t \mapsto t^2$。
事实上，这个 $X$ 就是上面的 $X$。
\end{remark}








\section{加厚}
\label{spaces-more-morphisms-section-thickenings}

\noindent
下述术语未必完全标准，但使用起来很方便。

\begin{definition}
\label{spaces-more-morphisms-definition-thickening}
加厚。设 $S$ 为概形。
\begin{enumerate}
\item 称代数空间 $X'$ 是代数空间 $X$ 的一个{\it 加厚}，
若 $X$ 是 $X'$ 的闭子空间且相伴的拓扑空间相等。
\item 称 $X'$ 是 $X$ 的一个{\it 一阶加厚}，
若 $X$ 是 $X'$ 的闭子空间，并且拟凝聚理想层
$\mathcal{I} \subset \mathcal{O}_{X'}$ 定义 $X$ 且平方为零。
\item 称 $X'$ 是 $X$ 的一个{\it 有限阶加厚}，
若 $X$ 是 $X$ 的闭子空间，% P09-ERR-0522
并且拟凝聚理想层
$\mathcal{I} \subset \mathcal{O}_{X'}$ 定义 $X$ 且幂零；也就是说，
存在整数 $n \geq 0$ 使得 $\mathcal{I}^{n + 1} = 0$。
\item 称 $X'$ 为一个{\it $n$ 阶加厚}（关于 $X$），
若 $X$ 是 $X$ 的闭子空间，且
$\mathcal{I}^{n + 1} = 0$，其中
$\mathcal{I} \subset \mathcal{O}_{X'}$ 是定义 $X$ 的拟凝聚理想层。
% P09-ERR-0522
\item 给定两个加厚 $X \subset X'$ 和 $Y \subset Y'$，
若态射 $f' : X' \to Y'$ 满足
$f(X) \subset Y$，则称它为{\it 加厚的态射}；% P09-ERR-0523
也就是说，$f'|_X$ 经闭子空间 $Y$ 分解。在此情形令
$f = f'|_X : X \to Y$，并称
$(f, f') : (X \subset X') \to (Y \subset Y')$
为加厚的态射。
\item 设 $B$ 为代数空间。类似地定义
{\it $B$ 上的加厚}和{\it $B$ 上加厚的态射}。
这表示上述空间 $X, X', Y, Y'$ 都是赋予了到 $B$
的结构态射的代数空间，而态射
$X \to X'$、$Y \to Y'$ 和 $f' : X' \to Y'$
都是 $B$ 上的态射。
\end{enumerate}
\end{definition}

\noindent
基本等价。注意，若 $X \subset X'$ 是加厚，则
$X \to X'$ 整且普遍双射。这蕴含
\begin{equation}
\label{spaces-more-morphisms-equation-equivalence-etale-spaces}
X_{spaces, \etale} = X'_{spaces, \etale}
\end{equation}
，其中等价由拉回函子给出；见定理
\ref{spaces-more-morphisms-theorem-topological-invariance}。
因此可把 $\mathcal{O}_{X'}$ 看作
$X_{spaces, \etale}$ 上的层。于是有局部带环拓扑斯的典范等价
\begin{equation}
\label{spaces-more-morphisms-equation-fundamental-equivalence}
(\Sh(X'_{spaces, \etale}), \mathcal{O}_{X'})
\cong
(\Sh(X_{spaces, \etale}), \mathcal{O}_{X'})
\end{equation}
。下文将经常把它同《空间的性质》定理
\ref{spaces-properties-theorem-fully-faithful}
的全忠实性结果结合使用。% P09-ERR-0525
例如，闭浸入 $i_X : X \to X'$ 对应于满射
$i_X^\sharp : \mathcal{O}_{X'} \to \mathcal{O}_X$。

\medskip\noindent
设 $S$ 为概形，并设 $B$ 为 $S$ 上的代数空间。
设 $(f, f') : (X \subset X') \to (Y \subset Y')$
为 $B$ 上加厚的态射。注意，连续函子图
$$
\xymatrix{
X_{spaces, \etale} &
Y_{spaces, \etale} \ar[l] \\
X'_{spaces, \etale} \ar[u] &
Y'_{spaces, \etale} \ar[u] \ar[l]
}
$$
交换，且竖直箭头是等价。因此
$f_{spaces, \etale}$、$f_{small}$、
$f'_{spaces, \etale}$ 和 $f'_{small}$
都定义同一个拓扑斯态射。因此可把
$$
(f')^\sharp :
f_{spaces, \etale}^{-1}\mathcal{O}_{Y'}
\longrightarrow
\mathcal{O}_{X'}
$$
看作 $\mathcal{O}_B$-代数层的映射，并且它嵌入交换图
$$
\xymatrix{
f_{spaces, \etale}^{-1}\mathcal{O}_Y \ar[r]_-{f^\sharp} \ar[r] &
\mathcal{O}_X \\
f_{spaces, \etale}^{-1}\mathcal{O}_{Y'} \ar[r]^-{(f')^\sharp}
\ar[u]^{i_Y^\sharp} &
\mathcal{O}_{X'} \ar[u]_{i_X^\sharp}
}
$$
。这里 $i_X : X \to X'$ 和 $i_Y : Y \to Y'$
是给定闭浸入的名称。

\begin{lemma}
\label{spaces-more-morphisms-lemma-first-order-thickening-maps}
设 $S$ 为概形。设 $B$ 为 $S$ 上的代数空间。
设 $X \subset X'$ 和 $Y \subset Y'$ 为 $B$ 上代数空间的加厚。
设 $f : X \to Y$ 为 $B$ 上代数空间的态射。给定任意
$\mathcal{O}_B$-代数映射
$$
\alpha : f_{spaces, \etale}^{-1}\mathcal{O}_{Y'} \to \mathcal{O}_{X'}
$$
，使得图
$$
\xymatrix{
f_{spaces, \etale}^{-1}\mathcal{O}_Y \ar[r]_-{f^\sharp} \ar[r] &
\mathcal{O}_X \\
f_{spaces, \etale}^{-1}\mathcal{O}_{Y'} \ar[r]^-\alpha
\ar[u]^{i_Y^\sharp} &
\mathcal{O}_{X'} \ar[u]_{i_X^\sharp}
}
$$
交换，则存在唯一的 $(f, f')$ 的态射% P09-ERR-0524
作为 $B$ 上加厚的态射，使得 $\alpha = (f')^\sharp$。
\end{lemma}

\begin{proof}
为求得 $f'$，由《空间的性质》定理
\ref{spaces-properties-theorem-fully-faithful}，
只需证明带环拓扑斯态射
$$
(f_{spaces, \etale}, \alpha) :
(\Sh(X_{spaces, \etale}), \mathcal{O}_{X'})
\longrightarrow
(\Sh(Y_{spaces, \etale}), \mathcal{O}_{Y'})
$$
是局部带环拓扑斯的态射。这直接来自局部带环拓扑斯态射的定义
（《位点上的模》定义
\ref{sites-modules-definition-morphism-locally-ringed-topoi}）、
$(f, f^\sharp)$ 是局部带环拓扑斯的态射这一事实
（《空间的性质》引理
\ref{spaces-properties-lemma-morphism-locally-ringed}）、
$\alpha$ 嵌入给定交换图这一事实，以及
$i_X^\sharp$ 与 $i_Y^\sharp$ 的核局部幂零这一事实。
最后，$f' \circ i_X = i_Y \circ f$ 来自该图的交换性和
再次应用《空间的性质》定理
\ref{spaces-properties-theorem-fully-faithful}。
关于 $f'$ 是 $B$ 上态射的验证略去。
\end{proof}

\begin{lemma}
\label{spaces-more-morphisms-lemma-open-subspace-thickening}
设 $S$ 为概形。设 $X \subset X'$ 为 $S$ 上代数空间的加厚。
对任意开子空间 $U \subset X$，存在唯一的开子空间
$U' \subset X'$，使得
$U = X \times_{X'} U'$。
\end{lemma}

\begin{proof}
令 $U' \to X'$ 为 $X'_{spaces, \etale}$ 中的对象，
它经由 (\ref{spaces-more-morphisms-equation-equivalence-etale-spaces}) 对应于对象
$U \to X$，后者属于 $X_{spaces, \etale}$。
态射 $U' \to X'$ 平展且普遍单射，因而为开浸入，见
《空间的态射》引理
\ref{spaces-morphisms-lemma-etale-universally-injective-open}。
\end{proof}

\noindent
设 $i_X : X \to X'$ 为代数空间的加厚。
核 $\mathcal{I} = \Ker(i_X^\sharp) \subset \mathcal{O}_{X'}$
的任意局部截面都局部幂零。这对结构层
$\mathcal{O}_{X'}$ 给出一定控制；但从技术上说，
有限阶加厚这一类 $X \subset X'$ 容易处理得多。
具体而言，若 $X'$ 是一个 $n$ 阶加厚，则有过滤
$$
0 \subset \mathcal{I}^n \subset \mathcal{I}^{n - 1} \subset
\ldots \subset \mathcal{I} \subset \mathcal{O}_{X'}
$$
，并可见 $X'$ 被闭子空间过滤：
$$
X = X_0 \subset X_1 \subset \ldots \subset X_{n - 1} \subset X_{n + 1} = X'
$$
，其中每一对 $X_i \subset X_{i + 1}$ 都是 $B$ 上的一阶加厚。
% P09-ERR-0526 P09-ERR-0527
利用简单的归纳论证，许多关于一阶加厚证明的结果都可改述为
有限阶加厚的结果。

\begin{lemma}
\label{spaces-more-morphisms-lemma-first-order-thickening-surjective}
设 $S$ 为概形。设 $X \subset X'$ 为 $S$ 上代数空间的加厚。
设 $U$ 为 $X_{spaces, \etale}$ 的仿射对象。则
$$
\Gamma(U, \mathcal{O}_{X'}) \to \Gamma(U, \mathcal{O}_X)
$$
满；这里通过 (\ref{spaces-more-morphisms-equation-fundamental-equivalence})
把 $\mathcal{O}_{X'}$ 看作 $X_{spaces, \etale}$ 上的层。
\end{lemma}

\begin{proof}
设 $U' \to X'$ 为代数空间的平展态射，使得
$U = X \times_{X'} U'$；见定理
\ref{spaces-more-morphisms-theorem-topological-invariance}。
由《空间的极限》引理 \ref{spaces-limits-lemma-affine}，
$U'$ 是仿射概形。因此
$\Gamma(U, \mathcal{O}_{X'}) = \Gamma(U', \mathcal{O}_{U'}) \to
\Gamma(U, \mathcal{O}_U)$
为满射，因为 $U \to U'$ 是仿射概形的闭浸入。
下面给出有限阶加厚的直接证明；这是实践中最常使用的情形。
\end{proof}

\begin{proof}[有限阶加厚情形的证明]
由上述原理，可假设 $X \subset X'$ 是一阶加厚。
以 $\mathcal{I}$ 表示满射
$\mathcal{O}_{X'} \to \mathcal{O}_X$ 的核。
由于 $\mathcal{I}$ 是拟凝聚
$\mathcal{O}_{X'}$-模，且一阶加厚的定义给出
$\mathcal{I}^2 = 0$，可应用《空间的态射》引理
\ref{spaces-morphisms-lemma-i-star-equivalence}，
得到 $\mathcal{I}$ 是拟凝聚 $\mathcal{O}_X$-模。
因此，引理由短正合列
$$
0 \to \mathcal{I} \to \mathcal{O}_{X'} \to \mathcal{O}_X \to 0
$$
所关联的上同调长正合列，以及
$H^1_\etale(U, \mathcal{I}) = 0$ 这一事实得出；
后者成立是因为 $\mathcal{I}$ 拟凝聚，见《下降》命题
\ref{descent-proposition-same-cohomology-quasi-coherent}
和《概形的上同调》引理
\ref{coherent-lemma-quasi-coherent-affine-cohomology-zero}。
\end{proof}

\begin{lemma}
\label{spaces-more-morphisms-lemma-thickening-scheme}
设 $S$ 为概形。设 $X \subset X'$ 为 $S$ 上代数空间的加厚。
若 $X$（可由）概形表示，则 $X'$ 亦然。
\end{lemma}

\begin{proof}
注意 $X'_{red} = X_{red}$。因此若 $X$ 是概形，
则 $X'_{red}$ 是概形。于是结论由
《空间的极限》引理
\ref{spaces-limits-lemma-reduction-scheme} 得出。
下面给出有限阶加厚的直接证明；这是实践中最常使用的情形。
\end{proof}

\begin{proof}[有限阶加厚情形的证明]
只需在 $X'$ 是 $X$ 的一阶加厚时证明。
由《空间的性质》引理
\ref{spaces-properties-lemma-subscheme}，
$X'$ 有一个最大的概形开子空间。因此只需证明，
每个点 $x$（在 $|X'| = |X|$ 中）都包含在 $X'$ 的某个概形开子空间中。
使用引理 \ref{spaces-more-morphisms-lemma-open-subspace-thickening}，
可把 $X \subset X'$ 替换为 $U \subset U'$，其中
$x \in U$ 且 $U$ 是仿射概形。因此可假设 $X$ 仿射。
这样就化归到下一段讨论的情形。

\medskip\noindent
假设 $X \subset X'$ 是一阶加厚，其中 $X$ 是仿射概形。
令 $A = \Gamma(X, \mathcal{O}_X)$ 且
$A' = \Gamma(X', \mathcal{O}_{X'})$。
由引理 \ref{spaces-more-morphisms-lemma-first-order-thickening-surjective}，
映射 $A' \to A$ 满。核 $I$ 是平方零理想。
由《空间的性质》引理
\ref{spaces-properties-lemma-morphism-to-affine-scheme}，
得到典范态射 $f : X' \to \Spec(A')$，它嵌入交换图
$$
\xymatrix{
X \ar@{=}[d] \ar[r] &  X' \ar[d]^f \\
\Spec(A) \ar[r] & \Spec(A')
}
$$
。因为水平箭头都是加厚，显然 $f$ 普遍单射且满。
因此只需证明 $f$ 平展；因为此时《空间的态射》引理
\ref{spaces-morphisms-lemma-etale-universally-injective-open}
将蕴含 $f$ 为同构。

\medskip\noindent
为证明 $f$ 平展，选择仿射概形 $U'$ 以及平展态射
$U' \to X'$。只需证明复合
$U' \to X' \to \Spec(A')$ 平展，见《空间的性质》定义
\ref{spaces-properties-definition-etale}。
写 $U' = \Spec(B')$。令 $U = X \times_{X'} U'$。
由于 $U$ 是 $U'$ 的闭子空间，它是闭子概形；故
$U = \Spec(B)$，其中 $B' \to B$ 满。
记 $J = \Ker(B' \to B)$，并注意
$J = \Gamma(U, \mathcal{I})$，其中
$\mathcal{I} = \Ker(\mathcal{O}_{X'} \to \mathcal{O}_X)$
是 $X_{spaces, \etale}$ 上的层，如引理
\ref{spaces-more-morphisms-lemma-first-order-thickening-surjective} 的证明中所述。
态射 $U' \to X' \to \Spec(A')$ 诱导交换图
$$
\xymatrix{
0 \ar[r] &
J \ar[r] &
B' \ar[r] &
B \ar[r] & 0 \\
0 \ar[r] &
I \ar[r] \ar[u] &
A' \ar[r] \ar[u] &
A \ar[r] \ar[u] & 0
}
$$
。现在，由于 $\mathcal{I}$ 是拟凝聚 $\mathcal{O}_X$-模，
有 $\mathcal{I} = (\widetilde I)^a$；记号见《下降》定义
\ref{descent-definition-structure-sheaf}，
这一事实的理由见《下降》命题
\ref{descent-proposition-equivalence-quasi-coherent}。
因此 $J = I \otimes_A B$。
最后，$A \to B$ 平展，因为 $U \to X$ 平展，
而后者是平展态射 $U' \to X'$ 的基变换。
由《代数》引理
\ref{algebra-lemma-lift-etale-infinitesimal}，
可断定 $A' \to B'$ 平展。
\end{proof}

\begin{lemma}
\label{spaces-more-morphisms-lemma-thickening-equivalence}
设 $S$ 为概形。设 $X \subset X'$ 为 $S$ 上代数空间的加厚。
函子
$$
V' \longmapsto V = X \times_{X'} V'
$$
定义范畴等价
$X'_\etale \to X_\etale$。
\end{lemma}

\begin{proof}
函子 $V' \mapsto V$ 定义范畴等价
$X'_{spaces, \etale} \to X_{spaces, \etale}$，见定理
\ref{spaces-more-morphisms-theorem-topological-invariance}。
因此只需证明 $V$ 是概形当且仅当 $V'$ 是概形。
这正是引理 \ref{spaces-more-morphisms-lemma-thickening-scheme} 的内容。
\end{proof}

\noindent
一阶加厚可描述如下。% P09-ERR-0528

\begin{lemma}
\label{spaces-more-morphisms-lemma-first-order-thickening}
设 $S$ 为概形。
设 $f : X \to B$ 为 $S$ 上代数空间的态射。
考虑短正合列
$$
0 \to \mathcal{I} \to \mathcal{A} \to \mathcal{O}_X \to 0
$$
，它由 $X_\etale$ 上的层组成，其中 $\mathcal{A}$ 是
$f^{-1}\mathcal{O}_B$-代数层，
$\mathcal{A} \to \mathcal{O}_X$ 是
$f^{-1}\mathcal{O}_B$-代数层的满射，而
$\mathcal{I}$ 是其核。若
\begin{enumerate}
\item $\mathcal{I}$ 是 $\mathcal{A}$ 中的平方零理想；且
\item $\mathcal{I}$ 作为 $\mathcal{O}_X$-模是拟凝聚的，
\end{enumerate}
则存在一阶加厚 $X \subset X'$（在 $B$ 上），以及同构
$\mathcal{O}_{X'} \to \mathcal{A}$；它是
$f^{-1}\mathcal{O}_B$-代数同构，并与到 $\mathcal{O}_X$ 的满射相容。
\end{lemma}

\begin{proof}
在本证明中，我们重做引理
\ref{spaces-more-morphisms-lemma-first-order-thickening-surjective} 和
\ref{spaces-more-morphisms-lemma-thickening-scheme} 的证明里使用过的一些论证。
先处理 $B = S = \Spec(\mathbf{Z})$ 的情形。
设 $U$ 为仿射概形，并设 $U \to X$ 平展。则
$$
0 \to \mathcal{I}(U) \to \mathcal{A}(U) \to \mathcal{O}_X(U) \to 0
$$
正合，因为 $H^1(U_\etale, \mathcal{I}) = 0$；
这是由于 $\mathcal{I}$ 拟凝聚，见《下降》命题
\ref{descent-proposition-same-cohomology-quasi-coherent}
和《概形的上同调》引理
\ref{coherent-lemma-quasi-coherent-affine-cohomology-zero}。
若 $V \to U$ 是 $X_{spaces, \etale}$ 的仿射对象之间的态射，则
$$
\mathcal{I}(V) = \mathcal{I}(U) \otimes_{\mathcal{O}_X(U)} \mathcal{O}_X(V)
$$
，因为 $\mathcal{I}$ 是拟凝聚 $\mathcal{O}_X$-模，见
《下降》命题
\ref{descent-proposition-equivalence-quasi-coherent}。
因此 $\mathcal{A}(U) \to \mathcal{A}(V)$ 是平展环同态，见
《代数》引理 \ref{algebra-lemma-lift-etale-infinitesimal}。
故可见
$$
U \longmapsto U' = \Spec(\mathcal{A}(U))
$$
是从 $X_{affine, \etale}$ 到以仿射概形为对象、平展态射为态射的范畴的函子。
事实上，我们断言该函子可扩张为函子 $U \mapsto U'$，
其定义域为整个 $X_\etale$。为说明这一点，若 $U$ 是 $X_\etale$ 的对象，注意
$$
0 \to \mathcal{I}|_{U_{Zar}} \to \mathcal{A}|_{U_{Zar}} \to
\mathcal{O}_X|_{U_{Zar}} \to 0
$$
且 $\mathcal{I}|_{U_{Zar}}$ 是 $U$ 上的拟凝聚层，见《下降》命题
\ref{descent-proposition-equivalence-quasi-coherent-functorial}。
因此由《关于态射的更多内容》引理
\ref{more-morphisms-lemma-first-order-thickening}，
得到概形的一阶加厚 $U \subset U'$，使得
$\mathcal{O}_{U'}$ 同构于 $\mathcal{A}|_{U_{Zar}}$。
显然，此构造与上面对仿射对象的构造相容。

\medskip\noindent
选择表示 $X = U/R$，见《空间》定义
\ref{spaces-definition-presentation}，使
$s, t : R \to U$ 定义平展等价关系。
应用上述函子，得到概形中的平展等价关系
$s', t' : R' \to U'$。考虑代数空间
$X' = U'/R'$（见《空间》定理
\ref{spaces-theorem-presentation}）。
态射 $X = U/R \to U'/R' = X'$ 是一阶加厚。
把 $\mathcal{O}_{X'}$ 看作 $X_\etale$ 上的层。
由构造，有同构
$$
\gamma :
\mathcal{O}_{X'}|_{U_\etale}
\longrightarrow
\mathcal{A}|_{U_\etale}
$$
，使得 $s^{-1}\gamma$ 与 $t^{-1}\gamma$ 在 $R_\etale$ 上一致。
因此由《空间的性质》引理
\ref{spaces-properties-lemma-descent-sheaf}，
$\gamma$ 来自唯一的所需同构
$\mathcal{O}_{X'} \to \mathcal{A}$。

\medskip\noindent
为处理一般基代数空间 $B$ 的情形，先按上法把
$X'$ 构造为 $\mathbf{Z}$ 上的代数空间。
然后用同构 $\mathcal{O}_{X'} \to \mathcal{A}$ 定义
$f^{-1}\mathcal{O}_B \to \mathcal{O}_{X'}$。
依引理 \ref{spaces-more-morphisms-lemma-first-order-thickening-maps}，
这定义了态射 $X' \to B$，它与给定态射
$X \to B$ 相容，证明完毕。
\end{proof}

\begin{lemma}
\label{spaces-more-morphisms-lemma-base-change-thickening}
设 $S$ 为概形。设 $Y \subset Y'$ 为 $S$ 上代数空间的加厚。
设 $X' \to Y'$ 为态射，并令 $X = Y \times_{Y'} X'$。
则
$(X \subset X') \to (Y \subset Y')$
是加厚的态射。若 $Y \subset Y'$ 是一阶
（相应地，有限阶）加厚，则 $X \subset X'$ 是一阶
（相应地，有限阶）加厚。
\end{lemma}

\begin{proof}
略。
\end{proof}

\begin{lemma}
\label{spaces-more-morphisms-lemma-composition-thickening}
设 $S$ 为概形。若 $X \subset X'$ 和 $X' \subset X''$
是 $S$ 上代数空间的加厚，则 $X \subset X''$ 亦然。
\end{lemma}

\begin{proof}
略。
\end{proof}

\begin{lemma}
\label{spaces-more-morphisms-lemma-descending-property-thickening}
“是加厚”这一性质是 fpqc 局部的。
一阶加厚亦然。
\end{lemma}

\begin{proof}
该陈述的含义如下：设 $S$ 为概形，并设
$X \to X'$ 为 $S$ 上代数空间的态射。
设 $\{g_i : X'_i \to X'\}$ 为代数空间的 fpqc 覆盖，
使基变换 $X_i \to X'_i$ 对所有 $i$ 都是加厚。
则 $X \to X'$ 是加厚。由于各态射 $g_i$ 联合满，
可知 $X \to X'$ 满。由《空间上的下降》引理
\ref{spaces-descent-lemma-descending-property-closed-immersion}，
可知 $X \to X'$ 是闭浸入。
因此 $X \to X'$ 是加厚。一阶加厚情形的证明略去。
\end{proof}


\section{加厚的态射}
\label{spaces-more-morphisms-section-morphisms-thickenings}

\noindent
若 $(f, f') : (X \subset X') \to (Y \subset Y')$ 是代数空间加厚的态射，
则态射 $f$ 的性质往往也由 $f'$ 继承。这里有若干变体。

\begin{lemma}
\label{spaces-more-morphisms-lemma-thicken-property-morphisms}
设 $S$ 为概形。设
$(f, f') : (X \subset X') \to (Y \subset Y')$
为 $S$ 上代数空间加厚的态射。则
\begin{enumerate}
\item $f$ 是仿射态射当且仅当 $f'$ 是仿射态射；
\item $f$ 是满态射当且仅当 $f'$ 是满态射；
\item $f$ 拟紧当且仅当 $f'$ 拟紧；% P09-ERR-0529
\item $f$ 普遍闭当且仅当 $f'$ 普遍闭；
\item $f$ 是整态射当且仅当 $f'$ 是整态射；
\item $f$（拟）分离当且仅当 $f'$（拟）分离；
\item $f$ 普遍单射当且仅当 $f'$ 普遍单射；
\item $f$ 普遍开当且仅当 $f'$ 普遍开；
\item $f$ 可表当且仅当 $f'$ 可表；且
\item 在此添加更多内容。% P09-ERR-0530
\end{enumerate}
\end{lemma}

\begin{proof}
注意 $Y \to Y'$ 和 $X \to X'$ 都是整的普遍同胚。
这立即推出第 (2)、(3)、(4)、(7) 和 (8) 项。
第 (1) 项由《空间的极限》命题
\ref{spaces-limits-proposition-affine} 得出；该命题告诉我们，
平展于 $X$ 与 $X'$ 的仿射概形之间有一一对应，
平展于 $Y$ 与 $Y'$ 的仿射概形之间亦然。
第 (5) 项由第 (1)、(4) 项及《空间的态射》引理
\ref{spaces-morphisms-lemma-integral-universally-closed} 得出。
最后，注意
$$
X \times_Y X = X \times_{Y'} X \to X \times_{Y'} X' \to X' \times_{Y'} X'
$$
是加厚（由引理 \ref{spaces-more-morphisms-lemma-base-change-thickening}，两个箭头都是加厚）。
因此，对加厚的态射
$(X \subset X') \to (X \times_Y X \to X' \times_{Y'} X')$
应用第 (3)、(4) 项，即得第 (6) 项。% P09-ERR-0531
最后，第 (9) 项来自如下事实：概形的代数空间加厚仍为概形，见引理
\ref{spaces-more-morphisms-lemma-thickening-scheme}。
\end{proof}

\begin{lemma}
\label{spaces-more-morphisms-lemma-thicken-property-morphisms-cartesian}
设 $S$ 为概形。设
$(f, f') : (X \subset X') \to (Y \subset Y')$
为 $S$ 上代数空间加厚的态射，并满足
$X = Y \times_{Y'} X'$。若 $X \subset X'$ 是有限阶加厚，则
\begin{enumerate}
\item $f$ 是闭浸入当且仅当 $f'$ 是闭浸入；
\item $f$ 局部有限型当且仅当 $f'$ 局部有限型；
\item $f$ 局部拟有限当且仅当 $f'$ 局部拟有限；
\item $f$ 局部有限型且相对维数为 $d$ 当且仅当
$f'$ 局部有限型且相对维数为 $d$；
\item $\Omega_{X/Y} = 0$ 当且仅当 $\Omega_{X'/Y'} = 0$；
\item $f$ 非分歧当且仅当 $f'$ 非分歧；
\item $f$ 固有当且仅当 $f'$ 固有；
\item $f$ 是有限态射当且仅当 $f'$ 是有限态射；% P09-ERR-0532
\item $f$ 是单态射当且仅当 $f'$ 是单态射；
\item $f$ 是浸入当且仅当 $f'$ 是浸入；且
\item 在此添加更多内容。% P09-ERR-0533
\end{enumerate}
\end{lemma}

\begin{proof}
选择概形 $V'$ 及满平展态射 $V' \to Y'$。
选择概形 $U'$ 及满平展态射
$U' \to X' \times_{Y'} V'$。令 $V = Y \times_{Y'} V'$ 且
$U = X \times_{X'} U'$。于是，对态射的平展局部性质，
可化归到概形加厚的态射
$(U \subset U') \to (V \subset V')$，并应用《关于态射的更多内容》引理
\ref{more-morphisms-lemma-thicken-property-morphisms-cartesian}。
这证明了第 (2)、(3)、(4)、(5) 和 (6) 项。

\medskip\noindent
第 (1)、(7)、(8)、(9)、(10) 项中的态射性质在基变换下稳定；
因此，若 $f'$ 具有性质 $\mathcal{P}$，则 $f$ 也具有它。见
《空间》引理 \ref{spaces-lemma-base-change-immersions}，以及
《空间的态射》引理
\ref{spaces-morphisms-lemma-base-change-proper}、
\ref{spaces-morphisms-lemma-base-change-integral} 和
\ref{spaces-morphisms-lemma-base-change-monomorphism}。

\medskip\noindent
第 (1)、(7)、(8)、(9)、(10) 项中值得讨论的方向，是假设
$f$ 具有该性质并推出 $f'$ 也具有它。
对加厚的阶数作归纳，可假设 $Y \subset Y'$ 是一阶加厚，
见上面对有限阶加厚的讨论。

\medskip\noindent
第 (1) 项的证明。选择概形 $V'$ 及满平展态射
$V' \to Y'$。令 $V = Y \times_{Y'} V'$、
$U' = X' \times_{Y'} V'$ 且 $U = X \times_Y V$。
于是 $U \to V$ 是闭浸入，这推出 $U$ 是概形，
进而推出 $U'$ 是概形（引理 \ref{spaces-more-morphisms-lemma-thickening-scheme}）。
故可对
$(U \subset U') \to (V \subset V')$
应用概形情形的引理
（《关于态射的更多内容》引理
\ref{more-morphisms-lemma-thicken-property-morphisms-cartesian}），
从而得到结论。

\medskip\noindent
第 (7) 项的证明。把第 (2) 项与引理
\ref{spaces-more-morphisms-lemma-thicken-property-morphisms} 的结果结合，
并使用如下事实：固有等于拟紧 $+$ 分离 $+$ 局部有限型
$+$ 普遍闭。

\medskip\noindent
第 (8) 项的证明。把第 (2) 项与引理
\ref{spaces-more-morphisms-lemma-thicken-property-morphisms} 的结果结合，
并使用如下事实：有限等于整 $+$ 局部有限型
（《态射》引理 \ref{morphisms-lemma-finite-integral}）。

\medskip\noindent
第 (9) 项的证明。由于 $f$ 是单态射，有 $X = X \times_Y X$。
可把目前已证的结果应用于加厚的态射
$(X \subset X') \to (X \times_Y X \subset X' \times_{Y'} X')$。
由第 (1) 项可知 $X' \to X' \times_{Y'} X'$ 是闭浸入。
事实上，它是一阶加厚，因为定义闭浸入
$X' \to X' \times_{Y'} X'$ 的理想包含在理想
$\mathcal{I} \subset \mathcal{O}_{Y'}$ 的拉回中；该理想切出
$Y$（在 $Y'$ 中）。确实，
$X = X \times_Y X = (X' \times_{Y'} X') \times_{Y'} Y$
包含在 $X'$ 中。闭浸入
$\Delta : X' \to X' \times_{Y'} X'$ 的余法层等于
$\Omega_{X'/Y'}$（这是《态射》引理
\ref{morphisms-lemma-differentials-diagonal} 对代数空间的类似陈述；
它可由平展局部化推出，也可由引理
\ref{spaces-more-morphisms-lemma-differentials-relative-immersion-section} 和
\ref{spaces-more-morphisms-lemma-differential-product} 结合推出；略去若干细节）。
因此只需证明 $\Omega_{X'/Y'} = 0$；这由第 (5) 项及
$X/Y$ 的相应陈述得出。

\medskip\noindent
第 (10) 项的证明。若 $f : X \to Y$ 是浸入，则它分解为
$X \to V \to Y$，其中 $V \to Y$ 是开子空间，而
$X \to V$ 是闭浸入，见《空间的态射》评注
\ref{spaces-morphisms-remark-immersion}。
设 $V' \subset Y'$ 为开子空间，使其底拓扑空间
$|V'|$ 与 $|V| \subset |Y| = |Y'|$ 相同。
于是 $X' \to Y'$ 通过 $V'$ 分解；由第 (1) 项可知
$X' \to V'$ 是闭浸入。证明完毕。
\end{proof}

\noindent
下面的引理是前一引理的一个变体。我们不再假设所涉及的加厚为
有限阶加厚（该假设使局部有限型这一性质可从 $f$ 传递到 $f'$），
而是直接假设 $f$ 和 $f'$ 都局部有限型。

\begin{lemma}
\label{spaces-more-morphisms-lemma-properties-that-extend-over-thickenings}
设 $S$ 为概形。设
$(f, f') : (X \subset X') \to (Y \to Y')$
为 $S$ 上代数空间加厚的态射。% P09-ERR-0534
假设 $f$ 和 $f'$ 局部有限型，并且
$X = Y \times_{Y'} X'$。则
\begin{enumerate}
\item $f$ 局部拟有限当且仅当 $f'$ 局部拟有限；
\item $f$ 有限当且仅当 $f'$ 有限；
\item $f$ 是闭浸入当且仅当 $f'$ 是闭浸入；
\item $\Omega_{X/Y} = 0$ 当且仅当 $\Omega_{X'/Y'} = 0$；
\item $f$ 非分歧当且仅当 $f'$ 非分歧；
\item $f$ 是单态射当且仅当 $f'$ 是单态射；
\item $f$ 是浸入当且仅当 $f'$ 是浸入；
\item $f$ 固有当且仅当 $f'$ 固有；且
\item 在此添加更多内容。% P09-ERR-0535
\end{enumerate}
\end{lemma}

\begin{proof}
选择概形 $V'$ 及满平展态射 $V' \to Y'$。
选择概形 $U'$ 及满平展态射
$U' \to X' \times_{Y'} V'$。令 $V = Y \times_{Y'} V'$ 且
$U = X \times_{X'} U'$。于是，对态射的平展局部性质，
可化归到概形加厚的态射
$(U \subset U') \to (V \subset V')$，并应用
《关于态射的更多内容》引理
\ref{more-morphisms-lemma-properties-that-extend-over-thickenings}。
这证明了第 (1)、(4) 和 (5) 项。

\medskip\noindent
第 (2)、(3)、(6)、(7)、(8) 项中的性质在基变换下稳定；
因此，若 $f'$ 具有性质 $\mathcal{P}$，则 $f$ 也具有它。见
《空间》引理 \ref{spaces-lemma-base-change-immersions}，以及
《空间的态射》引理
\ref{spaces-morphisms-lemma-base-change-proper}、
\ref{spaces-morphisms-lemma-base-change-integral} 和
\ref{spaces-morphisms-lemma-base-change-monomorphism}。
因此在每一种情形中，只需证明：若 $f$ 具有所需性质，
则 $f'$ 也具有它。

\medskip\noindent
第 (2) 种情形由引理 \ref{spaces-more-morphisms-lemma-thicken-property-morphisms}
的第 (5) 种情形，以及如下事实得出：有限态射恰为
局部有限型的整态射
（《空间的态射》引理 \ref{spaces-morphisms-lemma-finite-integral}）。

\medskip\noindent
第 (3) 种情形。这立即由《空间的极限》引理
\ref{spaces-limits-lemma-check-closed-infinitesimally} 得出。

\medskip\noindent
第 (6) 项的证明。由于 $f$ 是单态射，有 $X = X \times_Y X$。
可把目前已证的结果应用于加厚的态射
$(X \subset X') \to (X \times_Y X \subset X' \times_{Y'} X')$。
由第 (3) 项可知
$\Delta_{X'/Y'} : X' \to X' \times_{Y'} X'$ 是闭浸入。
事实上，$\Delta_{X'/Y'}$ 诱导双射
$|X'| \to |X' \times_{Y'} X'|$，故 $\Delta_{X'/Y'}$ 是加厚。
另一方面，由《空间的态射》引理
\ref{spaces-morphisms-lemma-diagonal-morphism-finite-type}，
$\Delta_{X'/Y'}$ 局部有限表示。换言之，
$\Delta_{X'/Y'}(X')$ 由有限型拟凝聚理想层
$\mathcal{J} \subset \mathcal{O}_{X' \times_{Y'} X'}$ 切出。
由于第 (5) 项给出 $\Omega_{X'/Y'} = 0$，可知闭浸入
$X' \to X' \times_{Y'} X'$ 的余法层为零。
（闭浸入 $\Delta_{X'/Y'}$ 的余法层等于
$\Omega_{X'/Y'}$；这是《态射》引理
\ref{morphisms-lemma-differentials-diagonal} 对代数空间的类似陈述；
它可由平展局部化推出，也可由引理
\ref{spaces-more-morphisms-lemma-differentials-relative-immersion-section} 和
\ref{spaces-more-morphisms-lemma-differential-product} 结合推出；略去若干细节。）
换言之，$\mathcal{J}/\mathcal{J}^2 = 0$。
这推出 $\Delta_{X'/Y'}$ 是同构，例如可用《代数》引理
\ref{algebra-lemma-ideal-is-squared-union-connected}。

\medskip\noindent
第 (7) 项的证明。若 $f : X \to Y$ 是浸入，则它分解为
$X \to V \to Y$，其中 $V \to Y$ 是开子空间，而
$X \to V$ 是闭浸入，见《空间的态射》评注
\ref{spaces-morphisms-remark-immersion}。
设 $V' \subset Y'$ 为开子空间，使其底拓扑空间
$|V'|$ 与 $|V| \subset |Y| = |Y'|$ 相同。
于是 $X' \to Y'$ 通过 $V'$ 分解；由第 (3) 项可知
$X' \to V'$ 是闭浸入。

\medskip\noindent
第 (8) 种情形由引理 \ref{spaces-more-morphisms-lemma-thicken-property-morphisms}
以及固有态射的定义得出：它们是局部有限型、拟紧、
普遍闭且分离的态射。
\end{proof}


\section{加厚的 Picard 群}
\label{spaces-more-morphisms-section-picard-group-thickening}

\noindent
这里给出一些关于加厚的 Picard 群的材料。

\begin{lemma}
\label{spaces-more-morphisms-lemma-picard-group-first-order-thickening}
设 $S$ 为概形。设 $X \subset X'$ 为 $S$ 上代数空间的一阶加厚，
其理想层为 $\mathcal{I}$。则有典范正合列
$$
\xymatrix{
0 \ar[r] &
H^0(X, \mathcal{I}) \ar[r] &
H^0(X', \mathcal{O}_{X'}^*) \ar[r] &
H^0(X, \mathcal{O}^*_X) \ar `r[d] `d[l] `l[llld] `d[dll] [dll] \\
& H^1(X, \mathcal{I}) \ar[r] &
\Pic(X') \ar[r] &
\Pic(X) \ar `r[d] `d[l] `l[llld] `d[dll] [dll] \\
& H^2(X, \mathcal{I}) \ar[r] & \ldots \ar[r] & \ldots
}
$$
，其中各项均为阿贝尔群。
\end{lemma}

\begin{proof}
回忆 $X_\etale = X'_\etale$，见引理
\ref{spaces-more-morphisms-lemma-thickening-equivalence}，以及更一般地见第
\ref{spaces-more-morphisms-section-thickenings} 节中的讨论。
引理中的序列是与阿贝尔群层的短正合列
$$
0 \to \mathcal{I} \to \mathcal{O}_{X'}^* \to \mathcal{O}_X^* \to 0
$$
相伴的长正合上同调列；该短正合列位于 $X_\etale$ 上，
其中第一个映射把局部截面 $f$（属于 $\mathcal{I}$）
送到可逆截面 $1 + f$（属于 $\mathcal{O}_{X'}$）。
还要使用带环位点的 Picard 群与可逆函数层的一阶上同调群之间的同一，
见《位点上的上同调》引理
\ref{sites-cohomology-lemma-h1-invertible}。
\end{proof}


\section{无穷小邻域}
\label{spaces-more-morphisms-section-first-order-infinitesimal-neighbourhood}

\noindent
有限阶加厚的一种自然构造如下。假设
$i : Z \to X$ 是代数空间的浸入。% P09-ERR-0536
选择开子空间 $U \subset X$，使 $i$ 把 $Z$ 同一为闭子空间
$Z \subset U$（见《空间的态射》评注
\ref{spaces-morphisms-remark-immersion}）。
设 $\mathcal{I} \subset \mathcal{O}_U$ 为定义 $Z$（在 $U$ 中）的
拟凝聚理想层，见《空间的态射》引理
\ref{spaces-morphisms-lemma-closed-immersion-ideals}。
对 $n \geq 1$，可考虑闭子空间 $Z_n \subset U$，
它由拟凝聚理想层 $\mathcal{I}^{n + 1}$ 定义。

\begin{definition}
\label{spaces-more-morphisms-definition-first-order-infinitesimal-neighbourhood}
设 $i : Z \to X$ 为代数空间的浸入。
\begin{enumerate}
\item $Z$ 在 $X$ 中的{\it 一阶无穷小邻域}，是上面描述的
一阶加厚 $Z \subset Z_1$（在 $X$ 上）。
\item {\it $n$ 阶无穷小邻域}，即 $Z$ 在 $X$ 中的这一邻域，
是上面描述的 $n$ 阶加厚 $Z \subset Z_n$（在 $X$ 上）。
\end{enumerate}
\end{definition}

\noindent
该加厚具有如下泛性质（这也消除了上述构造依赖于开集
$U$ 的任何疑虑）。

\begin{lemma}
\label{spaces-more-morphisms-lemma-first-order-infinitesimal-neighbourhood}
设 $i : Z \to X$ 为代数空间的浸入。
\begin{enumerate}
\item 一阶无穷小邻域 $Z'$（即 $Z$ 在 $X$ 中的这一邻域）
具有如下泛性质：
给定任意交换图
$$
\xymatrix{
Z \ar[d]_i & T \ar[l]^a \ar[d] \\
X & T' \ar[l]_b
}
$$
其中 $T \subset T'$ 是 $X$ 上的一阶加厚，则存在唯一的加厚的态射
$(a', a) : (T \subset T') \to (Z \subset Z')$（在 $X$ 上）。% P09-ERR-0537
\item 对 $n \geq 1$，$n$ 阶无穷小邻域 $Z_n$
（即 $Z$ 在 $X$ 中的这一邻域）具有如下泛性质：给定任意交换图
$$
\xymatrix{
Z \ar[d]_i & T \ar[l]^a \ar[d] \\
X & T' \ar[l]_b
}
$$
其中 $T \subset T'$ 是 $n$ 阶加厚（在 $X$ 上），则存在唯一的加厚的态射
$(a', a) : (T \subset T') \to (Z \subset Z_n)$（在 $X$ 上）。% P09-ERR-0537
\end{enumerate}
\end{lemma}

\begin{proof}
只证明第 (1) 项。
设 $U \subset X$ 为构造 $Z'$ 时使用的开子空间；换言之，
该开集把 $Z$ 同一为 $U$ 中由拟凝聚理想层
$\mathcal{I}$ 切出的闭子空间。
由于 $|T| = |T'|$，可见 $|b|(|T'|) \subset |U|$。
因此可把 $b$ 看作到 $U$ 的态射，见《空间的性质》引理
\ref{spaces-properties-lemma-factor-through-open-subspace}。
设 $\mathcal{J} \subset \mathcal{O}_{T'}$ 为切出 $T$ 的
平方零拟凝聚理想层。由图的交换性，有
$b|_T = i \circ a$，其中 $i : Z \to U$ 是闭浸入。
由《空间的态射》引理
\ref{spaces-morphisms-lemma-closed-immersion-ideals}，可知
$b^\sharp(b^{-1}\mathcal{I}) \subset \mathcal{J}$。
由于 $T'$ 是 $T$ 的一阶加厚，有 $\mathcal{J}^2 = 0$，
从而 $b^\sharp(b^{-1}(\mathcal{I}^2)) = 0$。
由《空间的态射》引理
\ref{spaces-morphisms-lemma-closed-immersion-ideals}，
这意味着 $b$ 通过 $Z'$ 分解。令
$a' : T' \to Z'$ 为该分解，即得结论。
\end{proof}

\begin{lemma}
\label{spaces-more-morphisms-lemma-infinitesimal-neighbourhood-conormal}
设 $i : Z \to X$ 为代数空间的浸入。
设 $Z \subset Z'$ 为 $Z$ 在 $X$ 中的一阶无穷小邻域。则图
$$
\xymatrix{
Z \ar[r] \ar[d] & Z' \ar[d] \\
Z \ar[r] & X
}
$$
由引理 \ref{spaces-more-morphisms-lemma-conormal-functorial} 诱导余法层之间的映射
$\mathcal{C}_{Z/X} \to \mathcal{C}_{Z/Z'}$。
该映射是同构。
\end{lemma}

\begin{proof}
这由上述 $Z'$ 的构造显然可见。
\end{proof}


\section{形式光滑、形式平展与形式非分歧的变换}
\label{spaces-more-morphisms-section-formally-smooth-etale-unramified}

\noindent
回忆：环同态 $R \to A$ 称为{\it 形式光滑}、相应地称为
{\it 形式平展}、相应地称为{\it 形式非分歧}
（见《代数》定义 \ref{algebra-definition-formally-smooth}、
相应地定义 \ref{algebra-definition-formally-etale}、
相应地定义 \ref{algebra-definition-formally-unramified}），
如果对每个实线部分交换的图
$$
\xymatrix{
A \ar[r] \ar@{-->}[rd] & B/I \\
R \ar[r] \ar[u] & B \ar[u]
}
$$
（其中 $I \subset B$ 为平方零理想），存在、相应地存在唯一、
相应地至多存在一个虚线箭头使该图交换。
这启发了代数空间态射以及更一般的函子的如下类似概念。

\begin{definition}
\label{spaces-more-morphisms-definition-formally-smooth-etale-unramified}
设 $S$ 为概形。设 $a : F \to G$ 为函子
$F, G : (\Sch/S)_{fppf}^{opp} \to \textit{Sets}$ 之间的变换。
考虑如下实线部分交换的图：
$$
\xymatrix{
F \ar[d]_a & T \ar[d]^i \ar[l] \\
G & T' \ar[l] \ar@{-->}[lu]
}
$$
其中 $T$ 和 $T'$ 为仿射概形，而 $i$ 是由平方零理想定义的闭浸入。
\begin{enumerate}
\item 若对任意上述实线图，均存在使图交换的虚线箭头，
则称 $a$ {\it 形式光滑}\footnote{这只是这里可以采用的一种定义。
另一稍弱的条件，是要求虚线箭头在 $T'$ 上 fppf 局部存在。
在某种意义下，这一较弱概念具有更好的形式性质。}。
\item 若对任意上述实线图，恰有一个使图交换的虚线箭头，
则称 $a$ {\it 形式平展}。
\item 若对任意上述实线图，至多有一个使图交换的虚线箭头，
则称 $a$ {\it 形式非分歧}。
\end{enumerate}
\end{definition}

\begin{lemma}
\label{spaces-more-morphisms-lemma-formally-etale-is-combination}
设 $S$ 为概形。设 $a : F \to G$ 为函子
$F, G : (\Sch/S)_{fppf}^{opp} \to \textit{Sets}$ 之间的变换。
则 $a$ 形式平展当且仅当 $a$ 既形式光滑又形式非分歧。
\end{lemma}

\begin{proof}
这由定义形式地得出。
\end{proof}

\begin{lemma}
\label{spaces-more-morphisms-lemma-composition-formally-smooth-etale-unramified}
关于复合，有：
\begin{enumerate}
\item 函子之间的形式光滑变换的复合仍形式光滑；
\item 函子之间的形式平展变换的复合仍形式平展；
\item 函子之间的形式非分歧变换的复合仍形式非分歧。
\end{enumerate}
\end{lemma}

\begin{proof}
这是形式的。
\end{proof}

\begin{lemma}
\label{spaces-more-morphisms-lemma-base-change-formally-smooth-etale-unramified}
设 $S$ 为包含在 $\Sch_{fppf}$ 中的概形。
设 $F, G, H : (\Sch/S)_{fppf}^{opp} \to \textit{Sets}$。
设 $a : F \to G$、$b : H \to G$ 为函子之间的变换。
考虑纤维积图
$$
\xymatrix{
H \times_{b, G, a} F \ar[r]_-{b'} \ar[d]_{a'} & F \ar[d]^a \\
H \ar[r]^b & G
}
$$
\begin{enumerate}
\item 若 $a$ 形式光滑，则基变换 $a'$ 形式光滑；
\item 若 $a$ 形式平展，则基变换 $a'$ 形式平展；
\item 若 $a$ 形式非分歧，则基变换 $a'$ 形式非分歧。
\end{enumerate}
\end{lemma}

\begin{proof}
这是形式的。
\end{proof}

\begin{lemma}
\label{spaces-more-morphisms-lemma-representable-property-formally-property}
设 $S$ 为概形。
设 $F, G : (\Sch/S)_{fppf}^{opp} \to \textit{Sets}$。
设 $a : F \to G$ 为可表的函子之间的变换。
\begin{enumerate}
\item 若 $a$ 光滑，则 $a$ 形式光滑；
\item 若 $a$ 平展，则 $a$ 形式平展；
\item 若 $a$ 非分歧，则 $a$ 形式非分歧。
\end{enumerate}
\end{lemma}

\begin{proof}
考虑实线部分交换的图
$$
\xymatrix{
F \ar[d]_a & T \ar[d]^i \ar[l] \\
G & T' \ar[l] \ar@{-->}[lu]
}
$$
，如定义 \ref{spaces-more-morphisms-definition-formally-smooth-etale-unramified} 中所述。
则 $F \times_G T'$ 是 $T'$ 上光滑（相应地，平展；相应地，非分歧）的概形。
因此由《关于态射的更多内容》引理
\ref{more-morphisms-lemma-smooth-formally-smooth}
（相应地，引理 \ref{more-morphisms-lemma-etale-formally-etale}；
相应地，引理 \ref{more-morphisms-lemma-unramified-formally-unramified}），
可以补入（相应地，唯一地补入；相应地，至多以一种方式补入）
下图中的虚线箭头：
$$
\xymatrix{
F \times_G T' \ar[d] & T \ar[d]^i \ar[l] \\
T' & T' \ar[l] \ar@{-->}[lu]
}
$$
，并因而也得到第一个图中的相应断言。% P09-ERR-0538
\end{proof}

\begin{lemma}
\label{spaces-more-morphisms-lemma-etale-on-top}
设 $S$ 为包含在 $\Sch_{fppf}$ 中的概形。
设 $F, G, H : (\Sch/S)_{fppf}^{opp} \to \textit{Sets}$。
设 $a : F \to G$、$b : G \to H$ 为函子之间的变换。
假设 $a$ 可表、满且平展。
\begin{enumerate}
\item 若 $b$ 形式光滑，则 $b \circ a$ 形式光滑；
\item 若 $b$ 形式平展，则 $b \circ a$ 形式平展；
\item 若 $b$ 形式非分歧，则 $b \circ a$ 形式非分歧。
\end{enumerate}
反过来，考虑实线部分交换的图
$$
\xymatrix{
G \ar[d]_b & T \ar[d]^i \ar[l] \\
H & T' \ar[l] \ar@{-->}[lu]
}
$$
，其中 $T'$ 是 $S$ 上的仿射概形，而
$i : T \to T'$ 是由平方零理想定义的闭浸入。
\begin{enumerate}
\item[(4)] 若 $b \circ a$ 形式光滑，则对每个 $t \in T$，
存在仿射概形之间的平展态射 $U' \to T'$ 及态射
$U' \to G$，使得图
$$
\xymatrix{
G \ar[d]_b & T \ar[l] & T \times_{T'} U' \ar[d] \ar[l]\\
H & T' \ar[l] & U' \ar[llu] \ar[l]
}
$$
交换，且 $t$ 位于 $U' \to T'$ 的像中。
\item[(5)] 若 $b \circ a$ 形式非分歧，则上述图中至多有一个
虚线箭头；即 $b$ 形式非分歧。
\item[(6)] 若 $b \circ a$ 形式平展，则上述图中恰有一个
虚线箭头；即 $b$ 形式平展。
\end{enumerate}
\end{lemma}

\begin{proof}
假设 $b$ 形式光滑（相应地，形式平展；相应地，形式非分歧）。
由于平展态射既光滑又非分歧，可知 $a$ 可表且光滑
（相应地，平展；相应地，非分歧）。因此第 (1)、(2)、(3) 项
由引理 \ref{spaces-more-morphisms-lemma-representable-property-formally-property} 与引理
\ref{spaces-more-morphisms-lemma-composition-formally-smooth-etale-unramified} 结合得出。

\medskip\noindent
假设 $b \circ a$ 形式光滑。考虑引理陈述中的图。令
$W = F \times_G T$。由假设，$W$ 是 $T$ 上满平展的概形。
由《平展态射》定理 \ref{etale-theorem-remarkable-equivalence}，
存在概形 $W'$，它平展于 $T'$，使得
$W = T \times_{T'} W'$。选择仿射开子概形
$U' \subset W'$，使 $t$ 位于 $U' \to T'$ 的像中。
因为 $b \circ a$ 形式光滑，可知存在态射 $U' \to F$，% P09-ERR-0539
使图
$$
\xymatrix{
F \ar[d]_{b \circ a} & W \ar[l] & T \times_{T'} U' \ar[d] \ar[l]\\
H & T' \ar[l] & U' \ar[llu] \ar[l]
}
$$
交换。取复合 $U' \to F \to G$，便得到引理第 (5) 项中的映射。
% P09-ERR-0540

\medskip\noindent
假设 $f, g : T' \to G$ 是两个适合引理中图的虚线箭头。
令 $W = F \times_G T$。由假设，$W$ 是 $T$ 上满平展的概形。
由《平展态射》定理 \ref{etale-theorem-remarkable-equivalence}，
存在概形 $W'$，它平展于 $T'$，使得
$W = T \times_{T'} W'$。由于 $a$ 形式平展，复合
$$
W' \to T' \xrightarrow{f} G
\quad\text{和}\quad
W' \to T' \xrightarrow{g} G
$$
可提升为态射 $f', g' : W' \to F$
（在仿射开集上提升，并由唯一性粘合）。
现在若 $b \circ a : F \to H$ 形式非分歧，则
$f' = g'$，故 $f = g$，因为 $W' \to T'$ 是平展覆盖。
这证明了引理的第 (6) 项。% P09-ERR-0541

\medskip\noindent
假设 $b \circ a$ 形式平展。由第 (4) 项，可在 $T'$ 上平展局部地
找到适合该图的虚线箭头；由第 (5) 项，该虚线箭头唯一。
因此可粘合局部解，得到断言 (6)。略去若干细节。
\end{proof}

\begin{remark}
\label{spaces-more-morphisms-remark-tempting}
在引理 \ref{spaces-more-morphisms-lemma-etale-on-top} 的情形中，很容易猜想
``$b$ 形式光滑'' $\Leftrightarrow$ ``$b \circ a$ 形式光滑''。
然而，这在一般情形下很可能不成立。
\end{remark}

\begin{lemma}
\label{spaces-more-morphisms-lemma-formally-permanence}
设 $S$ 为概形。
设 $F, G, H : (\Sch/S)_{fppf}^{opp} \to \textit{Sets}$。
设 $a : F \to G$、$b : G \to H$ 为函子之间的变换。
假设 $b$ 形式非分歧。
\begin{enumerate}
\item 若 $b \circ a$ 形式非分歧，则 $a$ 形式非分歧；
\item 若 $b \circ a$ 形式平展，则 $a$ 形式平展；
\item 若 $b \circ a$ 形式光滑，则 $a$ 形式光滑。
\end{enumerate}
\end{lemma}

\begin{proof}
设 $T \subset T'$ 为由平方零理想定义的仿射概形闭浸入。
给定 $g' : T' \to G$ 和 $f : T \to F$，使得
$g'|_T = a \circ f$。因为 $b$ 形式非分歧，下列两个集合之间
存在一一对应：
$$
\{f' : T' \to F \mid f = f'|_T\text{ 且 }a \circ f' = g'\}
$$
与
$$
\{f' : T' \to F \mid f = f'|_T\text{ 且 }b \circ a \circ f' = b \circ g'\}.
$$
引理由此形式地得出。
\end{proof}








\section{形式非分歧态射}
\label{spaces-more-morphisms-section-formally-unramified}

\noindent
本节具体说明代数空间态射形式非分歧的含义。

\begin{definition}
\label{spaces-more-morphisms-definition-formally-unramified}
设 $S$ 为概形。设 $f : X \to Y$ 为代数空间的态射，并且它在 $S$ 上。
若它作为函子间的变换是形式非分歧的，则称它为
{\it 形式非分歧的}；其含义见定义
\ref{spaces-more-morphisms-definition-formally-smooth-etale-unramified}。
\end{definition}

\noindent
我们不再重述在函子间形式非分歧变换这一更一般的框架中证明的结果，
参见第 \ref{spaces-more-morphisms-section-formally-smooth-etale-unramified} 节。
这一性质可用相对微分模的消失刻画，见引理
\ref{spaces-more-morphisms-lemma-characterize-formally-unramified}。

\begin{lemma}
\label{spaces-more-morphisms-lemma-formally-unramified}
设 $S$ 为概形。设 $f : X \to Y$ 为 $S$ 上代数空间的态射。
下列条件等价：
\begin{enumerate}
\item $f$ 形式非分歧；
\item 对每个图
$$
\xymatrix{
U \ar[d] \ar[r]_\psi & V \ar[d] \\
X \ar[r]^f & Y
}
$$
其中 $U$、$V$ 为概形且竖直箭头平展，概形态射 $\psi$
是形式非分歧的（参见《态射进阶》定义
\ref{more-morphisms-definition-formally-unramified}）；
\item 对于一个这样的图，若竖直箭头满，则 $\psi$ 形式非分歧。
\end{enumerate}
\end{lemma}

\begin{proof}
假设 $f$ 形式非分歧。由引理
\ref{spaces-more-morphisms-lemma-representable-property-formally-property}，
态射 $U \to X$ 和 $V \to Y$ 形式非分歧。因此，由引理
\ref{spaces-more-morphisms-lemma-composition-formally-smooth-etale-unramified}，
复合 $U \to Y$ 形式非分歧。再由引理
\ref{spaces-more-morphisms-lemma-formally-permanence}，可知 $U \to V$ 形式非分歧。
故 (1) 蕴含 (2)，而 (2) 显然蕴含 (3)。% P09-ERR-0542

\medskip\noindent
假设给定 (3) 中的图。由引理
\ref{spaces-more-morphisms-lemma-representable-property-formally-property}，
态射 $V \to Y$ 形式非分歧。因此，由引理
\ref{spaces-more-morphisms-lemma-composition-formally-smooth-etale-unramified}，
复合 $U \to Y$ 形式非分歧。再由引理
\ref{spaces-more-morphisms-lemma-etale-on-top}，可知 $X \to Y$ 形式非分歧，
即 (1) 成立。
\end{proof}

\begin{lemma}
\label{spaces-more-morphisms-lemma-formally-unramified-not-affine}
设 $S$ 为概形。
若 $f : X \to Y$ 是 $S$ 上代数空间的形式非分歧态射，则对于任意实线交换图
$$
\xymatrix{
X \ar[d]_f & T \ar[d]^i \ar[l] \\
Y & T' \ar[l] \ar@{-->}[lu]
}
$$
其中 $T \subset T'$ 是 $S$ 上代数空间的一阶加厚，至多存在一个
使该图交换的虚线箭头。换言之，在定义
\ref{spaces-more-morphisms-definition-formally-unramified} 中，可去掉 $T$ 为仿射概形的条件。
\end{lemma}

\begin{proof}
这是因为存在满平展态射 $U' \to T'$，其中 $U'$ 是仿射概形的不交并
（见《空间的性质》引理
\ref{spaces-properties-lemma-cover-by-union-affines}），
而态射 $T' \to X$ 由其在 $U'$ 上的限制唯一确定。
\end{proof}

\begin{lemma}
\label{spaces-more-morphisms-lemma-composition-formally-unramified}
形式非分歧态射的复合仍形式非分歧。
\end{lemma}

\begin{proof}
这是形式地成立的。
\end{proof}

\begin{lemma}
\label{spaces-more-morphisms-lemma-base-change-formally-unramified}
形式非分歧态射的基变换仍形式非分歧。
\end{lemma}

\begin{proof}
这是形式地成立的。
\end{proof}

\begin{lemma}
\label{spaces-more-morphisms-lemma-characterize-formally-unramified}
设 $S$ 为概形。设 $f : X \to Y$ 为 $S$ 上代数空间的态射。
下列条件等价：
\begin{enumerate}
\item $f$ 形式非分歧；
\item $\Omega_{X/Y} = 0$。
\end{enumerate}
\end{lemma}

\begin{proof}
这是下列结果的组合：引理
\ref{spaces-more-morphisms-lemma-formally-unramified}、《态射进阶》引理
\ref{more-morphisms-lemma-formally-unramified-differentials}
以及引理 \ref{spaces-more-morphisms-lemma-localize-differentials}。
\end{proof}

\begin{lemma}
\label{spaces-more-morphisms-lemma-unramified-formally-unramified}
设 $S$ 为概形。
设 $f : X \to Y$ 为 $S$ 上代数空间的态射。
下列条件等价：
\begin{enumerate}
\item 态射 $f$ 非分歧；
\item 态射 $f$ 局部有限型且 $\Omega_{X/Y} = 0$；
\item 态射 $f$ 局部有限型且形式非分歧。
\end{enumerate}
\end{lemma}

\begin{proof}
选择图
$$
\xymatrix{
U \ar[d] \ar[r]_\psi & V \ar[d] \\
X \ar[r]^f & Y
}
$$
其中 $U$、$V$ 为概形，且竖直箭头平展并且满。于是有
\begin{align*}
f\text{ 非分歧}
& \Leftrightarrow
\psi\text{ 非分歧} \\
& \Leftrightarrow
\psi\text{ 局部有限型且 }\Omega_{U/V} = 0 \\
& \Leftrightarrow
f\text{ 局部有限型且 }\Omega_{X/Y} = 0 \\
& \Leftrightarrow
f\text{ 局部有限型且形式非分歧}
\end{align*}
% P09-ERR-0543
这里使用了《态射》引理
\ref{morphisms-lemma-unramified-omega-zero}
以及引理 \ref{spaces-more-morphisms-lemma-characterize-formally-unramified}。
\end{proof}

\begin{lemma}
\label{spaces-more-morphisms-lemma-universally-injective-unramified}
设 $S$ 为概形。
设 $f : X \to Y$ 为 $S$ 上代数空间的态射。
下列条件等价：
\begin{enumerate}
\item $f$ 非分歧且为单态射；
\item $f$ 非分歧且普遍单射；
\item $f$ 局部有限型且为单态射；
\item $f$ 普遍单射、局部有限型且形式非分歧。
\end{enumerate}
此外，在这种情形下，$f$ 还可表、分离且局部拟有限。
\end{lemma}

\begin{proof}
由引理 \ref{spaces-more-morphisms-lemma-unramified-formally-unramified}，
形式非分歧且局部有限型等价于非分歧。因此 (4) 等价于 (2)。
单态射当然形式非分歧，故 (3) 蕴含 (4)。% P09-ERR-0544
显然 (1) 蕴含 (3)。最后，若 (2) 成立，则
$\Delta : X \to X \times_Y X$ 既是开浸入
（《空间的态射》引理
\ref{spaces-morphisms-lemma-diagonal-unramified-morphism}），
又是满的（《空间的态射》引理
\ref{spaces-morphisms-lemma-universally-injective}），
因而是同构，即 $f$ 为单态射。由此可见 (2) 蕴含 (1)。
最后，由《空间的态射》诸引理
\ref{spaces-morphisms-lemma-monomorphism-loc-finite-type-loc-quasi-finite}
和
\ref{spaces-morphisms-lemma-locally-quasi-finite-separated-representable}，
可知 $f$ 可表、分离且局部拟有限。
\end{proof}

\begin{lemma}
\label{spaces-more-morphisms-lemma-characterize-closed-immersion}
设 $S$ 为概形。
设 $f : X \to Y$ 为 $S$ 上代数空间的态射。
下列条件等价：
\begin{enumerate}
\item $f$ 是闭浸入；
\item $f$ 普遍闭、非分歧且为单态射；
\item $f$ 普遍闭、非分歧且普遍单射；
\item $f$ 普遍闭、局部有限型且为单态射；
\item $f$ 普遍闭、普遍单射、局部有限型且形式非分歧。
\end{enumerate}
\end{lemma}

\begin{proof}
(2) -- (5) 的等价性立即由引理
\ref{spaces-more-morphisms-lemma-universally-injective-unramified} 得出。
此外，若 (2) -- (5) 成立，则 $f$ 可表。
同样，若 (1) 成立，则 $f$ 可表。
故结论归结为概形情形；见《平展态射》引理
\ref{etale-lemma-characterize-closed-immersion}。
\end{proof}

\section{泛一阶加厚}
\label{spaces-more-morphisms-section-universal-thickening}

\noindent
设 $S$ 为概形。
设 $h : Z \to X$ 为 $S$ 上代数空间的态射。
$Z$ 在 $X$ 上的一个{\it 泛一阶加厚}，是一个一阶加厚
$Z \subset Z'$（在 $X$ 上），使得给定任意一阶加厚
$T \subset T'$（在 $X$ 上）及实线交换图
\begin{equation}
\label{spaces-more-morphisms-equation-universal-first-order-thickening}
\vcenter{
\xymatrix{
& Z \ar[ld] & & T \ar[rd] \ar[ll]^a \\
Z' \ar[rrd] & & & & T' \ar@{..>}[llll]_{a'} \ar[lld]^b \\
 & & X
}
}
\end{equation}
都存在唯一的虚线箭头使该图交换。注意，在这种情形下，
$(a, a') : (T \subset T') \to (Z \subset Z')$
是 $X$ 上加厚之间的态射。因此，若泛一阶加厚存在，
则它在唯一同构意义下唯一。一般而言，泛一阶加厚未必存在；
但若 $h$ 形式非分歧，则它确实存在。在证明这一点之前，
先说明概形范畴中的泛一阶加厚也是代数空间范畴中的泛一阶加厚。

\begin{lemma}
\label{spaces-more-morphisms-lemma-check-universal-first-order-thickening}
设 $S$ 为概形。
设 $h : Z \to X$ 为 $S$ 上代数空间的态射。
设 $Z \subset Z'$ 为 $X$ 上的一阶加厚。
下列条件等价：% P09-ERR-0545
\begin{enumerate}
\item $Z \subset Z'$ 是泛一阶加厚；
\item 对任意图 (\ref{spaces-more-morphisms-equation-universal-first-order-thickening})，
若 $T'$ 为概形，则存在唯一的虚线箭头使该图交换；
\item 对任意图 (\ref{spaces-more-morphisms-equation-universal-first-order-thickening})，
若 $T'$ 为仿射概形，则存在唯一的虚线箭头使该图交换。
\end{enumerate}
\end{lemma}

\begin{proof}
蕴含关系 (1) $\Rightarrow$ (2) $\Rightarrow$ (3) 是形式的。
假设 (3)，并给定任意一个图
(\ref{spaces-more-morphisms-equation-universal-first-order-thickening})。% P09-ERR-0546
选择展示 $T' = U'/R'$，见《空间》定义
\ref{spaces-definition-presentation}。
可假设 $U' = \coprod U'_i$ 是仿射概形的不交并，于是
$R' = U' \times_{T'} U' = \coprod_{i, j} U'_i \times_T' U'_j$。% P09-ERR-0547
对每一对 $(i, j)$，选择仿射开覆盖
$U'_i \times_T' U'_j = \bigcup_k R'_{ijk}$。
用 $U_i, R_{ijk}$ 表示它们与 $T$ 在 $T'$ 上的纤维积。
于是每个 $U_i \subset U'_i$ 和 $R_{ijk} \subset R'_{ijk}$
都是仿射概形的一阶加厚。用 $a_i : U_i \to Z$，相应地用
$a_{ijk} : R_{ijk} \to Z$，表示 $a : T \to Z$ 分别与态射
$U_i \to T$、$R_{ijk} \to T$ 的复合。
将 (3) 应用于 $a_i : U_i \to Z$，得到唯一态射
$a'_i : U'_i \to Z'$。将 (3) 应用于 $a_{ijk}$，可知两个复合
$R'_{ijk} \to R'_i \to Z'$ 和 $R'_{ijk} \to R'_j \to Z'$ 相等。% P09-ERR-0548
因此 $a' = \coprod a'_i : U' = \coprod U'_i \to Z'$
下降到商层 $T' = U'/R'$，结论成立。
\end{proof}

\begin{lemma}
\label{spaces-more-morphisms-lemma-universal-thickening-over-formally-etale}
设 $S$ 为概形。
设 $Z \to Y \to X$ 为 $S$ 上代数空间的态射。
若 $Z \subset Z'$ 是 $Z$ 在 $Y$ 上的泛一阶加厚，且
$Y \to X$ 形式平展，则 $Z \subset Z'$ 是 $Z$ 在 $X$ 上的
泛一阶加厚。
\end{lemma}

\begin{proof}
这是形式的。事实上，由引理
\ref{spaces-more-morphisms-lemma-check-universal-first-order-thickening}，
只需考虑 $T'$ 为仿射概形的实线交换图
(\ref{spaces-more-morphisms-equation-universal-first-order-thickening})。
复合 $T \to Z \to Y$ 唯一地提升为 $T' \to Y$，
因为假设 $Y \to X$ 形式平展。于是 $Z \subset Z'$ 是 $Y$ 上的
泛一阶加厚这一事实给出所需态射 $a' : T' \to Z'$。
\end{proof}

\begin{lemma}
\label{spaces-more-morphisms-lemma-etale-morphism-of-universal-thickenings}
设 $S$ 为概形。
设 $Z \to Y \to X$ 为 $S$ 上代数空间的态射。
假设 $Z \to Y$ 平展。
\begin{enumerate}
\item 若 $Y \subset Y'$ 是 $Y$ 在 $X$ 上的泛一阶加厚，
则唯一平展态射 $Z' \to Y'$，即满足 $Z = Y \times_{Y'} Z'$ 的那个态射，
（见定理 \ref{spaces-more-morphisms-theorem-topological-invariance}）
是 $Z$ 在 $X$ 上的泛一阶加厚。
\item 若 $Z \to Y$ 满，并且
$(Z \subset Z') \to (Y \subset Y')$ 是 $X$ 上一阶加厚之间的
平展态射，而 $Z'$ 是 $Z$ 在 $X$ 上的泛一阶加厚，
则 $Y'$ 是 $Y$ 在 $X$ 上的泛一阶加厚。
\end{enumerate}
\end{lemma}

\begin{proof}
先证 (1)。由引理
\ref{spaces-more-morphisms-lemma-check-universal-first-order-thickening}，
只需考虑 $T'$ 为仿射概形的实线交换图
(\ref{spaces-more-morphisms-equation-universal-first-order-thickening})。
复合 $T \to Z \to Y$ 唯一地提升为 $T' \to Y'$，
因为 $Y'$ 是泛一阶加厚。然后 $Z' \to Y'$ 平展这一事实
意味着（见引理 \ref{spaces-more-morphisms-lemma-representable-property-formally-property}）
$T' \to Y'$ 可提升为所需态射 $a' : T' \to Z'$。

\medskip\noindent
再证 (2)。设 $T \subset T'$ 为 $X$ 上的一阶加厚，
并设 $a : T \to Y$ 为态射。令 $W = T \times_Y Z$，
并以 $c : W \to Z$ 表示投影。% P09-ERR-0549
设 $W' \to T'$ 为满足 $W = T \times_{T'} W'$ 的唯一平展态射，
见定理 \ref{spaces-more-morphisms-theorem-topological-invariance}。
注意，$W' \to T'$ 满，因为 $Z \to Y$ 满。
由假设，得到唯一态射 $c' : W' \to Z'$，它在 $X$ 上，
且其限制是 $c$（在 $W$ 上）。由唯一性，$c'$ 在 $W' \times_{T'} W'$ 上的
两个限制相等（因为 $c$ 在 $W \times_T W$ 上的两个限制相等）。
因此 $c'$ 下降为唯一态射 $a' : T' \to Y'$，结论成立。
\end{proof}

\begin{lemma}
\label{spaces-more-morphisms-lemma-universal-thickening}
设 $S$ 为概形。
设 $h : Z \to X$ 为 $S$ 上代数空间的形式非分歧态射。
存在泛一阶加厚 $Z \subset Z'$，它是 $Z$ 在 $X$ 上的加厚。
\end{lemma}

\begin{proof}
任取交换图
$$
\xymatrix{
V \ar[d] \ar[r] & U \ar[d] \\
Z \ar[r] & X
}
$$
其中 $V$、$U$ 为概形，且竖直箭头平展。
注意，$V \to U$ 是概形的形式非分歧态射，见引理
\ref{spaces-more-morphisms-lemma-formally-unramified}。
合并引理 \ref{spaces-more-morphisms-lemma-check-universal-first-order-thickening}
与《态射进阶》引理
\ref{more-morphisms-lemma-universal-thickening}，
可知存在泛一阶加厚 $V \subset V'$，它是 $V$ 在 $U$ 上的加厚。
由引理 \ref{spaces-more-morphisms-lemma-universal-thickening-over-formally-etale} 第 (1) 项，% P09-ERR-0550
$V'$ 是 $V$ 在 $X$ 上的泛一阶加厚。

\medskip\noindent
固定概形 $U$ 及满平展态射 $U \to X$。上述论证表明，
对于任意平展态射 $V \to Z$，其中 $V$ 为概形，
若 $V \to Z \to X$ 分解通过 $U$，则存在泛一阶加厚
$V \subset V'$，它是 $V$ 在 $X$ 上的加厚（但它不依赖于所选的
$V \to X$ 通过 $U$ 的分解）。现在可以选择 $V$，使得
$V \to Z$ 满平展（见《空间》引理
\ref{spaces-lemma-lift-morphism-presentations}）。
于是 $R = V \times_Z V$ 是一个平展于 $Z$ 的概形，
且 $R \to X$ 也分解通过 $U$。% P09-ERR-0551
因此得到泛一阶加厚 $V \subset V'$ 和 $R \subset R'$，
它们都在 $X$ 上。由于 $V \subset V'$
是泛一阶加厚，两个投影 $s, t : R \to V$ 提升为态射
$s', t': R' \to V'$。由引理
\ref{spaces-more-morphisms-lemma-etale-morphism-of-universal-thickenings}，
因为 $R'$ 是 $R$ 在 $X$ 上的泛一阶加厚，这些态射平展。
于是 $(t', s') : R' \to V'$ 是平展等价关系，% P09-ERR-0552
可以令 $Z' = V'/R'$。由于 $V' \to Z'$ 满平展，且
$v'$ 是 $V$ 在 $X$ 上的泛一阶加厚，% P09-ERR-0553
由引理 \ref{spaces-more-morphisms-lemma-universal-thickening-over-formally-etale} 第 (2) 项，% P09-ERR-0554
可知 $Z'$ 是 $Z$ 在 $X$ 上的泛一阶加厚。
\end{proof}

\begin{definition}
\label{spaces-more-morphisms-definition-universal-thickening}
设 $S$ 为概形。
设 $h : Z \to X$ 为 $S$ 上代数空间的形式非分歧态射。
\begin{enumerate}
\item $Z$ 在 $X$ 上的{\it 泛一阶加厚}，是引理
\ref{spaces-more-morphisms-lemma-universal-thickening} 中构造的加厚 $Z \subset Z'$。
\item $Z$ 在 $X$ 上的{\it 余法层}，是 $Z$ 在其泛一阶加厚
$Z'$ 中的余法层；该加厚定义在 $X$ 上。
\end{enumerate}
在这种情形下，我们常将余法层记为 $\mathcal{C}_{Z/X}$。
\end{definition}

\noindent
因此，有层的短正合列
$$
0 \to \mathcal{C}_{Z/X} \to \mathcal{O}_{Z'} \to \mathcal{O}_Z \to 0
$$
定义在 $Z_\etale$ 上，且 $\mathcal{C}_{Z/X}$ 是拟凝聚
$\mathcal{O}_Z$-模。
下列引理证明，当 $Z \to X$ 是浸入时，本定义与相应定义并不冲突。

\begin{lemma}
\label{spaces-more-morphisms-lemma-immersion-universal-thickening}
设 $S$ 为概形。
设 $i : Z \to X$ 为 $S$ 上代数空间的浸入。则
\begin{enumerate}
\item $i$ 形式非分歧；
\item $Z$ 在 $X$ 上的泛一阶加厚是定义
\ref{spaces-more-morphisms-definition-first-order-infinitesimal-neighbourhood}
中的 $Z$ 在 $X$ 中的一阶无穷小邻域；
\item 定义 \ref{spaces-more-morphisms-definition-conormal-sheaf} 意义下 $i$ 的余法层
与定义 \ref{spaces-more-morphisms-definition-universal-thickening} 意义下 $i$ 的余法层相同。
\end{enumerate}
\end{lemma}

\begin{proof}
代数空间的浸入按定义是可表态射。因此，由《态射》诸引理
\ref{morphisms-lemma-open-immersion-unramified} 和
\ref{morphisms-lemma-closed-immersion-unramified}，
通过《空间》引理
\ref{spaces-lemma-representable-transformations-property-implication}
这一抽象原理，浸入是非分歧的。故由引理
\ref{spaces-more-morphisms-lemma-unramified-formally-unramified}，它形式非分歧。
其余断言由引理
\ref{spaces-more-morphisms-lemma-first-order-infinitesimal-neighbourhood}、
\ref{spaces-more-morphisms-lemma-infinitesimal-neighbourhood-conormal} 与各定义合并得出。
\end{proof}

\begin{lemma}
\label{spaces-more-morphisms-lemma-universal-thickening-unramified}
设 $S$ 为概形。
设 $Z \to X$ 为 $S$ 上代数空间的形式非分歧态射。
则泛一阶加厚 $Z'$ 在 $X$ 上形式非分歧。
\end{lemma}

\begin{proof}
设 $T \subset T'$ 为 $X$ 上仿射概形的一阶加厚。设
$$
\xymatrix{
Z' \ar[d] & T \ar[l]^c \ar[d] \\
X & T' \ar[l] \ar[lu]^{a, b}
}
$$
为交换图。令 $T_0 = c^{-1}(Z) \subset T$，并令
$T'_a = a^{-1}(Z)$（概形论意义下）。
由于 $Z'$ 是 $Z$ 的一阶加厚，可知 $T'$ 是 $T'_a$ 的一阶加厚。
又因 $c = a|_T$，可知 $T_0 = T \cap T'_a$（概形论意义下）。
由于 $T'$ 是 $T$ 的一阶加厚，故 $T'_a$ 是 $T_0$ 的一阶加厚。
现在 $a|_{T'_a}$ 和 $b|_{T'_a}$ 是从 $T'_a$ 到 $Z'$
的态射，并且在 $X$ 上；它们在 $T_0$ 上相同（作为到 $Z$ 的态射）。
因此，由 $Z'$ 的泛性质，得到
$a|_{T'_a} = b|_{T'_a}$。于是 $a$、$b$ 是从一阶加厚
$T'$ 出发的态射；这个加厚是 $T'_a$ 的，并且它们在 $T'_a$
上的限制作为到 $Z$ 的态射相同。% P09-ERR-0555
再次使用 $Z'$ 的泛性质，得到 $a = b$。换言之，
形式非分歧态射的定义性质对 $Z' \to X$ 成立，正如所需。
\end{proof}

\begin{lemma}
\label{spaces-more-morphisms-lemma-universal-thickening-functorial}
设 $S$ 为概形。% P09-ERR-0556
考虑 $S$ 上代数空间的交换图
$$
\xymatrix{
Z \ar[r]_h \ar[d]_f & X \ar[d]^g \\
W \ar[r]^{h'} & Y
}
$$
其中 $h$ 和 $h'$ 形式非分歧。设 $Z \subset Z'$ 是 $Z$ 在 $X$
上的泛一阶加厚，设 $W \subset W'$ 是 $W$ 在 $Y$ 上的泛一阶加厚。
存在加厚之间的典范态射
$(f, f') : (Z, Z') \to (W, W')$，它在 $Y$ 上，
并嵌入下列交换图：
$$
\xymatrix{
& & & Z' \ar[ld] \ar[d]^{f'} \\
Z \ar[rr] \ar[d]_f \ar[rrru] & & X \ar[d] & W' \ar[ld] \\
W \ar[rrru]|!{[rr];[rruu]}\hole \ar[rr] & & Y
}
$$
特别地，加厚之间的态射 $(f, f')$ 诱导余法层之间的态射
$f^*\mathcal{C}_{W/Y} \to \mathcal{C}_{Z/X}$。
\end{lemma}

\begin{proof}
第一项断言由 $W'$ 的泛性质显然成立。
余法层上的诱导映射，是将引理
\ref{spaces-more-morphisms-lemma-conormal-functorial} 应用于
$(Z \subset Z') \to (W \subset W')$ 所得的映射。
\end{proof}

\begin{lemma}
\label{spaces-more-morphisms-lemma-universal-thickening-fibre-product}
设 $S$ 为概形。设
$$
\xymatrix{
Z \ar[r]_h \ar[d]_f & X \ar[d]^g \\
W \ar[r]^{h'} & Y
}
$$
为 $S$ 上代数空间的纤维积图，其中 $h'$ 形式非分歧。
则 $h$ 形式非分歧；并且若 $W \subset W'$ 是 $W$ 在 $Y$ 上的
泛一阶加厚，则
$Z = X \times_Y W \subset X \times_Y W'$ 是 $Z$ 在 $X$ 上的
泛一阶加厚。特别地，引理
\ref{spaces-more-morphisms-lemma-universal-thickening-functorial} 中的典范映射
$f^*\mathcal{C}_{W/Y} \to \mathcal{C}_{Z/X}$ 是满射。
\end{lemma}

\begin{proof}
由引理 \ref{spaces-more-morphisms-lemma-base-change-formally-unramified}，态射 $h$
形式非分歧。显然，$X \times_Y W'$ 是一阶加厚。
由 $W'$ 的泛性质（以及纤维积的映射性质），直接验证可知
它具有该泛性质。关于这为何蕴含余法层之间的映射为满射，
见引理 \ref{spaces-more-morphisms-lemma-conormal-functorial-flat}。
\end{proof}

\begin{lemma}
\label{spaces-more-morphisms-lemma-universal-thickening-fibre-product-flat}
设 $S$ 为概形。设
$$
\xymatrix{
Z \ar[r]_h \ar[d]_f & X \ar[d]^g \\
W \ar[r]^{h'} & Y
}
$$
为 $S$ 上代数空间的纤维积图，其中 $h'$ 形式非分歧且 $g$ 平坦。
在这种情形下，相应的泛一阶加厚之间的映射 $Z' \to W'$ 平坦，
且 $f^*\mathcal{C}_{W/Y} \to \mathcal{C}_{Z/X}$ 是同构。
\end{lemma}

\begin{proof}
平坦性在基变换下保持，见《空间的态射》引理
\ref{spaces-morphisms-lemma-base-change-flat}。
因此，第一项断言由引理
\ref{spaces-more-morphisms-lemma-universal-thickening-fibre-product} 中对 $W'$ 的描述得出。
显然，$X \times_Y W'$ 是一阶加厚。
由 $W'$ 的泛性质（以及纤维积的映射性质），直接验证可知
它具有该泛性质。关于这为何蕴含余法层之间的映射为同构，
见引理 \ref{spaces-more-morphisms-lemma-conormal-functorial-flat}。
\end{proof}

\begin{lemma}
\label{spaces-more-morphisms-lemma-universal-thickening-localize}
取泛一阶加厚的操作与平展局部化交换。更准确地说，
设 $h : Z \to X$ 为基概形 $S$ 上代数空间的形式非分歧态射。
设
$$
\xymatrix{
V \ar[d] \ar[r] & U \ar[d] \\
Z \ar[r] & X
}
$$
为竖直箭头平展的交换图。
设 $Z'$ 为 $Z$ 在 $X$ 上的泛一阶加厚。
则 $V \to U$ 形式非分歧，并且泛一阶加厚 $V'$（即 $V$
在 $U$ 上的那个）平展于 $Z'$。特别地，
$\mathcal{C}_{Z/X}|_V = \mathcal{C}_{V/U}$。
\end{lemma}

\begin{proof}
第一项断言就是引理 \ref{spaces-more-morphisms-lemma-formally-unramified}。
泛一阶加厚的相容性由引理
\ref{spaces-more-morphisms-lemma-universal-thickening-over-formally-etale} 和
\ref{spaces-more-morphisms-lemma-etale-morphism-of-universal-thickenings} 得出。
\end{proof}

\begin{lemma}
\label{spaces-more-morphisms-lemma-differentials-universally-unramified}
设 $S$ 为概形，$B$ 为 $S$ 上的代数空间。
设 $h : Z \to X$ 为 $B$ 上代数空间的形式非分歧态射。
设 $Z \subset Z'$ 为 $Z$ 在 $X$ 上的泛一阶加厚，
其结构态射为 $h' : Z' \to X$。典范映射
$$
\text{d}h' : (h')^*\Omega_{X/B} \to \Omega_{Z'/B}
$$
% P09-ERR-0557
诱导同构
$h^*\Omega_{X/B} \to \Omega_{Z'/B} \otimes \mathcal{O}_Z$。
\end{lemma}

\begin{proof}
映射 $c_{h'}$ 是引理
\ref{spaces-more-morphisms-lemma-functoriality-differentials} 中定义的映射。
若 $i : Z \to Z'$ 是给定闭浸入，则 $i^*c_{h'}$ 是映射
$h^*\Omega_{X/S} \to \Omega_{Z'/S} \otimes \mathcal{O}_Z$。% P09-ERR-0558
通过平展局部化，验证它是同构可归结为概形情形，见引理
\ref{spaces-more-morphisms-lemma-universal-thickening-localize} 和
\ref{spaces-more-morphisms-lemma-localize-differentials}。
在这种情形下，结论就是《态射进阶》引理
\ref{more-morphisms-lemma-differentials-universally-unramified}。
\end{proof}

\begin{lemma}
\label{spaces-more-morphisms-lemma-universally-unramified-differentials-sequence}
设 $S$ 为概形，$B$ 为 $S$ 上的代数空间。
设 $h : Z \to X$ 为 $B$ 上代数空间的形式非分歧态射。
有典范正合列
$$
\mathcal{C}_{Z/X} \to h^*\Omega_{X/B} \to \Omega_{Z/B} \to 0.
$$
第一支箭头由 $\text{d}_{Z'/B}$ 诱导，其中 $Z'$ 是 $Z$ 在 $X$
上的泛一阶邻域。% P09-ERR-0559
\end{lemma}

\begin{proof}
已知有典范正合列
$$
\mathcal{C}_{Z/Z'} \to
\Omega_{Z'/S} \otimes \mathcal{O}_Z \to
\Omega_{Z/S} \to 0.
$$
% P09-ERR-0560
见引理 \ref{spaces-more-morphisms-lemma-differentials-relative-immersion}。
因此，应用引理 \ref{spaces-more-morphisms-lemma-differentials-universally-unramified}
即得结论。
\end{proof}

\begin{lemma}
\label{spaces-more-morphisms-lemma-two-unramified-morphisms}
设 $S$ 为概形。设
$$
\xymatrix{
Z \ar[r]_i \ar[rd]_j & X \ar[d] \\
& Y
}
$$
为 $S$ 上代数空间的交换图，其中 $i$、$j$ 形式非分歧。
则有典范正合列
$$
\mathcal{C}_{Z/Y} \to
\mathcal{C}_{Z/X} \to
i^*\Omega_{X/Y} \to 0
$$
其中第一支箭头来自引理
\ref{spaces-more-morphisms-lemma-universal-thickening-functorial}，
第二支箭头来自引理
\ref{spaces-more-morphisms-lemma-universally-unramified-differentials-sequence}。
\end{lemma}

\begin{proof}
这些映射已经定义，故通过平展局部化，验证该序列正合
可归结为概形情形，见引理
\ref{spaces-more-morphisms-lemma-universal-thickening-localize} 和
\ref{spaces-more-morphisms-lemma-localize-differentials}。
在这种情形下，结论就是《态射进阶》引理
\ref{more-morphisms-lemma-two-unramified-morphisms}。
\end{proof}

\begin{lemma}
\label{spaces-more-morphisms-lemma-transitivity-conormal-unramified}
设 $S$ 为概形。
设 $Z \to Y \to X$ 为 $S$ 上代数空间的形式非分歧态射。
\begin{enumerate}
\item 若 $Z \subset Z'$ 是 $Z$ 在 $X$ 上的泛一阶加厚，
$Y \subset Y'$ 是 $Y$ 在 $X$ 上的泛一阶加厚，则存在态射
$Z' \to Y'$，并且 $Y \times_{Y'} Z'$ 是 $Z$ 在 $Y$ 上的
泛一阶加厚。
\item 有典范正合列
$$
i^*\mathcal{C}_{Y/X} \to
\mathcal{C}_{Z/X} \to
\mathcal{C}_{Z/Y} \to 0
$$
其中映射来自引理 \ref{spaces-more-morphisms-lemma-universal-thickening-functorial}，
而 $i : Z \to Y$ 是第一个态射。
\end{enumerate}
\end{lemma}

\begin{proof}
(1) 中的映射 $h : Z' \to Y'$ 来自引理
\ref{spaces-more-morphisms-lemma-universal-thickening-functorial}。
$Y \times_{Y'} Z'$ 是 $Z$ 在 $Y$ 上的泛一阶加厚这一断言，
由 $Z'$ 和 $Y'$ 的泛性质显然成立。由引理
\ref{spaces-more-morphisms-lemma-transitivity-conormal}，有正合列
$$
(i')^*\mathcal{C}_{Y \times_{Y'} Z'/Z'} \to
\mathcal{C}_{Z/Z'} \to
\mathcal{C}_{Z/Y \times_{Y'} Z'} \to 0
$$
其中 $i' : Z \to Y \times_{Y'} Z'$ 是给定态射。由引理
\ref{spaces-more-morphisms-lemma-conormal-functorial-flat}，存在满射
$h^*\mathcal{C}_{Y/Y'} \to \mathcal{C}_{Y \times_{Y'} Z'/Z'}$。
结合等式
$\mathcal{C}_{Y/Y'} = \mathcal{C}_{Y/X}$、
$\mathcal{C}_{Z/Z'} = \mathcal{C}_{Z/X}$ 以及
$\mathcal{C}_{Z/Y \times_{Y'} Z'} = \mathcal{C}_{Z/Y}$，
即得引理。
\end{proof}

\section{形式平展态射}
\label{spaces-more-morphisms-section-formally-etale}

\noindent
本节具体说明代数空间态射形式平展的含义。

\begin{definition}
\label{spaces-more-morphisms-definition-formally-etale}
设 $S$ 为概形。设 $f : X \to Y$ 为代数空间的态射，并且它在 $S$ 上。
若它作为函子间的变换是形式平展的，则称它为{\it 形式平展的}；
其含义见定义 \ref{spaces-more-morphisms-definition-formally-smooth-etale-unramified}。
\end{definition}

\noindent
我们不再重述在函子间形式平展变换这一更一般的框架中证明的结果，
参见第 \ref{spaces-more-morphisms-section-formally-smooth-etale-unramified} 节。

\begin{lemma}
\label{spaces-more-morphisms-lemma-formally-etale}
设 $S$ 为概形。设 $f : X \to Y$ 为代数空间的态射，并且它在
$S$ 上。下列条件等价：
\begin{enumerate}
\item $f$ 形式平展；
\item 对每个图
$$
\xymatrix{
U \ar[d] \ar[r]_\psi & V \ar[d] \\
X \ar[r]^f & Y
}
$$
其中 $U$、$V$ 为概形，且竖直箭头平展，概形态射 $\psi$
是形式平展的（参见《态射进阶》定义
\ref{more-morphisms-definition-formally-etale}）；
\item 对于一个这样的图，若竖直箭头满，则 $\psi$ 形式平展。
\end{enumerate}
\end{lemma}

\begin{proof}
假设 $f$ 形式平展。由引理
\ref{spaces-more-morphisms-lemma-representable-property-formally-property}，
态射 $U \to X$ 和 $V \to Y$ 形式平展。因此，由引理
\ref{spaces-more-morphisms-lemma-composition-formally-smooth-etale-unramified}，
复合 $U \to Y$ 形式平展。再由引理
\ref{spaces-more-morphisms-lemma-formally-permanence}，可知 $U \to V$ 形式平展。
故 (1) 蕴含 (2)，而 (2) 显然蕴含 (3)。% P09-ERR-0561

\medskip\noindent
假设给定 (3) 中的图。由引理
\ref{spaces-more-morphisms-lemma-representable-property-formally-property}，
态射 $V \to Y$ 形式平展。因此，由引理
\ref{spaces-more-morphisms-lemma-composition-formally-smooth-etale-unramified}，
复合 $U \to Y$ 形式平展。再由引理
\ref{spaces-more-morphisms-lemma-etale-on-top}，可知 $X \to Y$ 形式平展，即 (1) 成立。
\end{proof}

\begin{lemma}
\label{spaces-more-morphisms-lemma-formally-etale-not-affine}
设 $S$ 为概形。
设 $f : X \to Y$ 为代数空间的形式平展态射，并且它在 $S$ 上。
则给定任意实线交换图
$$
\xymatrix{
X \ar[d]_f & T \ar[d]^i \ar[l]_a \\
Y & T' \ar[l] \ar@{-->}[lu]
}
$$
其中 $T \subset T'$ 是 $Y$ 上代数空间的一阶加厚，
恰有一个虚线箭头使该图交换。换言之，在定义
\ref{spaces-more-morphisms-definition-formally-etale} 中，可去掉 $T$ 为仿射概形的条件。
\end{lemma}

\begin{proof}
设 $U' \to T'$ 为满平展态射，其中
$U' = \coprod U'_i$ 是仿射概形的不交并。令
$U_i = T \times_{T'} U'_i$。于是得到态射
$a'_i : U'_i \to X$，使得 $a'_i|_{U_i}$ 等于复合
$U_i \to T \to X$。由唯一性（见引理
\ref{spaces-more-morphisms-lemma-formally-unramified-not-affine}），
可知 $a'_i$ 和 $a'_j$ 在纤维积
$U'_i \times_{T'} U'_j$ 上相同。因此
$\coprod a'_i : U' \to X$ 下降为唯一态射 $a' : T' \to X$。
\end{proof}

\begin{lemma}
\label{spaces-more-morphisms-lemma-composition-formally-etale}
形式平展态射的复合仍形式平展。
\end{lemma}

\begin{proof}
这是形式地成立的。
\end{proof}

\begin{lemma}
\label{spaces-more-morphisms-lemma-base-change-formally-etale}
形式平展态射的基变换仍形式平展。
\end{lemma}

\begin{proof}
这是形式地成立的。
\end{proof}

\begin{lemma}
\label{spaces-more-morphisms-lemma-characterize-formally-etale}
设 $S$ 为概形。
设 $f : X \to Y$ 为代数空间的态射，并且它在 $S$ 上。% P09-ERR-0562
下列条件等价：
\begin{enumerate}
\item $f$ 形式平展；
\item $f$ 形式非分歧，且 $X$ 在 $Y$ 上的泛一阶加厚等于 $X$；
\item $f$ 形式非分歧，且 $\mathcal{C}_{X/Y} = 0$；
\item $\Omega_{X/Y} = 0$ 且 $\mathcal{C}_{X/Y} = 0$。
\end{enumerate}
\end{lemma}

\begin{proof}
事实上，最后一项断言之所以有意义，只是因为
$\Omega_{X/Y} = 0$ 蕴含 $\mathcal{C}_{X/Y}$ 可通过引理
\ref{spaces-more-morphisms-lemma-characterize-formally-unramified} 和定义
\ref{spaces-more-morphisms-definition-universal-thickening} 加以定义。% P09-ERR-0563
这也说明 (3) 和 (4) 等价。

\medskip\noindent
(1)、(2)、(3) 中任一假设都蕴含 $f$ 形式非分歧。% P09-ERR-0564
故可假设 $f$ 形式非分歧。(1)、(2)、(3) 的等价性由
泛一阶加厚 $X'$（即 $X$ 在 $S$ 上的那个）的泛性质，% P09-ERR-0565 % P09-ERR-0566
以及下述事实得出：
$X = X' \Leftrightarrow \mathcal{C}_{X/Y} = 0$；毕竟按定义，
$\mathcal{C}_{X/Y} = \mathcal{C}_{X/X'}$ 是 $X$ 在 $X'$ 中的理想层。
\end{proof}

\begin{lemma}
\label{spaces-more-morphisms-lemma-unramified-flat-formally-etale}
非分歧平坦态射是形式平展的。
\end{lemma}

\begin{proof}
这由概形情形、《态射进阶》引理
\ref{more-morphisms-lemma-unramified-flat-formally-etale}，
以及平展局部化得出；对于后者，见引理
\ref{spaces-more-morphisms-lemma-formally-unramified}、\ref{spaces-more-morphisms-lemma-formally-etale}
以及《空间的态射》引理
\ref{spaces-morphisms-lemma-flat-local}。
\end{proof}

\begin{lemma}
\label{spaces-more-morphisms-lemma-etale-formally-etale}
设 $S$ 为概形。
设 $f : X \to Y$ 为 $S$ 上代数空间的态射。
下列条件等价：
\begin{enumerate}
\item 态射 $f$ 平展；
\item 态射 $f$ 局部有限表示且形式平展。
\end{enumerate}
\end{lemma}

\begin{proof}
这由概形情形、《态射进阶》引理
\ref{more-morphisms-lemma-etale-formally-etale}，
以及平展局部化得出；对于后者，见引理
\ref{spaces-more-morphisms-lemma-formally-etale}
以及《空间的态射》诸引理
\ref{spaces-morphisms-lemma-finite-presentation-local} 和
\ref{spaces-morphisms-lemma-etale-local}。
\end{proof}

\section{映射的无穷小形变}
\label{spaces-more-morphisms-section-action-by-derivations}

\noindent
本节说明如何用导子无穷小地移动一个映射。全节都使用如下事实：
加厚 $X'$（它是 $X$ 的加厚）上的层可视为 $X$ 上的层，见等式
(\ref{spaces-more-morphisms-equation-equivalence-etale-spaces}) 和
(\ref{spaces-more-morphisms-equation-fundamental-equivalence})。

\begin{lemma}
\label{spaces-more-morphisms-lemma-difference-derivation}
设 $S$ 为概形，$B$ 为 $S$ 上的代数空间。
设 $X \subset X'$ 和 $Y \subset Y'$ 是 $B$ 上代数空间的两个一阶加厚。
设
$(a, a'), (b, b') : (X \subset X') \to (Y \subset Y')$
是 $B$ 上加厚之间的两个态射。假设
\begin{enumerate}
\item $a = b$；
\item 两个映射 $a^*\mathcal{C}_{Y/Y'} \to \mathcal{C}_{X/X'}$
（引理 \ref{spaces-more-morphisms-lemma-conormal-functorial}）相等。
\end{enumerate}
则映射 $(a')^\sharp - (b')^\sharp$ 分解为
$$
\mathcal{O}_{Y'} \to \mathcal{O}_Y \xrightarrow{D}
a_*\mathcal{C}_{X/X'} \to a_*\mathcal{O}_{X'}
$$
其中 $D$ 是 $\mathcal{O}_B$-导子。
\end{lemma}

\begin{proof}
我们不在 $Y$ 上工作，而在 $X$ 上工作，其优点是拉回函子
$a^{-1}$ 正合。使用 (1) 和 (2)，得到行正合的交换图
$$
\xymatrix{
0 \ar[r] &
\mathcal{C}_{X/X'} \ar[r] &
\mathcal{O}_{X'} \ar[r] &
\mathcal{O}_X \ar[r] & 0 \\
0 \ar[r] &
a^{-1}\mathcal{C}_{Y/Y'} \ar[r] \ar[u] &
a^{-1}\mathcal{O}_{Y'}
\ar[r] \ar@<1ex>[u]^{(a')^\sharp} \ar@<-1ex>[u]_{(b')^\sharp} &
a^{-1}\mathcal{O}_Y \ar[r] \ar[u] & 0
}
$$
一般而言，在这种情形下，$\mathcal{O}_B$-代数映射
$(a')^\sharp$ 与 $(b')^\sharp$ 之差，是一个 $\mathcal{O}_B$-导子，
即从 $a^{-1}\mathcal{O}_Y$ 到 $\mathcal{C}_{X/X'}$ 的导子。
由函子 $a^{-1}$ 与 $a_*$ 的伴随性，这等同于一个
$\mathcal{O}_B$-导子，即从 $\mathcal{O}_Y$ 到
$a_*\mathcal{C}_{X/X'}$ 的导子。略去若干细节。
\end{proof}

\noindent
注意，在上述引理的情形中，可将 $D$ 写成
\begin{equation}
\label{spaces-more-morphisms-equation-D}
D = \text{d}_{Y/B} \circ \theta
\end{equation}
% P09-ERR-0567
其中 $\theta$ 是 $\mathcal{O}_Y$-线性映射
$\theta : \Omega_{Y/B} \to a_*\mathcal{C}_{X/X'}$。
当然，再次利用伴随性，也可将 $\theta$ 视为 $\mathcal{O}_X$-线性映射
$\theta : a^*\Omega_{Y/B} \to \mathcal{C}_{X/X'}$。

\begin{lemma}
\label{spaces-more-morphisms-lemma-action-by-derivations}
设 $S$ 为概形，$B$ 为 $S$ 上的代数空间。
设 $(a, a') : (X \subset X') \to (Y \subset Y')$
为 $B$ 上一阶加厚之间的态射。设
$$
\theta : a^*\Omega_{Y/B} \to \mathcal{C}_{X/X'}
$$
为 $\mathcal{O}_X$-线性映射。则存在唯一的序对态射
$(b, b') : (X \subset X') \to (Y \subset Y')$，
使得引理 \ref{spaces-more-morphisms-lemma-difference-derivation} 的 (1)、(2) 成立，
并且导子 $D$ 与 $\theta$ 由等式 (\ref{spaces-more-morphisms-equation-D}) 联系。
\end{lemma}

\begin{proof}
考虑映射
$$
\alpha = (a')^\sharp + D : a^{-1}\mathcal{O}_{Y'} \to \mathcal{O}_{X'}
$$
其中 $D$ 如等式 (\ref{spaces-more-morphisms-equation-D}) 所示。因为 $D$ 是
$\mathcal{O}_B$-导子，所以 $\alpha$ 是 $\mathcal{O}_B$-代数层之间的映射。
由构造，$i_X^\sharp \circ \alpha = a^\sharp \circ i_Y^\sharp$，其中
$i_X : X \to X'$ 和 $i_Y : Y \to Y'$ 是给定闭浸入。由引理
\ref{spaces-more-morphisms-lemma-first-order-thickening-maps}，得到加厚之间的唯一态射
$(a, b') : (X  \subset X') \to (Y \subset Y')$，它在 $B$ 上，
使得 $\alpha = (b')^\sharp$。令 $b = a$ 即得结论。
\end{proof}

\begin{remark}
\label{spaces-more-morphisms-remark-action-by-derivations}
沿用引理 \ref{spaces-more-morphisms-lemma-action-by-derivations} 的假设与记号。
局部截面 $\theta$ 在 $a'$ 上的作用有时记作
$\theta \cdot a'$。这无非意味着 $(a')^\sharp$ 与
$(\theta \cdot a')^\sharp$ 相差一个导子 $D$，
而该导子与 $\theta$ 由等式 (\ref{spaces-more-morphisms-equation-D}) 联系。
\end{remark}

\begin{lemma}
\label{spaces-more-morphisms-lemma-sheaf}
设 $S$ 为概形，$B$ 为 $S$ 上的代数空间。
设 $X \subset X'$ 和 $Y \subset Y'$ 是 $B$ 上的一阶加厚。
假设给定态射 $a : X \to Y$ 以及映射
$A : a^*\mathcal{C}_{Y/Y'} \to \mathcal{C}_{X/X'}$；
这是一个 $\mathcal{O}_X$-模的映射。对对象 $U'$（它属于
$(X')_{spaces, \etale}$），令
$U = X \times_{X'} U'$，并考虑满足下列条件的态射
$a' : U' \to Y'$：
\begin{enumerate}
\item $a'$ 是 $B$ 上的态射；
\item $a'|_U = a|_U$；
\item 诱导映射
$a^*\mathcal{C}_{Y/Y'}|_U \to \mathcal{C}_{X/X'}|_U$
是 $A$ 在 $U$ 上的限制。
\end{enumerate}
则规则
\begin{equation}
\label{spaces-more-morphisms-equation-sheaf}
U' \mapsto
\{a' : U' \to Y'\text{ 满足 (1)、(2)、(3)。}\}
\end{equation}
定义 $(X')_{spaces, \etale}$ 上的集合层。
\end{lemma}

\begin{proof}
用 $\mathcal{F}$ 表示引理中的规则。映射
$\mathcal{F}(U') \to \mathcal{F}(V')$（其中
$V' \subset U' \subset X'$）作为 $\mathcal{F}$ 的限制映射，实际上就是
$a' \mapsto a'|_{V'}$。在此定义下，$\mathcal{F}$ 显然是层，
因为代数空间的态射满足平展下降；见《空间上的下降》引理
\ref{spaces-descent-lemma-fpqc-universal-effective-epimorphisms}。
\end{proof}

\begin{lemma}
\label{spaces-more-morphisms-lemma-action-sheaf}
沿用引理 \ref{spaces-more-morphisms-lemma-sheaf} 的记号与假设。
我们识别 $X$ 和 $X'$ 上的层，所用等价见
(\ref{spaces-more-morphisms-equation-equivalence-etale-spaces})。
层
$$
\SheafHom_{\mathcal{O}_X}(a^*\Omega_{Y/B}, \mathcal{C}_{X/X'})
$$
作用在层 (\ref{spaces-more-morphisms-equation-sheaf}) 上。此外，对于任意对象
$U'$（它属于 $(X')_{spaces, \etale}$），若层
(\ref{spaces-more-morphisms-equation-sheaf}) 在其上有截面，则该作用简便传递。
\end{lemma}

\begin{proof}
这是下列引理的组合：
\ref{spaces-more-morphisms-lemma-difference-derivation}、
\ref{spaces-more-morphisms-lemma-action-by-derivations} 和 \ref{spaces-more-morphisms-lemma-sheaf}。
\end{proof}

\begin{remark}
\label{spaces-more-morphisms-remark-special-case}
引理 \ref{spaces-more-morphisms-lemma-difference-derivation}、
\ref{spaces-more-morphisms-lemma-action-by-derivations}、
\ref{spaces-more-morphisms-lemma-sheaf} 和
\ref{spaces-more-morphisms-lemma-action-sheaf} 的一个特殊情形是 $Y = Y'$。
此时映射 $A$ 总为零。引理 \ref{spaces-more-morphisms-lemma-sheaf} 中的层就是由规则
$$
U' \mapsto
\{a' : U' \to Y\text{ 在 }B\text{ 上，且 } a'|_U = a|_U\}
$$
给出的层，并由层
$\SheafHom_{\mathcal{O}_X}(a^*\Omega_{Y/B}, \mathcal{C}_{X/X'})$
作用于其上。
\end{remark}

\begin{remark}
\label{spaces-more-morphisms-remark-another-special-case}
引理 \ref{spaces-more-morphisms-lemma-difference-derivation}、
\ref{spaces-more-morphisms-lemma-action-by-derivations}、
\ref{spaces-more-morphisms-lemma-sheaf} 和
\ref{spaces-more-morphisms-lemma-action-sheaf} 的另一个特殊情形是：$B$ 本身是加厚
$Z \subset Z' = B$，且 $Y = Z \times_{Z'} Y'$。图示为
$$
\xymatrix{
(X \subset X') \ar@{..>}[rr]_{(a, ?)} \ar[rd]_{(g, g')} & &
(Y \subset Y') \ar[ld]^{(h, h')} \\
& (Z \subset Z')
}
$$
此时映射 $A : a^*\mathcal{C}_{Y/Y'} \to \mathcal{C}_{X/X'}$
由 $a$ 决定：映射
$h^*\mathcal{C}_{Z/Z'} \to \mathcal{C}_{Y/Y'}$ 是满射（因为假设
$Y = Z \times_{Z'} Y'$），所以拉回
$g^*\mathcal{C}_{Z/Z'} = a^*h^*\mathcal{C}_{Z/Z'} \to
a^*\mathcal{C}_{Y/Y'}$ 是满射，而复合
$g^*\mathcal{C}_{Z/Z'} \to a^*\mathcal{C}_{Y/Y'} \to \mathcal{C}_{X/X'}$
必须是由 $g'$ 诱导的典范映射。因此，引理
\ref{spaces-more-morphisms-lemma-sheaf} 中的层就是由规则
$$
U' \mapsto
\{a' : U' \to Y'\text{ 在 }Z'\text{ 上，且 } a'|_U = a|_U\}
$$
给出的层，并由层
$\SheafHom_{\mathcal{O}_X}(a^*\Omega_{Y/Z}, \mathcal{C}_{X/X'})$
作用于其上。
\end{remark}

\begin{lemma}
\label{spaces-more-morphisms-lemma-action-by-derivations-etale-localization}
设 $S$ 为概形。考虑一阶加厚的交换图
$$
\vcenter{
\xymatrix{
(T_2 \subset T_2') \ar[d]_{(h, h')} \ar[rr]_{(a_2, a_2')} & &
(X_2 \subset X_2') \ar[d]^{(f, f')} \\
(T_1 \subset T_1') \ar[rr]^{(a_1, a_1')} & &
(X_1 \subset X_1')
}
}
\quad
\begin{matrix}
\text{以及交换} \\
\text{图}
\end{matrix}
\quad
\vcenter{
\xymatrix{
X_2' \ar[r] \ar[d] & B_2 \ar[d] \\
X_1' \ar[r] & B_1
}
}
$$
它们都是 $S$ 上代数空间的图，并且 $X_2 \to X_1$ 和
$B_2 \to B_1$ 平展。对任意 $\mathcal{O}_{T_1}$-线性映射
$\theta_1 : a_1^*\Omega_{X_1/B_1} \to \mathcal{C}_{T_1/T'_1}$，
令 $\theta_2$ 为复合
$$
\xymatrix{
a_2^*\Omega_{X_2/B_2} \ar@{=}[r] &
h^*a_1^*\Omega_{X_1/B_1} \ar[r]^-{h^*\theta_1} &
h^*\mathcal{C}_{T_1/T'_1} \ar[r] &
\mathcal{C}_{T_2/T'_2}
}
$$
（等号在证明中说明）。则图
$$
\xymatrix{
T_2' \ar[rr]_{\theta_2 \cdot a_2'} \ar[d] & & X'_2 \ar[d] \\
T_1' \ar[rr]^{\theta_1 \cdot a_1'} & & X'_1
}
$$
交换，其中作用 $\theta_2 \cdot a_2'$ 和 $\theta_1 \cdot a_1'$
如评注 \ref{spaces-more-morphisms-remark-action-by-derivations} 所述。
\end{lemma}

\begin{proof}
等号来自下述识别：
$f^*\Omega_{X_1/S_1} = \Omega_{X_2/S_2}$；该识别之所以成立，
是因为微分层的构造与平展局部化（在源与靶两边）相容，见引理
\ref{spaces-more-morphisms-lemma-localize-differentials}。% P09-ERR-0568
事实上，由此有
$a_2^*\Omega_{X_2/S_2} = a_2^*f^*\Omega_{X_1/S_1} =
h^*a_1^*\Omega_{X_1/S_1}$，因为 $f \circ a_2 = a_1 \circ h$。% P09-ERR-0569
说明这一点后，可平展局部地检验该图的交换性。% P09-ERR-0570
因此可假设 $T'_i$、$X'_i$、$B_2$ 和 $B_1$ 都是概形；
在这种情形下，引理由《态射进阶》引理
\ref{more-morphisms-lemma-action-by-derivations-etale-localization} 得出。
另一种证明：由引理 \ref{spaces-more-morphisms-lemma-first-order-thickening-maps}，
只需证明 $X_1'$ 上环层的某个图交换；然后完全仿照上述
《态射进阶》引理
\ref{more-morphisms-lemma-action-by-derivations-etale-localization}
的证明，即可看出确实如此。
\end{proof}
\section{代数空间的无穷小形变}
\label{spaces-more-morphisms-section-deform}

\noindent
下面这个简单引理往往是检验映射的无穷小形变是否平坦的便利工具。

\begin{lemma}
\label{spaces-more-morphisms-lemma-deform}
设 $S$ 为概形。设
$(f, f') : (X \subset X') \to (Y \subset Y')$
是 $S$ 上代数空间一阶加厚的态射，并假设 $f$ 平坦。则下列条件等价：
\begin{enumerate}
\item $f'$ 平坦且 $X = Y \times_{Y'} X'$；
\item 典范映射 $f^*\mathcal{C}_{Y/Y'} \to \mathcal{C}_{X/X'}$
是同构。
\end{enumerate}
\end{lemma}

\begin{proof}
选取概形 $V'$ 和满平展态射 $V' \to Y'$。
选取概形 $U'$ 和满平展态射
$U' \to X' \times_{Y'} V'$。令 $U = X \times_{X'} U'$ 且
$V = Y \times_{Y'} V'$。由代数空间平坦态射的定义可知，诱导映射
$g : U \to V$ 是概形的平坦态射，并且 $f'$ 平坦当且仅当相应态射
$g' : U' \to V'$ 平坦。此外，$X = Y \times_{Y'} X'$ 当且仅当
$U = V \times_{V'} V'$。% P09-ERR-0571
最后，映射
$f^*\mathcal{C}_{Y/Y'} \to \mathcal{C}_{X/X'}$
是同构当且仅当
$g^*\mathcal{C}_{V/V'} \to \mathcal{C}_{U/U'}$ 是同构。
因此，引理由概形态射的相应结论得出；见《态射进阶》引理
\ref{more-morphisms-lemma-deform}。
\end{proof}

\noindent
下面这个引理是“按纤维判定平坦性的准则”的“幂零”版本；见
第 \ref{spaces-more-morphisms-section-criterion-flat-fibres} 节。

\begin{lemma}
\label{spaces-more-morphisms-lemma-flatness-morphism-thickenings}
设 $S$ 为概形。考虑交换图
$$
\xymatrix{
(X \subset X') \ar[rr]_{(f, f')} \ar[rd] & & (Y \subset Y') \ar[ld] \\
& (B \subset B')
}
$$
，它是 $S$ 上代数空间加厚的交换图。假设
\begin{enumerate}
\item $X'$ 在 $B'$ 上平坦；
\item $f$ 平坦；
\item $B \subset B'$ 是有限阶加厚；
\item $X = B \times_{B'} X'$ 且 $Y = B \times_{B'} Y'$。
\end{enumerate}
则 $f'$ 平坦，并且 $Y'$ 在 $B'$ 上平坦，至少在 $f'$ 的像中每一点处如此。
\end{lemma}

\begin{proof}
选取概形 $U'$ 和满平展态射 $U' \to B'$。
选取概形 $V'$ 和满平展态射
$V' \to U' \times_{B'} Y'$。
选取概形 $W'$ 和满平展态射
$W' \to V' \times_{Y'} X'$。令 $U, V, W$ 分别为
$U', V', W'$ 沿 $B \to B'$ 的基变换。于是 $f'$ 的平坦性等价于
$W' \to V'$ 的平坦性，而已知 $W \to V$ 平坦。因此可以把概形情形的
引理应用于概形加厚的图
$$
\xymatrix{
(W \subset W') \ar[rr] \ar[rd] & & (V \subset V') \ar[ld] \\
& (U \subset U')
}
$$
；见《态射进阶》引理
\ref{more-morphisms-lemma-flatness-morphism-thickenings}。
关于 $Y'/B'$ 在 $f'$ 像中各点处平坦的断言也以同样方式得出。
\end{proof}

\noindent
概形态射的许多性质在平坦形变下保持不变。

\begin{lemma}
\label{spaces-more-morphisms-lemma-deform-property}
设 $S$ 为概形。考虑交换图
$$
\xymatrix{
(X \subset X') \ar[rr]_{(f, f')} \ar[rd] & & (Y \subset Y') \ar[ld] \\
& (B \subset B')
}
$$
，它是 $S$ 上代数空间加厚的交换图。假设
$B \subset B'$ 是有限阶加厚，$X'$ 在 $B'$ 上平坦，
$X = B \times_{B'} X'$，且 $Y = B \times_{B'} Y'$。则
\begin{enumerate}
\item $f$ 可表当且仅当 $f'$ 可表，
\label{spaces-more-morphisms-item-representable}
\item $f$ 平坦当且仅当 $f'$ 平坦，
\label{spaces-more-morphisms-item-flat}
\item $f$ 是同构当且仅当 $f'$ 是同构，
\label{spaces-more-morphisms-item-isomorphism}
\item $f$ 是开浸入当且仅当 $f'$ 是开浸入，
\label{spaces-more-morphisms-item-open-immersion}
\item $f$ 拟紧当且仅当 $f'$ 拟紧，
\label{spaces-more-morphisms-item-quasi-compact}
\item $f$ 普遍闭当且仅当 $f'$ 普遍闭，
\label{spaces-more-morphisms-item-universally-closed}
\item $f$（拟）分离当且仅当 $f'$（拟）分离，
\label{spaces-more-morphisms-item-separated}
\item $f$ 是单态射当且仅当 $f'$ 是单态射，
\label{spaces-more-morphisms-item-monomorphism}
\item $f$ 满射当且仅当 $f'$ 满射，
\label{spaces-more-morphisms-item-surjective}
\item $f$ 普遍单射当且仅当 $f'$ 普遍单射，
\label{spaces-more-morphisms-item-universally-injective}
\item $f$ 仿射当且仅当 $f'$ 仿射，
\label{spaces-more-morphisms-item-affine}
\item
\label{spaces-more-morphisms-item-finite-type}
$f$ 局部有限型当且仅当 $f'$ 局部有限型，
\item $f$ 局部拟有限当且仅当 $f'$ 局部拟有限，
\label{spaces-more-morphisms-item-quasi-finite}
\item
\label{spaces-more-morphisms-item-finite-presentation}
$f$ 局部有限表示当且仅当 $f'$ 局部有限表示，
\item
\label{spaces-more-morphisms-item-relative-dimension-d}
$f$ 是相对维数 $d$ 的局部有限型态射当且仅当
$f'$ 是相对维数 $d$ 的局部有限型态射，
\item $f$ 普遍开当且仅当 $f'$ 普遍开，
\label{spaces-more-morphisms-item-universally-open}
\item $f$ 是 syntomic 态射当且仅当 $f'$ 是 syntomic 态射，
\label{spaces-more-morphisms-item-syntomic}
\item $f$ 光滑当且仅当 $f'$ 光滑，
\label{spaces-more-morphisms-item-smooth}
\item $f$ 非分歧当且仅当 $f'$ 非分歧，
\label{spaces-more-morphisms-item-unramified}
\item $f$ 平展当且仅当 $f'$ 平展，
\label{spaces-more-morphisms-item-etale}
\item $f$ 固有当且仅当 $f'$ 固有，
\label{spaces-more-morphisms-item-proper}
\item $f$ 整当且仅当 $f'$ 整，
\label{spaces-more-morphisms-item-integral}
\item $f$ 有限当且仅当 $f'$ 有限，
\label{spaces-more-morphisms-item-finite}
\item
\label{spaces-more-morphisms-item-finite-locally-free}
$f$ 有限局部自由（秩为 $d$）当且仅当 $f'$ 有限局部自由
（秩为 $d$）；
\item 在此添加更多情形。% P09-ERR-0572
\end{enumerate}
\end{lemma}

\begin{proof}
情形 (\ref{spaces-more-morphisms-item-representable}) 由引理
\ref{spaces-more-morphisms-lemma-thicken-property-morphisms} 得出。

\medskip\noindent
选取概形 $U'$ 和满平展态射 $U' \to B'$。
选取概形 $V'$ 和满平展态射
$V' \to U' \times_{B'} Y'$。
选取概形 $W'$ 和满平展态射
$W' \to V' \times_{Y'} X'$。令 $U, V, W$ 分别为
$U', V', W'$ 沿 $B \to B'$ 的基变换。考虑概形加厚的图
$$
\xymatrix{
(W \subset W') \ar[rr] \ar[rd] & & (V \subset V') \ar[ld] \\
& (U \subset U')
}
$$
。对于任何对源与靶平展局部的性质，把概形加厚态射的相应结论
应用于上图便立即得到所需结论。因此，情形
(\ref{spaces-more-morphisms-item-flat})、(\ref{spaces-more-morphisms-item-finite-type})、
(\ref{spaces-more-morphisms-item-quasi-finite})、(\ref{spaces-more-morphisms-item-finite-presentation})、
(\ref{spaces-more-morphisms-item-relative-dimension-d})、(\ref{spaces-more-morphisms-item-syntomic})、
(\ref{spaces-more-morphisms-item-smooth})、(\ref{spaces-more-morphisms-item-unramified})、(\ref{spaces-more-morphisms-item-etale})
由《态射进阶》引理
\ref{more-morphisms-lemma-deform-property} 的相应情形得出。

\medskip\noindent
由于 $X \to X'$ 和 $Y \to Y'$ 都是普遍同胚，关于映射
$X \to Y$ 和 $X' \to Y'$ 的任何拓扑问题都有相同答案。因此情形
(\ref{spaces-more-morphisms-item-quasi-compact})、(\ref{spaces-more-morphisms-item-universally-closed})、
(\ref{spaces-more-morphisms-item-surjective})、(\ref{spaces-more-morphisms-item-universally-injective}) 和
(\ref{spaces-more-morphisms-item-universally-open}) 成立。

\medskip\noindent
在其余每种情形中，我们只证明蕴含
$f\text{ 具有 }P \Rightarrow f'\text{ 具有 }P$，因为反向蕴含由
$P$ 在基变换下稳定得出；见《空间》引理
\ref{spaces-lemma-base-change-immersions} 以及《空间的态射》引理
\ref{spaces-morphisms-lemma-base-change-separated}、
\ref{spaces-morphisms-lemma-base-change-monomorphism}、
\ref{spaces-morphisms-lemma-base-change-affine}、
\ref{spaces-morphisms-lemma-base-change-proper}、
\ref{spaces-morphisms-lemma-base-change-integral} 和
\ref{spaces-morphisms-lemma-base-change-finite-locally-free}。

\medskip\noindent
情形 (\ref{spaces-more-morphisms-item-open-immersion})。假设 $f$ 是开浸入。
由 (\ref{spaces-more-morphisms-item-etale})，$f'$ 平展；又由
(\ref{spaces-more-morphisms-item-universally-injective}) 普遍单射，故 $f'$ 是开浸入；
见《空间的态射》引理
\ref{spaces-morphisms-lemma-etale-universally-injective-open}。
若先用 (\ref{spaces-more-morphisms-item-representable}) 约化到概形情形，则也可以不用这个引理。

\medskip\noindent
情形 (\ref{spaces-more-morphisms-item-isomorphism})。它由情形
(\ref{spaces-more-morphisms-item-open-immersion}) 和 (\ref{spaces-more-morphisms-item-surjective}) 得出。

\medskip\noindent
情形 (\ref{spaces-more-morphisms-item-separated})。见引理
\ref{spaces-more-morphisms-lemma-thicken-property-morphisms}。

\medskip\noindent
情形 (\ref{spaces-more-morphisms-item-monomorphism})。假设 $f$ 是单态射。
考虑对角态射 $\Delta_{X'/Y'} : X' \to X' \times_{Y'} X'$。
$\Delta_{X'/Y'}$ 沿 $B \to B'$ 的基变换是 $\Delta_{X/Y}$，
后者依假设是同构。由 (\ref{spaces-more-morphisms-item-isomorphism}) 可知
$\Delta_{X'/Y'}$ 是同构，因而 $f'$ 是单态射。

\medskip\noindent
情形 (\ref{spaces-more-morphisms-item-affine})。见引理 \ref{spaces-more-morphisms-lemma-thicken-property-morphisms}。

\medskip\noindent
情形 (\ref{spaces-more-morphisms-item-proper})。见引理
\ref{spaces-more-morphisms-lemma-thicken-property-morphisms-cartesian}。

\medskip\noindent
情形 (\ref{spaces-more-morphisms-item-integral})。见引理
\ref{spaces-more-morphisms-lemma-thicken-property-morphisms}。

\medskip\noindent
情形 (\ref{spaces-more-morphisms-item-finite})。见引理
\ref{spaces-more-morphisms-lemma-thicken-property-morphisms-cartesian}。

\medskip\noindent
情形 (\ref{spaces-more-morphisms-item-finite-locally-free})。假设 $f$ 有限局部自由。
由 (\ref{spaces-more-morphisms-item-finite}) 可知 $f'$ 有限。
由 (\ref{spaces-more-morphisms-item-flat}) 可知 $f'$ 平坦。
由 (\ref{spaces-more-morphisms-item-finite-presentation})，$f'$ 局部有限表示。
因此，由《空间的态射》引理
\ref{spaces-morphisms-lemma-finite-flat}，$f'$ 有限局部自由。
\end{proof}

\noindent
下面这个引理是“按纤维判定平坦性的准则”的“局部幂零”版本；见
第 \ref{spaces-more-morphisms-section-criterion-flat-fibres} 节。

\begin{lemma}
\label{spaces-more-morphisms-lemma-flatness-morphism-thickenings-fp-over-ft}
设 $S$ 为概形。考虑交换图
$$
\xymatrix{
(X \subset X') \ar[rr]_{(f, f')} \ar[rd] & & (Y \subset Y') \ar[ld] \\
& (B \subset B')
}
$$
，它是 $S$ 上代数空间加厚的交换图。假设
\begin{enumerate}
\item $Y' \to B'$ 局部有限型；
\item $X' \to B'$ 平坦且局部有限表示；
\item $f$ 平坦；
\item $X = B \times_{B'} X'$ 且 $Y = B \times_{B'} Y'$。
\end{enumerate}
则 $f'$ 平坦，并且对于所有 $y' \in |Y'|$，只要它在
$|f'|$ 的像中，态射 $Y' \to B'$ 就在 $y'$ 处平坦。
\end{lemma}

\begin{proof}
选取概形 $U'$ 和满平展态射 $U' \to B'$。
选取概形 $V'$ 和满平展态射
$V' \to U' \times_{B'} Y'$。
选取概形 $W'$ 和满平展态射
$W' \to V' \times_{Y'} X'$。令 $U, V, W$ 分别为
$U', V', W'$ 沿 $B \to B'$ 的基变换。于是 $f'$ 的平坦性等价于
$W' \to V'$ 的平坦性，而已知 $W \to V$ 平坦。因此可以把概形情形的
引理应用于概形加厚的图
$$
\xymatrix{
(W \subset W') \ar[rr] \ar[rd] & & (V \subset V') \ar[ld] \\
& (U \subset U')
}
$$
；见《态射进阶》引理
\ref{more-morphisms-lemma-flatness-morphism-thickenings-fp-over-ft}。
关于 $Y'/B'$ 在 $f'$ 像中各点处平坦的断言也以同样方式得出。
\end{proof}

\noindent
概形态射的许多性质在上述引理中的平坦形变下保持不变。

\begin{lemma}
\label{spaces-more-morphisms-lemma-deform-property-fp-over-ft}
设 $S$ 为概形。考虑交换图
$$
\xymatrix{
(X \subset X') \ar[rr]_{(f, f')} \ar[rd] & & (Y \subset Y') \ar[ld] \\
& (B \subset B')
}
$$
，它是 $S$ 上代数空间加厚的交换图。假设
$Y' \to B'$ 局部有限型，
$X' \to B'$ 平坦且局部有限表示，
$X = B \times_{B'} X'$，且 $Y = B \times_{B'} Y'$。则
\begin{enumerate}
\item $f$ 可表当且仅当 $f'$ 可表，
\label{spaces-more-morphisms-item-representable-fp-over-ft}
\item $f$ 平坦当且仅当 $f'$ 平坦，
\label{spaces-more-morphisms-item-flat-fp-over-ft}
\item $f$ 是同构当且仅当 $f'$ 是同构，
\label{spaces-more-morphisms-item-isomorphism-fp-over-ft}
\item $f$ 是开浸入当且仅当 $f'$ 是开浸入，
\label{spaces-more-morphisms-item-open-immersion-fp-over-ft}
\item $f$ 拟紧当且仅当 $f'$ 拟紧，
\label{spaces-more-morphisms-item-quasi-compact-fp-over-ft}
\item $f$ 普遍闭当且仅当 $f'$ 普遍闭，
\label{spaces-more-morphisms-item-universally-closed-fp-over-ft}
\item $f$（拟）分离当且仅当 $f'$（拟）分离，
\label{spaces-more-morphisms-item-separated-fp-over-ft}
\item $f$ 是单态射当且仅当 $f'$ 是单态射，
\label{spaces-more-morphisms-item-monomorphism-fp-over-ft}
\item $f$ 满射当且仅当 $f'$ 满射，
\label{spaces-more-morphisms-item-surjective-fp-over-ft}
\item $f$ 普遍单射当且仅当 $f'$ 普遍单射，
\label{spaces-more-morphisms-item-universally-injective-fp-over-ft}
\item $f$ 仿射当且仅当 $f'$ 仿射，
\label{spaces-more-morphisms-item-affine-fp-over-ft}
\item $f$ 局部拟有限当且仅当 $f'$ 局部拟有限，
\label{spaces-more-morphisms-item-quasi-finite-fp-over-ft}
\item
\label{spaces-more-morphisms-item-relative-dimension-d-fp-over-ft}
$f$ 是相对维数 $d$ 的局部有限型态射当且仅当
$f'$ 是相对维数 $d$ 的局部有限型态射，
\item $f$ 普遍开当且仅当 $f'$ 普遍开，
\label{spaces-more-morphisms-item-universally-open-fp-over-ft}
\item $f$ 是 syntomic 态射当且仅当 $f'$ 是 syntomic 态射，
\label{spaces-more-morphisms-item-syntomic-fp-over-ft}
\item $f$ 光滑当且仅当 $f'$ 光滑，
\label{spaces-more-morphisms-item-smooth-fp-over-ft}
\item $f$ 非分歧当且仅当 $f'$ 非分歧，
\label{spaces-more-morphisms-item-unramified-fp-over-ft}
\item $f$ 平展当且仅当 $f'$ 平展，
\label{spaces-more-morphisms-item-etale-fp-over-ft}
\item $f$ 固有当且仅当 $f'$ 固有，
\label{spaces-more-morphisms-item-proper-fp-over-ft}
\item $f$ 有限当且仅当 $f'$ 有限，
\label{spaces-more-morphisms-item-finite-fp-over-ft}
\item
\label{spaces-more-morphisms-item-finite-locally-free-fp-over-ft}
$f$ 有限局部自由（秩为 $d$）当且仅当 $f'$ 有限局部自由
（秩为 $d$）；
\item 在此添加更多情形。% P09-ERR-0573
\end{enumerate}
\end{lemma}

\begin{proof}
情形 (\ref{spaces-more-morphisms-item-representable-fp-over-ft}) 由引理
\ref{spaces-more-morphisms-lemma-thicken-property-morphisms} 得出。

\medskip\noindent
选取概形 $U'$ 和满平展态射 $U' \to B'$。
选取概形 $V'$ 和满平展态射
$V' \to U' \times_{B'} Y'$。
选取概形 $W'$ 和满平展态射
$W' \to V' \times_{Y'} X'$。令 $U, V, W$ 分别为
$U', V', W'$ 沿 $B \to B'$ 的基变换。考虑概形加厚的图
$$
\xymatrix{
(W \subset W') \ar[rr] \ar[rd] & & (V \subset V') \ar[ld] \\
& (U \subset U')
}
$$
。对于任何对源与靶平展局部的性质，把概形加厚态射的相应结论
应用于上图便立即得到所需结论。因此，情形
(\ref{spaces-more-morphisms-item-flat-fp-over-ft})、
(\ref{spaces-more-morphisms-item-quasi-finite-fp-over-ft})、
(\ref{spaces-more-morphisms-item-relative-dimension-d-fp-over-ft})、
(\ref{spaces-more-morphisms-item-syntomic-fp-over-ft})、
(\ref{spaces-more-morphisms-item-smooth-fp-over-ft})、
(\ref{spaces-more-morphisms-item-unramified-fp-over-ft})、
(\ref{spaces-more-morphisms-item-etale-fp-over-ft})
由《态射进阶》引理
\ref{more-morphisms-lemma-deform-property-fp-over-ft} 的相应情形得出。

\medskip\noindent
由于 $X \to X'$ 和 $Y \to Y'$ 都是普遍同胚，关于映射
$X \to Y$ 和 $X' \to Y'$ 的任何拓扑问题都有相同答案。因此情形
(\ref{spaces-more-morphisms-item-quasi-compact-fp-over-ft})、
(\ref{spaces-more-morphisms-item-universally-closed-fp-over-ft})、
(\ref{spaces-more-morphisms-item-surjective-fp-over-ft})、
(\ref{spaces-more-morphisms-item-universally-injective-fp-over-ft}) 和
(\ref{spaces-more-morphisms-item-universally-open-fp-over-ft}) 成立。

\medskip\noindent
在其余每种情形中，我们只证明蕴含
$f\text{ 具有 }P \Rightarrow f'\text{ 具有 }P$，因为反向蕴含由
$P$ 在基变换下稳定得出；见《空间》引理
\ref{spaces-lemma-base-change-immersions} 以及《空间的态射》引理
\ref{spaces-morphisms-lemma-base-change-separated}、
\ref{spaces-morphisms-lemma-base-change-monomorphism}、
\ref{spaces-morphisms-lemma-base-change-affine}、
\ref{spaces-morphisms-lemma-base-change-proper}、
\ref{spaces-morphisms-lemma-base-change-integral} 和
\ref{spaces-morphisms-lemma-base-change-finite-locally-free}。

\medskip\noindent
情形 (\ref{spaces-more-morphisms-item-open-immersion-fp-over-ft})。
假设 $f$ 是开浸入。
由 (\ref{spaces-more-morphisms-item-etale-fp-over-ft})，$f'$ 平展；又由
(\ref{spaces-more-morphisms-item-universally-injective-fp-over-ft}) 普遍单射，故
$f'$ 是开浸入；见《空间的态射》引理
\ref{spaces-morphisms-lemma-etale-universally-injective-open}。
若先用 (\ref{spaces-more-morphisms-item-representable-fp-over-ft}) 约化到概形情形，
则也可以不用这个引理。

\medskip\noindent
情形 (\ref{spaces-more-morphisms-item-isomorphism-fp-over-ft})。它由情形
(\ref{spaces-more-morphisms-item-open-immersion-fp-over-ft}) 和
(\ref{spaces-more-morphisms-item-surjective-fp-over-ft}) 得出。

\medskip\noindent
情形 (\ref{spaces-more-morphisms-item-separated-fp-over-ft})。见引理
\ref{spaces-more-morphisms-lemma-thicken-property-morphisms}。

\medskip\noindent
情形 (\ref{spaces-more-morphisms-item-monomorphism-fp-over-ft})。假设 $f$ 是单态射。
考虑对角态射 $\Delta_{X'/Y'} : X' \to X' \times_{Y'} X'$。
$\Delta_{X'/Y'}$ 沿 $B \to B'$ 的基变换是 $\Delta_{X/Y}$，
后者依假设是同构。由 (\ref{spaces-more-morphisms-item-isomorphism-fp-over-ft}) 可知
$\Delta_{X'/Y'}$ 是同构，因而 $f'$ 是单态射。

\medskip\noindent
情形 (\ref{spaces-more-morphisms-item-affine-fp-over-ft})。
见引理 \ref{spaces-more-morphisms-lemma-thicken-property-morphisms}。

\medskip\noindent
情形 (\ref{spaces-more-morphisms-item-proper-fp-over-ft})。见引理
\ref{spaces-more-morphisms-lemma-properties-that-extend-over-thickenings}。

\medskip\noindent
情形 (\ref{spaces-more-morphisms-item-finite-fp-over-ft})。见引理
\ref{spaces-more-morphisms-lemma-properties-that-extend-over-thickenings}。

\medskip\noindent
情形 (\ref{spaces-more-morphisms-item-finite-locally-free-fp-over-ft})。
假设 $f$ 有限局部自由。
由 (\ref{spaces-more-morphisms-item-finite-fp-over-ft}) 可知 $f'$ 有限。
由 (\ref{spaces-more-morphisms-item-flat-fp-over-ft}) 可知 $f'$ 平坦。
此外，由《空间的态射》引理% P09-ERR-0574
\ref{spaces-morphisms-lemma-finite-presentation-permanence}，
$f'$ 局部有限表示。
因此，由《空间的态射》引理
\ref{spaces-morphisms-lemma-finite-flat}，$f'$ 有限局部自由。
\end{proof}













\section{形式光滑态射}
\label{spaces-more-morphisms-section-formally-smooth}

\noindent
本节引入代数空间形式光滑态射 $X \to Y$ 的概念。这类态射的特征是：
只要 $T$ 仿射，$X$ 的 $T$-值点就能提升到 $T$ 的无穷小加厚。
主要结论是，形式光滑且局部有限表示的态射是光滑的；见引理
\ref{spaces-more-morphisms-lemma-smooth-formally-smooth}。事实证明，这个判据往往比雅可比判据
更便于使用。

\begin{definition}
\label{spaces-more-morphisms-definition-formally-smooth}
设 $S$ 为概形。若代数空间的态射 $f : X \to Y$（它在 $S$ 上）作为函子变换
在定义 \ref{spaces-more-morphisms-definition-formally-smooth-etale-unramified} 的意义下形式光滑，
则称它为{\it 形式光滑的}。
\end{definition}

\noindent
对于形式非分歧和形式平展态射，可以去掉 $T'$ 仿射这一条件；见引理
\ref{spaces-more-morphisms-lemma-formally-unramified-not-affine} 和
\ref{spaces-more-morphisms-lemma-formally-etale-not-affine}。对于形式光滑态射，这不再成立。
事实上，稍微自然一些的条件是要求能够在 $T'$ 上平展局部地补出虚线箭头。
进一步分析引理 \ref{spaces-more-morphisms-lemma-smooth-formally-smooth} 的证明可知，这与目前的
定义等价。要求虚线箭头在 $T'$ 上 fppf 局部地存在也足够，不过证明起来
稍微困难一些。

\medskip\noindent
第 \ref{spaces-more-morphisms-section-formally-smooth-etale-unramified} 节已经在形式光滑函子变换的
更一般背景下证明了若干结果，这里不再重述。

\begin{lemma}
\label{spaces-more-morphisms-lemma-composition-formally-smooth}
形式光滑态射的复合仍形式光滑。
\end{lemma}

\begin{proof}
略。
\end{proof}

\begin{lemma}
\label{spaces-more-morphisms-lemma-base-change-formally-smooth}
形式光滑态射的基变换仍形式光滑。
\end{lemma}

\begin{proof}
略；代数版本见《代数》引理
\ref{algebra-lemma-base-change-fs}。
\end{proof}

\begin{lemma}
\label{spaces-more-morphisms-lemma-formally-etale-unramified-smooth}
设 $f : X \to S$ 为概形态射。则 $f$ 形式平展当且仅当
$f$ 形式光滑且形式非分歧。
\end{lemma}

\begin{proof}
略。
\end{proof}

\noindent
下面是一个辅助引理，稍后将由引理 \ref{spaces-more-morphisms-lemma-formally-smooth} 取代。

\begin{lemma}
\label{spaces-more-morphisms-lemma-helper-formally-smooth}
设 $S$ 为概形，并设
$$
\xymatrix{
U \ar[d] \ar[r]_\psi & V \ar[d] \\
X \ar[r]^f & Y
}
$$
为 $S$ 上代数空间态射的交换图。若竖直箭头平展且 $f$ 形式光滑，
则 $\psi$ 形式光滑。
\end{lemma}

\begin{proof}
由引理 \ref{spaces-more-morphisms-lemma-representable-property-formally-property}，态射
$U \to X$ 和 $V \to Y$ 形式平展。由引理
\ref{spaces-more-morphisms-lemma-composition-formally-smooth-etale-unramified}，复合
$U \to Y$ 形式光滑。再由引理 \ref{spaces-more-morphisms-lemma-formally-permanence} 可知
$\psi : U \to V$ 形式光滑。
\end{proof}

\noindent
下面的引理是本节的主要结果。结合《空间的极限》命题
\ref{spaces-limits-proposition-characterize-locally-finite-presentation}，
它说明可以判定代数空间态射 $f : X \to Y$ 是否光滑，而只需考察函子变换
$X \to Y$ 的“简单”性质。

\begin{lemma}[无穷小提升判据]
\label{spaces-more-morphisms-lemma-smooth-formally-smooth}
设 $S$ 为概形。设 $f : X \to Y$ 为 $S$ 上代数空间的态射。
下列条件等价：
\begin{enumerate}
\item 态射 $f$ 光滑。
\item 态射 $f$ 局部有限表示且形式光滑。
\end{enumerate}
\end{lemma}

\begin{proof}
假设 $f : X \to S$ 局部有限表示且形式光滑。% P09-ERR-0575
考虑交换图
$$
\xymatrix{
U \ar[d] \ar[r]_\psi & V \ar[d] \\
X \ar[r]^f & Y
}
$$
，其中 $U$ 和 $V$ 是概形，竖直箭头平展且满。由引理
\ref{spaces-more-morphisms-lemma-helper-formally-smooth} 可知 $\psi : U \to V$ 形式光滑。
由《空间的态射》引理
\ref{spaces-morphisms-lemma-finite-presentation-local}，态射 $\psi$ 局部有限表示。
因此，由概形情形可知态射 $\psi$ 光滑；见《态射进阶》引理
\ref{more-morphisms-lemma-smooth-formally-smooth}。于是 $f$ 光滑；见
《空间的态射》引理
\ref{spaces-morphisms-lemma-smooth-local}。

\medskip\noindent
反过来，假设 $f : X \to Y$ 光滑。考虑定义
\ref{spaces-more-morphisms-definition-formally-smooth} 中实线部分交换的图
$$
\xymatrix{
X \ar[d]_f & T \ar[d]^i \ar[l]^a \\
Y & T' \ar[l] \ar@{-->}[lu]
}
$$
。我们将证明虚线箭头存在，从而证明 $f$ 形式光滑。
令 $\mathcal{F}$ 为引理 \ref{spaces-more-morphisms-lemma-sheaf} 在评注
\ref{spaces-more-morphisms-remark-special-case} 所述特殊情形下定义于
$(T')_{spaces, \etale}$ 上的集合层。令
$$
\mathcal{H} =
\SheafHom_{\mathcal{O}_T}(a^*\Omega_{X/Y}, \mathcal{C}_{T/T'})
$$
为 $\mathcal{O}_T$-模层（定义于 $T_{spaces, \etale}$ 上），并按引理
\ref{spaces-more-morphisms-lemma-action-sheaf} 取作用
$\mathcal{H} \times \mathcal{F} \to \mathcal{F}$。
这个作用 $\mathcal{H} \times \mathcal{F} \to \mathcal{F}$ 使
$\mathcal{F}$ 成为伪 $\mathcal{H}$-挠子；见《位点上的上同调》定义
\ref{sites-cohomology-definition-torsor}。目标是证明 $\mathcal{F}$ 是平凡的
$\mathcal{H}$-挠子。分两步：(I) 为证明 $\mathcal{F}$ 是挠子，必须证明
$\mathcal{F}$ 在平展局部有截面；(II) 为证明 $\mathcal{F}$ 是平凡挠子，
只需证明 $H^1(T_\etale, \mathcal{H}) = 0$；见《位点上的上同调》引理
\ref{sites-cohomology-lemma-torsors-h1}。

\medskip\noindent
先证明 (I)。为此，选取交换图
$$
\xymatrix{
U \ar[d] \ar[r]_\psi & V \ar[d] \\
X \ar[r]^f & Y
}
$$
，其中 $U$ 和 $V$ 是概形，竖直箭头平展且满。由于假设 $f$ 光滑，
可知 $\psi$ 光滑，因而由引理
\ref{spaces-more-morphisms-lemma-representable-property-formally-property} 可知它形式光滑。
由同一引理，态射 $V \to Y$ 形式平展。因此，由引理
\ref{spaces-more-morphisms-lemma-composition-formally-smooth-etale-unramified}，复合
$U \to Y$ 形式光滑。于是 (I) 由引理 \ref{spaces-more-morphisms-lemma-etale-on-top} 的第 (4) 部分得出。

\medskip\noindent
最后证明 (II)。由引理
\ref{spaces-more-morphisms-lemma-finite-presentation-differentials} 可知
$\Omega_{X/S}$ 有限表示。% P09-ERR-0576
因此 $a^*\Omega_{X/S}$ 有限表示（见《空间的性质》第
\ref{spaces-properties-section-properties-modules} 节）。故层
$\mathcal{H} =
\SheafHom_{\mathcal{O}_T}(a^*\Omega_{X/Y}, \mathcal{C}_{T/T'})$
拟凝聚；见《空间的性质》引理
\ref{spaces-properties-lemma-properties-quasi-coherent}。因而由《下降》命题
\ref{descent-proposition-same-cohomology-quasi-coherent} 和《概形的上同调》引理
\ref{coherent-lemma-quasi-coherent-affine-cohomology-zero}，有
$$
H^1(T_{spaces, \etale}, \mathcal{H}) =
H^1(T_\etale, \mathcal{H}) =
H^1(T, \mathcal{H}) = 0
$$
，正合所需。
\end{proof}

\noindent
光滑态射满足强局部提升性质；见引理 \ref{spaces-more-morphisms-lemma-smooth-strong-lift}。
若在该引理中假设 $T'$ 仿射，我们不知道是否有必要取平展覆盖。
更确切地说，若有代数空间的交换图
$$
\xymatrix{
X \ar[d] & T \ar[l] \ar[d] \\
Y & T' \ar[l] \ar@{..>}[lu]
}
$$
，其中 $X \to Y$ 光滑且 $T \to T'$ 是仿射概形的加厚，那么使图交换的
虚线箭头是否总存在？% P09-ERR-0577
若你知道答案或掌握参考文献，请发送邮件至
\href{mailto:stacks.project@gmail.com}{stacks.project@gmail.com}。

\begin{lemma}
\label{spaces-more-morphisms-lemma-smooth-strong-lift}
设 $S$ 为概形。考虑交换图
$$
\xymatrix{
X \ar[d] & T \ar[l] \ar[d] \\
Y & T' \ar[l]
}
$$
，它是 $S$ 上代数空间的交换图，其中 $X \to Y$ 光滑且
$T \to T'$ 是加厚。则存在平展覆盖
$\{T'_i \to T'\}$，使得能在图
$$
\xymatrix{
X \ar[d] & T \ar[l] \ar[d] & T \times_{T'} T'_i \ar[l] \ar[d] \\
Y & T' \ar[l] & T'_i \ar[l] \ar@{..>}[llu]
}
$$
中找到使图交换的虚线箭头（对所有 $i$）。
\end{lemma}

\begin{proof}
选取平展覆盖 $\{Y_i \to Y\}$，其中每个 $Y_i$ 都仿射。
用诱导的平展覆盖替换 $T'$ 后，可假设 $Y$ 仿射。

\medskip\noindent
假设 $Y$ 仿射。选取平展覆盖 $\{X_i \to X\}$。
这给出 $T$ 的平展覆盖。由定理 \ref{spaces-more-morphisms-theorem-topological-invariance}
（另见第 \ref{spaces-more-morphisms-section-thickenings} 节的讨论），$T$ 的这个平展覆盖来自
$T'$ 的某个平展覆盖。因此可假设 $X$ 仿射。

\medskip\noindent
假设 $X$ 和 $Y$ 都仿射。可再对 $T'$ 取一次平展覆盖并假设 $T'$ 仿射。
在这种情形下，引理由《代数》引理
\ref{algebra-lemma-smooth-strong-lift} 得出。
\end{proof}







\noindent
下面再做一些工作，证明形式光滑性对源平展局部。首先证明形式光滑态射
具有良好的微分层。局部投射拟凝聚模的概念定义于《空间的性质》第
\ref{spaces-properties-section-locally-projective} 节。

\begin{lemma}
\label{spaces-more-morphisms-lemma-formally-smooth-sheaf-differentials}
设 $S$ 为概形。设 $f : X \to Y$ 为 $S$ 上代数空间的形式光滑态射。
则 $\Omega_{X/Y}$ 在 $X$ 上局部投射。
\end{lemma}

\begin{proof}
选取图
$$
\xymatrix{
U \ar[d] \ar[r]_\psi & V \ar[d] \\
X \ar[r]^f & Y
}
$$
，其中 $U$ 和 $V$ 是仿射（！）概形，竖直箭头平展。由引理
\ref{spaces-more-morphisms-lemma-helper-formally-smooth} 可知 $\psi : U \to V$ 形式光滑。
因此
$\Gamma(V, \mathcal{O}_V) \to \Gamma(U, \mathcal{O}_U)$
是形式光滑环同态；见《态射进阶》引理
\ref{more-morphisms-lemma-affine-formally-smooth}。于是由《代数》引理
\ref{algebra-lemma-characterize-formally-smooth-again}，
$\Gamma(U, \mathcal{O}_U)$-模
$\Omega_{\Gamma(U, \mathcal{O}_U)/\Gamma(V, \mathcal{O}_V)}$
是投射模。因此 $\Omega_{U/V}$ 局部投射；见《性质》第
\ref{properties-section-locally-projective} 节。由于
$\Omega_{X/Y}|_U = \Omega_{U/V}$，可知 $\Omega_{X/Y}$ 也局部投射。
（因为可以找到 $X$ 的平展覆盖，其成员是上述图中的仿射 $U$；细节从略。）
\end{proof}

\begin{lemma}
\label{spaces-more-morphisms-lemma-h1-is-zero}
设 $T$ 为仿射概形。设 $\mathcal{F}$、$\mathcal{G}$ 是拟凝聚
$\mathcal{O}_T$-模（定义于 $T_\etale$ 上）。考虑内 Hom 层
$\mathcal{H} = \SheafHom_{\mathcal{O}_T}(\mathcal{F}, \mathcal{G})$，
它定义于 $T_\etale$ 上。
若 $\mathcal{F}$ 局部投射，则 $H^1(T_\etale, \mathcal{H}) = 0$。
\end{lemma}

\begin{proof}
由代数空间上局部投射层的定义（见《空间的性质》定义
\ref{spaces-properties-definition-locally-projective}），可知
$\mathcal{F}_{Zar} = \mathcal{F}|_{T_{Zar}}$ 是概形 $T$ 上的局部投射层。
因此 $\mathcal{F}_{Zar}$ 是某个自由 $\mathcal{O}_{T_{Zar}}$-模的直和项。
进而（因为 $\mathcal{F} = (\mathcal{F}_{Zar})^a$，见《下降》命题
\ref{descent-proposition-equivalence-quasi-coherent}），可知 $\mathcal{F}$ 是
某个自由 $\mathcal{O}_T$-模的直和项（定义于 $T_\etale$ 上）。故可假设
$\mathcal{F} = \bigoplus_{i \in I} \mathcal{O}_T$ 是自由模。
在这种情形下，$\mathcal{H} = \prod_{i \in I} \mathcal{G}$ 是拟凝聚模的积。
由《位点上的上同调》引理
\ref{sites-cohomology-lemma-cohomology-products} 可得 $H^1 = 0$，因为
仿射概形上拟凝聚层的上同调为零；见《下降》命题
\ref{descent-proposition-same-cohomology-quasi-coherent} 和《概形的上同调》引理
\ref{coherent-lemma-quasi-coherent-affine-cohomology-zero}。
\end{proof}

\begin{lemma}
\label{spaces-more-morphisms-lemma-formally-smooth}
设 $S$ 为概形。设 $f : X \to Y$ 为 $S$ 上代数空间的态射。
下列条件等价：
\begin{enumerate}
\item $f$ 形式光滑；
\item 对于每个图
$$
\xymatrix{
U \ar[d] \ar[r]_\psi & V \ar[d] \\
X \ar[r]^f & Y
}
$$
，其中 $U$ 和 $V$ 是概形且竖直箭头平展，概形态射 $\psi$ 都形式光滑
（即《态射进阶》定义
\ref{more-morphisms-definition-formally-unramified} 的意义）；% P09-ERR-0578
\item 对于一个竖直箭头为满射的这类图，态射 $\psi$ 形式光滑。
\end{enumerate}
\end{lemma}

\begin{proof}
引理 \ref{spaces-more-morphisms-lemma-helper-formally-smooth} 已经说明 (1) 蕴含 (2) 和 (3)。
假设 (3)。证明 $f$ 形式光滑的过程与引理
\ref{spaces-more-morphisms-lemma-smooth-formally-smooth} 中 (1) $\Rightarrow$ (2) 的证明完全类似。

\medskip\noindent
考虑定义 \ref{spaces-more-morphisms-definition-formally-smooth} 中实线部分交换的图
$$
\xymatrix{
X \ar[d]_f & T \ar[d]^i \ar[l]^a \\
Y & T' \ar[l] \ar@{-->}[lu]
}
$$
。我们将证明虚线箭头存在，从而证明 $f$ 形式光滑。
令 $\mathcal{F}$ 为引理 \ref{spaces-more-morphisms-lemma-sheaf} 在评注
\ref{spaces-more-morphisms-remark-special-case} 所述特殊情形下定义于
$(T')_{spaces, \etale}$ 上的集合层。令
$$
\mathcal{H} =
\SheafHom_{\mathcal{O}_T}(a^*\Omega_{X/Y}, \mathcal{C}_{T/T'})
$$
为 $\mathcal{O}_T$-模层（定义于 $T_{spaces, \etale}$ 上），并按引理
\ref{spaces-more-morphisms-lemma-action-sheaf} 取作用
$\mathcal{H} \times \mathcal{F} \to \mathcal{F}$。
这个作用 $\mathcal{H} \times \mathcal{F} \to \mathcal{F}$ 使
$\mathcal{F}$ 成为伪 $\mathcal{H}$-挠子；见《位点上的上同调》定义
\ref{sites-cohomology-definition-torsor}。目标是证明 $\mathcal{F}$ 是平凡的
$\mathcal{H}$-挠子。分两步：(I) 为证明 $\mathcal{F}$ 是挠子，必须证明
$\mathcal{F}$ 在平展局部有截面；(II) 为证明 $\mathcal{F}$ 是平凡挠子，
只需证明 $H^1(T_\etale, \mathcal{H}) = 0$；见《位点上的上同调》引理
\ref{sites-cohomology-lemma-torsors-h1}。

\medskip\noindent
先证明 (I)。为此考虑图（它因假设 (3) 而存在）
$$
\xymatrix{
U \ar[d] \ar[r]_\psi & V \ar[d] \\
X \ar[r]^f & Y
}
$$
，其中 $U$ 和 $V$ 是概形，竖直箭头平展且满，并且 $\psi$ 形式光滑。
由引理 \ref{spaces-more-morphisms-lemma-representable-property-formally-property}，态射
$V \to Y$ 形式平展。因此，由引理
\ref{spaces-more-morphisms-lemma-composition-formally-smooth-etale-unramified}，复合
$U \to Y$ 形式光滑。于是 (I) 由引理 \ref{spaces-more-morphisms-lemma-etale-on-top} 的第 (4) 部分得出。

\medskip\noindent
最后证明 (II)。由引理
\ref{spaces-more-morphisms-lemma-formally-smooth-sheaf-differentials} 可知
$\Omega_{U/V}$ 局部投射。% P09-ERR-0579
因此 $\Omega_{X/Y}$ 局部投射；见《空间上的下降》引理
\ref{spaces-descent-lemma-locally-projective-descends}。
故 $a^*\Omega_{X/Y}$ 局部投射；见《空间的性质》引理
\ref{spaces-properties-lemma-locally-projective-pullback}。因此
$$
H^1(T_\etale, \mathcal{H}) =
H^1(T_\etale,
\SheafHom_{\mathcal{O}_T}(a^*\Omega_{X/Y}, \mathcal{C}_{T/T'}) = 0
$$
% P09-ERR-0580
由引理 \ref{spaces-more-morphisms-lemma-h1-is-zero} 得到，正合所需。
\end{proof}

\begin{lemma}
\label{spaces-more-morphisms-lemma-descending-property-formally-smooth}
性质 $\mathcal{P}(f) =$“$f$ 形式光滑”在基底上是 fpqc 局部的。
\end{lemma}

\begin{proof}
设 $f : X \to Y$ 是某概形 $S$ 上代数空间的态射。
选取指标集 $I$ 和图
$$
\xymatrix{
U_i \ar[d] \ar[r]_{\psi_i} & V_i \ar[d] \\
X \ar[r]^f & Y
}
$$
，其中竖直箭头平展，且 $U_i$、$V_i$ 是仿射概形。此外，假设
$\coprod U_i \to X$ 和 $\coprod V_i \to Y$ 满射；见《空间的性质》引理
\ref{spaces-properties-lemma-cover-by-union-affines}。由引理
\ref{spaces-more-morphisms-lemma-formally-smooth} 可知，$f$ 形式光滑当且仅当每个态射
$\psi_i$ 都形式光滑。% P09-ERR-0581
因此约化到仿射概形的态射。在这种情形下，结论由《代数》引理
\ref{algebra-lemma-descent-formally-smooth} 得出。部分细节从略。
\end{proof}

\begin{lemma}
\label{spaces-more-morphisms-lemma-triangle-differentials-formally-smooth}
设 $S$ 为概形。设 $f : X \to Y$、$g : Y \to Z$ 是 $S$ 上代数空间的态射。
假设 $f$ 形式光滑。则
$$
0 \to f^*\Omega_{Y/Z} \to \Omega_{X/Z} \to \Omega_{X/Y} \to 0
$$
在引理 \ref{spaces-more-morphisms-lemma-triangle-differentials} 中是短正合的。% P09-ERR-0582
\end{lemma}

\begin{proof}
由概形情形和平展局部化得出；前者见《态射进阶》引理
\ref{more-morphisms-lemma-triangle-differentials-formally-smooth}，
后者见引理 \ref{spaces-more-morphisms-lemma-formally-smooth} 和
\ref{spaces-more-morphisms-lemma-localize-differentials}。
\end{proof}

\begin{lemma}
\label{spaces-more-morphisms-lemma-differentials-formally-unramified-formally-smooth}
设 $S$ 为概形，$B$ 为 $S$ 上的代数空间。设 $h : Z \to X$ 是
$B$ 上代数空间的形式非分歧态射。假设 $Z$ 在 $B$ 上形式光滑。
则引理 \ref{spaces-more-morphisms-lemma-universally-unramified-differentials-sequence} 中的
典范正合序列
$$
0 \to \mathcal{C}_{Z/X} \to h^*\Omega_{X/B} \to \Omega_{Z/B} \to 0
$$
是短正合的。
\end{lemma}

\begin{proof}
设 $Z \to Z'$ 是 $Z$ 在 $X$ 上的泛一阶加厚。由引理
\ref{spaces-more-morphisms-lemma-universally-unramified-differentials-sequence} 的证明可知，
我们的序列与下列序列相识别：
$$
\mathcal{C}_{Z/Z'} \to \Omega_{Z'/B} \otimes \mathcal{O}_Z \to
\Omega_{Z/B} \to 0.
$$
由于 $Z \to S$ 形式光滑，% P09-ERR-0583
可以在 $Z'$ 上平展局部地找到左逆 $Z' \to Z$；它在 $B$ 上左逆于
含入映射 $Z \to Z'$。因此该序列平展局部分裂；见引理
\ref{spaces-more-morphisms-lemma-differentials-relative-immersion-section}。
\end{proof}

\begin{lemma}
\label{spaces-more-morphisms-lemma-two-unramified-morphisms-formally-smooth}
设 $S$ 为概形。设
$$
\xymatrix{
Z \ar[r]_i \ar[rd]_j & X \ar[d]^f \\
& Y
}
$$
为 $S$ 上代数空间的交换图，其中 $i$ 和 $j$ 形式非分歧，且 $f$ 形式光滑。
则引理 \ref{spaces-more-morphisms-lemma-two-unramified-morphisms} 中的典范正合序列
$$
0 \to
\mathcal{C}_{Z/Y} \to
\mathcal{C}_{Z/X} \to
i^*\Omega_{X/Y} \to 0
$$
正合且局部分裂。
\end{lemma}

\begin{proof}
记 $Z \to Z'$ 为 $Z$ 在 $X$ 上的泛一阶加厚。
记 $Z \to Z''$ 为 $Z$ 在 $Y$ 上的泛一阶加厚。由引理
\ref{spaces-more-morphisms-lemma-universally-unramified-differentials-sequence}，% P09-ERR-0584
这里有典范态射 $Z' \to Z''$，% P09-ERR-0585
从而得到交换图
$$
\xymatrix{
Z \ar[r]_{i'} \ar[rd]_{j'} & Z' \ar[r]_a \ar[d]^k & X \ar[d]^f \\
& Z'' \ar[r]^b & Y
}
$$
。上述序列借助形式非分歧态射余法层的定义与序列
$$
\mathcal{C}_{Z/Z''} \to
\mathcal{C}_{Z/Z'} \to
(i')^*\Omega_{Z'/Z''} \to 0
$$
相识别。设 $U'' \to Z''$ 是平展态射且 $U''$ 仿射。
记 $U \to Z$ 和 $U' \to Z'$ 为相应的仿射概形，它们分别在
$Z$ 和 $Z'$ 上平展。由于 $f$ 形式光滑，存在态射 $h : U'' \to X$，
它与 $i$ 一致（限制到 $U$ 后），并且 $f \circ h$ 等于 $b|_{U''}$。
由于 $Z'$ 是泛一阶加厚，得到唯一态射 $g : U'' \to Z'$，使得
$g = a \circ h$。% P09-ERR-0586
$Z''$ 的泛性质蕴含 $k \circ g$ 是含入映射 $U'' \to Z''$。
因此 $g$ 是 $k$ 的左逆。图示为
$$
\xymatrix{
U \ar[d] \ar[r] & Z' \ar[d]^k \\
U'' \ar[r] \ar[ru]^g & Z''
}
$$
。于是 $g$ 诱导映射
$\mathcal{C}_{Z/Z'}|_U \to \mathcal{C}_{Z/Z''}|_U$，它是映射
$\mathcal{C}_{Z/Z''} \to \mathcal{C}_{Z/Z'}$ 在 $U$ 上的左逆。
\end{proof}











\section{诺特基底上的光滑性}
\label{spaces-more-morphisms-section-smooth-Noetherian}

\noindent
本节对应于《态射进阶》第
\ref{more-morphisms-section-smooth-Noetherian} 节。

\begin{lemma}
\label{spaces-more-morphisms-lemma-lifting-along-artinian-at-point}
设 $S$ 为概形。设 $f : X \to Y$ 是 $S$ 上代数空间的态射。
设 $x \in |X|$。假设 $Y$ 局部诺特且 $f$ 局部有限型。
下列条件等价：
\begin{enumerate}
\item $f$ 在 $x$ 处光滑；
\item 对于每个实线部分交换的图
$$
\xymatrix{
X \ar[d]_f & \Spec(B) \ar[d]^i \ar[l]^-\alpha \\
Y & \Spec(B') \ar[l]_-{\beta} \ar@{-->}[lu]
}
$$
，其中 $B' \to B$ 是局部环的满射，$\Ker(B' \to B)$ 的平方为零，
并且 $\alpha$ 把 $\Spec(B)$ 的闭点映到 $x$，都存在使图交换的虚线箭头；
\item 与 (2) 相同，但让 $B' \to B$ 遍历小扩张（见《代数》定义
\ref{algebra-definition-small-extension}）。
\end{enumerate}
\end{lemma}

\begin{proof}
条件 (1) 意味着存在开子空间 $X' \subset X$，使得 $X' \to Y$ 光滑。
因此，由引理 \ref{spaces-more-morphisms-lemma-smooth-formally-smooth}，(1) 蕴含 (2) 和 (3)。
条件 (2) 显然蕴含条件 (3)。假设 (3)。选取交换图
$$
\xymatrix{
X \ar[d] & U \ar[l] \ar[d] \\
Y & V \ar[l]
}
$$
，其中 $U$ 和 $V$ 仿射，水平箭头平展，并且存在点
$u \in U$ 映到 $x$。接着，考虑 (3) 中的图
$$
\xymatrix{
X \ar[d] & U \ar[l] \ar[d] & \Spec(B) \ar[d]^i \ar[l]^-\alpha  \\
Y & V \ar[l] & \Spec(B') \ar[l]_-{\beta}
}
$$
，不过现针对 $u \in U \to V$。令 $\gamma : \Spec(B') \to X$
为假设 (3) 对 $X$ 成立所给出的箭头。因为 $U \to X$ 平展，故形式平展
（引理 \ref{spaces-more-morphisms-lemma-etale-formally-etale}），态射 $\gamma$ 有唯一提升到 $U$，
并与 $\alpha$ 相容。又因为 $V \to Y$ 平展，因而形式平展，该提升与
$\beta$ 相容。因此 (3) 对 $u \in U \to V$ 成立，由《态射进阶》引理
\ref{more-morphisms-lemma-lifting-along-artinian-at-point} 可知
$U \to V$ 在 $u$ 处光滑。这证明 $X \to Y$ 在 $x$ 处光滑，证明完毕。
\end{proof}

\noindent
有时，只需对以有限型点为“中心”的小扩张检验提升判据这一事实很有用
（见《空间的态射》第
\ref{spaces-morphisms-section-points-finite-type} 节）。

\begin{lemma}
\label{spaces-more-morphisms-lemma-lifting-along-artinian}
设 $S$ 为概形。设 $f : X \to Y$ 是 $S$ 上代数空间的态射。
假设 $Y$ 局部诺特且 $f$ 局部有限型。下列条件等价：
\begin{enumerate}
\item $f$ 光滑；
\item 对于每个实线部分交换的图
$$
\xymatrix{
X \ar[d]_f & \Spec(B) \ar[d]^i \ar[l]^-\alpha \\
Y & \Spec(B') \ar[l]_-{\beta} \ar@{-->}[lu]
}
$$
，其中 $B' \to B$ 是阿廷局部环的小扩张，且 $\beta$ 有限型（！），
都存在使图交换的虚线箭头。
\end{enumerate}
\end{lemma}

\begin{proof}
若 $f$ 光滑，则无穷小提升判据（引理
\ref{spaces-more-morphisms-lemma-smooth-formally-smooth}）说明 $f$ 形式光滑，从而 (2) 成立。

\medskip\noindent
假设 $f$ 不光滑。点 $x \in X$ 中使 $f$ 不光滑者构成闭子集
$T$，它位于 $|X|$ 中。由《空间的态射》引理
\ref{spaces-morphisms-lemma-enough-finite-type-points}，存在点
$x \in T \subset X$，且 $x \in X_{\text{ft-pts}}$。选取交换图
$$
\xymatrix{
X \ar[d] & U \ar[l] \ar[d] & u \ar@{|->}[d] \\
Y & V \ar[l] & v
}
$$
，其中 $U$ 和 $V$ 仿射，水平箭头平展，并且存在点
$u \in U$ 映到 $x$。于是 $u$ 是 $U$ 的有限型点。因为
$U \to V$ 在点 $u$ 处不光滑，由《态射进阶》引理
\ref{more-morphisms-lemma-lifting-along-artinian-at-point}，存在图
$$
\xymatrix{
X \ar[d] & U \ar[l] \ar[d] & \Spec(B) \ar[d]^i \ar[l]^-\alpha  \\
Y & V \ar[l] & \Spec(B') \ar[l]_-{\beta} \ar@{-->}[lu]
}
$$
，其中 $B' \to B$ 是（阿廷）局部环的小扩张，$B$ 的剩余域等于
$\kappa(v)$，并且虚线箭头不存在。由于 $U \to V$ 有限型，
可知 $v$ 是 $V$ 的有限型点。由《态射》引理
\ref{morphisms-lemma-artinian-finite-type}，态射 $\beta$ 有限型，
因而复合 $\Spec(B) \to Y$ 也有限型。% P09-ERR-0587
完全仿照引理 \ref{spaces-more-morphisms-lemma-lifting-along-artinian-at-point} 的证明
（利用 $U \to X$ 和 $V \to Y$ 平展，因而形式平展），可知不存在
箭头 $\Spec(B) \to X$ 嵌入最后一个图的外矩形。% P09-ERR-0588
换言之，(2) 不成立，证明完毕。
\end{proof}

\noindent
下面是一个有用的应用。

\begin{lemma}
\label{spaces-more-morphisms-lemma-check-smoothness-on-infinitesimal-nbhds}
设 $S$ 为概形。设 $f : X \to Y$ 是 $S$ 上代数空间的态射。
假设 $f$ 局部有限型且 $Y$ 局部诺特。设 $Z \subset Y$ 为闭子空间，
其 $n$ 阶无穷小邻域为 $Z_n \subset Y$。令
$X_n = Z_n \times_Y X$。
\begin{enumerate}
\item 若 $X_n \to Z_n$ 对所有 $n$ 都光滑，则 $f$ 在
$f^{-1}(Z)$ 的每一点处光滑。
\item 若 $X_n \to Z_n$ 对所有 $n$ 都平展，则 $f$ 在
$f^{-1}(Z)$ 的每一点处平展。
\end{enumerate}
\end{lemma}

\begin{proof}
假设 $X_n \to Z_n$ 对所有 $n$ 都光滑。设 $x \in X$ 位于 $Z$ 的某点上方。
给定引理 \ref{spaces-more-morphisms-lemma-lifting-along-artinian-at-point} 第 (3) 部分中的
小扩张 $B' \to B$ 和态射 $\alpha$、$\beta$，$B'$ 的极大理想幂零
（因为 $B'$ 是阿廷环），所以态射 $\beta$ 通过 $Z_n$ 分解，而 $\alpha$
通过 $X_n$ 分解；对某个适当的 $n$ 会出现这种情形。于是
$X_n \to Z_n$ 的提升性质
给出图中所需的虚线箭头。这证明 (1)。第 (2) 部分由 (1) 以及下述事实得出：
态射平展当且仅当它光滑且相对维数为 $0$。
\end{proof}











\section{朴素余切复形}
\label{spaces-more-morphisms-section-netherlander}

\noindent
本节接续《位点上的模》第 \ref{sites-modules-section-netherlander} 节，
后者又接续《代数》第 \ref{algebra-section-netherlander} 节中的讨论。

\begin{definition}
\label{spaces-more-morphisms-definition-netherlander}
设 $S$ 为概形。设 $f : X \to Y$ 是 $S$ 上代数空间的态射。
{\it $f$ 的朴素余切复形}是《位点上的模》定义
\ref{sites-modules-definition-cotangent-complex-morphism-ringed-topoi}
对带环拓扑斯态射 $f_{small}$ 所定义的复形；该态射位于 $X$ 与 $Y$ 的
小平展位点之间；
见《空间的性质》引理
\ref{spaces-properties-lemma-morphism-ringed-topoi}。记作 $\NL_f$ 或
$\NL_{X/Y}$。
\end{definition}

\noindent
下面几个引理说明，这一定义与环同态和概形的定义相容，并且
$\NL_{X/Y}$ 是 $D_\QCoh(\mathcal{O}_X)$ 的对象。

\begin{lemma}
\label{spaces-more-morphisms-lemma-NL-etale-localization}
设 $S$ 为概形。考虑交换图
$$
\xymatrix{
U \ar[d]_p \ar[r]_g & V \ar[d]^q \\
X \ar[r]^f & Y
}
$$
，它是 $S$ 上代数空间的图，且 $p$ 和 $q$ 平展。则
有典范识别 $\NL_{X/Y}|_{U_\etale} = \NL_{U/V}$，它位于
$D(\mathcal{O}_U)$ 中。
\end{lemma}

\begin{proof}
朴素余切复形的形成与拉回交换（《位点上的模》引理
\ref{sites-modules-lemma-pullback-NL}），并且有
$p_{small}^{-1}\mathcal{O}_X = \mathcal{O}_U$ 以及
$g_{small}^{-1}\mathcal{O}_{V_\etale} =
p_{small}^{-1}f_{small}^{-1}\mathcal{O}_{Y_\etale}$，
因为由《空间的性质》引理
\ref{spaces-properties-lemma-etale-exact-pullback} 有
$q_{small}^{-1}\mathcal{O}_{Y_\etale} =
\mathcal{O}_{V_\etale}$。逐一定义追踪可得
$\NL_{X/Y}|_{U_\etale} = \NL_{U/V}$。
\end{proof}

\begin{lemma}
\label{spaces-more-morphisms-lemma-NL-compare-spaces-schemes}
设 $S$ 为概形。设 $f : X \to Y$ 是 $S$ 上代数空间的态射。
假设 $X$ 和 $Y$ 分别由概形 $X_0$ 和 $Y_0$ 表示。% P09-ERR-0589
则有典范识别 $\NL_{X/Y} = \epsilon^*\NL_{X_0/Y_0}$；它位于
$D(\mathcal{O}_X)$ 中，其中 $\epsilon$ 如《空间的导出范畴》
第 \ref{spaces-perfect-section-derived-quasi-coherent-etale} 节所述，
而 $\NL_{X_0/Y_0}$ 如《态射进阶》定义
\ref{more-morphisms-definition-netherlander} 所述。
\end{lemma}

\begin{proof}
设 $f_0 : X_0 \to Y_0$ 为与 $f$ 对应的概形态射。有典范映射
$\epsilon^{-1}f_0^{-1}\mathcal{O}_{Y_0} \to f_{small}^{-1}\mathcal{O}_Y$，
它与
$\epsilon^\sharp : \epsilon^{-1}\mathcal{O}_{X_0} \to \mathcal{O}_X$
相容，因为有交换图
$$
\xymatrix{
X_{0, Zar} \ar[d]_{f_0} & X_\etale \ar[l]^\epsilon \ar[d]^f \\
Y_{0, Zar} & Y_\etale \ar[l]_\epsilon
}
$$
；见《空间的导出范畴》评注
\ref{spaces-perfect-remark-match-total-direct-images}。因此由朴素余切复形的
函子性得到典范映射
$$
\epsilon^{-1}\NL_{X_0/Y_0} =
\epsilon^{-1}\NL_{\mathcal{O}_{X_0}/f_0^{-1}\mathcal{O}_{Y_0}} =
\NL_{\epsilon^{-1}\mathcal{O}_{X_0}/\epsilon^{-1}f_0^{-1}\mathcal{O}_{Y_0}}
\to
\NL_{\mathcal{O}_X/f^{-1}_{small}\mathcal{O}_Y} = \NL_{X/Y}
$$
。
为看出诱导映射 $\epsilon^*\NL_{X_0/Y_0} \to \NL_{X/Y}$ 在
$D(\mathcal{O}_X)$ 中是同构，可在几何点的茎上检验（《空间的性质》定理
\ref{spaces-properties-theorem-exactness-stalks}）。设
$\overline{x} : \Spec(k) \to X_0$ 是位于 $x \in X_0$ 上方的几何点，
并令 $\overline{y} = f \circ \overline{x}$ 位于 $y \in Y_0$ 上方。则
$$
\NL_{X/Y, \overline{x}} =
\NL_{\mathcal{O}_{X, \overline{x}}/\mathcal{O}_{Y, \overline{y}}}
$$
。
这是因为在 $\overline{x}$ 处取茎等同于经
$\overline{x} : \Spec(k) \to X$ 取逆像，因而可以应用《位点上的模》引理
\ref{sites-modules-lemma-pullback-NL}。另一方面，有
$$
(\epsilon^*\NL_{X_0/Y_0})_{\overline{x}} =
\NL_{X_0/Y_0, x} \otimes_{\mathcal{O}_{X_0, x}}
\mathcal{O}_{X, \overline{x}} =
\NL_{\mathcal{O}_{X_0, x}/\mathcal{O}_{Y_0, y}}
\otimes_{\mathcal{O}_{X_0, x}} \mathcal{O}_{X, \overline{x}}
$$
。
略去一些细节（提示：使用拉回的茎等于像点处的茎这一事实，见《位点》引理
\ref{sites-lemma-point-morphism-sites}；也使用相应的模版本，见《位点上的模》
引理 \ref{sites-modules-lemma-pullback-stalk}）。注意
$\mathcal{O}_{X, \overline{x}}$ 是 $\mathcal{O}_{X_0, x}$ 的严格亨泽尔化，
$\mathcal{O}_{Y, \overline{y}}$ 也类似（《空间的性质》引理
\ref{spaces-properties-lemma-describe-etale-local-ring}）。因此结论由《代数进阶》
引理 \ref{more-algebra-lemma-henselization-NL} 得出。
\end{proof}

\begin{lemma}
\label{spaces-more-morphisms-lemma-netherlander-quasi-coherent}
设 $S$ 为概形。设 $f : X \to Y$ 是 $S$ 上代数空间的态射。
复形 $\NL_{X/Y}$ 的上同调层拟凝聚，在次数 $-1$、$0$ 之外为零，
并且等于 $\Omega_{X/Y}$（在次数 $0$）。
\end{lemma}

\begin{proof}
由《位点上的模》第 \ref{sites-modules-section-netherlander} 节中朴素余切复形的
构造可知，$\NL_{X/Y}$ 是位于次数 $-1$、$0$ 的复形，其次数 $0$ 上同调为
$\Omega_{X/Y}$（根据第 \ref{spaces-more-morphisms-section-sheaf-differentials} 节中
$\Omega_{X/Y}$ 的构造）。微分层拟凝聚（引理
\ref{spaces-more-morphisms-lemma-module-differentials-quasi-coherent}）。为完成证明，只需说明
$H^{-1}(\NL_{X/Y})$ 拟凝聚。这可通过平展局部检验得出：引理
\ref{spaces-more-morphisms-lemma-NL-etale-localization} 和《空间的性质》引理
\ref{spaces-properties-lemma-characterize-quasi-coherent} 允许这样做；
再由引理 \ref{spaces-more-morphisms-lemma-NL-compare-spaces-schemes} 约化到概形情形；
最后使用概形情形的结果，即《态射进阶》引理
\ref{more-morphisms-lemma-netherlander-quasi-coherent}。
\end{proof}

\begin{lemma}
\label{spaces-more-morphisms-lemma-netherlander-fp}
设 $S$ 为概形。设 $f : X \to Y$ 是 $S$ 上代数空间的态射。
若 $f$ 局部有限表示，则 $\NL_{X/Y}$ 在 $X$ 上平展局部拟同构于复形
$$
\ldots \to 0 \to \mathcal{F}^{-1} \to \mathcal{F}^0 \to 0 \to \ldots
$$
，其各项为拟凝聚 $\mathcal{O}_X$-模，并且 $\mathcal{F}^0$ 有限表示，
$\mathcal{F}^{-1}$ 有限型。
\end{lemma}

\begin{proof}
由引理 \ref{spaces-more-morphisms-lemma-NL-etale-localization}，朴素余切复形的形成与平展局部化交换。
再由引理 \ref{spaces-more-morphisms-lemma-NL-compare-spaces-schemes} 约化到概形情形。
概形情形的结论是《态射进阶》引理
\ref{more-morphisms-lemma-netherlander-fp}。
\end{proof}

\begin{lemma}
\label{spaces-more-morphisms-lemma-NL-formally-smooth}
设 $S$ 为概形。设 $f : X \to Y$ 是 $S$ 上代数空间的态射。
下列条件等价：
\begin{enumerate}
\item $f$ 形式光滑；
\item $H^{-1}(\NL_{X/Y}) = 0$，并且
$H^0(\NL_{X/Y}) = \Omega_{X/Y}$ 局部投射。
\end{enumerate}
\end{lemma}

\begin{proof}
这由引理 \ref{spaces-more-morphisms-lemma-formally-smooth}、
引理 \ref{spaces-more-morphisms-lemma-NL-etale-localization}、
引理 \ref{spaces-more-morphisms-lemma-NL-compare-spaces-schemes} 以及概形情形得出；
后者即《态射进阶》引理
\ref{more-morphisms-lemma-NL-formally-smooth}。
\end{proof}

\begin{lemma}
\label{spaces-more-morphisms-lemma-NL-formally-etale}
设 $f : X \to Y$ 为概形态射。% P09-ERR-0590
下列条件等价：
\begin{enumerate}
\item $f$ 形式平展；
\item $H^{-1}(\NL_{X/Y}) = H^0(\NL_{X/Y}) = 0$。
\end{enumerate}
\end{lemma}

\begin{proof}
假设 (1)。形式平展态射是形式光滑态射。因此，由引理
\ref{spaces-more-morphisms-lemma-NL-formally-smooth}，$H^{-1}(\NL_{X/Y}) = 0$。
另一方面，形式平展态射是形式非分歧的，% P09-ERR-0591
故由引理 \ref{spaces-more-morphisms-lemma-characterize-formally-unramified} 有
$\Omega_{X/Y} = 0$。反过来，若 (2) 成立，则由引理
\ref{spaces-more-morphisms-lemma-NL-formally-smooth}，$f$ 形式光滑；由引理
\ref{spaces-more-morphisms-lemma-characterize-formally-unramified}，它形式非分歧；
因而由引理 \ref{spaces-more-morphisms-lemma-formally-etale-unramified-smooth}，它形式平展。% P09-ERR-0592
\end{proof}

\begin{lemma}
\label{spaces-more-morphisms-lemma-NL-smooth}
设 $f : X \to Y$ 为概形态射。% P09-ERR-0593
下列条件等价：
\begin{enumerate}
\item $f$ 光滑；
\item $f$ 局部有限表示，$H^{-1}(\NL_{X/Y}) = 0$，并且
$H^0(\NL_{X/Y}) = \Omega_{X/Y}$ 有限局部自由。
\end{enumerate}
\end{lemma}

\begin{proof}
这由引理 \ref{spaces-more-morphisms-lemma-formally-smooth}、
引理 \ref{spaces-more-morphisms-lemma-NL-etale-localization}、
引理 \ref{spaces-more-morphisms-lemma-NL-compare-spaces-schemes} 以及概形情形得出；
后者即《态射进阶》引理
\ref{more-morphisms-lemma-NL-smooth}。
\end{proof}











\section{平坦轨迹的开性}
\label{spaces-more-morphisms-section-open-flat}

\noindent
本节对应于《态射进阶》第
\ref{more-morphisms-section-open-flat} 节。% P09-ERR-0594
注意，《空间的态射》第 \ref{spaces-morphisms-section-flat-modules} 节已经定义了
代数空间上拟凝聚模的平坦性。

\begin{theorem}
\label{spaces-more-morphisms-theorem-openness-flatness}
设 $S$ 为概形。设 $f : X \to Y$ 是 $S$ 上代数空间的态射。
设 $\mathcal{F}$ 为 $X$ 上的拟凝聚层。假设 $f$ 局部有限表示，
且 $\mathcal{F}$ 是局部有限表示的 $\mathcal{O}_X$-模。则
$$
\{x \in |X| : \mathcal{F}\text{ is flat over }Y\text{ 在 }x\}
$$
在 $|X|$ 中开。
\end{theorem}

\begin{proof}
选取交换图
$$
\xymatrix{
U \ar[d]_p \ar[r]_\alpha &
V \ar[d]^q \\
X \ar[r]^a & Y
}
$$
，其中 $U$、$V$ 是概形，而 $p$、$q$ 如《空间》引理
\ref{spaces-lemma-lift-morphism-presentations} 中那样满且平展。
由《态射进阶》定理
\ref{more-morphisms-theorem-openness-flatness}，集合
$U' = \{u \in |U| : p^*\mathcal{F}\text{ is flat over }V\text{ 在 }u\}$
在 $U$ 中开。由《空间的态射》定义
\ref{spaces-morphisms-definition-flat-module}，$U'$ 在 $|X|$ 中的像就是定理中的
集合。最后，映射 $|U| \to |X|$ 是开映射，故结论成立；见《空间的性质》引理
\ref{spaces-properties-lemma-topology-points}。
\end{proof}

\begin{lemma}
\label{spaces-more-morphisms-lemma-flat-locus-base-change}
设 $S$ 为概形。设
$$
\xymatrix{
X' \ar[r]_{g'} \ar[d]_{f'} & X \ar[d]^f \\
Y' \ar[r]^g & Y
}
$$
是 $S$ 上代数空间的笛卡尔图。设 $\mathcal{F}$ 是拟凝聚
$\mathcal{O}_X$-模。假设 $g$ 平坦，$f$ 局部有限表示，并且
$\mathcal{F}$ 局部有限表示。则
$$
\{x' \in |X'| : (g')^*\mathcal{F}\text{ is flat over }Y'\text{ 在 }x'\}
$$
是定理 \ref{spaces-more-morphisms-theorem-openness-flatness} 中开子集在连续映射
$|g'| : |X'| \to |X|$ 下的逆像。
\end{lemma}

\begin{proof}
这由《空间的态射》引理
\ref{spaces-morphisms-lemma-base-change-module-flat} 得出。
\end{proof}












\section{平坦性的纤维判据}
\label{spaces-more-morphisms-section-criterion-flat-fibres}

\noindent
设 $S$ 为概形。考虑 $S$ 上代数空间的交换图
$$
\xymatrix{
X \ar[rr]_f \ar[dr]_g & & Y \ar[dl]^h \\
& Z
}
$$
以及拟凝聚 $\mathcal{O}_X$-模 $\mathcal{F}$。
给定一点 $x \in |X|$，我们考察 $\mathcal{F}$ 是否相对于
$Y$ 在 $x$ 处平坦。若 $\mathcal{F}$ 相对于 $Z$ 在 $x$ 处平坦，
则下述定理说明，
这个问题与如下问题密切相关：$\mathcal{F}$ 在 $X \to Z$ 位于
$g(x)$ 上的纤维上的限制，是否相对于 $Y \to Z$ 位于 $g(x)$ 上的纤维
平坦。为使这一说法有意义，我们先给出如下引理。

\begin{lemma}
\label{spaces-more-morphisms-lemma-flat-on-fibres-at-point}
在上述情形中，下列条件等价：
\begin{enumerate}
\item 选取几何点 $\overline{x}$；它是 $X$ 的几何点并位于 $x$ 上方。
令 $\overline{y} = f \circ \overline{x}$ 且
$\overline{z} = g \circ \overline{x}$。则模
$\mathcal{F}_{\overline{x}}/
\mathfrak m_{\overline{z}}\mathcal{F}_{\overline{x}}$
在
$\mathcal{O}_{Y, \overline{y}}/
\mathfrak m_{\overline{z}}\mathcal{O}_{Y, \overline{y}}$
上平坦。
\item 选取态射 $x : \Spec(K) \to X$；它位于 $x$ 的等价类中。
令 $z = g \circ x$、$X_z = \Spec(K) \times_{z, Z} X$、
$Y_z = \Spec(K) \times_{z, Z} Y$，并令 $\mathcal{F}_z$ 为
$\mathcal{F}$ 到 $X_z$ 的拉回。则 $\mathcal{F}_z$ 在 $x$ 处相对于
$Y_z$ 平坦（如《空间的态射》定义
\ref{spaces-morphisms-definition-flat-module} 所定义）。
\item 选取交换图
% P09-ERR-0595
$$
\xymatrix{
& & & U \ar[llld]_a \ar[rr] \ar[dr] & & V \ar[llld]_>>>>>>>b \ar[dl] \\
X \ar[rr]_f \ar[dr]_g & & Y \ar[dl]^h &  & W \ar[llld]_c \\
& Z
}
$$
，其中 $U, V, W$ 是概形，$a, b, c$ 是平展的；并选取一点
$u \in U$，它映到 $x$。令 $w \in W$ 为 $u$ 的像。令
$\mathcal{F}_w$ 为 $\mathcal{F}$ 到纤维 $U_w$ 上的拉回；这是
$U \to W$ 在 $w$ 处的纤维。则 $\mathcal{F}_w$ 相对于 $V_w$
在 $u$ 处平坦。
\end{enumerate}
\end{lemma}

\begin{proof}
注意，在 (2) 中，态射 $x : \Spec(K) \to X$ 给出一个
$K$-有理点，它属于 $X_z$，故该陈述有意义。此外，(2) 中的条件与
$\Spec(K) \to X$ 的选择无关，只要该选择属于 $x$ 的等价类
（略去细节；这也将由下述论证推出，因为
其余条件不依赖于这个选择）。还要注意，我们总能如下选取 (3) 中的图：
首先选取概形 $W$ 以及满平展态射 $W \to Z$，再选取概形 $V$ 以及
满平展态射 $V \to W \times_Z Y$，最后选取概形 $U$ 以及满平展态射
$U \to V \times_Y X$。作出这些选择后，令 $U \to W$ 等于复合
$U \to V \to W$；可以选取一点 $u \in U$ 映到 $x$，因为态射
$U \to X$ 是满的。

\medskip\noindent
假设同时给定 (3) 中的图以及 (1) 中的几何点
$\overline{x} : \Spec(k) \to X$。由《空间的性质》引理
\ref{spaces-properties-lemma-geometric-lift-to-usual}，
可选取几何点 $\overline{u} : \Spec(k) \to U$，它位于 $u$ 上方，
使得 $\overline{x} = a \circ \overline{u}$。
用 $\overline{v} : \Spec(k) \to V$ 和
$\overline{w} : \Spec(k) \to W$ 表示由此诱导的 $V$ 和 $W$ 的几何点。
在此情形下，我们知道
$\mathcal{O}_{X, \overline{x}} = \mathcal{O}_{U, u}^{sh}$，
对 $Y$ 和 $Z$ 亦然；见《空间的性质》引理
\ref{spaces-properties-lemma-describe-etale-local-ring}。
同样有
$$
\mathcal{F}_{\overline{x}} =
(a^*\mathcal{F})_u \otimes_{\mathcal{O}_{U, u}}
\mathcal{O}_{U, u}^{sh}
$$
；见《空间的性质》引理
\ref{spaces-properties-lemma-stalk-quasi-coherent}。
注意，$\mathcal{F}_w$ 在 $u$ 处的茎为
$$
(\mathcal{F}_w)_u = (a^*\mathcal{F})_u/\mathfrak m_w(a^*\mathcal{F})_u
$$
，而 $V_w$ 在 $v$ 处的局部环为
$$
\mathcal{O}_{V_w, v} = \mathcal{O}_{V, v}/\mathfrak m_w\mathcal{O}_{V, v}.
$$
由于 $\mathfrak m_{\overline{z}} =
\mathfrak m_w \mathcal{O}_{Z, \overline{z}} =
\mathfrak m_w \mathcal{O}_{W, w}^{sh}$，可知
\begin{align*}
\mathcal{F}_{\overline{x}}/
\mathfrak m_{\overline{z}}\mathcal{F}_{\overline{x}} & =
(a^*\mathcal{F})_u \otimes_{\mathcal{O}_{U, u}}
\mathcal{O}_{X, \overline{x}}/
\mathfrak m_{\overline{z}}\mathcal{O}_{X, \overline{x}} \\
& =
(\mathcal{F}_w)_u \otimes_{\mathcal{O}_{U_w, u}}
\mathcal{O}_{U, u}^{sh}/\mathfrak m_w\mathcal{O}_{U, u}^{sh} \\
& = (\mathcal{F}_w)_u \otimes_{\mathcal{O}_{U_w, u}}
\mathcal{O}_{U_w, \overline{u}}^{sh} \\
& = (\mathcal{F}_w)_{\overline{u}}
\end{align*}
；倒数第二个等号由《代数》引理
\ref{algebra-lemma-quotient-strict-henselization} 得出，
最后一个等号由《空间的性质》引理
\ref{spaces-properties-lemma-stalk-quasi-coherent} 得出。
对 $V$ 和 $Y$ 的结构层应用同样的论证，得到
$$
\mathcal{O}_{V_w, \overline{v}}^{sh} =
\mathcal{O}_{V, v}^{sh}/\mathfrak m_w \mathcal{O}_{V, v}^{sh} =
\mathcal{O}_{Y, \overline{y}}/
\mathfrak m_{\overline{z}}\mathcal{O}_{Y, \overline{y}}.
$$
现在应用《空间的态射》引理
\ref{spaces-morphisms-lemma-flat-at-point}，可知 (1) 与 (3) 等价。

\medskip\noindent
最后证明 (2) 与 (3) 等价。
为此，选取域扩张 $\tilde K$；它扩张 $K$。再选取态射
$\tilde x : \Spec(\tilde K) \to U$，它位于 $u$ 上方（这是可能的，因为
$u \times_{X, x} \Spec(K)$ 是非空概形）。令
$\tilde z : \Spec(\tilde K) \to U \to W$ 为该复合。% P09-ERR-0596
我们得到交换图
% P09-ERR-0595
$$
\xymatrix{
& & & U_w \times_w \tilde z \ar[llld]_a \ar[rr] \ar[dr] & &
V_w \times_w \tilde z \ar[llld]_>>>>>>>b \ar[dl] \\
X_z \ar[rr]_f \ar[dr]_g & & Y_z \ar[dl]^h &  & \tilde z \ar[llld]_c \\
& z
}
$$
，其中 $z = \Spec(K)$ 且 $w = \Spec(\kappa(w))$。显然，
$\mathcal{F}_w$ 与 $\mathcal{F}_z$ 在 $U_w \times_w \tilde z$ 上的拉回
是同一个模。由此得到交换图
$$
\xymatrix{
X_z \ar[d] & U_w \times_w \tilde z \ar[l] \ar[d] \ar[r] & U_w \ar[d] \\
Y_z & V_w \times_w \tilde z \ar[l] \ar[r] & V_w
}
$$
；它的两个方块都是笛卡尔的，且底部两个水平箭头都是平坦的：
左下水平箭头是态射
$Y \times_Z \tilde z \to Y \times_Z z = Y_z$（平坦态射的基变换）、
平展态射 $V \times_Z \tilde z \to Y \times_Z \tilde z$ 以及
平展态射 $V \times_W \tilde z \to V \times_Z \tilde z$ 的复合。
因此，由《空间的态射》引理
\ref{spaces-morphisms-lemma-base-change-module-flat}，有
$$
\mathcal{F}_z\text{ flat at }x\text{ 相对于 }Y_z
\Leftrightarrow
\mathcal{F}|_{U_w \times_w \tilde z}
\text{ flat at }\tilde x\text{ 相对于 }V_w \times_w \tilde z
\Leftrightarrow
\mathcal{F}_w\text{ flat at }u\text{ 相对于 }V_w
$$
，于是结论成立。
\end{proof}

\begin{definition}
\label{spaces-more-morphisms-definition-module-flat-on-fibre}
设 $S$ 为概形。设 $X \to Y \to Z$ 是 $S$ 上代数空间的态射。
设 $\mathcal{F}$ 为拟凝聚 $\mathcal{O}_X$-模。
设 $x \in |X|$ 为一点，并用 $z \in |Z|$ 表示它的像。
\begin{enumerate}
\item 我们称 {\it $\mathcal{F}$ 到其位于 $z$ 上的纤维的限制在 $x$ 处
相对于 $Y$ 位于 $z$ 上的纤维平坦}，如果引理
\ref{spaces-more-morphisms-lemma-flat-on-fibres-at-point} 的等价条件成立。
\item 我们称 {\it $X$ 位于 $z$ 上的纤维在 $x$ 处相对于 $Y$ 位于
$z$ 上的纤维平坦}，如果引理 \ref{spaces-more-morphisms-lemma-flat-on-fibres-at-point}
的等价条件在 $\mathcal{F} = \mathcal{O}_X$ 时成立。
\item 我们称 {\it $X$ 位于 $z$ 上的纤维相对于 $Y$ 位于 $z$ 上的纤维
平坦}，如果对所有 $x \in |X|$，只要它位于 $z$ 上方，$X$ 位于
$z$ 上的纤维在 $x$ 处相对于 $Y$ 位于 $z$ 上的纤维就平坦
% P09-ERR-0597
\end{enumerate}
\end{definition}

\noindent
有了这个定义，可以如下陈述该判据的一个版本。诺特情形见第
\ref{spaces-more-morphisms-section-flat-over-Noetherian} 节。

\begin{theorem}
\label{spaces-more-morphisms-theorem-criterion-flatness-fibre}
设 $S$ 为概形。
设 $f : X \to Y$ 和 $Y \to Z$ 是 $S$ 上代数空间的态射。
设 $\mathcal{F}$ 为拟凝聚 $\mathcal{O}_X$-模。
假设
\begin{enumerate}
\item $X$ 在 $Z$ 上局部有限表示；
\item $\mathcal{F}$ 是有限表示的 $\mathcal{O}_X$-模；并且 % P09-ERR-0598
\item $Y$ 在 $Z$ 上局部有限型。
\end{enumerate}
设 $x \in |X|$，并令 $y \in |Y|$ 和 $z \in |Z|$ 为 $x$ 的像。
若 $\mathcal{F}_{\overline{x}} \not = 0$，则下列条件等价：% P09-ERR-0599
\begin{enumerate}
\item $\mathcal{F}$ 相对于 $Z$ 在 $x$ 处平坦，且
$\mathcal{F}$ 到其位于 $z$ 上的纤维的限制在 $x$ 处相对于
$Y$ 位于 $z$ 上的纤维平坦；
\item $Y$ 相对于 $Z$ 在 $y$ 处平坦，且 $\mathcal{F}$ 相对于 $Y$
在 $x$ 处平坦。
\end{enumerate}
此外，(1) 和 (2) 成立的点 $x$ 所成的集合在
$\text{Supp}(\mathcal{F})$ 中开。
\end{theorem}

\begin{proof}
选取引理 \ref{spaces-more-morphisms-lemma-flat-on-fibres-at-point} 第 (3) 项中的图。
由定义可知，这把问题化归为概形态射 $U \to V \to W$、
拟凝聚层 $a^*\mathcal{F}$ 以及点 $u \in U$ 的相应定理。
因此，该定理由概形的相应结果，即《态射进阶》定理
\ref{more-morphisms-theorem-criterion-flatness-fibre} 得出。
\end{proof}

\begin{lemma}
\label{spaces-more-morphisms-lemma-morphism-between-flat}
设 $S$ 为概形。
设 $f : X \to Y$ 和 $Y \to Z$ 是 $S$ 上代数空间的态射。% P09-ERR-0600
假设
\begin{enumerate}
\item $X$ 在 $Z$ 上局部有限表示；
\item $X$ 在 $Z$ 上平坦；
\item 对每个 $z \in |Z|$，$X$ 位于 $z$ 上的纤维相对于
$Y$ 位于 $z$ 上的纤维平坦；并且
\item $Y$ 在 $Z$ 上局部有限型。
\end{enumerate}
则 $f$ 平坦。若 $f$ 还是满的，则 $Y$ 在 $Z$ 上平坦。
\end{lemma}

\begin{proof}
这是定理 \ref{spaces-more-morphisms-theorem-criterion-flatness-fibre} 的特殊情形。
\end{proof}

\begin{lemma}
\label{spaces-more-morphisms-lemma-base-change-criterion-flatness-fibre}
设 $S$ 为概形。设 $f : X \to Y$ 和 $Y \to Z$ 是 $S$ 上代数空间的
态射。设 $\mathcal{F}$ 为拟凝聚 $\mathcal{O}_X$-模。
假设
\begin{enumerate}
\item $X$ 在 $Z$ 上局部有限表示；
\item $\mathcal{F}$ 是有限表示的 $\mathcal{O}_X$-模；% P09-ERR-0598
\item $\mathcal{F}$ 在 $Z$ 上平坦；并且
\item $Y$ 在 $Z$ 上局部有限型。
\end{enumerate}
则集合
% P09-ERR-0601
$$
A = \{x \in |X| : \mathcal{F} \text{ flat at }x \text{ 相对于 }Y\}.
$$
在 $|X|$ 中开，且它的形成与任意基变换可交换：
若 $Z' \to Z$ 是代数空间的态射，并且 $A'$ 是
$X' = X \times_Z Z'$ 中使
$\mathcal{F}' = \mathcal{F} \times_Z Z'$ 在
$Y' = Y \times_Z Z'$ 上平坦的点所成的集合，% P09-ERR-0602
则 $A'$ 是 $A$ 在连续映射 $|X'| \to |X|$ 下的逆像。
\end{lemma}

\begin{proof}
一种证明方式是把《态射进阶》引理
\ref{more-morphisms-lemma-morphism-between-flat} 中的证明
翻译到代数空间的范畴中。这里我们改为把它化归到概形情形。
具体而言，选取引理 \ref{spaces-more-morphisms-lemma-flat-on-fibres-at-point} 第 (3) 项中的图，
使得 $a$、$b$、$c$ 都是满的。
由定义可知，这把问题化归为概形态射 $U \to V \to W$、
拟凝聚层 $a^*\mathcal{F}$ 以及点 $u \in U$ 的相应定理。
唯一需要说明的小点是：给定代数空间的态射 $Z' \to Z$，
选取概形 $W'$ 以及满平展态射 $W' \to W \times_Z Z'$。
然后令 $U' = W' \times_W U$ 且 $V' = W' \times_W V$。
用 $a', b', c'$ 表示从 $U', V', W'$ 到 $X', Y', Z'$ 的态射。
在这种情形下，$A$ 和 $A'$ 分别是 $U$ 和 $U'$ 中与
$a^*\mathcal{F}$ 和 $(a')^*\mathcal{F}'$ 相关的开子集的像。
这确实把引理化归为《态射进阶》引理
\ref{more-morphisms-lemma-morphism-between-flat}。
\end{proof}

\begin{lemma}
\label{spaces-more-morphisms-lemma-base-change-flatness-fibres}
设 $S$ 为概形。
设 $f : X \to Y$ 和 $Y \to Z$ 是 $S$ 上代数空间的态射。% P09-ERR-0600
假设
\begin{enumerate}
\item $X$ 在 $Z$ 上局部有限表示；
\item $X$ 在 $Z$ 上平坦；并且
\item $Y$ 在 $Z$ 上局部有限型。
\end{enumerate}
则集合
% P09-ERR-0601
$$
\{x \in |X| : X\text{ flat at }x \text{ 相对于 }Y\}.
$$
在 $|X|$ 中开，且它的形成与任意基变换 $Z' \to Z$ 可交换。
\end{lemma}

\begin{proof}
这是引理 \ref{spaces-more-morphisms-lemma-base-change-criterion-flatness-fibre} 的特殊情形。
\end{proof}

\begin{lemma}
\label{spaces-more-morphisms-lemma-flat-and-free-at-point-fibre}
设 $S$ 为概形。设 $f : X \to Y$ 是 $S$ 上代数空间的态射，
并且局部有限表示。
设 $\mathcal{F}$ 是有限表示的 $\mathcal{O}_X$-模。
设 $x \in |X|$，其像为 $y \in |Y|$。若 $\mathcal{F}$ 在 $x$ 处
相对于 $Y$ 平坦，则下列条件等价：% P09-ERR-0603
\begin{enumerate}
\item $(\mathcal{F}_{\overline{y}})_{\overline{x}}$ 是平坦的
$\mathcal{O}_{X_{\overline{y}}, \overline{x}}$-模；
\item $(\mathcal{F}_{\overline{y}})_{\overline{x}}$ 是自由的
$\mathcal{O}_{X_{\overline{y}}, \overline{x}}$-模；
\item $\mathcal{F}_{\overline{y}}$ 在 $\overline{x}$ 的某个平展邻域上
有限自由；该邻域位于 $X_{\overline{y}}$ 中；并且
\item $\mathcal{F}$ 在 $x$ 的某个平展邻域上有限自由；该邻域位于
$X$ 中。
\end{enumerate}
这里 $\overline{x}$ 是 $X$ 的一个位于 $x$ 上方的几何点，
且 $\overline{y} = f \circ \overline{x}$。
\end{lemma}

\begin{proof}
选取交换图
$$
\xymatrix{
U \ar[d] \ar[r] & V \ar[d] \\
X \ar[r] & Y
}
$$
，其中 $U$ 和 $V$ 是概形，竖直箭头是平展的，且存在一点
$u \in U$ 映到 $x$。令 $v \in V$ 为 $u$ 的像。
把引理 \ref{spaces-more-morphisms-lemma-flat-on-fibres-at-point} 应用于
$\text{id} : X \to X$（底为 $Y$），可知 (1) 转化为条件
``$\mathcal{F}|_{U_v}$ 相对于 $U_v$ 在 $u$ 处平坦''。
换言之，(1) 等价于 $(\mathcal{F}|_{U_v})_u$ 是平坦的
$\mathcal{O}_{U_v, u}$-模。
由概形情形（《态射进阶》引理
\ref{more-morphisms-lemma-flat-and-free-at-point-fibre}），
这推出 $\mathcal{F}|_U$ 在 $u$ 的某个开邻域上有限自由。
由此可见，(1) 推出 (4)。
(4) $\Rightarrow$ (3) 和 (2) $\Rightarrow$ (1) 是立即的。
为证明 (3) $\Rightarrow$ (2)，使用《空间的性质》诸引理
\ref{spaces-properties-lemma-describe-etale-local-ring} 和
\ref{spaces-properties-lemma-stalk-quasi-coherent}
对局部环和茎的描述。
\end{proof}

\begin{lemma}
\label{spaces-more-morphisms-lemma-finite-free-open}
设 $S$ 为概形。设 $f : X \to Y$ 是 $S$ 上代数空间的态射，
并且局部有限表示。
设 $\mathcal{F}$ 是有限表示的 $\mathcal{O}_X$-模，并且在 $Y$ 上平坦。
则集合
$$
\{x \in |X| : \mathcal{F}\text{ free in an \'etale neighbourhood of }x\}
$$
在 $|X|$ 中开，且它的形成与任意基变换 $Y' \to Y$ 可交换。
\end{lemma}

\begin{proof}
开性是显然的。设 $Y' \to Y$ 是代数空间的态射，
令 $X' = Y' \times_Y X$，
并设 $x' \in |X'|$ 是位于 $x \in |X|$ 上方的一点。
由引理 \ref{spaces-more-morphisms-lemma-flat-and-free-at-point-fibre}，
$x$ 属于我们的集合，当且仅当
$(\mathcal{F}_{\overline{y}})_{\overline{x}}$ 是平坦的
$\mathcal{O}_{X_{\overline{y}}, \overline{x}}$-模。
同样，$x'$ 属于拉回 $\mathcal{F}'$ 所对应的类似集合；这里该模是
$\mathcal{F}$ 到 $X'$ 的拉回，当且仅当
$(\mathcal{F}'_{\overline{y}'})_{\overline{x}'}$ 是平坦的
$\mathcal{O}_{X'_{\overline{y}'}, \overline{x}'}$-模
（符号含义显然）。把引理 \ref{spaces-more-morphisms-lemma-flat-on-fibres-at-point} 应用于
态射 $\text{id} : X \to X$（底为 $Y$），可知这两个断言等价。
因此，关于基变换的陈述成立。
\end{proof}







\section{诺特基底上的平坦性}
\label{spaces-more-morphisms-section-flat-over-Noetherian}

\noindent
下面给出诺特情形的``平坦性的纤维判据''。

\begin{theorem}
\label{spaces-more-morphisms-theorem-criterion-flatness-fibre-Noetherian}
设 $S$ 为概形。
设 $f : X \to Y$ 和 $Y \to Z$ 是 $S$ 上代数空间的态射。
设 $\mathcal{F}$ 为拟凝聚 $\mathcal{O}_X$-模。
假设
\begin{enumerate}
\item $X$、$Y$、$Z$ 都局部诺特；并且 % P09-ERR-0604
\item $\mathcal{F}$ 是凝聚 $\mathcal{O}_X$-模。% P09-ERR-0604
\end{enumerate}
设 $x \in |X|$，并令 $y \in |Y|$ 和 $z \in |Z|$ 为 $x$ 的像。
若 $\mathcal{F}_{\overline{x}} \not = 0$，则下列条件等价：% P09-ERR-0605
\begin{enumerate}
\item $\mathcal{F}$ 相对于 $Z$ 在 $x$ 处平坦，且
$\mathcal{F}$ 到其位于 $z$ 上的纤维的限制在 $x$ 处相对于
$Y$ 位于 $z$ 上的纤维平坦；
\item $Y$ 相对于 $Z$ 在 $y$ 处平坦，且 $\mathcal{F}$ 相对于 $Y$
在 $x$ 处平坦。
\end{enumerate}
\end{theorem}

\begin{proof}
选取引理 \ref{spaces-more-morphisms-lemma-flat-on-fibres-at-point} 第 (3) 项中的图。
由定义可知，这把问题化归为概形态射 $U \to V \to W$、
拟凝聚层 $a^*\mathcal{F}$ 以及点 $u \in U$ 的相应定理。
因此，该定理由概形的相应结果，即《态射进阶》定理
\ref{more-morphisms-theorem-criterion-flatness-fibre-Noetherian} 得出。
\end{proof}

\begin{lemma}
\label{spaces-more-morphisms-lemma-morphism-between-flat-Noetherian}
设 $S$ 为概形。
设 $f : X \to Y$ 和 $Y \to Z$ 是 $S$ 上代数空间的态射。% P09-ERR-0606
假设
\begin{enumerate}
\item $X$、$Y$、$Z$ 都局部诺特；% P09-ERR-0604
\item $X$ 在 $Z$ 上平坦；
\item 对每个 $z \in |Z|$，$X$ 位于 $z$ 上的纤维相对于
$Y$ 位于 $z$ 上的纤维平坦。
\end{enumerate}
则 $f$ 平坦。若 $f$ 还是满的，则 $Y$ 在 $Z$ 上平坦。
\end{lemma}

\begin{proof}
这是定理 \ref{spaces-more-morphisms-theorem-criterion-flatness-fibre-Noetherian} 的特殊情形。
\end{proof}

\noindent
正如检验光滑性一样，若基底诺特，只需在阿廷环上检验平坦性。
下面给出一个示例性陈述。

\begin{lemma}
\label{spaces-more-morphisms-lemma-flatness-over-Noetherian-ring}
设 $A$ 为诺特环。设 $I \subset A$ 为理想。
设 $X$ 为在 $S = \Spec(A)$ 上局部有限表示的代数空间。
对 $n \geq 1$，令 $S_n = \Spec(A/I^n)$ 且
$X_n = S_n \times_S X$。设 $\mathcal{F}$ 为凝聚
$\mathcal{O}_X$-模。% P09-ERR-0607
若对每个 $n \geq 1$，$\mathcal{F}_n$ 是 $\mathcal{F}$ 到 $X$ 的拉回，% P09-ERR-0608
并且在 $S_n$ 上平坦，则 $\mathcal{F}$ 在 $X$ 上平坦的（开）轨迹 % P09-ERR-0609
包含 $V(I)$ 在 $X \to S$ 下的逆像。
\end{lemma}

\begin{proof}
$\mathcal{F}$ 在 $S$ 上平坦的轨迹由定理
\ref{spaces-more-morphisms-theorem-openness-flatness} 在 $|X|$ 中开。
把 $X$ 替换为一个平展覆盖的成员不影响该陈述，故可假设 $X$ 是仿射概形。
在这种情形下，结论立即由《代数》引理
\ref{algebra-lemma-flat-module-powers} 得出。
略去一些细节。
\end{proof}







\section{再论正规化}
\label{spaces-more-morphisms-section-normalization}

\noindent
正规化与光滑基变换可交换。

\begin{lemma}
\label{spaces-more-morphisms-lemma-integral-closure-smooth-pullback}
设 $S$ 为概形。设 $f : Y \to X$ 为 $S$ 上代数空间的光滑态射。
设 $\mathcal{A}$ 为拟凝聚 $\mathcal{O}_X$-代数层。则
$\mathcal{O}_Y$ 在 $f^*\mathcal{A}$ 中的整闭包等于 $f^*\mathcal{A}'$，
其中 $\mathcal{A}' \subset \mathcal{A}$ 是 $\mathcal{O}_X$ 在
$\mathcal{A}$ 中的整闭包。
\end{lemma}

\begin{proof}
由整闭包的构造（见《空间的态射》定义
\ref{spaces-morphisms-definition-integral-closure}），
这立即化归为 $X$ 和 $Y$ 均仿射的情形。
在这种情形下，结论就是《代数》引理
\ref{algebra-lemma-integral-closure-commutes-smooth}。
\end{proof}

\begin{lemma}[正规化与光滑基变换可交换]
\label{spaces-more-morphisms-lemma-normalization-smooth-localization}
设 $S$ 为概形。设
$$
\xymatrix{
Y_2 \ar[r] \ar[d] & Y_1 \ar[d]^f \\
X_2 \ar[r]^\varphi & X_1
}
$$
是 $S$ 上代数空间的笛卡尔方块。假设 $f$ 拟紧且拟分离，
并且 $\varphi$ 光滑。
设 $Y_i \to X_i' \to X_i$ 为 $X_i$ 在 $Y_i$ 中的正规化。
则 $X_2' \cong X_2 \times_{X_1} X_1'$。
\end{lemma}

\begin{proof}
把分解 $Y_1 \to X_1' \to X_1$ 基变换到 $X_2$，得到分解
$Y_2 \to X_2 \times_{X_1} X_1' \to X_1$，且
$X_2 \times_{X_1} X_1' \to X_1$ 是整态射 % P09-ERR-0610
（《空间的态射》引理 \ref{spaces-morphisms-lemma-base-change-integral}）。
因此，由《空间的态射》引理
\ref{spaces-morphisms-lemma-characterize-normalization} 的泛性质，
得到态射 $h : X_2' \to X_2 \times_{X_1} X_1'$。
注意，$X_2'$ 是 $\mathcal{O}_{X_2}$ 在
$f_{2, *}\mathcal{O}_{Y_2}$ 中的整闭包的相对谱。
若 $\mathcal{A}' \subset f_{1, *}\mathcal{O}_{Y_1}$ 表示
$\mathcal{O}_{X_2}$ 的整闭包，% P09-ERR-0611
则 $X_2 \times_{X_1} X_1'$ 是 $\varphi^*\mathcal{A}'$ 的相对谱，
因为相对谱的构造与任意基变换可交换。由《空间的上同调》引理
\ref{spaces-cohomology-lemma-flat-base-change-cohomology}，
我们知道 $f_{2, *}\mathcal{O}_{Y_2} = \varphi^*f_{1, *}\mathcal{O}_{Y_1}$。
因此，结论由引理 \ref{spaces-more-morphisms-lemma-integral-closure-smooth-pullback} 得出。
\end{proof}









\section{Cohen--Macaulay 态射}
\label{spaces-more-morphisms-section-CM}

\noindent
本节对应于《态射进阶》第 \ref{more-morphisms-section-CM} 节。

\begin{lemma}
\label{spaces-more-morphisms-lemma-CM-local-ring-fibre}
概形芽的态射的性质
\begin{align*}
& \mathcal{P}((X, x) \to (S, s)) = \\
& \text{纤维的局部环 }
\mathcal{O}_{X_s, x}
\text{ 是诺特且 Cohen--Macaulay 的}
\end{align*}
在源和目标上都是平展局部的（《下降》定义
\ref{descent-definition-local-source-target-at-point}）。
\end{lemma}

\begin{proof}
给定《下降》定义
\ref{descent-definition-local-source-target-at-point} 中的图，
我们得到纤维的平展态射 $U'_{v'} \to U_v$，它把 $u'$ 映到 $u$；
见《下降》引理 \ref{descent-lemma-etale-on-fiber}。
因此，局部环 $\mathcal{O}_{U'_{v'}, u'}$ 和
$\mathcal{O}_{U_v, u}$ 的严格亨泽尔化相同。结论由《代数进阶》引理
\ref{more-algebra-lemma-henselization-CM} 得出。
\end{proof}

\begin{definition}
\label{spaces-more-morphisms-definition-CM}
设 $S$ 为概形。
设 $f : X \to Y$ 为 $S$ 上代数空间的态射。
假设 $f$ 的纤维局部诺特（《空间上的除子》定义
\ref{spaces-divisors-definition-locally-Noetherian-fibre}）。
\begin{enumerate}
\item 设 $x \in |X|$ 且 $y = f(x)$。我们称 $f$
{\it 在 $x$ 处为 Cohen--Macaulay}，如果 $f$ 在 $x$ 处平坦，并且
《空间的态射》引理
\ref{spaces-morphisms-lemma-local-source-target-at-point} 中的等价条件
对于引理 \ref{spaces-more-morphisms-lemma-CM-local-ring-fibre} 所述性质 $\mathcal{P}$ 成立。
\item 我们称 $f$ 为 {\it Cohen--Macaulay 态射}，如果 $f$ 在 $X$
的每一点处都是 Cohen--Macaulay 的。
\end{enumerate}
\end{definition}

\noindent
下面给出相应的转述。

\begin{lemma}
\label{spaces-more-morphisms-lemma-CM}
设 $S$ 为概形。设 $f : X \to Y$ 为 $S$ 上代数空间的态射。
假设 $f$ 的纤维局部诺特。下列条件等价：% P09-ERR-0612
\begin{enumerate}
\item $f$ 是 Cohen--Macaulay 的；
\item $f$ 平坦，并且存在满平展态射 $V \to Y$，其中 $V$ 是概形，
使得 $X_V \to V$ 的纤维都是 Cohen--Macaulay 代数空间；并且
\item $f$ 平坦，并且对任意平展态射 $V \to Y$，其中 $V$ 是概形，
$X_V \to V$ 的纤维都是 Cohen--Macaulay 代数空间。
\end{enumerate}
给定 $x \in |X|$，其像为 $y \in |Y|$，下列条件等价：% P09-ERR-0612
\begin{enumerate}
\item[(a)] $f$ 在 $x$ 处为 Cohen--Macaulay；
\item[(b)] $\mathcal{O}_{Y, \overline{y}} \to \mathcal{O}_{X, \overline{x}}$
平坦，且
$\mathcal{O}_{X, \overline{x}}/
\mathfrak m_{\overline{y}}\mathcal{O}_{X, \overline{x}}$
是 Cohen--Macaulay 的。% P09-ERR-0613
\end{enumerate}
\end{lemma}

\begin{proof}
给定平展态射 $V \to Y$，其中 $V$ 是概形；选取概形 $U$ 以及
满平展态射 $U \to X \times_Y V$。考虑交换图
$$
\xymatrix{
U \ar[d] \ar[r] & V \ar[d] \\
X \ar[r] & Y
}
$$
设 $u \in U$，其像为 $x \in |X|$、$y \in |Y|$ 和 $v \in V$。
则 $f$ 在 $x$ 处为 Cohen--Macaulay，当且仅当 $U \to V$ 在 $u$ 处为
Cohen--Macaulay（由定义）。此外，态射
$U_v \to X_v = (X_V)_v$ 是满平展的。因此，概形 $U_v$ 是
Cohen--Macaulay 的，当且仅当代数空间 $X_v$ 是 Cohen--Macaulay 的。
于是，(1)、(2)、(3) 的等价性由概形态射的相应等价性通过形式论证得出；
见《态射进阶》引理 \ref{more-morphisms-lemma-CM}。

\medskip\noindent
证明 (a) 与 (b) 等价。平坦性的相应等价性即《空间的态射》引理
\ref{spaces-morphisms-lemma-flat-at-point-etale-local-rings}。
因此，在证明该等价性时，可以假设 $f$ 在 $x$ 处平坦。
考虑上述图以及 $x, y, u, v$。则
$\mathcal{O}_{Y, \overline{y}} \to \mathcal{O}_{X, \overline{x}}$
等于局部环严格亨泽尔化之间的映射
$\mathcal{O}_{V, v}^{sh} \to \mathcal{O}_{U, u}^{sh}$；
见《空间的性质》引理
\ref{spaces-properties-lemma-describe-etale-local-ring}。
因此有
$$
\mathcal{O}_{X, \overline{x}}/
\mathfrak m_{\overline{y}}\mathcal{O}_{X, \overline{x}} =
(\mathcal{O}_{U, u}/\mathfrak m_v \mathcal{O}_{U, u})^{sh}
$$
；这是《代数》引理
\ref{algebra-lemma-quotient-strict-henselization}。
故只需证明诺特局部环
$\mathcal{O}_{U, u}/\mathfrak m_v \mathcal{O}_{U, u}$
为 Cohen--Macaulay，当且仅当其严格亨泽尔化为 Cohen--Macaulay。
这就是《代数进阶》引理 \ref{more-algebra-lemma-henselization-CM}。
\end{proof}

\begin{lemma}
\label{spaces-more-morphisms-lemma-composition-CM}
设 $S$ 为概形。
设 $f : X \to Y$ 和 $g : Y \to Z$ 是 $S$ 上代数空间的态射。
假设 $f$、$g$ 和 $g \circ f$ 的纤维局部诺特。
设 $x \in |X|$，其像为 $y \in |Y|$ 和 $z \in |Z|$。
\begin{enumerate}
\item 若 $f$ 在 $x$ 处为 Cohen--Macaulay，且 $g$ 在 $f(x)$ 处为
Cohen--Macaulay，则 $g \circ f$ 在 $x$ 处为 Cohen--Macaulay。
\item 若 $f$ 和 $g$ 都是 Cohen--Macaulay 的，则 $g \circ f$ 也是
Cohen--Macaulay 的。
\item 若 $g \circ f$ 在 $x$ 处为 Cohen--Macaulay，且 $f$ 在 $x$ 处平坦，
则 $f$ 在 $x$ 处为 Cohen--Macaulay，且 $g$ 在 $f(x)$ 处为
Cohen--Macaulay。
\item 若 $f \circ g$ 是 Cohen--Macaulay 的，且 $f$ 平坦，% P09-ERR-0614
则 $f$ 是 Cohen--Macaulay 的，且 $g$ 在 $f$ 的像中的每一点处都是
Cohen--Macaulay 的。
\end{enumerate}
\end{lemma}

\begin{proof}
在平展局部上工作，结论由概形的相应结果得出；见《态射进阶》引理
\ref{more-morphisms-lemma-composition-CM}。
或者，可以使用引理 \ref{spaces-more-morphisms-lemma-CM} 中 (a) 与 (b) 的等价性。
因此，考虑诺特局部环的局部同态
$$
\mathcal{O}_{Y, \overline{y}}/
\mathfrak m_{\overline{z}}\mathcal{O}_{Y, \overline{y}}
\longrightarrow
\mathcal{O}_{X, \overline{x}}/
\mathfrak m_{\overline{z}}\mathcal{O}_{X, \overline{x}}
$$
，其纤维为
$$
\mathcal{O}_{X, \overline{x}}/
\mathfrak m_{\overline{y}}\mathcal{O}_{X, \overline{x}}
$$
；然后应用《代数》引理 \ref{algebra-lemma-CM-goes-up}。
\end{proof}

\begin{lemma}
\label{spaces-more-morphisms-lemma-flat-morphism-from-CM}
设 $S$ 为概形。
设 $f : X \to Y$ 为 $S$ 上局部诺特代数空间之间的平坦态射。
若 $X$ 是 Cohen--Macaulay 的，则 $f$ 是 Cohen--Macaulay 的，并且
$\mathcal{O}_{Y, f(\overline{x})}$ 对所有 $x \in |X|$ 都是
Cohen--Macaulay 的。% P09-ERR-0615
\end{lemma}

\begin{proof}
利用引理 \ref{spaces-more-morphisms-lemma-CM} 转译为代数（与引理
\ref{spaces-more-morphisms-lemma-composition-CM} 的证明比较）后，结论由《代数》引理
\ref{algebra-lemma-CM-goes-up} 得出。
\end{proof}

\begin{lemma}
\label{spaces-more-morphisms-lemma-base-change-CM}
设 $S$ 为概形。
设 $f : X \to Y$ 为 $S$ 上代数空间的态射。
假设 $f$ 的纤维局部诺特。
设 $Y' \to Y$ 局部有限型。设 $f' : X' \to Y'$ 为 $f$ 的基变换。
设 $x' \in |X'|$ 为一点，其像为 $x \in |X|$。
\begin{enumerate}
\item 若 $f$ 在 $x$ 处为 Cohen--Macaulay，则
$f' : X' \to Y'$ 在 $x'$ 处为 Cohen--Macaulay。
\item 若 $f$ 在 $x$ 处平坦，且 $f'$ 在 $x'$ 处为 Cohen--Macaulay，
则 $f$ 在 $x$ 处为 Cohen--Macaulay。
\item 若 $Y' \to Y$ 在 $f'(x')$ 处平坦，且 $f'$ 在 $x'$ 处为
Cohen--Macaulay，则 $f$ 在 $x$ 处为 Cohen--Macaulay。
\end{enumerate}
\end{lemma}

\begin{proof}
用 $y \in |Y|$ 和 $y' \in |Y'|$ 表示 $x'$ 的像。% P09-ERR-0616
选取满平展态射 $V \to Y$，其中 $V$ 是概形。
选取满平展态射 $U \to X \times_Y V$，其中 $U$ 是概形。
选取满平展态射 $V' \to Y' \times_Y V$，其中 $V'$ 是概形。% P09-ERR-0617
则 $U' = U \times_V V'$ 是概形，并带有满平展态射 $U' \to X'$。
选取 $u' \in U'$ 映到 $x'$。用 $u \in U$ 表示 $u'$ 的像。
于是，引理由 $U \to V$、其基变换 $U' \to V'$ 以及点 $u'$、$u$
的引理得出（这由定义可知）。
故引理由概形情形得出；见《态射进阶》引理
\ref{more-morphisms-lemma-base-change-CM}。
\end{proof}

\begin{lemma}
\label{spaces-more-morphisms-lemma-flat-finite-presentation-CM-open}
设 $S$ 为概形。
设 $f : X \to Y$ 为 $S$ 上代数空间的态射，并且平坦且局部有限表示。
令
$$
W = \{x \in |X| : f\text{ 为 Cohen--Macaulay 于点 }x\}
$$
则 $W$ 在 $|X|$ 中开，且 $W$ 的形成与 $f$ 的任意基变换可交换：
对任意态射 $g : Y' \to Y$，考虑基变换 $f' : X' \to Y'$（它是
$f$ 的基变换）以及投影 $g' : X' \to X$。
则相应集合 $W'$（针对态射 $f'$）等于 $W' = (g')^{-1}(W)$。% P09-ERR-0618
\end{lemma}

\begin{proof}
选取交换图
$$
\xymatrix{
U \ar[d] \ar[r] & V \ar[d] \\
X \ar[r] & Y
}
$$
，其中竖直箭头平展，而 $U$ 和 $V$ 是概形。
设 $u \in U$，其像为 $x \in |X|$。
则 $f$ 在 $x$ 处为 Cohen--Macaulay，当且仅当 $U \to V$ 在 $u$ 处为
Cohen--Macaulay（由定义）。
因而化归为态射 $U \to V$ 的情形。见《态射进阶》引理
\ref{more-morphisms-lemma-flat-finite-presentation-CM-open}。
\end{proof}

\begin{lemma}
\label{spaces-more-morphisms-lemma-lfp-CM-relative-dimension}
设 $S$ 为概形。设 $f : X \to Y$ 为 $S$ 上代数空间的态射。
假设 $f$ 局部有限表示且为 Cohen--Macaulay。
则存在开闭子概形 $X_d \subset X$，% P09-ERR-0619
使得 $X = \coprod_{d \geq 0} X_d$，且 $f|_{X_d} : X_d \to Y$
具有相对维数 $d$。
\end{lemma}

\begin{proof}
选取交换图
$$
\xymatrix{
U \ar[d] \ar[r] & V \ar[d] \\
X \ar[r] & Y
}
$$
，其中竖直箭头平展，而 $U$ 和 $V$ 是概形。
则 $U \to V$ 局部有限表示且为 Cohen--Macaulay（由定义立即得出）。
因此有分解 $U = \coprod_{d \geq 0} U_d$，其中每个分量都是开闭子概形，
且 $f|_{U_d} : U_d \to V$ 具有相对维数 $d$；% P09-ERR-0620
见《态射》引理
\ref{morphisms-lemma-flat-finite-presentation-CM-fibres-relative-dimension}。
设 $u \in U$，其像为 $x \in |X|$。则 $f$ 具有相对维数 $d$
在 $x$ 处成立，当且仅当 $U \to V$ 具有相对维数 $d$ 在 $u$ 处成立
（这由定义得出）。
由此可见，$U_d$ 是子集 $X_d \subset |X|$ 的逆像，
而该子集必为开闭的。用 $X_d$ 表示 $X$ 的相应开闭代数子空间，
即可看出引理成立。
\end{proof}







\section{Gorenstein 态射}
\label{spaces-more-morphisms-section-gorenstein}

\noindent
本节对应于《概形的对偶性》第
\ref{duality-section-gorenstein-morphisms} 节。

\begin{lemma}
\label{spaces-more-morphisms-lemma-gorenstein-local-ring-fibre}
概形芽的态射的下述性质
\begin{align*}
& \mathcal{P}((X, x) \to (S, s)) = \\
& \text{纤维的局部环 }
\mathcal{O}_{X_s, x}
\text{ 是诺特且 Gorenstein 的}
\end{align*}
在源和靶上为平展局部的（《下降》定义
\ref{descent-definition-local-source-target-at-point}）。
\end{lemma}

\begin{proof}
给定《下降》定义
\ref{descent-definition-local-source-target-at-point} 中的图，
我们得到纤维的平展态射 $U'_{v'} \to U_v$，它把 $u'$ 映到 $u$；
见《下降》引理 \ref{descent-lemma-etale-on-fiber}。
因此，$\mathcal{O}_{U_v, u} \to \mathcal{O}_{U'_{v'}, u'}$
是一个平展环映射的局部化。故前者是诺特的，当且仅当后者是诺特的；
见《代数进阶》引理 \ref{more-algebra-lemma-Noetherian-etale-extension}。
又因为
$\mathcal{O}_{U'_{v'}, u'}/\mathfrak m_u \mathcal{O}_{U'_{v'}, u'}
= \kappa(u')$（《代数》引理 \ref{algebra-lemma-etale-at-prime}）
是 Gorenstein 环，所以由《对偶化复形》引理
\ref{dualizing-lemma-flat-under-gorenstein} 可见，
$\mathcal{O}_{U_v, u}$ 是 Gorenstein 的，当且仅当
$\mathcal{O}_{U'_{v'}, u'}$ 是 Gorenstein 的。
\end{proof}

\begin{definition}
\label{spaces-more-morphisms-definition-gorenstein}
设 $S$ 为概形。
设 $f : X \to Y$ 为 $S$ 上代数空间的态射。
假设 $f$ 的纤维局部诺特（《空间上的除子》定义
\ref{spaces-divisors-definition-locally-Noetherian-fibre}）。
\begin{enumerate}
\item 设 $x \in |X|$ 且 $y = f(x)$。我们称 $f$
{\it 在 $x$ 处为 Gorenstein}，如果 $f$ 在 $x$ 处平坦，并且
《空间的态射》引理
\ref{spaces-morphisms-lemma-local-source-target-at-point} 中的等价条件
对于引理 \ref{spaces-more-morphisms-lemma-gorenstein-local-ring-fibre} 所述性质
$\mathcal{P}$ 成立。
\item 我们称 $f$ 为 {\it Gorenstein 态射}，如果 $f$ 在 $X$
的每一点处都是 Gorenstein 的。
\end{enumerate}
\end{definition}

\noindent
下面给出相应的转述。

\begin{lemma}
\label{spaces-more-morphisms-lemma-gorenstein}
设 $S$ 为概形。设 $f : X \to Y$ 为 $S$ 上代数空间的态射。
假设 $f$ 的纤维局部诺特。下列条件等价：% P09-ERR-0621
\begin{enumerate}
\item $f$ 是 Gorenstein 的；
\item $f$ 平坦，并且存在满平展态射 $V \to Y$，其中 $V$ 是概形，
使得 $X_V \to V$ 的纤维都是 Gorenstein 代数空间；并且
\item $f$ 平坦，并且对任意平展态射 $V \to Y$，其中 $V$ 是概形，
$X_V \to V$ 的纤维都是 Gorenstein 代数空间。
\end{enumerate}
给定 $x \in |X|$，其像为 $y \in |Y|$，下列条件等价：% P09-ERR-0621
\begin{enumerate}
\item[(a)] $f$ 在 $x$ 处为 Gorenstein；
\item[(b)] $\mathcal{O}_{Y, \overline{y}} \to \mathcal{O}_{X, \overline{x}}$
平坦，并且
$\mathcal{O}_{X, \overline{x}}/
\mathfrak m_{\overline{y}}\mathcal{O}_{X, \overline{x}}$
是 Gorenstein 的。% P09-ERR-0622
\end{enumerate}
\end{lemma}

\begin{proof}
给定平展态射 $V \to Y$，其中 $V$ 是概形；选取概形 $U$ 以及
满平展态射 $U \to X \times_Y V$。考虑交换图
$$
\xymatrix{
U \ar[d] \ar[r] & V \ar[d] \\
X \ar[r] & Y
}
$$
设 $u \in U$，其像为 $x \in |X|$、$y \in |Y|$ 和 $v \in V$。
则 $f$ 在 $x$ 处为 Gorenstein，当且仅当 $U \to V$ 在 $u$ 处为
Gorenstein（由定义）。此外，态射
$U_v \to X_v = (X_V)_v$ 是满平展的。因此，概形 $U_v$ 是
Gorenstein 的，当且仅当代数空间 $X_v$ 是 Gorenstein 的。
于是，(1)、(2) 与 (3) 的等价性由概形态射的相应等价性形式地得出；
见《概形的对偶性》引理 \ref{duality-lemma-gorenstein}。

\medskip\noindent
证明 (a) 与 (b) 等价。关于平坦性的相应等价性是《空间的态射》引理
\ref{spaces-morphisms-lemma-flat-at-point-etale-local-rings}。
因此，在证明该等价性时可以假设 $f$ 在 $x$ 处平坦。
考虑上述图以及 $x, y, u, v$。于是
$\mathcal{O}_{Y, \overline{y}} \to \mathcal{O}_{X, \overline{x}}$
等于局部环的严格亨泽尔化之间的映射
$\mathcal{O}_{V, v}^{sh} \to \mathcal{O}_{U, u}^{sh}$；
见《空间的性质》引理
\ref{spaces-properties-lemma-describe-etale-local-ring}。
故由《代数》引理 \ref{algebra-lemma-quotient-strict-henselization}，有
$$
\mathcal{O}_{X, \overline{x}}/
\mathfrak m_{\overline{y}}\mathcal{O}_{X, \overline{x}} =
(\mathcal{O}_{U, u}/\mathfrak m_v \mathcal{O}_{U, u})^{sh}
$$
因此，只须证明诺特局部环
$\mathcal{O}_{U, u}/\mathfrak m_v \mathcal{O}_{U, u}$
为 Gorenstein，当且仅当其严格亨泽尔化为 Gorenstein。
这立即由《对偶化复形》引理
\ref{dualizing-lemma-completion-henselization-dualizing}
以及 Gorenstein 局部环的定义得出：它是一个以自身为其对偶化复形的诺特局部环。
\end{proof}

\begin{lemma}
\label{spaces-more-morphisms-lemma-composition-gorenstein}
设 $S$ 为概形。
设 $f : X \to Y$ 和 $g : Y \to Z$ 是 $S$ 上代数空间的态射。
假设 $f$、$g$ 和 $g \circ f$ 的纤维局部诺特。
设 $x \in |X|$，其像为 $y \in |Y|$ 和 $z \in |Z|$。
\begin{enumerate}
\item 若 $f$ 在 $x$ 处为 Gorenstein，且 $g$ 在 $f(x)$ 处为
Gorenstein，则 $g \circ f$ 在 $x$ 处为 Gorenstein。
\item 若 $f$ 和 $g$ 都是 Gorenstein 的，则 $g \circ f$ 也是
Gorenstein 的。
\item 若 $g \circ f$ 在 $x$ 处为 Gorenstein，且 $f$ 在 $x$ 处平坦，
则 $f$ 在 $x$ 处为 Gorenstein，且 $g$ 在 $f(x)$ 处为 Gorenstein。
\item 若 $f \circ g$ 是 Gorenstein 的，且 $f$ 平坦，% P09-ERR-0623
则 $f$ 是 Gorenstein 的，且 $g$ 在 $f$ 的像中的每一点处都是
Gorenstein 的。
\end{enumerate}
\end{lemma}

\begin{proof}
在平展局部上工作，结论由概形的相应结果得出；见《概形的对偶性》引理
\ref{duality-lemma-composition-gorenstein}。
或者，可以使用引理 \ref{spaces-more-morphisms-lemma-gorenstein} 中 (a) 与 (b) 的等价性。
因此考虑诺特局部环的局部同态
$$
\mathcal{O}_{Y, \overline{y}}/
\mathfrak m_{\overline{z}}\mathcal{O}_{Y, \overline{y}}
\longrightarrow
\mathcal{O}_{X, \overline{x}}/
\mathfrak m_{\overline{z}}\mathcal{O}_{X, \overline{x}}
$$
其纤维为
$$
\mathcal{O}_{X, \overline{x}}/
\mathfrak m_{\overline{y}}\mathcal{O}_{X, \overline{x}}
$$
；然后应用《对偶化复形》引理
\ref{dualizing-lemma-flat-under-gorenstein}。
\end{proof}

\begin{lemma}
\label{spaces-more-morphisms-lemma-flat-morphism-from-gorenstein}
设 $S$ 为概形。
设 $f : X \to Y$ 为 $S$ 上局部诺特代数空间之间的平坦态射。
若 $X$ 是 Gorenstein 的，则 $f$ 是 Gorenstein 的，并且
$\mathcal{O}_{Y, f(\overline{x})}$ 对所有 $x \in |X|$ 都是
Gorenstein 的。% P09-ERR-0624
\end{lemma}

\begin{proof}
利用引理 \ref{spaces-more-morphisms-lemma-gorenstein} 转译为代数（与引理
\ref{spaces-more-morphisms-lemma-composition-gorenstein} 的证明比较）后，结论由
《对偶化复形》引理 \ref{dualizing-lemma-flat-under-gorenstein} 得出。
\end{proof}

\begin{lemma}
\label{spaces-more-morphisms-lemma-base-change-gorenstein}
设 $S$ 为概形。
设 $f : X \to Y$ 为 $S$ 上代数空间的态射。
假设 $f$ 的纤维局部诺特。
设 $Y' \to Y$ 局部有限型。设 $f' : X' \to Y'$ 为 $f$ 的基变换。
设 $x' \in |X'|$ 为一点，其像为 $x \in |X|$。
\begin{enumerate}
\item 若 $f$ 在 $x$ 处为 Gorenstein，则
$f' : X' \to Y'$ 在 $x'$ 处为 Gorenstein。
\item 若 $f$ 在 $x$ 处平坦，且 $f'$ 在 $x'$ 处为 Gorenstein，
则 $f$ 在 $x$ 处为 Gorenstein。
\item 若 $Y' \to Y$ 在 $f'(x')$ 处平坦，且 $f'$ 在 $x'$ 处为
Gorenstein，则 $f$ 在 $x$ 处为 Gorenstein。
\end{enumerate}
\end{lemma}

\begin{proof}
用 $y \in |Y|$ 和 $y' \in |Y'|$ 表示 $x'$ 的像。% P09-ERR-0625
选取满平展态射 $V \to Y$，其中 $V$ 是概形。
选取满平展态射 $U \to X \times_Y V$，其中 $U$ 是概形。
选取满平展态射 $V' \to Y' \times_Y V$，其中 $V'$ 是概形。% P09-ERR-0626
则 $U' = U \times_V V'$ 是概形，并带有满平展态射 $U' \to X'$。
选取 $u' \in U'$ 映到 $x'$。用 $u \in U$ 表示 $u'$ 的像。
于是，引理由 $U \to V$、其基变换 $U' \to V'$ 以及点 $u'$、$u$
的引理得出（这由定义可知）。
故引理由概形情形得出；见《概形的对偶性》引理
\ref{duality-lemma-base-change-gorenstein}。
\end{proof}

\begin{lemma}
\label{spaces-more-morphisms-lemma-flat-finite-presentation-gorenstein-open}
设 $S$ 为概形。
设 $f : X \to Y$ 为 $S$ 上代数空间的态射，并且平坦且局部有限表示。
令
$$
W = \{x \in |X| : f\text{ 为 Gorenstein 于点 }x\}
$$
则 $W$ 在 $|X|$ 中开，且 $W$ 的形成与 $f$ 的任意基变换可交换：
对任意态射 $g : Y' \to Y$，考虑基变换 $f' : X' \to Y'$（它是
$f$ 的基变换）以及投影 $g' : X' \to X$。
则相应集合 $W'$（针对态射 $f'$）等于 $W' = (g')^{-1}(W)$。% P09-ERR-0627
\end{lemma}

\begin{proof}
选取交换图% P09-ERR-0628
$$
\xymatrix{
U \ar[d] \ar[r] & V \ar[d] \\
X \ar[r] & Y
}
$$
设 $u \in U$，其像为 $x \in |X|$。
则 $f$ 在 $x$ 处为 Gorenstein，当且仅当 $U \to V$ 在 $u$ 处为
Gorenstein（由定义）。
因而化归为概形态射 $U \to V$ 的情形。开性由
《概形的对偶性》引理
\ref{duality-lemma-flat-finite-presentation-characterize-gorenstein} 证明，
与基变换的相容性由《概形的对偶性》引理
\ref{duality-lemma-flat-lft-base-change-gorenstein} 证明。
\end{proof}






\section{Cohen--Macaulay 态射的切片}
\label{spaces-more-morphisms-section-slice}

\noindent
设 $S$ 为概形。设 $X$ 为 $S$ 上的代数空间。
设 $f_1, \ldots, f_r \in \Gamma(X, \mathcal{O}_X)$。在此情形下，
用 $V(f_1, \ldots, f_r)$ 表示 {\it $X$ 中由 $f_1, \ldots, f_r$
切出的闭子空间}。更确切地，可以把 $V(f_1, \ldots, f_r)$
定义为 $X$ 的闭子空间，它对应于由 $f_1, \ldots, f_r$ 生成的
拟凝聚理想层；见《空间的态射》引理
\ref{spaces-morphisms-lemma-closed-immersion-ideals}。
也可以选取表示 $X = U/R$，并考虑闭子概形 $Z \subset U$，它由
$f_1|U, \ldots, f_r|_U$ 切出。显然，$Z$ 是 $R$-不变闭子概形
（见《广群》定义 \ref{groupoids-definition-invariant-open}），并可令
$V(f_1, \ldots, f_r) = Z/R_Z$。

\begin{lemma}
\label{spaces-more-morphisms-lemma-slice}
设 $S$ 为概形。考虑笛卡尔图
$$
\xymatrix{
X \ar[d] & F \ar[l]^p \ar[d] \\
Y & \Spec(k) \ar[l]
}
$$
其中 $X \to Y$ 是 $S$ 上代数空间的平坦且局部有限表示的态射，
而 $k$ 是 $S$ 上的域。设
$f_1, \ldots, f_r \in \Gamma(X, \mathcal{O}_X)$，且 $z \in |F|$，
使得 $f_1, \ldots, f_r$ 在局部环
$\mathcal{O}_{F, \overline{z}}$ 中的像构成正则序列。% P09-ERR-0629
则以 $X$ 的一个包含 $p(z)$ 的开子空间替换之以后，态射
$$
V(f_1, \ldots, f_r) \longrightarrow Y
$$
平坦且局部有限表示。
\end{lemma}

\begin{proof}
令 $Z = V(f_1, \ldots, f_r)$。显然，$Z \to X$ 局部有限表示，
故复合 $Z \to Y$ 局部有限表示；见《空间的态射》引理
\ref{spaces-morphisms-lemma-composition-finite-presentation}。
因此，只须证明 $Z \to Y$ 在 $p(z)$ 的一个邻域中平坦。
设 $k'/k$ 为域扩张。则
$F' = F \times_{\Spec(k)} \Spec(k')$ 在 $F$ 上满且平坦，% P09-ERR-0630
故可找到点 $z' \in |F'|$，它映到 $z$，且局部环映射
$\mathcal{O}_{F, \overline{z}} \to \mathcal{O}_{F', \overline{z}'}$
平坦；见《空间的态射》引理
\ref{spaces-morphisms-lemma-flat-at-point-etale-local-rings}。
因此，$f_1, \ldots, f_r$ 在
$\mathcal{O}_{F', \overline{z}'}$ 中的像也构成正则序列；
见《代数》引理 \ref{algebra-lemma-flat-increases-depth}。
所以在证明过程中可以用扩域替换 $k$。特别地，可以假设
$z \in |F|$ 来自截面 $z : \Spec(k) \to F$，它是结构态射
$F \to \Spec(k)$ 的截面。

\medskip\noindent
选取概形 $V$ 以及满平展态射 $V \to Y$。选取概形 $U$ 以及
满平展态射 $U \to X \times_Y V$。必要时再次扩大 $k$，可以假设
$\Spec(k) \to F \to X$ 通过 $U$ 分解（因为 $U \to X$ 是满的）。
设 $u : \Spec(k) \to U$ 为这样的分解，并用 $v \in V$ 表示 $u$ 的像。
注意，态射
$$
U_v \times_{\Spec(\kappa(v))} \Spec(k) =
U \times_V \Spec(k) \to U \times_Y \Spec(k) \to F
$$
都是平展的（第一个是 $V \to V \times_Y V$ 的基变换，
第二个是 $U \to X$ 的基变换）。此外，依构造，点
$u : \Spec(k) \to U$ 给出最左空间中的一点，它在右端映到 $z$。% P09-ERR-0631
因此，$f_1, \ldots, f_r$ 在下述映射右端的局部环中的像构成正则序列：
$$
\mathcal{O}_{U_v, u}
\longrightarrow
\mathcal{O}_{U_v \times_{\Spec(\kappa(v)} \Spec(k), \overline{u}}
=
\mathcal{O}_{U \times_V \Spec(k), \overline{u}}.
$$
% P09-ERR-0632
但是，该显示箭头平坦（合用《平坦性进阶》引理
\ref{flat-lemma-flat-up-down-henselization} 与《空间的态射》引理
\ref{spaces-morphisms-lemma-flat-at-point-etale-local-rings}），
故由《代数》引理 \ref{algebra-lemma-flat-increases-depth} 可见，
$f_1, \ldots, f_r$ 在 $\mathcal{O}_{U_v, u}$ 中映成正则序列。% P09-ERR-0633
由《态射进阶》引理 \ref{more-morphisms-lemma-slice-given-elements}，
概形态射
$$
V(f_1, \ldots, f_r) \times_X U = V(f_1|_U, \ldots, f_r|_U) \to V
$$
在开邻域 $U'$ 中平坦，该邻域含 $u$。设 $X' \subset X$ 为与
$|U'| \to |X|$ 的像对应的开子空间（见《空间的性质》引理
\ref{spaces-properties-lemma-topology-points} 和
\ref{spaces-properties-lemma-open-subspaces}）。
由此，$V(f_1, \ldots, f_r) \cap X' \to Y$ 平坦
（见《空间的态射》定义 \ref{spaces-morphisms-definition-flat}），
因为有交换图
$$
\xymatrix{
V(f_1, \ldots, f_r) \times_X U' \ar[d]_a \ar[r] & V \ar[d]^b \\
V(f_1, \ldots, f_r) \cap X' \ar[r] & Y
}
$$
，其中 $a, b$ 平展，且 $a$ 满。
\end{proof}






\section{约化纤维}
\label{spaces-more-morphisms-section-reduced}

\noindent
本节对应于《态射进阶》第
\ref{more-morphisms-section-reduced} 节。

\begin{lemma}
\label{spaces-more-morphisms-lemma-geometrically-reduced-fibre}
设 $S$ 为概形。设 $f : X \to Y$ 为 $S$ 上代数空间的态射。
设 $y \in |Y|$。下列条件等价：% P09-ERR-0634
\begin{enumerate}
\item 对于某个态射 $\Spec(k) \to Y$，它位于 $y$ 的等价类中，
代数空间 $X_k$ 在 $k$ 上几何约化；
\item 对于每个态射 $\Spec(k) \to Y$，它位于 $y$ 的等价类中，
代数空间 $X_k$ 在 $k$ 上几何约化；
\item 对于每个态射 $\Spec(k) \to Y$，它位于 $y$ 的等价类中，
代数空间 $X_k$ 约化。
\end{enumerate}
\end{lemma}

\begin{proof}
这立即由《域上的空间》引理
\ref{spaces-over-fields-lemma-geometrically-reduced-upstairs}
以及《空间的性质》第 \ref{spaces-properties-section-points} 节中给出的、
定义 $|X|$ 的等价关系得出。% P09-ERR-0635
\end{proof}

\begin{definition}
\label{spaces-more-morphisms-definition-geometrically-reduced-fibre}
设 $S$ 为概形。设 $f : X \to Y$ 为 $S$ 上代数空间的态射。
设 $y \in |Y|$。若引理
\ref{spaces-more-morphisms-lemma-geometrically-reduced-fibre} 的等价条件成立，
则称 {\it $f : X \to Y$ 在 $y$ 处的纤维几何约化}。
\end{definition}

\noindent
下面给出必需的引理。

\begin{lemma}
\label{spaces-more-morphisms-lemma-base-change-fibres-geometrically-reduced}
设 $S$ 为概形。设 $f : X \to Y$ 和 $g : Y' \to Y$ 为 $S$ 上
代数空间的态射。用 $f' : X' \to Y'$ 表示 $f$ 沿 $g$ 的基变换。则
\begin{align*}
\{y' \in |Y'| :
\text{态射 }f' : X' \to Y'\text{ 在点 }y'
\text{ 的纤维几何约化}\} \\
= g^{-1}(\{y \in |Y| :
\text{态射 }f : X \to Y\text{ 在点 }y
\text{ 的纤维几何约化}\}).
\end{align*}
\end{lemma}

\begin{proof}
对 $y' \in |Y'|$，选取态射 $\Spec(k) \to Y'$，它位于 $y'$
的等价类中。则 $g(y')$ 由复合
$\Spec(k) \to Y' \to Y$ 表示。于是
$X' \times_{Y'} \Spec(k) = X \times_Y \Spec(k)$，
结论由定义得出。
\end{proof}

\begin{lemma}
\label{spaces-more-morphisms-lemma-geometrically-reduced-constructible}
设 $S$ 为概形。设 $f : X \to Y$ 为 $S$ 上代数空间的拟紧且
局部有限表示的态射。则集合
$$
E = \{y \in |Y| : \text{态射 }f : X \to Y\text{ 在点 }y
\text{ 的纤维几何约化}\}
$$
平展局部可构造。
\end{lemma}

\begin{proof}
选取仿射概形 $V$ 以及平展态射 $V  \to Y$。
陈述的含义是 $E$ 在 $|V|$ 中的逆像可构造。由引理
\ref{spaces-more-morphisms-lemma-base-change-fibres-geometrically-reduced}，
可以替换 $Y$，以 $V$ 代之，即可以假设 $Y$ 为仿射概形。
于是 $X$ 拟紧。选取仿射概形 $U$ 以及满平展态射 $U \to X$。
对于态射 $\Spec(k) \to Y$，纤维之间的态射
$U_k \to X_k$ 满且平展。因此，$U_k$ 在 $k$ 上几何约化，
当且仅当 $X_k$ 在 $k$ 上几何约化；见《域上的空间》引理
\ref{spaces-over-fields-lemma-geometrically-reduced-etale-local}。
所以，集合 $E$（针对 $X \to Y$）与集合 $E$（针对 $U \to Y$）相同。
由此可见，引理由概形情形得出；见《态射进阶》引理
\ref{more-morphisms-lemma-geometrically-reduced-constructible}。
\end{proof}

\begin{lemma}
\label{spaces-more-morphisms-lemma-proper-flat-over-dvr-reduced-fibre}
设 $X$ 为离散赋值环 $R$ 上的代数空间，其结构态射
$X \to \Spec(R)$ 固有且平坦。若特殊纤维约化，则
$X$ 和一般纤维 $X_\eta$ 都约化。
\end{lemma}

\begin{proof}
选取平展态射 $U \to X$，其中 $U$ 为仿射概形。
则 $U$ 在 $R$ 上有限型。设 $u \in U$ 位于特殊纤维中。
局部环 $A = \mathcal{O}_{U, u}$ 在 $R$ 上本质有限型，因而是诺特的。
设 $\pi \in R$ 为一致化元。由于 $X$ 在 $R$ 上平坦，
可见 $\pi \in \mathfrak m_A$ 在 $A$ 上为非零因子；又由于 $X$
的特殊纤维约化，$A/\pi A$ 约化。
若 $a \in A$ 且 $a \not = 0$，则存在 $n \geq 0$ 以及元素
$a' \in A$，使得 $a = \pi^n a'$ 且 $a' \not \in \pi A$。
这由 Krull 交定理得出（《代数》引理
\ref{algebra-lemma-intersect-powers-ideal-module-zero}）。
若 $a$ 幂零，则 $a'$ 也幂零，因为 $\pi$ 是非零因子。
但 $a'$ 在约化环 $A/\pi A$ 中的像非零，这是不可能的。
故 $A$ 约化。因此，存在 $u$ 的一个约化开邻域，它位于 $U$ 中
（略去一个小细节；使用 $U$ 为诺特的这一事实）。% P09-ERR-0636
于是可以找到平展态射 $U \to X$，其中 $U$ 是约化概形，
使得 $X$ 的特殊纤维的每一点都在其像中。由于 $X$ 在 $R$ 上固有，
可知 $U \to X$ 满。因此 $X$ 约化。由于 $U \to \Spec(R)$ 的
一般纤维也约化（在仿射片上，它通过取局部化计算），
一般纤维也具有同样的性质。
\end{proof}

\begin{lemma}
\label{spaces-more-morphisms-lemma-geometrically-reduced-open}
设 $S$ 为概形。设 $f : X \to Y$ 为 $S$ 上代数空间的态射。
若 $f$ 平坦、固有且有限表示，则集合
$$
E = \{y \in |Y| : \text{态射 }f : X \to Y\text{ 在点 }y
\text{ 的纤维几何约化}\}
$$
在 $|Y|$ 中开。
\end{lemma}

\begin{proof}
由引理 \ref{spaces-more-morphisms-lemma-base-change-fibres-geometrically-reduced}，
$E$ 的形成与基变换可交换。为检验 $|Y|$ 的子集为开集，
可以用一个平展覆盖的各成员替换 $Y$。故可假设 $Y$ 仿射。
于是 $Y$ 是有限型仿射 $\mathbf{Z}$-概形的余滤极限。
因此，可以假设 $X \to Y$ 是 $X_0 \to Y_0$ 的基变换，其中
$Y_0$ 是有限型 $\mathbf{Z}$-代数的谱，而 $X_0 \to Y_0$ 平坦且固有。
见《空间的极限》引理
\ref{spaces-limits-lemma-descend-finite-presentation}、
\ref{spaces-limits-lemma-descend-flat} 和
\ref{spaces-limits-lemma-eventually-proper}。由于 $E$ 的形成与基变换
可交换（见上），可以假设基底为诺特的。

\medskip\noindent
假设 $Y$ 诺特。由引理
\ref{spaces-more-morphisms-lemma-geometrically-reduced-constructible}，该集合可构造。
因此，只须证明该集合在一般化下稳定（《拓扑》引理
\ref{topology-lemma-characterize-closed-Noetherian}）。由《性质》引理
\ref{properties-lemma-locally-Noetherian-specialization-dvr}，
约化到如下情形：$Y = \Spec(R)$，$R$ 为离散赋值环，且闭纤维
$X_y$ 几何约化。要证明：一般纤维 $X_\eta$ 几何约化。

\medskip\noindent
若不然，则存在有限扩张 $L$，它扩张 $R$ 的分式域，且使 $X_L$ 非约化；
见《域上的空间》引理
\ref{spaces-over-fields-lemma-perfect-reduced}（特征零）和
\ref{spaces-over-fields-lemma-geometrically-reduced-positive-characteristic}
（正特征）。存在离散赋值环 $R' \subset L$，其分式域为 $L$ 且支配
$R$；见《代数》引理 \ref{algebra-lemma-integral-closure-Dedekind}。
替换 $R$，以 $R'$ 代之以后，约化到引理
\ref{spaces-more-morphisms-lemma-proper-flat-over-dvr-reduced-fibre}。
\end{proof}







\section{纤维的连通分支}
\label{spaces-more-morphisms-section-connected}

\noindent
本节对应于《态射进阶》第
\ref{more-morphisms-section-connected} 节。

\begin{lemma}
\label{spaces-more-morphisms-lemma-base-change-fibres-nr-geometrically-connected-components}
设 $S$ 为概形。
设 $f : X \to Y$ 为 $S$ 上代数空间的态射。设
$$
n_{X/Y} : |Y| \to \{0, 1, 2, 3, \ldots, \infty\}
$$
为如下函数：它把 $y \in Y$ 映到 $X_k$ 的连通分支数，% P09-ERR-0637
其中 $\Spec(k) \to Y$ 位于 $y$ 的等价类中，且 $k$ 代数闭。
该函数良定义；若 $g : Y' \to Y$ 为态射，则
$$
n_{X'/Y'} = n_{X/Y} \circ g
% P09-ERR-0645
$$
，其中 $X' \to Y'$ 是 $f$ 的基变换。
\end{lemma}

\begin{proof}
假设 $y' \in Y'$ 的像为 $y \in Y$。% P09-ERR-0638
设 $\Spec(k') \to Y'$ 位于 $y'$ 的等价类中，且 $k'$ 代数闭。
则可以选取交换图
$$
\xymatrix{
\Spec(K) \ar[r] \ar[rd] &
\Spec(k') \ar[r] & Y' \ar[d] \\
& \Spec(k) \ar[r] & Y
}
$$
，其中 $K$ 为代数闭域。结论成立，因为概形态射
$$
\xymatrix{
X'_{k'} & (X'_{k'})_K = (X_k)_K \ar[l] \ar[r] & X_k
}
$$
在连通分支之间诱导双射；见《域上的空间》引理
\ref{spaces-over-fields-lemma-separably-closed-field-connected-components}。
为用此证明函数良定义，取 $Y' = Y$。
\end{proof}







\section{纤维的维数}
\label{spaces-more-morphisms-section-dimension}

\noindent
本节对应于《态射进阶》第
\ref{more-morphisms-section-dimension} 节。

\begin{lemma}
\label{spaces-more-morphisms-lemma-dimension-fibre}
设 $S$ 为概形。设 $f : X \to Y$ 为 $S$ 上代数空间的有限型态射。
设 $y \in |Y|$。下列量相同：% P09-ERR-0639
\begin{enumerate}
\item 令 $d = -\infty$（如果 $y$ 不在 $|f|$ 的像中）；否则，$d$ 是满足
如下条件的最小整数：$f$ 在所有下述点处的相对维数均 $\leq d$：
$x \in |X|$ 且其像为 $y$；
\item 代数空间 $X_k = \Spec(k) \times_Y X$ 的维数，其中
$\Spec(k) \to Y$ 是定义 $y$ 的等价类中的任意态射。
\end{enumerate}
\end{lemma}

\begin{proof}
为理解此陈述，须参阅《空间的态射》定义
\ref{spaces-morphisms-definition-dimension-fibre}、
《空间的性质》定义 \ref{spaces-properties-definition-dimension}，以及
《空间的性质》定义 \ref{spaces-properties-definition-dimension-at-point}。
对于一个固定态射 $\Spec(k) \to Y$，我们将证明 (1) 与 (2) 中的数相等。
选取平展态射 $V \to Y$，其中 $V$ 为仿射概形；并选取点
$v \in V$，其像为 $y$。由于 $V \times_Y \Spec(k) \to \Spec(k)$ 满且平展
（由《空间的性质》引理 \ref{spaces-properties-lemma-points-cartesian}），
可以找到有限可分扩张 $k'/k$（由《态射》引理
\ref{morphisms-lemma-etale-over-field}）以及交换图
$$
\xymatrix{
\Spec(k') \ar[r] \ar[d] & V \ar[d] \\
\Spec(k) \ar[r] & Y
}
$$
可以把 $X \to Y$ 替换为 $V \times_Y X \to V$，并把
$X_k$ 替换为 $X_{k'} = \Spec(k') \times_V (V \times_Y X)$，
因为由《空间的性质》引理
\ref{spaces-properties-lemma-dimension-decent-invariant-under-etale}
和《空间的态射》引理
\ref{spaces-morphisms-lemma-dimension-fibre-after-base-change}，
这不改变所论维数。因此，可以假设 $Y$ 为仿射概形。
在此情形下，可以假设 $k = \kappa(y)$，因为
$X_{\kappa(y)}$ 与 $X_k$ 的维数相同；这由前述《空间的态射》引理
\ref{spaces-morphisms-lemma-dimension-fibre-after-base-change}
以及如下事实得出：对于代数空间 $Z$ 及其底域 $K$，$Z$ 在点
$z \in |Z|$ 处的相对维数依定义等于 $\dim_z(Z)$。
假设 $Y$ 仿射且 $k = \kappa(y)$。则 $X$ 拟紧；% P09-ERR-0640
可以选取仿射概形 $U$ 以及满平展态射 $U \to X$。% P09-ERR-0641
于是 $\dim(X_k) = \dim(U_k) = \max \dim_u(U_k)$，% P09-ERR-0642
它依定义等于 (1) 中给出的数。
\end{proof}

\begin{lemma}
\label{spaces-more-morphisms-lemma-base-change-dimension-fibres}
设 $S$ 为概形。设 $f : X \to Y$ 为 $S$ 上代数空间的有限型态射。设
$$
n_{X/Y} : |Y| \to \{-\infty, 0, 1, 2, 3, \ldots\}
$$
为把 $y \in |Y|$ 映到引理 \ref{spaces-more-morphisms-lemma-dimension-fibre} 中所论整数的函数。
若 $g : Y' \to Y$ 为态射，则
$$
n_{X'/Y'} = n_{X/Y} \circ |g|
$$
，其中 $X' \to Y'$ 是 $f$ 的基变换。
\end{lemma}

\begin{proof}
这立即由引理 \ref{spaces-more-morphisms-lemma-dimension-fibre} 得出。
\end{proof}

\begin{lemma}
\label{spaces-more-morphisms-lemma-dimension-fibres-flat}
设 $S$ 为概形。
设 $f : X \to Y$ 为 $S$ 上代数空间的平坦有限表示态射。
设 $n_{X/Y}$ 为引理 \ref{spaces-more-morphisms-lemma-base-change-dimension-fibres}
中引入的、在 $Y$ 上给出 $f$ 的纤维维数的函数。
则 $n_{X/Y}$ 下半连续。
\end{lemma}

\begin{proof}
设 $V \to Y$ 为满平展态射，其中 $V$ 是概形。
若能证明复合 $n_{X/Y} \circ |g|$ 下半连续，% P09-ERR-0643
则因为 $|g|$ 开，引理成立。故可以假设 $Y$ 为概形。
在局部上工作，可以假设 $V$ 为仿射概形。% P09-ERR-0644
于是可以选取仿射概形 $U$ 以及满平展态射 $U \to X$。
于是 $n_{X/Y} = n_{U/Y}$。故可以假设 $X$ 和 $Y$ 都是概形。
在此情形下，引理由《态射进阶》引理
\ref{more-morphisms-lemma-dimension-fibres-flat} 得出。
\end{proof}

\begin{lemma}
\label{spaces-more-morphisms-lemma-dimension-fibres-proper}
设 $S$ 为概形。
设 $f : X \to Y$ 为 $S$ 上代数空间的固有态射。
设 $n_{X/Y}$ 为引理 \ref{spaces-more-morphisms-lemma-base-change-dimension-fibres}
中引入的、在 $Y$ 上给出 $f$ 的纤维维数的函数。
则 $n_{X/Y}$ 上半连续。
\end{lemma}

\begin{proof}
令 $Z_d = \{x \in |X| :
\text{态射 }f\text{ 在点 }x\text{ 的纤维维数 }> d\}$。
由《空间的态射》引理
\ref{spaces-morphisms-lemma-openness-bounded-dimension-fibres}，
$Z_d$ 是 $|X|$ 的闭子集。由于 $f$ 固有，% P09-ERR-0646
$f(Z_d)$ 在 $|Y|$ 中闭。
由于 $y \in f(Z_d) \Leftrightarrow n_{X/Y}(y) > d$，引理成立。
\end{proof}

\begin{lemma}
\label{spaces-more-morphisms-lemma-dimension-fibres-proper-flat}
设 $S$ 为概形。设 $f : X \to Y$ 为 $S$ 上代数空间的固有、平坦且
有限表示的态射。
设 $n_{X/Y}$ 为引理 \ref{spaces-more-morphisms-lemma-base-change-dimension-fibres}
中引入的、在 $Y$ 上给出 $f$ 的纤维维数的函数。
则 $n_{X/Y}$ 局部常值。
\end{lemma}

\begin{proof}
这是引理 \ref{spaces-more-morphisms-lemma-dimension-fibres-flat} 和
\ref{spaces-more-morphisms-lemma-dimension-fibres-proper} 的直接推论。
\end{proof}





\section{catenary 代数空间}
\label{spaces-more-morphisms-section-catenary}

\noindent
本节继续《良态空间》第 \ref{decent-spaces-section-catenary} 节中开始的讨论。
下一个引理的证明将使用如下引理。

\begin{lemma}
\label{spaces-more-morphisms-lemma-construct-glueing}
设 $S$ 为概形。设 $f : X \to Y$ 为 $S$ 上代数空间的整态射。
设 $y \in |Y|$ 为可由闭浸入 $y : \Spec(k) \to Y$ 表示的点。
则存在分解 $X \to X' \to Y$，其复合为 $f$，并且
\begin{enumerate}
\item $X' \to Y$ 为整态射；
\item $X \to X'$ 在 $X' \setminus X'_y$ 上为同构；
\item $X'_y$ 有唯一点 $x'$，且 $\kappa(x') = k$。
\end{enumerate}
此外，若 $f$ 有限且 $Y$ 局部 Noetherian，则 $X' \to Y$ 有限。
\end{lemma}

\begin{proof}
由《空间的态射》引理 \ref{spaces-morphisms-lemma-pushforward}，层
$f_*\mathcal{O}_X$、$(X_y \to Y)_*\mathcal{O}_{X_y}$ 和
$y_*\mathcal{O}_{\Spec(k)}$ 都是拟凝聚 $\mathcal{O}_Y$-代数层。考虑映射
$$
f_*\mathcal{O}_Y \longrightarrow
% P09-ERR-0647
(X_y \to Y)_*\mathcal{O}_{X_y} \longleftarrow
y_*\mathcal{O}_{\Spec(k)}
$$
其纤维积是拟凝聚 $\mathcal{O}_Y$-代数层 $\mathcal{A}'$。
可以把 $X' \to Y$ 定义为 $\mathcal{A}'$ 在 $Y$ 上的相对谱；见
《态射》引理 \ref{morphisms-lemma-affine-equivalence-algebras}。
此构造与任意基变换可交换。特别地，显然在开子空间
$|Y| \setminus \{y\}$ 上态射 $X \to X'$ 是同构，而在
$|Y| \setminus \{y\}$ 上态射 $X' \to Y$ 是整态射。
尚须在 $y$ 的一个小邻域内证明这些陈述。
选取仿射概形 $V = \Spec(R)$ 以及平展态射
$\varphi : V \to Y$，使得 $y$ 位于 $\varphi$ 的像中。
则 $V_y$ 是 $V$ 的闭子概形，并且在 $k$ 上平展，因而由有限多个点组成，
且每个点的剩余域都是 $k$ 的可分扩张
（见《良态空间》评注 \ref{decent-spaces-remark-recall}）。
缩小 $V$ 后，可以假设存在唯一的闭点
$v = \Spec(l) \to V$ 映到 $y$，其中 $l/k$ 有限可分。
可以写成 $V \times_Y X = \Spec(C)$，其中 $R \to C$ 是整环同态。
上述与基变换的相容性给出
$U \times_X Y' = \Spec(C')$，其中 % P09-ERR-0648
$$
C' = C \times_{C \otimes_R l} l
$$
由于 $R \to l$ 满，$C \to C \otimes_R l$ 也满，因而这是
《代数进阶》情形 \ref{more-algebra-situation-module-over-fibre-product}
中研究的那类纤维积
（其中 $A, A', B, B'$ 分别对应于 $C \otimes_R l, C, l, C'$）。
注意到 $C'$ 是一个 $R$-子代数，包含于 $C$ 中，因而在 $R$ 上整；
这证明了 (1)。
最后，《代数进阶》引理
\ref{more-algebra-lemma-points-of-fibre-product} 表明
$V \times_X Y' = \Spec(C')$ % P09-ERR-0649
有唯一的点 $y''$ 位于 $v$ 上方，其剩余域为 $l$
（它对应于显然的满映射 $C' \to l$）。
因此 $X_y \times_{\Spec(k)} \Spec(l)$ 有唯一的点，% P09-ERR-0650
其剩余域为 $l$。由于 $l/k$ 有限可分，这蕴含 $X'_y$
有唯一的点，且其剩余域为 $k$，即 (3) 成立。

\medskip\noindent
为证明最后一个陈述，注意到若 $Y$ 局部 Noetherian，则 $R$ 是
Noetherian 环；若 $f$ 有限，则 $R \to C$ 有限。于是由《代数进阶》引理
\ref{more-algebra-lemma-fibre-product-finite-type}，$C'$ 是有限型
$R$-代数。这证明了 $X' \to Y$ 有限。
\end{proof}

\begin{lemma}
\label{spaces-more-morphisms-lemma-universally-catenary-dimension-function}
设 $S$ 为概形。设 $B$ 为 $S$ 上的代数空间。设
$\delta : |B| \to \mathbf{Z}$ 为函数。假设 $B$ 良态、局部 Noetherian
且普遍 catenary，并且 $\delta$ 是维数函数。若 $X$ 是 $B$ 上的良态
代数空间，其结构态射 $f : X \to B$ 局部有限型，则按下述规则定义
$\delta_X : |X| \to \mathbf{Z}$：
$$
\delta_X(x) = \delta(f(x)) + \text{超越次数 }x/f(x)
% P09-ERR-0651
$$
（《空间的态射》定义
\ref{spaces-morphisms-definition-dimension-fibre}）。
则 $\delta_X$ 是维数函数。
\end{lemma}

\begin{proof}
问题在 $B$ 上是局部的，故可以假设 $B$ 拟紧。
由《良态空间》引理
\ref{decent-spaces-lemma-locally-Noetherian-decent-quasi-separated}，
$B$ 拟分离。由《空间的极限》命题
\ref{spaces-limits-proposition-there-is-a-scheme-finite-over}，
可以选取有限满态射 $\pi : Y \to X$，其中 $Y$ 是概形。% P09-ERR-0652
断言：$\delta_Y$ 是维数函数。

\medskip\noindent
该断言蕴含引理。对引理中的 $X \to B$，令
$Z = Y \times_B X$，并记投影为 $p : Z \to Y$ 和 $q : Z \to X$。
则
$$
\delta_Z(z) = \delta_Y(p(z)) + \text{超越次数 }z/p(z)
% P09-ERR-0651
$$
且 $\delta_Z(z) = \delta_X(q(z))$。这由《空间的态射》引理
\ref{spaces-morphisms-lemma-dimension-fibre-at-a-point-additive}
以及这些超越次数对于有限态射为零这一事实得出。由《良态空间》引理
\ref{decent-spaces-lemma-scheme-with-dimension-function} 和该断言，
可知 $\delta_Z$ 是维数函数。再由《良态空间》引理
\ref{decent-spaces-lemma-check-dimension-function-finite-cover}，
可知 $\delta_X$ 是维数函数。

\medskip\noindent
证明该断言。考虑特化 $y \leadsto y'$，其中 $y \not = y'$，
二者都是 Noetherian 概形 $Y$ 的点。于是
$\delta_Y(y) > \delta_Y(y')$，因为 $Y$ 的纤维中的点之间不存在特化
（见《良态空间》引理
\ref{decent-spaces-lemma-conditions-on-fibre-and-qf}）。
把它用于一条特化链，得到
$$
\delta_Y(y) - \delta_Y(y') \geq
\text{余维}(\overline{\{y'\}}, \overline{\{y\}})
$$
我们的任务是证明等号成立。由《性质》引理
\ref{properties-lemma-locally-Noetherian-closed-point}，
可以选取一个特化 $y' \leadsto y_0$。只须证明
$\delta_Y(y) - \delta_Y(y_0) =
\text{余维}(\overline{\{y_0\}}, \overline{\{y\}})$，
因为这将同时蕴含 $y \leadsto y'$ 和 $y' \leadsto y_0$ 的等式。

\medskip\noindent
选取极大特化链
$y = y_c \leadsto y_{c - 1} \leadsto \ldots \leadsto y_0$，它位于 $Y$ 中。
令 $b = \pi(y)$ 且 $b_0 = \pi(y_0)$。
选取极大特化链
$b = b_e \leadsto b_{e - 1} \leadsto \ldots \leadsto b_0$，它位于 $|B|$ 中。
须证明 $e = c$。由于 $\pi$ 闭
（《空间的态射》引理 \ref{spaces-morphisms-lemma-finite-proper}），
可以找到特化序列
$y = y'_e \leadsto y'_{e - 1} \leadsto \ldots \leadsto y'_0$，
它映到
$b = b_e \leadsto b_{e - 1} \leadsto \ldots \leadsto b_0$。
注意 $y'_e \leadsto y'_{e - 1} \leadsto \ldots \leadsto y'_0$
也是极大链。若 $y_0 = y'_0$，则因 $Y$ catenary，得到所需的
$e = c$。下一段用一个小技巧约化到此情形，并以同样方式得出结论。

\medskip\noindent
由于 $\pi$ 闭，$b_0$ 是 $|B|$ 的闭点。
由《良态空间》引理 \ref{decent-spaces-lemma-decent-space-closed-point}，
可以把 $b_0$ 表示为闭浸入 $b_0 : \Spec(k) \to B$。
由引理 \ref{spaces-more-morphisms-lemma-construct-glueing}，可以找到分解
$$
Y \to Y' \to X
% P09-ERR-0653
$$
，其中 $\pi' : Y' \to X$ 有限，而 $Y \to Y'$ 是把 % P09-ERR-0653
$y_0$ 和 $y'_0$ 映到同一点、且在 $b_0$ 的逆像之外为同构的态射。
% P09-ERR-0654
（当然，$Y'$ 不会是概形，但这不影响下面的论证。）
显然，极大特化链
$y_c \leadsto y_{c - 1} \leadsto \ldots \leadsto y_0$ 和
$y'_e \leadsto y'_{e - 1} \leadsto \ldots \leadsto y'_0$
映为 $Y'$ 中有相同起点和终点的极大特化链。
由于 $B$ 普遍 catenary，$|Y'|$ catenary，结论同前。
\end{proof}





\section{态射的平展局部化}
\label{spaces-more-morphisms-section-etale-localization}

\noindent
本节对应于《态射进阶》第
\ref{more-morphisms-section-etale-localization} 节。

\begin{lemma}
\label{spaces-more-morphisms-lemma-etale-splits-off-quasi-finite-part}
设 $S$ 为概形。
设 $f : X \to Y$ 为 $S$ 上代数空间的态射。
设 $y \in |Y|$。设 $x_1, \ldots, x_n \in |X|$ 均映到 $y$。
假设
\begin{enumerate}
\item $f$ 局部有限型；
\item $f$ 分离；
\item $f$ 在 $x_1, \ldots, x_n$ 处拟有限；并且
\item $f$ 拟紧，或 $Y$ 良态。
\end{enumerate}
则存在带基点代数空间的平展态射 $(U, u) \to (Y, y)$ 以及分解
$$
U \times_Y X = W \amalg V
$$
为既开又闭子空间，使得态射 $V \to U$ 有限；
$|V| \to |U|$ 在 $u$ 上方纤维的每个点都映到某个 $x_i$；
而 $|W| \to |U|$ 在 $u$ 上方的纤维不含映到任何 $x_i$ 的点。
\end{lemma}

\begin{proof}
设 $(U, u) \to (Y, y)$ 为代数空间的平展态射，并考虑满足如下条件的
$w \in |U \times_Y X|$ 之集合：它映到 $u \in |U|$ 以及某个
$x_i \in |X|$。由《良态空间》引理
\ref{decent-spaces-lemma-qf-and-qc-finite-fibre}（若 $f$ 有限型）
或《良态空间》引理
\ref{decent-spaces-lemma-decent-finite-fibre}（若 $Y$ 良态），
此集合有限。因此可以把 $f$ 替换为基变换 $U \times_Y X \to U$，
并把 $x_1, \ldots, x_n$ 替换为这些 $w$ 之集合。
特别地，可以并且确实假设 $Y$ 为仿射概形；于是 $X$ 是分离代数空间。

\medskip\noindent
选取仿射概形 $Z$ 以及平展态射 $Z \to X$，使
$x_1, \ldots, x_n$ 位于 $|Z| \to |X|$ 的像中。
$|Z| \to |X|$ 的纤维有限；见《空间的性质》引理
\ref{spaces-properties-lemma-finite-fibres-presentation}
（或《良态空间》第
\ref{decent-spaces-section-reasonable-decent} 节中更一般的讨论）。
令 $\{z_1, \ldots, z_m\} \subset |Z|$ 为
$\{x_1, \ldots, x_n\}$ 的逆像。由《态射进阶》引理
\ref{more-morphisms-lemma-etale-splits-off-quasi-finite-part-technical}，
存在平展态射 $(U, u) \to (Y, y)$，使得
$U \times_Y Z = Z_1 \amalg Z_2$，其中 $Z_1 \to U$ 有限，且
$(Z_1)_y = \{z_1, \ldots, z_m\}$。% P09-ERR-0655
可以假设 $U$ 仿射，因而 $Z_1$ 也仿射。

\medskip\noindent
由于 $f$ 分离，像 $V$（即 $Z_1 \to X$ 的像）既开又闭
（《空间的态射》引理
\ref{spaces-morphisms-lemma-universally-closed-permanence}）。% P09-ERR-0656
令 $W = X \setminus V$，便得到引理所述的分解。% P09-ERR-0656
为完成证明，须证明 $V \to U$ 有限。
由于 $Z_1 \to V$ 满且平展，$V$ 是 $Z_1$ 关于平展等价关系
$R = Z_1 \times_V Z_1$ 的商；见《空间》引理
\ref{spaces-lemma-space-presentation}。
由于 $f$ 分离，$V \to U$ 分离，且 $R$ 在
$Z_1 \times_U Z_1$ 中闭。由于 $Z_1 \to U$ 有限，投影
$s, t : R \to Z_1$ 有限。因此，由《群胚》命题
\ref{groupoids-proposition-finite-flat-equivalence}，$V$ 是仿射概形。
由《态射》引理 \ref{morphisms-lemma-image-proper-is-proper}，
可知 $V \to U$ 固有；再由《态射》引理
\ref{morphisms-lemma-finite-proper}，可知 $V \to U$ 有限，
从而完成证明。
\end{proof}

\begin{lemma}
\label{spaces-more-morphisms-lemma-etale-splits-off-just-one-quasi-finite-part}
设 $S$ 为概形。
设 $f : X \to Y$ 为 $S$ 上代数空间的态射。
设 $x \in |X|$，其像为 $y \in |Y|$。假设
\begin{enumerate}
\item $f$ 局部有限型；
\item $f$ 分离；并且
\item $f$ 在 $x$ 处拟有限。
\end{enumerate}
则存在带基点代数空间的平展态射 $(U, u) \to (Y, y)$ 以及分解
$$
U \times_Y X = W \amalg V
$$
为既开又闭子空间，使得态射 $V \to U$ 有限，并且存在点
$v \in |V|$，它映到 $x$（在 $|X|$ 中）以及 $u$（在 $|U|$ 中）。
\end{lemma}

\begin{proof}
选取概形 $U$、点 $u \in U$ 以及平展态射
$U \to Y$，该态射把 $u$ 映到 $y$。存在点
$x' \in |U \times_Y X|$，它映到 $x$（在 $|X|$ 中）以及
$u$（在 $|U|$ 中）（《空间的性质》引理
\ref{spaces-properties-lemma-points-cartesian}）。
为完成证明，把引理 \ref{spaces-more-morphisms-lemma-etale-splits-off-quasi-finite-part}
应用于态射 $U \times_Y X \to U$ 和点 $x'$。
它适用，因为 $U$ 是概形，因而 $u$ 来自单态射
$\Spec(\kappa(u)) \to U$。
\end{proof}





\section{Zariski 主定理}
\label{spaces-more-morphisms-section-structure-quasi-finite}

\noindent
本节应用上一节的结果，证明代数空间态射的 Zariski 主定理。
本节对应于《态射进阶》第
\ref{more-morphisms-section-application-etale-neighbourhoods} 节。

\begin{lemma}
\label{spaces-more-morphisms-lemma-finite-type-separated}
设 $S$ 为概形。设 $f : X \to Y$ 为 $S$ 上代数空间的有限型分离态射。
设 $Y'$ 为 $Y$ 在 $X$ 中的正规化。图示为：
$$
\xymatrix{
X \ar[rd]_f \ar[rr]_{f'} & & Y' \ar[ld]^\nu \\
& Y &
}
$$
则存在开子空间 $U' \subset Y'$，使得
\begin{enumerate}
\item $(f')^{-1}(U') \to U'$ 为同构；并且
\item $(f')^{-1}(U') \subset X$ 恰为 $f$ 拟有限的点之集合。
\end{enumerate}
\end{lemma}

\begin{proof}
由《空间的态射》引理
\ref{spaces-morphisms-lemma-locally-finite-type-quasi-finite-part}，
存在开子空间 $U \subset X$，它对应于 $|X|$ 中 $f$ 拟有限的点。
须证明
\begin{enumerate}
\item[(a)] $|U| \to |Y'|$ 的像为 $|U'|$，其中 $U'$ 是
$Y'$ 的某个开子空间；
\item[(b)] $U = f^{-1}(U')$；并且 % P09-ERR-0657
\item[(c)] $U \to U'$ 为同构。
\end{enumerate}
由于 $U$ 的形成与任意基变换可交换
（《空间的态射》引理
\ref{spaces-morphisms-lemma-locally-finite-type-quasi-finite-part}），
由于正规化 $Y'$ 的形成与光滑基变换可交换
（引理 \ref{spaces-more-morphisms-lemma-normalization-smooth-localization}），
由于平展态射是开态射，并且由于“为同构”在底上是 fpqc 局部的
（《空间上的下降》引理
\ref{spaces-descent-lemma-descending-property-isomorphism}），
只须在 $Y$ 上平展局部地证明 (a)、(b)、(c)
（省略若干细节）。% P09-ERR-0658
因此，可以假设 $Y$ 为仿射概形。这蕴含 $Y'$ 也是（仿射）概形。

\medskip\noindent
设 $x \in |U|$。断言：存在开邻域 $f'(x) \in V \subset Y'$，
使得 $(f')^{-1}V \to V$ 为同构。先证明该断言蕴含引理。
事实上，此时 $(f')^{-1}V \cong V$ 是概形（因为它是 $Y'$ 的开集），
在 $Y$ 上局部有限型（因为它是 $X$ 的开子空间）；而对于 $v \in V$，
剩余域扩张 $\kappa(v)/\kappa(\nu(v))$ 是代数的（因为
$V \subset Y'$ 且 $Y'$ 在 $Y$ 上整）。因此 $V \to Y$ 的纤维离散
（《态射》引理
\ref{morphisms-lemma-algebraic-residue-field-extension-closed-point-fibre}），
且 $(f')^{-1}V \to Y$ 局部拟有限
（《态射》引理 \ref{morphisms-lemma-locally-quasi-finite-fibres}）。
这蕴含 $(f')^{-1}V \subset U$ 且 $V \subset U'$。% P09-ERR-0659
由于 $x$ 任意，可知 (a)、(b)、(c) 成立。

\medskip\noindent
令 $y = f(x) \in |Y|$。设 $(T, t) \to (Y, y)$ 为带基点概形的
平展态射。用下标 $ {}_T$ 表示到 $T$ 的基变换。
设 $z \in X_T$ 为纤维 $X_t$ 中位于 $x$ 上方的点。% P09-ERR-0660
注意 $U_T \subset X_T$ 是 $f_T$ 拟有限的点之集合；见
《空间的态射》引理
\ref{spaces-morphisms-lemma-locally-finite-type-quasi-finite-part}。
注意
$$
X_T \xrightarrow{f'_T} Y'_T \xrightarrow{\nu_T} T
$$
是 $T$ 在 $X_T$ 中的正规化；见引理
\ref{spaces-more-morphisms-lemma-normalization-smooth-localization}。
假设断言对 $z \in U_T \subset X_T \to Y'_T \to T$ 成立，% P09-ERR-0660
即假设可以找到开邻域 $f'_T(z) \in V' \subset Y'_T$，使得
$(f'_T)^{-1}V' \to V'$ 为同构。态射 $Y'_T \to Y'$ 平展，
故像 $V \subset Y'$（即 $V'$ 的像）是开集。注意 $f'(x) \in V$，
因为 $f'_T(z) \in V'$。还要注意
$$
\xymatrix{
(f'_T)^{-1}V' \ar[r] \ar[d] & (f')^{-1}(V) \ar[d] \\
V' \ar[r] & V
}
$$
是纤维方块（因为 $Y'_T \times_{Y'} X = X_T$）。
由于左侧竖直箭头是同构，且 $\{V' \to V\}$ 是平展覆盖，
由《空间上的下降》引理
\ref{spaces-descent-lemma-descending-property-isomorphism}，
右侧竖直箭头也是同构。换言之，断言对
$x \in U \subset X \to Y' \to Y$ 成立。% P09-ERR-0660

\medskip\noindent
由上一段的结果，为证明断言对 $x \in |U|$ 成立，可以把 $Y$
替换为平展邻域 $T$，它是 $y = f(x)$ 的邻域；并把 $x$ 替换为
位于 $x$ 上方、属于 $T \times_Y X$ 的任意点。% P09-ERR-0660
因此，可以假设存在分解
$$
X = V \amalg W
$$
为既开又闭子空间，其中 $V \to Y$ 有限且 $x \in V$；
见引理 \ref{spaces-more-morphisms-lemma-etale-splits-off-quasi-finite-part}。
由于 $X$ 是 $V$ 与 $W$ 在 $Y$ 上的不交并，并且 $V \to Y$ 有限，
可知 $Y$ 在 $X$ 中的正规化是态射
$$
X = V \amalg W \longrightarrow V \amalg W' \longrightarrow S
% P09-ERR-0661
$$
，其中 $W'$ 是 $Y$ 在 $W$ 中的正规化；见《空间的态射》引理
\ref{spaces-morphisms-lemma-normalization-in-disjoint-union}、
\ref{spaces-morphisms-lemma-finite-integral} 和
\ref{spaces-morphisms-lemma-normalization-in-integral}。
断言成立，证明完成。
\end{proof}

\noindent
下一个引理与《空间的态射》引理
\ref{spaces-morphisms-lemma-quasi-finite-separated-quasi-affine}
重复。保留同一引理的两个副本，是因为证明有所不同。下面给出的证明
依赖于上面所述不可表示代数空间态射的 Zariski 主定理，而《空间的态射》
引理 \ref{spaces-morphisms-lemma-quasi-finite-separated-quasi-affine}
的证明则依赖于《空间的态射》命题
\ref{spaces-morphisms-proposition-locally-quasi-finite-separated-over-scheme}，
从而约化到概形态射的情形。

\begin{lemma}
\label{spaces-more-morphisms-lemma-quasi-finite-separated-quasi-affine}
设 $S$ 为概形。
设 $f : X \to Y$ 为 $S$ 上代数空间的态射。
假设 $f$ 拟有限且分离。设 $Y'$ 为 $Y$ 在 $X$ 中的正规化。
图示为：
$$
\xymatrix{
X \ar[rd]_f \ar[rr]_{f'} & & Y' \ar[ld]^\nu \\
& Y &
}
$$
则 $f'$ 是拟紧开浸入，且 $\nu$ 是整态射。特别地，$f$ 拟仿射。
\end{lemma}

\begin{proof}
这由引理 \ref{spaces-more-morphisms-lemma-finite-type-separated} 得出。事实上，由该引理，
存在开子空间 $U' \subset Y'$，使得
$(f')^{-1}(U') = X$（！）且 $X \to U'$ 是同构！换言之，$f'$ 是开浸入。
注意 $f'$ 拟紧，因为 $f$ 拟紧，且 $\nu : Y' \to Y$ 分离
（《空间的态射》引理
\ref{spaces-morphisms-lemma-quasi-compact-permanence}）。
因此，对于每个仿射概形 $Z$ 和态射 $Z \to Y$，纤维积
$Z \times_Y X$ 是仿射概形 $Z \times_Y Y'$ 的拟紧开子概形。
故依定义，$f$ 拟仿射。
\end{proof}

\begin{lemma}[Zariski 主定理]
\label{spaces-more-morphisms-lemma-quasi-finite-separated-pass-through-finite}
设 $S$ 为概形。
设 $f : X \to Y$ 为 $S$ 上代数空间的态射。
假设 $f$ 拟有限且分离，并假设 $Y$ 拟紧且拟分离。则存在分解
$$
\xymatrix{
X \ar[rd]_f \ar[rr]_j & & T \ar[ld]^\pi \\
& Y &
}
$$
，其中 $j$ 是拟紧开浸入，且 $\pi$ 有限。
\end{lemma}

\begin{proof}
设 $X \to Y' \to Y$ 如引理
\ref{spaces-more-morphisms-lemma-quasi-finite-separated-quasi-affine} 的结论所述。
由《空间的极限》引理
\ref{spaces-limits-lemma-integral-algebra-directed-colimit-finite}，
可以把 $\nu_*\mathcal{O}_{Y'} = \colim_{i \in I} \mathcal{A}_i$
写成有限拟凝聚 $\mathcal{O}_X$-代数 % P09-ERR-0662
$\mathcal{A}_i \subset \nu_*\mathcal{O}_{Y'}$ 的有向余极限。于是
$\pi_i : T_i = \underline{\Spec}_Y(\mathcal{A}_i) \to Y$
对于每个 $i$ 都是有限态射。注意转移态射 $T_{i'} \to T_i$
都是仿射态射，并且
$Y' = \lim T_i$。

\medskip\noindent
由《空间的极限》引理 \ref{spaces-limits-lemma-descend-opens}，
存在 $i$ 以及拟紧开集 $U_i \subset T_i$，其在 $Y'$ 中的逆像等于
$f'(X)$。对于 $i' \geq i$，令 $U_{i'}$ 为 $U_i$ 在 $T_{i'}$
中的逆像。于是
$X \cong f'(X) = \lim_{i' \geq i} U_{i'}$；见《空间的极限》引理
\ref{spaces-limits-lemma-directed-inverse-system-has-limit}。
由《空间的极限》引理
\ref{spaces-limits-lemma-finite-type-eventually-closed}，可知
$X \to U_{i'}$ 对于某个 $i' \geq i$ 是闭浸入。
（事实上，$X \cong U_{i'}$ 对充分大的 $i'$ 成立，但这里不需要。）
因此 $X \to T_{i'}$ 是浸入。由《空间的态射》引理
\ref{spaces-morphisms-lemma-factor-the-other-way}，可以将其分解为
$X \to T \to T_{i'}$，其中第一个箭头是开浸入，第二个是闭浸入。
证明完成。
\end{proof}

\begin{lemma}
\label{spaces-more-morphisms-lemma-quasi-finite-separated-pass-through-finite-addendum}
采用引理 \ref{spaces-more-morphisms-lemma-quasi-finite-separated-pass-through-finite}
的记号和假设。再假设 $f$ 局部有限表示。则可以选取分解，使
$T$ 在 $Y$ 上有限且有限表示。
\end{lemma}

\begin{proof}
由《空间的极限》引理
\ref{spaces-limits-lemma-finite-in-finite-and-finite-presentation}，
可以写成 $T = \lim T_i$，其中所有 $T_i$ 在 $Y$ 上有限且有限表示，
而转移态射 $T_{i'} \to T_i$ 都是闭浸入。由《空间的极限》引理
\ref{spaces-limits-lemma-descend-opens}，存在 $i$ 以及开子概形
$U_i \subset T_i$，其在 $T$ 中的逆像是 $X$。% P09-ERR-0663
由《空间的极限》引理
\ref{spaces-limits-lemma-finite-type-eventually-closed}，可知
$X \cong U_i$ 对充分大的 $i$ 成立。把 $T$ 替换为 $T_i$ 即完成证明。
\end{proof}





\section{Zariski 主定理的应用 I}
\label{spaces-more-morphisms-section-applications-zmt}

\noindent
第一个应用是把有限态射刻画为具有有限纤维的固有态射。

\begin{lemma}
\label{spaces-more-morphisms-lemma-characterize-finite}
设 $S$ 为概形。设 $f : X \to Y$ 为 $S$ 上代数空间的态射。
下列条件等价：
\begin{enumerate}
\item $f$ 有限；
\item $f$ 固有且局部拟有限；
\item $f$ 固有，并且 $|X_k|$ 对于每个态射
$\Spec(k) \to Y$ 都是离散空间，其中 $k$ 是域；
\item $f$ 万有闭、分离、局部有限型，并且 $|X_k|$ 对于每个态射
$\Spec(k) \to Y$ 都是离散空间，其中 $k$ 是域。
\end{enumerate}
\end{lemma}

\begin{proof}
由《空间的态射》引理
\ref{spaces-morphisms-lemma-finite-proper}、
\ref{spaces-morphisms-lemma-finite-quasi-finite}，有 (1) $\Rightarrow$ (2)。
由《空间的态射》引理
\ref{spaces-morphisms-lemma-locally-quasi-finite}，有 (2) $\Rightarrow$ (3)。
依定义，(3) 蕴含 (4)。

\medskip\noindent
假设 (4)。由于 $f$ 万有闭，它拟紧
（《空间的态射》引理
\ref{spaces-morphisms-lemma-universally-closed-quasi-compact}）。
选取点 $y$，它属于 $|Y|$。选取一个表示 $y$ 的态射
$\Spec(k) \to Y$。
注意 $|X_k|$ 作为拟紧离散空间是有限离散空间。
映射 $|X_k| \to |X|$ 满射到 $|X| \to |Y|$ 在 $y$ 上方的纤维
（《空间的性质》引理 \ref{spaces-properties-lemma-points-cartesian}）。
由《空间的态射》引理
\ref{spaces-morphisms-lemma-quasi-finite-at-point}，可知
$X \to Y$ 在 $|X| \to |Y|$ 位于 $y$ 上方的纤维的所有点处拟有限。
选取基本平展邻域 $(U, u) \to (Y, y)$ 以及分解
$X_U = V \amalg W$，如引理
\ref{spaces-more-morphisms-lemma-etale-splits-off-quasi-finite-part} 所述，并使其适合
$|X|$ 中位于 $y$ 上方的所有点。注意 $W_u = \emptyset$，
因为采用了 $|X| \to |Y|$ 位于 $y$ 上方纤维中的所有点。
由于 $f$ 万有闭，$|W|$ 在 $|U|$ 中的像是不含 $u$ 的闭集。
缩小 $U$ 后，可以假设 $W = \emptyset$。换言之，
$X_U = V$ 在 $U$ 上有限。由于 $y \in |Y|$ 任意，这意味着存在
平展态射族 $\{U_i \to Y\}$，其像覆盖 $Y$，并且基变换
$X_{U_i} \to U_i$ 都有限。由《空间的态射》引理
\ref{spaces-morphisms-lemma-integral-local}，可知 $f$ 有限。
\end{proof}

\noindent
作为推论，得到如下有用结果。

\begin{lemma}
\label{spaces-more-morphisms-lemma-proper-finite-fibre-finite-in-neighbourhood}
设 $S$ 为概形。设 $f : X \to Y$ 为 $S$ 上代数空间的态射。
设 $y \in |Y|$。假设
\begin{enumerate}
\item $f$ 固有；并且
\item $f$ 在所有 $x \in |X|$ 处拟有限，且这些点位于 $y$ 上方
（《良态空间》引理
\ref{decent-spaces-lemma-conditions-on-fibre-and-qf}）。
\end{enumerate}
则存在开邻域 $V \subset Y$，它是 $y$ 的邻域，并使得
$f|_{f^{-1}(V)} : f^{-1}(V) \to V$ 有限。
\end{lemma}

\begin{proof}
由《空间的态射》引理
\ref{spaces-morphisms-lemma-locally-finite-type-quasi-finite-part}，
$f$ 拟有限的点之集合是开集 $U \subset X$。
令 $Z = X \setminus U$。则 $y \not \in f(Z)$。由于 $f$ 固有，
$f(Z) \subset Y$ 是闭集。选取任意开邻域 $V \subset Y$，它含 $y$
且满足 $Z \cap V = \emptyset$。% P09-ERR-0664
于是 $f^{-1}(V) \to V$ 局部拟有限且固有。因此由引理
\ref{spaces-more-morphisms-lemma-characterize-finite}，$f^{-1}(V) \to V$ 有限。
\end{proof}

\begin{lemma}
\label{spaces-more-morphisms-lemma-flat-proper-family-cannot-collapse-fibre}
\begin{slogan}
压缩固有族的一个纤维，会迫使附近的纤维也被压缩。
\end{slogan}
设 $S$ 为概形。设
$$
\xymatrix{
X \ar[rr]_h \ar[rd]_f & & Y \ar[ld]^g \\
& B
}
$$
为 $S$ 上代数空间态射的交换图。% P09-ERR-0665
设 $b \in B$，并设 $\Spec(k) \to B$ 为 $b$ 的等价类中的态射。
% P09-ERR-0666
假设
\begin{enumerate}
\item $X \to B$ 是固有态射；
\item $Y \to B$ 分离且局部有限型；
\item 下列条件之一成立：
\begin{enumerate}
\item $|X_k| \to |Y_k|$ 的像有限；
\item $|f|^{-1}(\{b\})$ 在 $|Y|$ 中的像有限，且 $B$ 良态。
\end{enumerate}
\end{enumerate}
则存在开子空间 $B' \subset B$，它包含 $b$，并使得
$X_{B'} \to Y_{B'}$ 分解经过闭子空间 $Z \subset Y_{B'}$，
且后者在 $B'$ 上有限。
\end{lemma}

\begin{proof}
令 $Z \subset Y$ 为 $h$ 的概形论像；见《空间的态射》第
\ref{spaces-morphisms-section-scheme-theoretic-image} 节。
由《空间的态射》引理
\ref{spaces-morphisms-lemma-scheme-theoretic-image-is-proper}，
态射 $X \to Z$ 满，且 $Z \to B$ 固有。因此
$$
\{x \in |X|\text{ 映到 }b\} \to
\{z \in |Z|\text{ 映到 }b\}
$$
以及 $|X_k| \to |Z_k|$ 都满。可见，(3)(a) 或 (3)(b) 都蕴含
$Z \to B$ 在 $|Z|$ 中位于 $b$ 上方的所有点处拟有限。% P09-ERR-0667
这由《良态空间》引理
\ref{decent-spaces-lemma-conditions-on-fibre-and-qf} 得出。
因此，由引理
\ref{spaces-more-morphisms-lemma-proper-finite-fibre-finite-in-neighbourhood}，
$Z \to B$ 在 $b$ 的一个开邻域上有限。
\end{proof}







\section{Stein 分解}
\label{spaces-more-morphisms-section-stein-factorization}

\noindent
Stein 分解是如下陈述：若固有态射 $f : X \to S$ 满足
$f_*\mathcal{O}_X = \mathcal{O}_S$，则它的纤维连通。

\begin{lemma}
\label{spaces-more-morphisms-lemma-stein-universally-closed}
设 $S$ 为概形。设 $f : X \to Y$ 为 $S$ 上代数空间之间的普遍闭
拟分离态射。存在分解
$$
\xymatrix{
X \ar[rr]_{f'} \ar[rd]_f & & Y' \ar[dl]^\pi \\
& Y &
}
$$
满足下列性质：
\begin{enumerate}
\item 态射 $f'$ 普遍闭、拟紧、拟分离且满；
\item 态射 $\pi : Y' \to Y$ 为整态射；
\item 有 $f'_*\mathcal{O}_X = \mathcal{O}_{Y'}$；
\item 有 $Y' = \underline{\Spec}_Y(f_*\mathcal{O}_X)$；
\item $Y'$ 是《空间的态射》定义
\ref{spaces-morphisms-definition-normalization-X-in-Y}
所定义的 $Y$ 在 $X$ 中的正规化。
\end{enumerate}
分解 $f = \pi \circ f'$ 的形成与平坦基变换可交换。
\end{lemma}

\begin{proof}
由《空间的态射》引理
\ref{spaces-morphisms-lemma-universally-closed-quasi-compact}，
态射 $f$ 拟紧。我们直接把 $Y'$ 定义为 $Y$ 在 $X$ 中的正规化，
故 (5) 与 (2) 自动成立。由《空间的态射》引理
\ref{spaces-morphisms-lemma-normalization-in-universally-closed}，
可知 (4) 成立。由《空间的态射》引理
\ref{spaces-morphisms-lemma-universally-closed-permanence}，
态射 $f'$ 普遍闭。由《空间的态射》引理
\ref{spaces-morphisms-lemma-quasi-compact-permanence}，它拟紧；
由《空间的态射》引理
\ref{spaces-morphisms-lemma-compose-after-separated}，它拟分离。

\medskip\noindent
为证明其余陈述，可以假设基底 $Y$ 仿射
（因为取正规化与平展局部化可交换）。设 $Y = \Spec(R)$。
则 $Y' = \Spec(A)$，其中
$A = \Gamma(X, \mathcal{O}_X)$ 是整 $R$-代数。
于是显然 $f'_*\mathcal{O}_X$ 就是 $\mathcal{O}_{Y'}$
（因为由《空间的态射》引理
\ref{spaces-morphisms-lemma-pushforward}，
$f'_*\mathcal{O}_X$ 拟凝聚，从而等于 $\widetilde{A}$）。这证明了 (3)。

\medskip\noindent
下面证明 $f'$ 满。由于 $f'$ 普遍闭（见上文），$f'$ 的像是闭子集
$V(I) \subset Y' = \Spec(A)$。取 $h \in I$。则
$h|_X = f^\sharp(h)$ 是 $X$ 的结构层的整体截面，并且在每一点消失。
% P09-ERR-0684
由于 $X$ 拟紧，这意味着 $h|_X$ 是幂零截面，即
$h^n|X = 0$ 对某个 $n > 0$ 成立。% P09-ERR-0668
但 $A = \Gamma(X, \mathcal{O}_X)$，故 $h^n = 0$。
换言之，$I$ 包含于 $A$ 的 Jacobson 根中；由此得到
$V(I) = Y'$，如所需。% P09-ERR-0669
\end{proof}

\begin{lemma}
\label{spaces-more-morphisms-lemma-stein-universally-closed-residue-fields}
在引理 \ref{spaces-more-morphisms-lemma-stein-universally-closed} 中，进一步假设
$f$ 局部有限型且 $Y$ 仿射。则对 $y \in Y$，纤维
$\pi^{-1}(\{y\}) = \{y_1, \ldots, y_n\}$ 有限，并且域扩张
$\kappa(y_i)/\kappa(y)$ 有限。% P09-ERR-0670
\end{lemma}

\begin{proof}
回忆，$\pi^{-1}(\{y\})$ 的各点之间不存在特殊化；见《代数》引理
\ref{algebra-lemma-integral-no-inclusion}。
由于 $f'$ 满，可知 $|X_y| \to \pi^{-1}(\{y\})$ 满。
注意，$X_y$ 是域上有限型的拟分离代数空间
（拟紧性已在所引引理的证明中证得）。因此 $|X_y|$ 是诺特拓扑空间
（《空间的态射》引理
\ref{spaces-morphisms-lemma-finite-presentation-noetherian}）。
现在一个拓扑论证（略）表明 $\pi^{-1}(\{y\})$ 有限。% P09-ERR-0671
对每个 $i$，可以选取有限型点 $x_i \in |X_y|$，它映到 $y_i$
（《空间的态射》引理
\ref{spaces-morphisms-lemma-enough-finite-type-points}）。
由此可知 $\kappa(y_i)/\kappa(y)$ 有限：$x_i$ 可由有限型态射
$\Spec(k_i) \to X_y$ 表示（根据有限型点的定义），因而
$\Spec(k_i) \to y = \Spec(\kappa(y))$ 是有限型态射
（作为有限型态射的复合），% P09-ERR-0672
所以 $k_i/\kappa(y)$ 有限（《态射》引理
\ref{morphisms-lemma-point-finite-type}）。
\end{proof}

\noindent
设 $f : X \to Y$ 为代数空间的态射，并设
$\overline{y} : \Spec(k) \to Y$ 为几何点。则 $f$ 在
$\overline{y}$ 上方的{\it 纤维}是代数空间
$X_{\overline{y}} = X \times_{Y, \overline{y}} \Spec(k)$，它定义在 $k$ 上。
若 $Y$ 是概形且 $y \in Y$ 是一点，则仍照常记
$X_y = X \times_Y \Spec(\kappa(y))$ 为该纤维。% P09-ERR-0673

\begin{lemma}
\label{spaces-more-morphisms-lemma-characterize-geometrically-connected-fibres}
设 $S$ 为概形。设 $f : X \to Y$ 为 $S$ 上代数空间的态射。
设 $\overline{y}$ 为 $Y$ 的几何点。则 $X_{\overline{y}}$ 连通，当且仅当：
对每个平展邻域 $(V, \overline{v}) \to (Y, \overline{y})$，其中
$V$ 为概形，基变换 $X_V \to V$ 的纤维 $X_v$ 都连通。
\end{lemma}

\begin{proof}
由于 $\overline{y}$ 的平展邻域所成范畴是余滤过的，并包含由概形组成的
余终子集（《空间的性质》引理
\ref{spaces-properties-lemma-cofinal-etale}），
可以用其中一个邻域替换 $Y$，从而假设 $Y$ 为概形。
设 $y \in Y$ 为对应于 $\overline{y}$ 的点。则 $X_y$ 在
$\kappa(y)$ 上几何连通，当且仅当 $X_{\overline{y}}$ 连通；这又当且仅当
$(X_y)_{k'}$ 对每个有限可分扩张 $k'$（它扩张 $\kappa(y)$）都连通。
见《域上空间》第
\ref{spaces-over-fields-section-geometrically-connected} 节，尤其是引理
\ref{spaces-over-fields-lemma-characterize-geometrically-disconnected}。
由《态射进阶》引理
\ref{more-morphisms-lemma-realize-prescribed-residue-field-extension-etale}，
存在仿射平展邻域 $(V, v) \to (Y, y)$，使得
$\kappa(s) \subset \kappa(u)$ 与 $\kappa(s) \subset k'$
对任意给定的有限可分扩张相识别。% P09-ERR-0674
引理得证。
\end{proof}

\begin{theorem}[Stein 分解；诺特情形]
\label{spaces-more-morphisms-theorem-stein-factorization-Noetherian}
设 $S$ 为概形。设 $f : X \to Y$ 为 $S$ 上代数空间的固有态射，
并且 $Y$ 局部诺特。存在分解
$$
\xymatrix{
X \ar[rr]_{f'} \ar[rd]_f & & Y' \ar[dl]^\pi \\
& Y &
}
$$
满足下列性质：
\begin{enumerate}
\item 态射 $f'$ 固有且具有连通的几何纤维；
\item 态射 $\pi : Y' \to Y$ 有限；
\item 有 $f'_*\mathcal{O}_X = \mathcal{O}_{Y'}$；
\item 有 $Y' = \underline{\Spec}_Y(f_*\mathcal{O}_X)$；
\item $Y'$ 是 $Y$ 在 $X$ 中的正规化；见《态射》定义
\ref{morphisms-definition-normalization-X-in-Y}。
\end{enumerate}
\end{theorem}

\begin{proof}
令 $f = \pi \circ f'$ 为引理
\ref{spaces-more-morphisms-lemma-stein-universally-closed} 给出的分解。注意，除引理
\ref{spaces-more-morphisms-lemma-stein-universally-closed} 的结论外，$f'$ 还分离
（《空间的态射》引理
\ref{spaces-morphisms-lemma-compose-after-separated}）且有限型
（《空间的态射》引理
\ref{spaces-morphisms-lemma-permanence-finite-type}）。
因此 $f'$ 固有。由《空间的上同调》引理
\ref{spaces-cohomology-lemma-proper-pushforward-coherent}，
可知 $f_*\mathcal{O}_X$ 是凝聚 $\mathcal{O}_Y$-模。
因此 $\pi$ 有限，即 (2) 成立。

\medskip\noindent
这证明了除最有意义的断言外的所有断言；该断言是 $f'$ 的几何纤维连通。
由上面的讨论显然可以把 $Y$ 替换为 $Y'$。于是 $Y$ 局部诺特，
$f : X \to Y$ 固有，并且 $f_*\mathcal{O}_X = \mathcal{O}_Y$。
设 $\overline{y}$ 为 $Y$ 的一个几何点。此时应用形式函数定理，更确切地说，
应用《空间的上同调》引理
\ref{spaces-cohomology-lemma-formal-functions-stalk}。它告诉我们
$$
\mathcal{O}^\wedge_{Y, \overline{y}} =
\lim_n H^0(X_n, \mathcal{O}_{X_n})
$$
其中 $X_n =
\Spec(\mathcal{O}_{Y, \overline{y}}/\mathfrak m_{\overline{y}}^n) \times_Y X$。
注意，$X_1 = X_{\overline{y}} \to X_n$ 是（有限阶）加厚，
因而 $X_n$ 的底拓扑空间等于 $X_{\overline{y}}$ 的底拓扑空间。
于是，若 $X_{\overline{y}} = T_1 \amalg T_2$ 是两个非空开闭子空间的
不交并，则同样有 $X_n = T_{1, n} \amalg T_{2, n}$，且这对所有
$n$ 成立。这又意味着
$H^0(X_n, \mathcal{O}_{X_n})$ 含有非平凡幂等元 $e_{1, n}$，
即恒取值 $1$ 于 $T_{1, n}$、且恒取值 $0$ 于 $T_{2, n}$ 的函数。
显然，$e_{1, n + 1}$ 限制为 $e_{1, n}$ 于 $X_n$ 上。
因此 $e_1 = \lim e_{1, n}$ 是该极限中的非平凡幂等元。
这与 $\mathcal{O}^\wedge_{Y, \overline{y}}$ 是局部环矛盾。
故假设错误，即 $X_{\overline{y}}$ 连通，如所需。
\end{proof}

\begin{theorem}[Stein 分解；一般情形]
\label{spaces-more-morphisms-theorem-stein-factorization-general}
设 $S$ 为概形。设 $f : X \to Y$ 为 $S$ 上代数空间的固有态射。
存在分解
$$
\xymatrix{
X \ar[rr]_{f'} \ar[rd]_f & & Y' \ar[dl]^\pi \\
& Y &
}
$$
满足下列性质：
\begin{enumerate}
\item 态射 $f'$ 固有且具有连通的几何纤维；
\item 态射 $\pi : Y' \to Y$ 为整态射；
\item 有 $f'_*\mathcal{O}_X = \mathcal{O}_{Y'}$；
\item 有 $Y' = \underline{\Spec}_Y(f_*\mathcal{O}_X)$；
\item $Y'$ 是 $Y$ 在 $X$ 中的正规化（《空间的态射》定义
\ref{spaces-morphisms-definition-normalization-X-in-Y}）。
\end{enumerate}
\end{theorem}

\begin{proof}
可以应用引理 \ref{spaces-more-morphisms-lemma-stein-universally-closed} 得到态射
$f' : X \to Y'$。注意，除引理
\ref{spaces-more-morphisms-lemma-stein-universally-closed} 的结论外，$f'$ 还分离
（《空间的态射》引理
\ref{spaces-morphisms-lemma-compose-after-separated}）且有限型
（《空间的态射》引理
\ref{spaces-morphisms-lemma-permanence-finite-type}）。
因此 $f'$ 固有。至此，除 $f'$ 具有连通几何纤维这一陈述外，
所有陈述都已得到证明。

\medskip\noindent
由讨论可知，显然可以把 $Y$ 替换为 $Y'$。于是
$f : X \to Y$ 固有且 $f_*\mathcal{O}_X = \mathcal{O}_Y$。
注意，这些条件在平坦基变换下保持不变
（《空间的态射》引理 \ref{spaces-morphisms-lemma-base-change-proper}
以及《空间的上同调》引理
\ref{spaces-cohomology-lemma-flat-base-change-cohomology}）。
设 $\overline{y}$ 为 $Y$ 的几何点。由引理
\ref{spaces-more-morphisms-lemma-characterize-geometrically-connected-fibres}
及刚才的评注，问题约化到如下情形：$Y$ 是概形，$y \in Y$ 是一点，
% P09-ERR-0675
$f : X \to Y$ 是 $Y$ 上的固有代数空间，且
$f_*\mathcal{O}_X = \mathcal{O}_Y$；需要证明纤维 $X_y$ 连通。
把 $Y$ 替换为 $y$ 的仿射邻域，可以假设 $Y = \Spec(R)$ 仿射。
于是 $f_*\mathcal{O}_X = \mathcal{O}_Y$ 意味着环同态
$R \to \Gamma(X, \mathcal{O}_X)$ 是双射。

\medskip\noindent
由《空间的极限》引理
\ref{spaces-limits-lemma-proper-limit-of-proper-finite-presentation-noetherian}，
可以写成 $(X \to Y) = \lim (X_i \to Y_i)$，其中
$X_i \to Y_i$ 固有且有限表示，$Y_i$ 诺特。对充分大的 $i$，
$Y_i$ 仿射（《空间的极限》引理
\ref{spaces-limits-lemma-limit-is-affine}）。设 $Y_i = \Spec(R_i)$，
并令 $R'_i = \Gamma(X_i, \mathcal{O}_{X_i})$。
注意，有环同态 $R_i \to R_i' \to R$：第一个来自 $X_i$ 是
$R_i$ 上的代数空间；第二个来自 $X \to X_i$ 以及
$R = \Gamma(X, \mathcal{O}_X)$。由《空间的极限》引理
\ref{spaces-limits-lemma-descend-section}，注意到 $R = \colim R'_i$。
于是
$$
\xymatrix{
X \ar[d]  \ar[r] & X_i \ar[d] \\
Y \ar[r] & Y'_i \ar[r] & Y_i
}
$$
交换，其中 $Y'_i = \Spec(R'_i)$。设 $y'_i \in Y'_i$ 为 $y$ 的像。
% P09-ERR-0676
有 $X_y = \lim X_{i, y'_i}$，因为 $X = \lim X_i$、
$Y = \lim Y'_i$ 且 $\kappa(y) = \colim \kappa(y'_i)$。现在设
$X_y = U \amalg V$，其中 $U$ 与 $V$ 都既开又闭。
则 $U, V$ 是开集 $U_i, V_i$ 的逆像，而这些开集位于 $X_{i, y'_i}$ 中
（《空间的极限》引理 \ref{spaces-limits-lemma-descend-opens}）。
由定理 \ref{spaces-more-morphisms-theorem-stein-factorization-Noetherian}，
$X_i \to Y'_i$ 的纤维连通，故 $U$ 或 $V$ 之一为空。
证明完毕。
\end{proof}

\noindent
这里给出一个应用。

\begin{lemma}
\label{spaces-more-morphisms-lemma-geometrically-connected-fibres-towards-normal}
设 $S$ 为概形。设 $f : X \to Y$ 为 $S$ 上代数空间的态射。假设
\begin{enumerate}
\item $f$ 固有；
\item $Y$ 是整的（《域上空间》定义
\ref{spaces-over-fields-definition-integral-algebraic-space}），
其一般点为 $\xi$；
\item $Y$ 正规；
\item $X$ 约化；
\item $|X|$ 的每个不可约分支的一般点都映到 $\xi$；
\item 有 $H^0(X_\xi, \mathcal{O}) = \kappa(\xi)$。
\end{enumerate}
则 $f_*\mathcal{O}_X = \mathcal{O}_Y$，并且 $f$ 具有几何连通纤维。
\end{lemma}

\begin{proof}
应用定理 \ref{spaces-more-morphisms-theorem-stein-factorization-general} 得到分解
$X \to Y' \to Y$。只须证明 $Y' = Y$。
只须证明 $Y' \times_Y V \to V$ 是同构，其中 $V \to Y$ 为平展态射，
且 $V$ 是仿射整概形；见《域上空间》引理
\ref{spaces-over-fields-lemma-normal-integral-cover-by-affines}。
$Y'$ 的形成与平展基变换可交换；见《空间的态射》引理
\ref{spaces-morphisms-lemma-properties-normalization}。
$X \times_Y V$ 的一般点位于 $X$ 的一般点上方
（《良态空间》引理 \ref{decent-spaces-lemma-decent-generic-points}），
故由假设 (5)，它们映到 $V$ 的一般点。此外，条件 (6) 在由
$V \to Y$ 所作的基变换下保持不变，例如这可由平坦基变换得出
（《空间的上同调》引理
\ref{spaces-cohomology-lemma-flat-base-change-cohomology}）。
因此只须在 $Y$ 是正规整仿射概形的情形证明引理。

\medskip\noindent
假设 $Y$ 是正规整仿射概形。我们将应用《态射》引理
\ref{morphisms-lemma-finite-birational-over-normal}
证明 $Y' \to Y$ 是同构。事实上，$Y'$ 约化，因为 $X$ 约化
（《空间的态射》引理
\ref{spaces-morphisms-lemma-normalization-in-reduced}）。
由上面所引定理，态射 $Y' \to Y$ 为整态射。
由于 $Y$ 良态且 $X \to Y$ 分离，可知 $X$ 也良态；为此可用
《良态空间》引理
\ref{decent-spaces-lemma-properties-trivial-implications} 和
\ref{decent-spaces-lemma-property-over-property}。由假设 (5)、
《空间的态射》引理 \ref{spaces-morphisms-lemma-normalization-generic}
以及《良态空间》引理
\ref{decent-spaces-lemma-decent-generic-points}，可知
$|Y'|$ 的每个不可约分支的一般点都映到 $\xi$。另一方面，
由于 $Y'$ 是 $f_*\mathcal{O}_X$ 的相对谱，可知概形论纤维
$Y'_\xi$ 是 $H^0(X_\xi, \mathcal{O})$ 的谱；由假设，它等于
$\kappa(\xi)$。故 $Y'$ 是整概形，并且其函数域等于 $Y$ 的函数域。
证明完毕。
\end{proof}

\noindent
这里再给出一个应用。

\begin{lemma}
\label{spaces-more-morphisms-lemma-proper-flat-nr-geom-conn-comps-lower-semicontinuous}
设 $S$ 为概形。设 $X \to Y$ 为 $S$ 上代数空间的态射。
若 $f$ 固有、平坦且有限表示，则函数
% P09-ERR-0677
$n_{X/Y} : |Y| \to \mathbf{Z}$ 是下半连续的；该函数计数 $f$ 的纤维的
几何连通分支数（引理
\ref{spaces-more-morphisms-lemma-base-change-fibres-nr-geometrically-connected-components}）。
\end{lemma}

\begin{proof}
问题在 $Y$ 上是平展局部的，故可以且确实假设 $Y$ 是仿射概形。
设 $y \in Y$。令 $n = n_{X/S}(y)$。% P09-ERR-0678
注意 $n < \infty$，因为 $X \to Y$ 在 $y$ 处的几何纤维是域上的
固有代数空间，故诺特，因而只有有限个连通分支。
需要找出开邻域 $V$，它含 $y$，使得 $n_{X/S}|_V \geq n$。
% P09-ERR-0679
令 $X \to Y' \to Y$ 为定理
\ref{spaces-more-morphisms-theorem-stein-factorization-general} 中的 Stein 分解。
由引理 \ref{spaces-more-morphisms-lemma-stein-universally-closed-residue-fields}，
存在有限多个点 $y'_1, \ldots, y'_m \in Y'$，它们位于 $y$ 上方，
且扩张
$\kappa(y'_i)/\kappa(y)$ 有限。
《态射进阶》引理
\ref{more-morphisms-lemma-etale-makes-integral-split}
告诉我们，在把 $Y$ 替换为 $y$ 的一个平展邻域后，可以假设
$Y' = V_1 \amalg \ldots \amalg V_m$ 是概形，且
$y'_i \in V_i$ 以及 $\kappa(y'_i)/\kappa(y)$ 纯不可分。
于是代数空间 $X_{y_i'}$ 在 $\kappa(y)$ 上几何连通，故 $m = n$。
代数空间 $X_i = (f')^{-1}(V_i)$（$i = 1, \ldots, n$）
在 $Y$ 上平坦且有限表示。因此 $X_i \to Y$ 的像是开集
（《空间的态射》引理 \ref{spaces-morphisms-lemma-fppf-open}）。
故在 $y$ 的一个邻域内，$n_{X/Y}$ 至少为 $n$。
\end{proof}

\begin{lemma}
\label{spaces-more-morphisms-lemma-proper-flat-geom-red}
设 $S$ 为概形。设 $f : X \to Y$ 为 $S$ 上代数空间的态射。假设
\begin{enumerate}
\item $f$ 固有、平坦且有限表示；
\item $f$ 的几何纤维约化。
\end{enumerate}
则函数 $n_{X/S} : |Y| \to \mathbf{Z}$ 是局部常值的；该函数计数
% P09-ERR-0680
$f$ 的纤维的几何连通分支数
（引理 \ref{spaces-more-morphisms-lemma-base-change-fibres-nr-geometrically-connected-components}）。
\end{lemma}

\begin{proof}
由引理 \ref{spaces-more-morphisms-lemma-proper-flat-nr-geom-conn-comps-lower-semicontinuous}，
函数 $n_{X/Y}$ 下半连续。% P09-ERR-0681
故只须证明它上半连续。为此，可以在 $Y$ 上平展局部地工作，
因而可以假设 $Y$ 是仿射概形。对 $y \in Y$，考虑
$\kappa(y)$-代数
$$
A = H^0(X_y, \mathcal{O}_{X_y})
$$
由《域上空间》引理
\ref{spaces-over-fields-lemma-proper-geometrically-reduced-global-sections}
以及 $X_y$ 几何约化这一事实，$A$ 是 $\kappa(y)$ 的有限可分扩张的
有限乘积。因此 $A \otimes_{\kappa(y)} \kappa(\overline{y})$ 是由
$\beta_0(y) = \dim_{\kappa(y)} A$ 个 $\kappa(\overline{y})$ 副本组成的积。
故 $X_{\overline{y}}$ 有 $\beta_0(y)$ 个连通分支。
换言之，有 $n_{X/S} = \beta_0$，这是 $Y$ 上函数的等式。
% P09-ERR-0682
因此，由《空间的导出范畴》引理
\ref{spaces-perfect-lemma-jump-loci-geometric}，$n_{X/Y}$ 上半连续。
证明完毕。
\end{proof}

\begin{lemma}
\label{spaces-more-morphisms-lemma-stein-factorization-etale}
设 $S$ 为概形。设 $f : X \to Y$ 为 $S$ 上代数空间的固有态射。
设 $X \to Y' \to Y$ 为 $f$ 的 Stein 分解
（定理 \ref{spaces-more-morphisms-theorem-stein-factorization-general}）。
若 $f$ 有限表示、平坦且具有几何约化纤维
（定义 \ref{spaces-more-morphisms-definition-geometrically-reduced-fibre}），
则 $Y' \to Y$ 有限平展。
\end{lemma}

\begin{proof}
Stein 分解的形成与平坦基变换可交换；见引理
\ref{spaces-more-morphisms-lemma-stein-universally-closed}。因此可以在 $Y$ 上平展局部地工作，
并可假设 $Y$ 是仿射概形。此时 $Y'$ 是仿射概形，且
$Y' \to Y$ 为整态射。

\medskip\noindent
设 $y \in Y$。令 $n$ 为几何纤维 $X_{\overline{y}}$ 的连通分支数。
注意 $n < \infty$，因为 $X \to Y$ 在 $y$ 处的几何纤维是域上的
固有代数空间，故诺特，因而只有有限个连通分支。
由引理 \ref{spaces-more-morphisms-lemma-stein-universally-closed-residue-fields}，
存在有限多个点 $y'_1, \ldots, y'_m \in Y'$，它们位于 $y$ 上方；
并且对每个 $i$，可以选取有限型点
$x_i \in |X_y|$ 映到 $y'_i$，且扩张
$\kappa(y'_i)/\kappa(y)$ 有限。% P09-ERR-0683
因此《态射进阶》引理
\ref{more-morphisms-lemma-etale-makes-integral-split}
告诉我们，在把 $Y$ 替换为 $y$ 的一个平展邻域后，可以假设
$Y' = V_1 \amalg \ldots \amalg V_m$ 是概形，且
$y'_i \in V_i$ 以及 $\kappa(y'_i)/\kappa(y)$ 纯不可分。
此时，代数空间 $X_{y_i'}$ 在 $\kappa(y)$ 上几何连通，故 $m = n$。
代数空间 $X_i = (f')^{-1}(V_i)$（$i = 1, \ldots, n$）
在 $Y$ 上固有、平坦、有限表示，且具有几何约化纤维。
只须对每个态射 $X_i \to Y$ 证明引理。这把我们约化到
$X_{\overline{y}}$ 连通的情形。

\medskip\noindent
假设 $X_{\overline{y}}$ 连通。由引理
\ref{spaces-more-morphisms-lemma-proper-flat-geom-red}，可知 $X \to Y$ 在 $y$ 的一个邻域内
具有几何连通纤维。因此可以假设 $X \to Y$ 的纤维都几何连通。
于是由《空间的导出范畴》引理
\ref{spaces-perfect-lemma-proper-flat-geom-red-connected}，有
$f_*\mathcal{O}_X = \mathcal{O}_Y$，这就完成了证明。
\end{proof}

\noindent
下面引理的证明使用了概形的 Stein 分解，这就是它出现在本节的原因。

\begin{lemma}
\label{spaces-more-morphisms-lemma-split-off-proper-part-henselian}
设 $(A, I)$ 为亨泽尔对。设 $X$ 为 $A$ 上分离有限型的代数空间。令
$X_0 = X \times_{\Spec(A)} \Spec(A/I)$。
设 $Y \subset X_0$ 为既开又闭子空间，并且
$Y \to \Spec(A/I)$ 固有。则存在既开又闭子空间 $W \subset X$，
它在 $A$ 上固有，且
$W \times_{\Spec(A)} \Spec(A/I) = Y$。
\end{lemma}

\begin{proof}
我们用 $T \mapsto T_0$ 表示由
$\Spec(A/I) \to \Spec(A)$ 所作的基变换。
由 Chow 引理的一个弱版本（其形式见《空间的上同调》引理
\ref{spaces-cohomology-lemma-weak-chow}），存在满固有态射
$\varphi : X' \to X$，使得 $X'$ 可浸入 $\mathbf{P}^n_A$。
令 $Y' = \varphi^{-1}(Y)$。这是 $X'_0$ 的开闭子概形。
由《态射进阶》引理
\ref{more-morphisms-lemma-split-off-proper-part-henselian}，
引理对 $(X', Y')$ 成立。设 $W' \subset X'$ 为 $A$ 上固有的开闭子概形，
使得 $Y' = W'_0$。由《空间的态射》引理
\ref{spaces-morphisms-lemma-universally-closed-permanence}，
$Q_1 = \varphi(|W'|) \subset |X|$ 和
$Q_2 = \varphi(|X' \setminus W'|) \subset |X|$
是闭子集；又由《空间的态射》引理
\ref{spaces-morphisms-lemma-image-proper-is-proper}，
$Q_1$ 上的任意闭子空间结构在 $A$ 上固有。
$Q_1 \cap Q_2$ 在 $\Spec(A)$ 中的像是闭集。
由于 $(A, I)$ 是亨泽尔对，若 $Q_1 \cap Q_2$ 非空，则可知
$Q_1 \cap Q_2$ 有一点位于 $\Spec(A/I)$ 上方。
这是不可能的，因为 $W'_0 = Y' = \varphi^{-1}(Y)$。
由此得到 $Q_1$ 在 $|X|$ 中既开又闭。
令 $W \subset X$ 为对应的既开又闭子空间。则 $W$ 在 $A$ 上固有，
且 $W_0 = Y$。
\end{proof}










\section{从开子空间延拓性质}
\label{spaces-more-morphisms-section-extending-properties}

\noindent
本节汇集若干如下形式的结果：若 $f : X \to Y$ 是代数空间的平坦态射，
并且 $f$ 在 $Y$ 的一个稠密开子空间上满足某个性质，则 $f$ 在整个
$Y$ 上满足同一性质。

\begin{lemma}
\label{spaces-more-morphisms-lemma-flat-finite-type-finitely-presented-over-dense-open}
设 $S$ 为概形。设 $f : X \to Y$ 为 $S$ 上代数空间的态射。
设 $\mathcal{F}$ 为拟凝聚 $\mathcal{O}_X$-模。
设 $V \subset Y$ 为开子空间。假设
\begin{enumerate}
\item $f$ 局部有限表示；
\item $\mathcal{F}$ 有限型且在 $Y$ 上平坦；
\item $V \to Y$ 拟紧且概形论稠密；
\item $\mathcal{F}|_{f^{-1}V}$ 有限表示。
\end{enumerate}
则 $\mathcal{F}$ 有限表示。
\end{lemma}

\begin{proof}
只须证明 $\mathcal{F}$ 到某个概形上的拉回有限表示，该概形满且平展地
映到 $X$。故可假设 $X$ 为概形。% P09-ERR-0685
类似地，可以用一个概形替换 $Y$，该概形满且平展地映到 $Y$
（$V$ 在此概形中的逆像概形论稠密；见《空间的态射》第
\ref{spaces-morphisms-section-scheme-theoretic-closure} 节）。
% P09-ERR-0686
于是问题约化到概形情形，即《平坦性进阶》引理
\ref{flat-lemma-flat-finite-type-finitely-presented-over-dense-open}。
\end{proof}

\begin{lemma}
\label{spaces-more-morphisms-lemma-flat-finite-type-finitely-presented-over-dense-open-X}
设 $S$ 为概形。设 $f : X \to Y$ 为 $S$ 上代数空间的态射。
设 $V \subset Y$ 为开子空间。假设
\begin{enumerate}
\item $f$ 局部有限型且平坦；
\item $V \to Y$ 拟紧且概形论稠密；
\item $f|_{f^{-1}V} : f^{-1}V \to V$ 局部有限表示。
\end{enumerate}
则 $f$ 局部有限表示。% P09-ERR-0687
\end{lemma}

\begin{proof}
证明与引理
\ref{spaces-more-morphisms-lemma-flat-finite-type-finitely-presented-over-dense-open}
的证明相同，只须改用《平坦性进阶》引理
\ref{flat-lemma-flat-finite-type-finitely-presented-over-dense-open-X}。
\end{proof}

\begin{lemma}
\label{spaces-more-morphisms-lemma-flat-fp-dimension-over-dense-open}
设 $S$ 为概形。设 $f : X \to Y$ 为 $S$ 上代数空间的态射，
它平坦且局部有限型。设 $V \subset Y$ 为开子空间，满足
$|V| \subset |Y|$ 稠密且 $X_V \to V$ 的相对维数 $\leq d$。
若还满足下列条件之一：
\begin{enumerate}
\item $f$ 局部有限表示；
\item $V \to Y$ 拟紧，
\end{enumerate}
则 $f : X \to Y$ 的相对维数 $\leq d$。
\end{lemma}

\begin{proof}
可以把 $Y$ 替换为其约化，因而可以假设 $Y$ 约化。
于是 $V$ 在 $Y$ 中概形论稠密；见《空间的态射》引理
\ref{spaces-morphisms-lemma-quasi-compact-immersion}。
根据定义，相对维数 $\leq d$ 这一性质可以在平展覆盖上检验；见
《空间的态射》第 \ref{spaces-morphisms-section-relative-dimension} 节。
% P09-ERR-0688
因此，只须证明 $f$ 的相对维数 $\leq d$；为此可把 $X$ 替换为
满且平展地映到 $X$ 的一个概形。类似地，可以用一个概形替换 $Y$，
该概形满且平展地映到 $Y$。% P09-ERR-0689
$V$ 在此概形中的逆像概形论稠密；见《空间的态射》第
\ref{spaces-morphisms-section-scheme-theoretic-closure} 节。
由于概形的概形论稠密开子概形尤其稠密，问题约化到概形情形，即
《平坦性进阶》引理
\ref{flat-lemma-flat-finite-presentation-dimension-over-dense-open}。
\end{proof}

\begin{lemma}
\label{spaces-more-morphisms-lemma-proper-flat-finite-over-dense-open}
设 $S$ 为概形。设 $f : X \to Y$ 为 $S$ 上代数空间的态射，
它平坦且固有。设 $V \to Y$ 为开子空间，满足
$|V| \subset |Y|$ 稠密且 $X_V \to V$ 有限。% P09-ERR-0690
若还满足 $f$ 局部有限表示或 $V \to Y$ 拟紧之一，则 $f$ 有限。
% P09-ERR-0691
\end{lemma}

\begin{proof}
由引理 \ref{spaces-more-morphisms-lemma-flat-fp-dimension-over-dense-open}，$f$ 的纤维维数为零。
由《空间的态射》引理
\ref{spaces-morphisms-lemma-locally-quasi-finite-rel-dimension-0}，
这蕴含 $f$ 局部拟有限。由《空间的态射》引理
\ref{spaces-morphisms-lemma-locally-quasi-finite-separated-representable}，
这蕴含 $f$ 可表。可以检验 $f$ 是否有限，而这一检验在 $Y$ 上平展局部；
故可假设 $Y$ 是概形。由于 $f$ 可表，问题约化到概形情形，即
《平坦性进阶》引理 \ref{flat-lemma-proper-flat-finite-over-dense-open}。
\end{proof}

\begin{lemma}
\label{spaces-more-morphisms-lemma-zariski}
设 $S$ 为概形。设 $f : X \to Y$ 为 $S$ 上代数空间的态射。
设 $V \subset Y$ 为开子空间。若
\begin{enumerate}
\item $f$ 分离、局部有限型且平坦；
\item $f^{-1}(V) \to V$ 是同构；
\item $V \to Y$ 拟紧且概形论稠密，
\end{enumerate}
则 $f$ 是开浸入。
\end{lemma}

\begin{proof}
应用引理
\ref{spaces-more-morphisms-lemma-flat-finite-type-finitely-presented-over-dense-open-X}，
可知 $f$ 局部有限表示。应用引理
\ref{spaces-more-morphisms-lemma-flat-fp-dimension-over-dense-open}，可知 $f$ 的相对维数
$\leq 0$。由《空间的态射》引理
\ref{spaces-morphisms-lemma-locally-quasi-finite-rel-dimension-0}，
这蕴含 $f$ 局部拟有限。由《空间的态射》引理
\ref{spaces-morphisms-lemma-locally-quasi-finite-separated-representable}，
这蕴含 $f$ 可表。由《空间上的下降》引理
\ref{spaces-descent-lemma-descending-property-open-immersion}，
可以检验 $f$ 是否为开浸入，而这一检验在 $Y$ 上平展局部。
故可假设 $Y$ 是概形。由于 $f$ 可表，问题约化到概形情形，即
《平坦性进阶》引理 \ref{flat-lemma-zariski}。
\end{proof}

\section{爆破与平坦性}
\label{spaces-more-morphisms-section-blowup-flat}

\noindent
我们不再重复《平坦性进阶》第 \ref{flat-section-blowup-flat} 节中的工作，
而是证明《平坦性进阶》引理 \ref{flat-lemma-push-ideal} 的一个类似结果；
它说明，寻找合适爆破的问题通常在基底上是平展局部的。

\begin{lemma}
\label{spaces-more-morphisms-lemma-push-ideal}
设 $S$ 为概形。设 $X$ 为 $S$ 上拟紧且拟分离的代数空间。
设 $\varphi : W \to X$ 为拟紧、分离的平展态射。
设 $U \subset X$ 为拟紧开子空间。
设 $\mathcal{I} \subset \mathcal{O}_W$ 为有限型拟凝聚理想层，满足
$V(\mathcal{I}) \cap \varphi^{-1}(U) = \emptyset$。
则存在有限型拟凝聚理想层
$\mathcal{J} \subset \mathcal{O}_X$，使得
\begin{enumerate}
\item $V(\mathcal{J}) \cap U = \emptyset$；
\item $\varphi^{-1}(\mathcal{J})\mathcal{O}_W = \mathcal{I} \mathcal{I}'$，
其中 $\mathcal{I}' \subset \mathcal{O}_W$ 是某个有限型拟凝聚理想。
\end{enumerate}
\end{lemma}

\begin{proof}
选择分解 $W \to Y \to X$，其中 $j : W \to Y$ 是拟紧开浸入，
$\pi : Y \to X$ 是有限表示的有限态射
（引理 \ref{spaces-more-morphisms-lemma-quasi-finite-separated-pass-through-finite-addendum}）。
令 $V = j(W) \cup \pi^{-1}(U) \subset Y$。注意，
$\mathcal{I}$ 在 $W \cong j(W)$ 上，它与另一开子空间上的
$\mathcal{O}_{\pi^{-1}(U)}$ 可粘合成有限型拟凝聚理想层
$\mathcal{I}_1 \subset \mathcal{O}_V$。由《空间的极限》引理
\ref{spaces-limits-lemma-extend}，存在有限型拟凝聚理想层
$\mathcal{I}_2 \subset \mathcal{O}_Y$，使得
$\mathcal{I}_2|_V = \mathcal{I}_1$。换言之，
$\mathcal{I}_2 \subset \mathcal{O}_Y$ 是有限型拟凝聚理想层，满足
$V(\mathcal{I}_2)$ 与 $\pi^{-1}(U)$ 不交，并且
$j^{-1}\mathcal{I}_2 = \mathcal{I}$。以 $i : Z \to Y$
表示相应的闭浸入；它是有限表示的（《空间的态射》引理
\ref{spaces-morphisms-lemma-closed-immersion-finite-presentation}）。
% P09-ERR-0692
特别地，复合 $\tau = \pi \circ i : Z \to X$ 是有限且有限表示的
（《空间的态射》引理
\ref{spaces-morphisms-lemma-composition-finite-presentation} 和
\ref{spaces-morphisms-lemma-composition-integral}）。

\medskip\noindent
令 $\mathcal{F} = \tau_*\mathcal{O}_Z$，并把它视为拟凝聚
$\mathcal{O}_X$-模。由《空间上的下降》引理
\ref{spaces-descent-lemma-finite-finitely-presented-module}，
可知 $\mathcal{F}$ 是有限表示的 $\mathcal{O}_X$-模。
令 $\mathcal{J} = \text{Fit}_0(\mathcal{F})$。
（此处插入带环拓扑斯上模的 Fitting 理想的参考文献。）
% P09-ERR-0693
这是 $X$ 上的有限型拟凝聚理想层（因为 $\mathcal{F}$ 有限表示；见
《代数进阶》引理 \ref{more-algebra-lemma-fitting-ideal-basics}）。
$|\tau|(|Z|) \cap |U| = \emptyset$；这是由
$\mathcal{I}_2$ 的选取而得。而零模的第 $0$ 个 Fitting 理想
等于结构层，故引理的第 (1) 部分成立。为证明 (2)，注意
$\varphi^{-1}(\mathcal{J})\mathcal{O}_W = \text{Fit}_0(\varphi^*\mathcal{F})$，
因为取 Fitting 理想与基变换可交换。另一方面，由于
$\varphi : W \to X$ 分离且平展，可知
$(1, j) : W \to W \times_X Y$ 是开闭浸入。
故有 $W \times_Y Z = V(\mathcal{I}) \amalg Z'$，其中
$\tau' : Z' \to W$ 是某个有限且有限表示的代数空间态射。
% P09-ERR-0694
因此
\begin{align*}
\text{Fit}_0(\varphi^*\mathcal{F}) & =
\text{Fit}_0((W \times_Y Z \to W)_*\mathcal{O}_{W \times_Y Z}) \\
& =
\text{Fit}_0(\mathcal{O}_W/\mathcal{I}) \cdot
\text{Fit}_0(\tau'_*\mathcal{O}_{Z'}) \\
& =
\mathcal{I} \cdot \text{Fit}_0(\tau'_*\mathcal{O}_{Z'})
\end{align*}
其中第二个等式由《代数进阶》引理
\ref{more-algebra-lemma-fitting-ideal-basics}
在带环拓扑斯上的层中转述而得。% P09-ERR-0695
令 $\mathcal{I}' = \text{Fit}_0(\tau'_*\mathcal{O}_{Z'})$，
即完成引理的证明。
\end{proof}

\begin{theorem}
\label{spaces-more-morphisms-theorem-flatten-module}
设 $S$ 为概形。设 $B$ 为 $S$ 上拟紧且拟分离的代数空间。
设 $X$ 为 $B$ 上的代数空间。设 $\mathcal{F}$ 为 $X$ 上的拟凝聚模。
设 $U \subset B$ 为拟紧开子空间。假设
\begin{enumerate}
\item $X$ 拟紧；
\item $X$ 在 $B$ 上局部有限表示；
\item $\mathcal{F}$ 是有限型模；
\item $\mathcal{F}_U$ 有限表示；
\item $\mathcal{F}_U$ 在 $U$ 上平坦。
\end{enumerate}
则存在一个 $U$-容许爆破 $B' \to B$，使得严格变换
$\mathcal{F}'$（即 $\mathcal{F}$ 的严格变换）是有限表示的
$\mathcal{O}_{X \times_B B'}$-模，并且在 $B'$ 上平坦。
\end{theorem}

\begin{proof}
选择仿射概形 $V$ 以及满平展态射 $V \to X$。
由于严格变换与平展局部化可交换
（《空间上的除子》引理 \ref{spaces-divisors-lemma-strict-transform-local}），
只须证明把 $X$ 替换为 $V$ 后所得的结论。故可假设
$X \to B$ 可表（此外仍满足引理的各项假设）。% P09-ERR-0696

\medskip\noindent
假设 $X \to B$ 可表。选择仿射概形 $W$ 以及满平展态射
$\varphi : W \to B$。注意 $X \times_B W$ 是概形。
由概形情形（《平坦性进阶》定理 \ref{flat-theorem-flatten-module}），
可以找到有限型拟凝聚理想层
$\mathcal{I} \subset \mathcal{O}_W$，使得
(a) $|V(\mathcal{I})| \cap |\varphi^{-1}(U)| = \emptyset$，并且 (b)
$\mathcal{F}|_{X \times_B W}$ 关于爆破 $W' \to W$
（其中心为 $\mathcal{I}$）的严格变换在 $W'$ 上平坦且是有限表示模。
按引理 \ref{spaces-more-morphisms-lemma-push-ideal} 选择有限型理想层
$\mathcal{J} \subset \mathcal{O}_B$。令 $B' \to B$ 为
$\mathcal{J}$ 的爆破。我们断言这一爆破满足要求。事实上，显然
$B' \to B$ 是 $U$-容许的；这是由 $\mathcal{J}$ 的选取而得。
此外，基变换 $B' \times_B W \to W$ 是 $W$ 沿
$\varphi^{-1}\mathcal{J} = \mathcal{I}\mathcal{I}'$ 的爆破
（爆破与平坦基变换的相容性；见《空间上的除子》引理
\ref{spaces-divisors-lemma-flat-base-change-blowing-up}）。
因而存在分解
$$
W \times_B B' \to W' \to W
$$
其中第一个态射也是爆破；见《空间上的除子》引理
\ref{spaces-divisors-lemma-blowing-up-two-ideals}。
$\mathcal{F}'$（它位于 $B' \times_B X$ 上）到
$W \times_B B' \times_B X$ 的限制，是
$\mathcal{F}|_{X \times_B W}$ 的严格变换
（《空间上的除子》引理 \ref{spaces-divisors-lemma-strict-transform-local}），
因而也是 $\mathcal{F}|_{X \times_B W}$ 经上述两个爆破所得的
两次迭代严格变换（《空间上的除子》引理
\ref{spaces-divisors-lemma-strict-transform-composition-blowups}）。
第一次爆破后，所得层按构造已经在基底上平坦且有限表示。
因此第二次严格变换后仍然如此（因为由《空间上的除子》引理
\ref{spaces-divisors-lemma-strict-transform-flat}，该变换是拉回）。
于是可知，$\mathcal{F}'$ 到 $B' \times_B X$ 的一个平展覆盖上的限制
具有所需性质，定理得证。
\end{proof}

\section{应用}
\label{spaces-more-morphisms-section-applications-flattening-by-blowing-up}

\noindent
本节应用通过爆破实现平坦化的结果。

\begin{lemma}
\label{spaces-more-morphisms-lemma-flat-after-blowing-up}
设 $S$ 为概形。设 $f : X \to B$ 为 $S$ 上代数空间的态射。
设 $U \subset B$ 为开子空间。假设
\begin{enumerate}
\item $B$ 拟紧且拟分离；
\item $U$ 拟紧；
\item $f : X \to B$ 有限型且拟分离；
\item $f^{-1}(U) \to U$ 平坦且局部有限表示。
\end{enumerate}
则存在一个 $U$-容许爆破 $B' \to B$，使得严格变换
$X'$（即 $X$ 的严格变换）在 $B'$ 上平坦且有限表示。
\end{lemma}

\begin{proof}
设 $B' \to B$ 为一个 $U$-容许爆破。注意，$X$ 的严格变换在
$B'$ 上拟紧且拟分离，因为 $X$ 在 $B$ 上拟紧且拟分离。
故只须寻找一个 $U$-容许爆破，使得严格变换平坦且局部有限表示。
不能直接应用定理 \ref{spaces-more-morphisms-theorem-flatten-module}，因为 $X$ 在 $B$
上不局部有限表示。

\medskip\noindent
选择仿射概形 $V$ 以及满平展态射 $V \to X$。
（这是可能的，因为 $X$ 作为拟紧空间 $B$ 上的有限型空间而拟紧。）
于是只须对态射 $V \to B$ 证明结论（因为严格变换与平展局部化可交换；
见《空间上的除子》引理
\ref{spaces-divisors-lemma-strict-transform-local}）。
故可假设 $X \to B$ 除有限型外还分离。
在此情形下，可以找到闭浸入 $i : X \to Y$，其中 $Y \to B$
分离且有限表示；见《空间的极限》命题
\ref{spaces-limits-proposition-separated-closed-in-finite-presentation}。

\medskip\noindent
对 $\mathcal{F} = i_*\mathcal{O}_X$（视为 $Y/B$ 上的模）应用定理
\ref{spaces-more-morphisms-theorem-flatten-module}。得到一个 $U$-容许爆破
$B' \to B$，使得 $\mathcal{F}$ 的严格变换在 $B'$ 上平坦且有限表示。
设 $X'$ 为 $X$ 在爆破 $B' \to B$ 下的严格变换。
设 $i' : X' \to Y \times_B B'$ 为诱导态射。
由于取严格变换与沿仿射态射的推前可交换（《空间上的除子》引理
\ref{spaces-divisors-lemma-strict-transform-affine}），
可知 $i'_*\mathcal{O}_{X'}$ 在 $B'$ 上平坦，并且作为
$\mathcal{O}_{Y \times_B B'}$-模有限表示。
因此 $X' \to B'$ 平坦且局部有限表示。
结合前述讨论即得引理。
\end{proof}

\begin{lemma}
\label{spaces-more-morphisms-lemma-finite-after-blowing-up}
设 $S$ 为概形。设 $f : X \to B$ 为 $S$ 上代数空间的态射。
设 $U \subset B$ 为开子空间。假设
\begin{enumerate}
\item $B$ 拟紧且拟分离；
\item $U$ 拟紧；
\item $f : X \to B$ 固有；
\item $f^{-1}(U) \to U$ 有限局部自由。
\end{enumerate}
则存在一个 $U$-容许爆破 $B' \to B$，使得严格变换
$X'$（即 $X$ 的严格变换）在 $B'$ 上有限局部自由。
\end{lemma}

\begin{proof}
由引理 \ref{spaces-more-morphisms-lemma-flat-after-blowing-up}，可假设
$X \to B$ 平坦且有限表示。必要时把 $B$ 替换为一个
$U$-容许爆破，可以假设 $U \subset B$ 概形论稠密。
于是由引理 \ref{spaces-more-morphisms-lemma-proper-flat-finite-over-dense-open}，$f$ 有限。
故由《空间的态射》引理
\ref{spaces-morphisms-lemma-finite-flat}，$f$ 有限局部自由。
\end{proof}

\begin{lemma}
\label{spaces-more-morphisms-lemma-isomorphism-after-blowing-up}
设 $S$ 为概形。设 $f : X \to B$ 为 $S$ 上代数空间的态射。
设 $U \subset B$ 为开子空间。假设
\begin{enumerate}
\item $B$ 拟紧且拟分离；
\item $U$ 拟紧；
\item $f : X \to B$ 固有；
\item $f^{-1}(U) \to U$ 是同构。
\end{enumerate}
则存在一个 $U$-容许爆破 $B' \to B$，使得严格变换
$X'$（即 $X$ 的严格变换）同构地映到 $B'$。
\end{lemma}

\begin{proof}
由引理 \ref{spaces-more-morphisms-lemma-flat-after-blowing-up}，可假设
$X \to B$ 平坦且有限表示。必要时把 $B$ 替换为一个
$U$-容许爆破，可以假设 $U \subset B$ 概形论稠密。
于是由引理 \ref{spaces-more-morphisms-lemma-proper-flat-finite-over-dense-open}，$f$ 有限；
由引理 \ref{spaces-more-morphisms-lemma-zariski}，它是开浸入。
故 $f$ 是像为闭集且包含稠密开子空间 $U$ 的开浸入，
从而 $f$ 是同构。
\end{proof}

\begin{lemma}
\label{spaces-more-morphisms-lemma-zariski-after-blowup}
设 $S$ 为概形。设 $f : X \to B$ 为 $S$ 上代数空间的态射。
设 $U \subset B$ 为开子空间。假设
\begin{enumerate}
\item $B$ 拟紧且拟分离；% P09-ERR-0698
\item $U$ 拟紧；
\item $f$ 有限型；
\item $f^{-1}(U) \to U$ 是同构。
\end{enumerate}
则存在一个 $U$-容许爆破 $B' \to B$，使 $U$ 在 $B'$ 中概形论稠密，
并且严格变换 $X'$（即 $X$ 的严格变换）同构地映到 $B'$ 的一个开子空间。
\end{lemma}

\begin{proof}
本引理是引理 \ref{spaces-more-morphisms-lemma-isomorphism-after-blowing-up} 的推广。
由于 $U$-容许爆破的复合仍是 $U$-容许的（《空间上的除子》引理
\ref{spaces-divisors-lemma-composition-admissible-blowups}），
可以分阶段进行。选取有限型拟凝聚理想层
$\mathcal{I} \subset \mathcal{O}_B$，使得
$|B| \setminus |U| = |V(\mathcal{I})|$。把 $B$ 替换为 $B$ 沿
$\mathcal{I}$ 的爆破，并把 $X$ 替换为 $X$ 的严格变换。
替换后，$B \setminus U$ 是某个有效 Cartier 除子 $D$ 的支撑
（《空间上的除子》引理
\ref{spaces-divisors-lemma-blowing-up-gives-effective-Cartier-divisor}）。
特别地，$U$ 在 $B$ 中概形论稠密（《空间上的除子》引理
\ref{spaces-divisors-lemma-complement-effective-Cartier-divisor}）。
接着再作一次 $U$-容许爆破，使得 $X \to B$ 平坦且有限表示；
见引理 \ref{spaces-more-morphisms-lemma-flat-after-blowing-up}。
注意 $U$ 在 $B$ 中仍概形论稠密。因此由引理
\ref{spaces-more-morphisms-lemma-zariski}，$X \to B$ 是开浸入。
\end{proof}

\noindent
下述引理说明，一个修改可被某个爆破支配。

\begin{lemma}
\label{spaces-more-morphisms-lemma-dominate-modification-by-blowup}
设 $S$ 为概形。设 $f : X \to B$ 为 $S$ 上代数空间的态射。
设 $U \subset B$ 为开子空间。假设
\begin{enumerate}
\item $B$ 拟紧且拟分离；
\item $U$ 拟紧；
\item $f : X \to B$ 固有；
\item $f^{-1}(U) \to U$ 是同构。% P09-ERR-0699
\end{enumerate}
则存在一个 $U$-容许爆破 $B' \to B$，它支配 $X$；
换言之，爆破态射存在分解 $B' \to X \to B$。
\end{lemma}

\begin{proof}
由引理 \ref{spaces-more-morphisms-lemma-isomorphism-after-blowing-up}，可找到一个
$U$-容许爆破 $B' \to B$，使得严格变换 $X'$ 同构地映到
$B'$。于是可用 $B' = X' \to X$ 作为该分解。
\end{proof}

\begin{lemma}
\label{spaces-more-morphisms-lemma-get-section-after-blowup}
设 $S$ 为概形。设 $X$、$Y$ 为 $S$ 上的代数空间。
设 $U \subset W \subset Y$ 为开子空间。
设 $f : X \to W$ 与 $s : U \to X$ 为态射，并满足
$f \circ s = \text{id}_U$。假设
\begin{enumerate}
\item $f$ 固有；
\item $Y$ 拟紧且拟分离；
\item $U$ 与 $W$ 拟紧。
\end{enumerate}
则存在一个 $U$-容许爆破 $b : Y' \to Y$ 以及态射
$s' : b^{-1}(W) \to X$；它延拓 $s$，并满足
$f \circ s' = b|_{b^{-1}(W)}$。
\end{lemma}

\begin{proof}
可以并且将把 $X$ 替换为 $s$ 的概形论像。
此时 $X \to W$ 在 $U$ 上是同构；见《空间的态射》引理
\ref{spaces-morphisms-lemma-scheme-theoretic-image-of-partial-section}。
由引理 \ref{spaces-more-morphisms-lemma-dominate-modification-by-blowup}，存在一个
$U$-容许爆破 $W' \to W$ 以及态射 $W' \to X$；后者延拓 $s$。
最后应用《空间上的除子》引理
\ref{spaces-divisors-lemma-extend-admissible-blowups}，
把 $W' \to W$ 延拓为一个 $U$-容许爆破，其靶为 $Y$，即完成证明。
\end{proof}
 
\section{Chow 引理}
\label{spaces-more-morphisms-section-chow}

\noindent
本节证明 Chow 引理
（引理 \ref{spaces-more-morphisms-lemma-chow-noetherian-separated}）。
我们建议读者参阅《空间的上同调》第
\ref{spaces-cohomology-section-weak-chow} 节，那里给出了
Chow 引理的一个容易证明并足以用于许多应用的弱版本。

\medskip\noindent
由于我们尚未定义代数空间的射影态射，
各引理的陈述（例如见引理
\ref{spaces-more-morphisms-lemma-blowup-to-find-embedding}）将涉及可表固有态射，
而不是射影态射。

\begin{lemma}
\label{spaces-more-morphisms-lemma-find-common-blowups}
设 $S$ 为概形。设 $Y$ 为 $S$ 上拟紧且拟分离的代数空间。
设 $U \to X_1$ 与 $U \to X_2$ 是 $Y$ 上代数空间的开浸入，
并假设 $U$、$X_1$、$X_2$ 在 $Y$ 上都是有限型且分离的。% P09-ERR-0700
则存在代数空间的交换图
$$
\xymatrix{
X_1' \ar[d] \ar[r] & X & X_2' \ar[l] \ar[d] \\
X_1 & U \ar[l] \ar[lu] \ar[u] \ar[ru] \ar[r] & X_2
}
$$
它在 $Y$ 上，其中 $X_i' \to X_i$ 是 $U$-容许爆破，
$X_i' \to X$ 是开浸入，
且 $X$ 在 $Y$ 上分离且有限型。
\end{lemma}

\begin{proof}
在整个证明中，所有代数空间都在 $Y$ 上分离且有限型。% P09-ERR-0701
特别地，这蕴含这些代数空间拟紧且拟分离，并且它们之间的态射
拟紧且分离。见《空间的态射》第
\ref{spaces-morphisms-section-separation-axioms} 节和第
\ref{spaces-morphisms-section-quasi-compact} 节。
我们将使用如下事实：若 $U \to W$ 是这类 $Y$ 上空间的浸入，
则概形论像 $Z$，即 $U$ 在 $W$ 中的概形论像，是 $W$ 的闭子空间，
$U \to Z$ 是开浸入，$U \subset Z$ 概形论稠密，并且
$|U| \subset |Z|$ 稠密。见《空间的态射》引理
\ref{spaces-morphisms-lemma-quasi-compact-immersion}。

\medskip\noindent
设 $X_{12} \subset X_1 \times_Y X_2$ 为
$U \to X_1 \times_Y X_2$ 的概形论像。由《空间的态射》引理
\ref{spaces-morphisms-lemma-scheme-theoretic-image-of-partial-section}，
投影 $p_i : X_{12} \to X_i$ 诱导同构 $p_i^{-1}(U) \to U$。
选取一个 $U$-容许爆破 $X_i^i \to X_i$，使得严格变换
$X_{12}^i$，即 $X_{12}$ 的严格变换，同构于 $X_i^i$ 的一个开子空间；见引理
\ref{spaces-more-morphisms-lemma-zariski-after-blowup}。
设 $\mathcal{I}_i \subset \mathcal{O}_{X_i}$ 为相应的有限型拟凝聚理想层。
回忆 $X_{12}^i \to X_{12}$ 是沿
$p_i^{-1}\mathcal{I}_i \mathcal{O}_{X_{12}}$ 的爆破；见
《空间上的除子》引理 \ref{spaces-divisors-lemma-strict-transform}。
设 $X_{12}'$ 为 $X_{12}$ 沿
$p_1^{-1}\mathcal{I}_1 p_2^{-1}\mathcal{I}_2 \mathcal{O}_{X_{12}}$
的爆破；关于这意味着什么，见《空间上的除子》引理
\ref{spaces-divisors-lemma-blowing-up-two-ideals}。我们得到交换图
$$
\xymatrix{
X_{12}' \ar[d] \ar[r] & X_{12}^2 \ar[d] \\
X_{12}^1 \ar[r] & X_{12}
}
$$
其中所有态射都是 $U$-容许爆破。由于
$X_{12}^i \subset X_i^i$ 是开子空间，可以选取一个
$U$-容许爆破 $X_i' \to X_i^i$，其限制为
$X_{12}' \to X_{12}^i$；见《空间上的除子》引理
\ref{spaces-divisors-lemma-extend-admissible-blowups}。
于是 $X_{12}' \subset X_i'$ 是开子空间，并且图
$$
\xymatrix{
X_{12}' \ar[d] \ar[r] & X_i' \ar[d] \\
X_{12}^i \ar[r] & X_i^i
}
$$
交换，其中竖直箭头是爆破，水平箭头是开浸入。% P09-ERR-0702
注意 $X'_{12} \to X_1' \times_Y X_2'$ 是浸入且固有
（使用 $X'_{12} \to X_{12}$ 固有、
$X_{12} \to X_1 \times_Y X_2$ 为闭态射，且
$X_1' \times_Y X_2' \to X_1 \times_Y X_2$ 分离，再应用
《空间的态射》引理
\ref{spaces-morphisms-lemma-universally-closed-permanence}）。
因此 $X'_{12} \to  X_1' \times_Y X_2'$ 是闭浸入。
若把 $X$ 定义为将 $X_1'$ 与 $X_2'$ 沿公共开子空间
$X_{12}'$ 黏合所得的空间，则 $X \to Y$ 有限型且分离
\footnote{因为可以检验对角态射 $X \to X \times_Y X$ 在四个开部分
$X'_i \times_Y X'_j$ 上的闭性；这些开部分覆盖 $X \times_Y X$，
而在其上闭性是显然的。}。
由于 $U$-容许爆破的复合仍为 $U$-容许爆破
（《空间上的除子》引理
\ref{spaces-divisors-lemma-composition-admissible-blowups}），
引理得证。
\end{proof}

\begin{lemma}
\label{spaces-more-morphisms-lemma-blowup-to-find-embedding}
设 $S$ 为概形。设 $f : X \to Y$ 为 $S$ 上代数空间的态射。
设 $U \subset X$ 为开子空间。假设
\begin{enumerate}
\item $U$ 拟紧；
\item $Y$ 拟紧且拟分离；
\item 存在浸入 $U \to \mathbf{P}^n_Y$，它在 $Y$ 上；
\item $f$ 有限型且分离。
\end{enumerate}
则存在交换图
$$
\xymatrix{
& U \ar[ld] \ar[d] \ar[rd] \ar[rrd] \\
X \ar[rd] & X' \ar[l] \ar[d] \ar[r] & Z' \ar[ld] \ar[r] & Z \ar[ld] \\
& Y & \mathbf{P}^n_Y \ar[l]
}
$$
其中以 $U$ 为源的箭头是开浸入，
$X' \to X$ 是 $U$-容许爆破，
$X' \to Z'$ 是开浸入，
$Z' \to Y$ 是代数空间的固有可表态射。
更精确地，$Z' \to Z$ 是 $U$-容许爆破，
而 $Z \to \mathbf{P}^n_Y$ 是闭浸入。
\end{lemma}

\begin{proof}
设 $Z \subset \mathbf{P}^n_Y$ 为浸入
$U \to \mathbf{P}^n_Y$ 的概形论像。由于
$U \to \mathbf{P}^n_Y$ 拟紧，可知 $U \subset Z$ 是
（概形论）稠密开子空间（《空间的态射》引理
\ref{spaces-morphisms-lemma-quasi-compact-immersion}）。
应用引理 \ref{spaces-more-morphisms-lemma-find-common-blowups} 得到图
$$
\xymatrix{
X' \ar[d] \ar[r] & \overline{X}' & Z' \ar[l] \ar[d] \\
X & U \ar[l] \ar[lu] \ar[u] \ar[ru] \ar[r] & Z
}
$$
它具有该引理陈述中列出的性质。
由于 $X' \to X$ 与 $Z' \to Z$ 都是 $U$-容许爆破，
可知 $U$ 是 $X'$ 与 $Z'$ 二者的概形论稠密开子空间
（见《空间上的除子》引理
\ref{spaces-divisors-lemma-blowing-up-gives-effective-Cartier-divisor} 和
\ref{spaces-divisors-lemma-complement-effective-Cartier-divisor}）。
由于 $Z' \to Z \to Y$ 固有，可知
$Z' \subset \overline{X}'$ 是闭子空间（见《空间的态射》引理
\ref{spaces-morphisms-lemma-universally-closed-permanence}）。
于是 $X' \subset Z'$（在概形论意义下），因而 $X'$ 是 $Z'$ 的
开子空间（略去一个小细节），引理得证。% P09-ERR-0703
\end{proof}

\begin{lemma}
\label{spaces-more-morphisms-lemma-chow-noetherian}
设 $S$ 为概形。设 $f : X \to Y$ 为 $S$ 上代数空间的态射。
假设 $f$ 分离且有限型，并且 $Y$ 诺特。% P09-ERR-0704
则存在稠密开子空间 $U \subset X$ 以及交换图
$$
\xymatrix{
& U \ar[ld] \ar[d] \ar[rd] \ar[rrd] \\
X \ar[rd] & X' \ar[l] \ar[d] \ar[r] & Z' \ar[ld] \ar[r] & Z \ar[ld] \\
& Y & \mathbf{P}^n_Y \ar[l]
}
$$
其中以 $U$ 为源的箭头是开浸入，
$X' \to X$ 是 $U$-容许爆破，
$X' \to Z'$ 是开浸入，
$Z' \to Y$ 是代数空间的固有可表态射。
更精确地，$Z' \to Z$ 是 $U$-容许爆破，
而 $Z \to \mathbf{P}^n_Y$ 是闭浸入。
\end{lemma}

\begin{proof}
由《空间的极限》引理
\ref{spaces-limits-lemma-embedding-into-affine-over-qs}，存在稠密开子空间
$U \subset X$ 以及浸入 $U \to \mathbf{A}^n_Y$，它在 $Y$ 上。
将其与开浸入 $\mathbf{A}^n_Y \to \mathbf{P}^n_Y$ 复合，
便得到引理 \ref{spaces-more-morphisms-lemma-blowup-to-find-embedding} 的情形，结论随之成立。
\end{proof}

\begin{remark}
\label{spaces-more-morphisms-remark-chow-Noetherian}
在引理 \ref{spaces-more-morphisms-lemma-blowup-to-find-embedding} 和
\ref{spaces-more-morphisms-lemma-chow-noetherian} 中，态射 $g : Z' \to Y$ 是射影态射的复合。
推测（由《态射》引理 \ref{morphisms-lemma-ample-composition}
对代数空间的类似结论），存在一个 $g$-丰沛可逆层，它位于 $Z'$ 上。
若以后用到这一点，我们将在此处陈述并证明它。
\end{remark}

\noindent
下述结果是 \cite[IV Theorem 3.1]{Kn}。注意浸入
$X' \to \mathbf{P}^n_Y$ 拟紧，因而可分解为
$X' \to Z' \to \mathbf{P}^n_Y$，其中第一个态射为开浸入，
第二个态射为闭浸入（《空间的态射》引理
\ref{spaces-morphisms-lemma-quasi-compact-immersion}）。

\begin{lemma}[Chow 引理]
\label{spaces-more-morphisms-lemma-chow-noetherian-separated}
\begin{reference}
\cite[IV Theorem 3.1]{Kn}
\end{reference}
设 $S$ 为概形。设 $f : X \to Y$ 为 $S$ 上代数空间的态射。
假设 $f$ 分离且有限型，并且 $Y$ 分离且诺特。% P09-ERR-0705
则存在交换图
$$
\xymatrix{
X \ar[rd] & X' \ar[l] \ar[d] \ar[r] & \mathbf{P}^n_Y \ar[ld] \\
& Y
}
$$
其中 $X' \to X$ 是 $U$-容许爆破，这里 $U \subset X$
是某个稠密开子空间，并且态射 $X' \to \mathbf{P}^n_Y$ 是浸入。
\end{lemma}

\begin{proof}
在证明的第一段中，我们把引理约化到 $Y$ 在
$\Spec(\mathbf{Z})$ 上有限型的情形。
我们可以并且将把基概形 $S$ 替换为 $\Spec(\mathbf{Z})$。
可将 $Y = \lim Y_i$ 写成 $\Spec(\mathbf{Z})$ 上有限型分离代数空间的
有向逆极限；见《空间的极限》命题
\ref{spaces-limits-proposition-approximate} 和引理
\ref{spaces-limits-lemma-descend-separated}。
对所有充分大的 $i$，可以找到有限型分离态射
$X_i \to Y_i$，使得 $X = Y \times_{Y_i} X_i$；见《空间的极限》引理
\ref{spaces-limits-lemma-descend-finite-presentation} 和
\ref{spaces-limits-lemma-descend-separated-morphism}。
设 $\eta_1, \ldots, \eta_n$ 为 $|X|$ 的不可约分支的一般点
（$X$ 是诺特代数空间 $Y$ 上有限型分离的代数空间，因而诺特，
所以 $|X|$ 是诺特拓扑空间）。
由《空间的极限》引理 \ref{spaces-limits-lemma-topology-limit}，
可知 $\eta_1, \ldots, \eta_n$ 在 $|X_i|$ 中的像彼此不同，
只要 $i$ 足够大。可以把 $X_i$ 替换为态射 $X \to X_i$
（它拟紧，事实上是仿射的）的概形论像。
替换后可知，$\eta_1, \ldots, \eta_n$ 在 $|X_i|$ 中的像是
$|X_i|$ 的不可约分支的一般点；见《空间的态射》引理
\ref{spaces-morphisms-lemma-quasi-compact-scheme-theoretic-image}。
在作出这些说明后，假设能找到图
$$
\xymatrix{
X_i \ar[rd] & X_i' \ar[l] \ar[d] \ar[r] & \mathbf{P}^n_{Y_i} \ar[ld] \\
& Y % P09-ERR-0706
}
$$
其中 $X_i' \to X_i$ 是 $U_i$-容许爆破，这里
$U_i \subset X_i$ 是某个稠密开子空间，并且态射
$X_i' \to \mathbf{P}^n_{Y_i}$ 是浸入。
则严格变换 $X' \to X$，即 $X$ 关于 $X_i' \to X_i$ 的严格变换，
是 $U$-容许爆破，其中 $U \subset X$ 是 $U_i$ 在 $X$ 中的逆像。
由于我们精心选取了指标 $i$，可得
$\eta_1, \ldots, \eta_n \in |U|$，且 $U \subset X$ 稠密。
此外，$X' \to \mathbf{P}^n_Y$ 是浸入，因为 $X'$ 在
$X_i' \times_{X_i} X = X_i' \times_{Y_i} Y$ 中为闭子空间，
而后者带有到 $\mathbf{P}^n_Y$ 的浸入。因此我们已约化到下一段的情形。

\medskip\noindent
假设 $Y$ 在 $\Spec(\mathbf{Z})$ 上分离且有限型。
于是 $X \to \Spec(\mathbf{Z})$ 也分离且有限型。
把引理 \ref{spaces-more-morphisms-lemma-chow-noetherian} 应用于
$X \to \Spec(\mathbf{Z})$，得到稠密开子空间 $U \subset X$
以及交换图
$$
\xymatrix{
& U \ar[ld] \ar[d] \ar[rd] \ar[rrd] \\
X \ar[rd] & X' \ar[l] \ar[d] \ar[r] & Z' \ar[ld] \ar[r] & Z \ar[ld] \\
& \Spec(\mathbf{Z}) & \mathbf{P}^n_\mathbf{Z} \ar[l]
}
$$
它具有该引理中列出的全部性质。
注意 $Z$ 有丰沛可逆层，即
$\mathcal{O}_{\mathbf{P}^n}(1)|_Z$。因此由《态射》引理
\ref{morphisms-lemma-projective-over-quasi-projective-is-H-projective}，
$Z' \to Z$ 是 H-射影态射。
由《态射》引理 \ref{morphisms-lemma-H-projective-composition}，
可知 $Z' \to \Spec(\mathbf{Z})$ 是 H-射影态射。
因此存在闭浸入 $Z' \to \mathbf{P}^m_{\Spec(\mathbf{Z})}$，
其中 $m \geq 0$。
于是对角态射
$$
X' \to Y \times \mathbf{P}^m_\mathbf{Z} = \mathbf{P}^m_Y
$$
是浸入（因为它与到 $\mathbf{P}^m_\mathbf{Z}$ 的投影之复合是浸入），
证明即告完成。
\end{proof}

\section{Chow 引理的变体}
\label{spaces-more-morphisms-section-chows-lemma}

\noindent
本节证明 Chow 引理的若干变体，它们处理非诺特代数空间之间的态射。
相应的诺特版本是引理 \ref{spaces-more-morphisms-lemma-chow-noetherian} 和引理
\ref{spaces-more-morphisms-lemma-chow-noetherian-separated}。

\begin{lemma}
\label{spaces-more-morphisms-lemma-chow-finite-type}
设 $S$ 为概形。设 $Y$ 为 $S$ 上拟紧且拟分离的代数空间。
设 $f : X \to Y$ 为有限型分离态射。则存在交换图
$$
\xymatrix{
X \ar[rd] & X' \ar[l] \ar[d] \ar[r] & \overline{X}' \ar[ld] \\
& Y
}
$$
其中 $X' \to X$ 固有且满，
$X' \to \overline{X}'$ 是开浸入，并且
$\overline{X}' \to Y$ 是代数空间的固有可表态射。% P09-ERR-0707
\end{lemma}

\begin{proof}
由《空间的极限》命题
\ref{spaces-limits-proposition-separated-closed-in-finite-presentation}，
可以找到闭浸入 $X \to X_1$，其中 $X_1$ 在 $Y$ 上分离且有限表示。
显然，若对 $X_1 \to Y$ 证明了断言，则 $X$ 的结论随之成立。
故可假设 $X$ 在 $Y$ 上有限表示。

\medskip\noindent
我们可以并且将把基概形 $S$ 替换为 $\Spec(\mathbf{Z})$。
把 $Y = \lim_i Y_i$ 写成 $\Spec(\mathbf{Z})$ 上有限型拟分离
代数空间的有向逆极限；见《空间的极限》命题
\ref{spaces-limits-proposition-approximate}。
由《空间的极限》引理
\ref{spaces-limits-lemma-descend-finite-presentation}，可以找到指标
$i \in I$ 以及有限表示的概形 $X_i \to Y_i$，% P09-ERR-0708
使得 $X = Y \times_{Y_i} X_i$。
由《空间的极限》引理
\ref{spaces-limits-lemma-descend-separated-morphism}，
可以假设 $X_i \to Y_i$ 分离。
显然，若对 $X_i$（在 $Y_i$ 上）证明了断言，则 $X$ 的断言也成立。
$X_i \to Y_i$ 的情形由引理 \ref{spaces-more-morphisms-lemma-chow-noetherian} 处理。
\end{proof}

\begin{lemma}
\label{spaces-more-morphisms-lemma-chow-finite-type-separated}
设 $S$ 为概形。设 $f : X \to Y$ 为 $S$ 上代数空间的态射。
假设 $f$ 分离且有限型，并且 $Y$ 分离且拟紧。% P09-ERR-0709
则存在交换图
$$
\xymatrix{
X \ar[rd] & X' \ar[l] \ar[d] \ar[r] & \mathbf{P}^n_Y \ar[ld] \\
& Y
}
$$
其中 $X' \to X$ 是固有满态射，% P09-ERR-0710
并且态射 $X' \to \mathbf{P}^n_Y$ 是浸入。
\end{lemma}

\begin{proof}
由《空间的极限》命题
\ref{spaces-limits-proposition-separated-closed-in-finite-presentation}，
可以找到闭浸入 $X \to X_1$，其中 $X_1$ 在 $Y$ 上分离且有限表示。
显然，若对 $X_1 \to Y$ 证明了断言，则 $X$ 的结论随之成立。
故可假设 $X$ 在 $Y$ 上有限表示。

\medskip\noindent
我们可以并且将把基概形 $S$ 替换为 $\Spec(\mathbf{Z})$。
把 $Y = \lim_i Y_i$ 写成 $\Spec(\mathbf{Z})$ 上有限型拟分离
代数空间的有向逆极限；见《空间的极限》命题
\ref{spaces-limits-proposition-approximate}。
由《空间的极限》引理 \ref{spaces-limits-lemma-descend-separated}，
可以假设 $Y_i$ 对所有 $i$ 都分离。
由《空间的极限》引理
\ref{spaces-limits-lemma-descend-finite-presentation}，可以找到指标
$i \in I$ 以及有限表示的概形 $X_i \to Y_i$，% P09-ERR-0711
使得 $X = Y \times_{Y_i} X_i$。
由《空间的极限》引理
\ref{spaces-limits-lemma-descend-separated-morphism}，
可以假设 $X_i \to Y_i$ 分离。
显然，若对 $X_i$（在 $Y_i$ 上）证明了断言，则 $X$ 的断言也成立。
$X_i \to Y_i$ 的情形由引理
\ref{spaces-more-morphisms-lemma-chow-noetherian-separated} 处理。
\end{proof}

\section{Grothendieck 存在定理}
\label{spaces-more-morphisms-section-existence-proper}

\noindent
本节讨论代数空间的 Grothendieck 存在定理。我们不发展“形式代数空间”
的理论，而是暂时建立一些术语，以替代“诺特 adic 形式空间上的凝聚模”
这一概念。

\medskip\noindent
设 $S$ 为概形。设 $X$ 为 $S$ 上的诺特代数空间。
设 $\mathcal{I} \subset \mathcal{O}_X$ 为拟凝聚理想层。
下面考虑逆系统 $(\mathcal{F}_n)$，其各项为凝聚
$\mathcal{O}_X$-模，并且满足
\begin{enumerate}
\item $\mathcal{F}_n$ 被 $\mathcal{I}^n$ 零化；
\item 转移态射诱导同构
$\mathcal{F}_{n + 1}/\mathcal{I}^n\mathcal{F}_{n + 1} \to \mathcal{F}_n$。
\end{enumerate}
这类逆系统之间的态射
$\alpha : (\mathcal{F}_n) \to (\mathcal{G}_n)$，就是相容的态射系统
$\alpha_n : \mathcal{F}_n \to \mathcal{G}_n$。
把这些逆系统所成的范畴记为 $\textit{Coh}(X, \mathcal{I})$。
我们将建立关于这些系统的一些理论，它与概形情形中的相应结果平行；
见《概形的上同调》第 \ref{coherent-section-existence}、
\ref{coherent-section-existence-proper}、
\ref{coherent-section-existence-proper-support} 和
\ref{coherent-section-algebraization} 节。% P09-ERR-0712

\medskip\noindent
函子性。设 $f : X \to Y$ 为某概形 $S$ 上诺特代数空间的态射，
并设 $\mathcal{J} \subset \mathcal{O}_Y$ 为拟凝聚理想层。
令 $\mathcal{I} = f^{-1}\mathcal{J}\mathcal{O}_X$。
在此情形下有函子
$$
f^* : \textit{Coh}(Y, \mathcal{J}) \longrightarrow \textit{Coh}(X, \mathcal{I})
$$
它把 $(\mathcal{G}_n)$ 映到 $(f^*\mathcal{G}_n)$。
参见《概形的上同调》引理
\ref{coherent-lemma-inverse-systems-pullback}。
若 $f$ 平展，则可把它简单地看作将系统限制到 $X$；
见《空间的性质》公式
\ref{spaces-properties-equation-restrict-modules}。

\medskip\noindent
平展下降。设 $S$ 为概形。设 $U_0 \to X$ 为诺特代数空间的满平展态射。
令 $U_1 = U_0 \times_X U_0$ 以及
$U_2 = U_0 \times_X U_0 \times_X U_0$。
设 $\mathcal{I} \subset \mathcal{O}_{X}$ 为拟凝聚理想层。
令 $\mathcal{I}_i = \mathcal{I}|_{U_i}$。在此情形下得到范畴图
$$
\xymatrix{
\textit{Coh}(X, \mathcal{I}) \ar[r] &
\textit{Coh}(U_0, \mathcal{I}_0) \ar@<0.5ex>[r] \ar@<-0.5ex>[r] &
\textit{Coh}(U_1, \mathcal{I}_1) \ar@<1ex>[r] \ar[r] \ar@<-1ex>[r] &
\textit{Coh}(U_2, \mathcal{I}_2)
}
$$
并且第一个箭头把 $\textit{Coh}(X, \mathcal{I})$ 展示为该图右部的
同伦极限。% P09-ERR-0713
更精确地，给定一个{\it 下降资料}，即一对
$((\mathcal{G}_n), \varphi)$，其中 $(\mathcal{G}_n)$ 是
$\textit{Coh}(U_0, \mathcal{I}_0)$ 的对象，且
$\varphi : \text{pr}_0^*(\mathcal{G}_n) \to \text{pr}_1^*(\mathcal{G}_n)$
是 $\textit{Coh}(U_1, \mathcal{I}_1)$ 中的同构，并在
$\textit{Coh}(U_2, \mathcal{I}_2)$ 中满足余圈条件，则存在唯一对象
$(\mathcal{F}_n)$ 属于 $\textit{Coh}(X, \mathcal{I})$，
其相伴的典范下降资料同构于 $((\mathcal{G}_n), \varphi)$。
参见《空间上的下降》定义
\ref{spaces-descent-definition-descent-datum-effective-quasi-coherent}。
对下降资料 $(\mathcal{G}_n, \varphi_n)$（其中 $n$ 变化）应用
《空间上的下降》命题
\ref{spaces-descent-proposition-fpqc-descent-quasi-coherent}，
便立即得到这一陈述的证明。

\begin{lemma}
\label{spaces-more-morphisms-lemma-inverse-systems-abelian}
设 $S$ 为概形。设 $X$ 为 $S$ 上的诺特代数空间，
并设 $\mathcal{I} \subset \mathcal{O}_X$ 为拟凝聚理想层。
\begin{enumerate}
\item 范畴 $\textit{Coh}(X, \mathcal{I})$ 是阿贝尔范畴。
\item $\textit{Coh}(X, \mathcal{I})$ 中的正合性可在平展局部检验。
\item 对诺特代数空间的任意平坦态射 $f : X' \to X$，函子
$f^* : \textit{Coh}(X, \mathcal{I}) \to
\textit{Coh}(X', f^{-1}\mathcal{I}\mathcal{O}_{X'})$ 是正合的。
\end{enumerate}
\end{lemma}

\begin{proof}
先证 (1)。选择仿射概形 $U_0$ 以及满平展态射
$U_0 \to X$。如同上述平展下降的讨论，令
$U_1 = U_0 \times_X U_0$ 以及
$U_2 = U_0 \times_X U_0 \times_X U_0$。
各范畴 $\textit{Coh}(U_i, \mathcal{I}_i)$ 都是阿贝尔范畴
（《概形的上同调》引理
\ref{coherent-lemma-inverse-systems-abelian}），并且拉回函子
$\textit{Coh}(U_0, \mathcal{I}_0) \to \textit{Coh}(U_1, \mathcal{I}_1)$
以及
$\textit{Coh}(U_1, \mathcal{I}_1) \to \textit{Coh}(U_2, \mathcal{I}_2)$
都是正合函子（《概形的上同调》引理
\ref{coherent-lemma-inverse-systems-pullback}）。
于是，由把 $\textit{Coh}(X, \mathcal{I})$ 描述成下降资料范畴，
形式地得到引理。略去一些细节；参见《群胚》引理
\ref{groupoids-lemma-abelian} 的证明。% P09-ERR-0714

\medskip\noindent
第 (2) 部分立即由上一段的讨论得到。在 (3) 的情形中，
选择交换图
$$
\xymatrix{
U' \ar[d] \ar[r] & U \ar[d] \\
X' \ar[r] & X
}
$$
其中 $U'$ 与 $U$ 是仿射概形，竖直态射满且平展。
于是 $U' \to U$ 是诺特概形的平坦态射（《空间的态射》引理
\ref{spaces-morphisms-lemma-flat-local}），故拉回函子
$\textit{Coh}(U, \mathcal{I}\mathcal{O}_U) \to
\textit{Coh}(U', \mathcal{I}\mathcal{O}_{U'})$
由《概形的上同调》引理
\ref{coherent-lemma-inverse-systems-pullback} 而正合。
由于可以检验 $\textit{Coh}(X, \mathcal{O}_X)$ 中的正合性，
而这一检验可在 $U$ 上进行，并且对 $X', U'$ 也类似，
断言成立。% P09-ERR-0715
\end{proof}

\begin{lemma}
\label{spaces-more-morphisms-lemma-inverse-systems-surjective}
设 $S$ 为概形。设 $X$ 为 $S$ 上的诺特代数空间，
并设 $\mathcal{I} \subset \mathcal{O}_X$ 为拟凝聚理想层。
态射 $(\mathcal{F}_n) \to (\mathcal{G}_n)$ 在
$\textit{Coh}(X, \mathcal{I})$ 中为满，当且仅当
$\mathcal{F}_1 \to \mathcal{G}_1$ 为满。
\end{lemma}

\begin{proof}
由引理 \ref{spaces-more-morphisms-lemma-inverse-systems-abelian}，可以在 $X$ 的某个仿射
平展覆盖上检验。因而问题约化到概形情形，即《概形的上同调》引理
\ref{coherent-lemma-inverse-systems-surjective}。
\end{proof}

\noindent
设 $S$ 为概形。设 $X$ 为 $S$ 上的诺特代数空间，
并设 $\mathcal{I} \subset \mathcal{O}_X$ 为拟凝聚理想层。
存在函子
\begin{equation}
\label{spaces-more-morphisms-equation-completion-functor}
\textit{Coh}(\mathcal{O}_X) \longrightarrow \textit{Coh}(X, \mathcal{I}), \quad
\mathcal{F} \longmapsto \mathcal{F}^\wedge
\end{equation}
它把凝聚 $\mathcal{O}_X$-模 $\mathcal{F}$ 映到对象
$\mathcal{F}^\wedge = (\mathcal{F}/\mathcal{I}^n\mathcal{F})$，
该对象属于 $\textit{Coh}(X, \mathcal{I})$。

\begin{lemma}
\label{spaces-more-morphisms-lemma-exact}
函子 (\ref{spaces-more-morphisms-equation-completion-functor}) 是正合的。
\end{lemma}

\begin{proof}
只须在 $X$ 上平展局部检验这一点；见引理
\ref{spaces-more-morphisms-lemma-inverse-systems-abelian}。
因而问题约化到概形情形，即《概形的上同调》引理
\ref{coherent-lemma-exact}。
\end{proof}

\begin{lemma}
\label{spaces-more-morphisms-lemma-completion-internal-hom}
设 $S$ 为概形。设 $X$ 为 $S$ 上的诺特代数空间，
并设 $\mathcal{I} \subset \mathcal{O}_X$ 为拟凝聚理想层。
设 $\mathcal{F}$、$\mathcal{G}$ 为凝聚 $\mathcal{O}_X$-模。令
$\mathcal{H} = \SheafHom_{\mathcal{O}_X}(\mathcal{F}, \mathcal{G})$。
则
$$
\lim H^0(X, \mathcal{H}/\mathcal{I}^n\mathcal{H}) =
\Mor_{\textit{Coh}(X, \mathcal{I})}
(\mathcal{F}^\wedge, \mathcal{G}^\wedge).
$$
\end{lemma}

\begin{proof}
由于 $\mathcal{H}$ 是 $X_\etale$ 上的层，而且
$\textit{Coh}(X, \mathcal{I})$ 的对象具有平展下降，
只须在平展局部证明。因而问题约化到概形情形，即《概形的上同调》引理
\ref{coherent-lemma-completion-internal-hom}。
\end{proof}

\noindent
现在引入本节其余部分始终关注的情形。

\begin{situation}
\label{spaces-more-morphisms-situation-existence}
这里 $A$ 是关于理想 $I$ 完备的诺特环。此外，
$f : X \to \Spec(A)$ 是代数空间的有限型分离态射，且
$\mathcal{I} = I\mathcal{O}_X$。
\end{situation}

\noindent
在此情形中，记
$$
\textit{Coh}_{\text{支撑固有于 } A}(\mathcal{O}_X)
$$
为 $\textit{Coh}(\mathcal{O}_X)$ 的全子范畴，% P09-ERR-0716
其对象是凝聚 $\mathcal{O}_X$-模，而且其支撑在 $\Spec(A)$ 上固有；
等价地，其概形论支撑在 $\Spec(A)$ 上固有。见
《空间的导出范畴》引理
\ref{spaces-perfect-lemma-module-support-proper-over-base}。
类似地，令
$$
\textit{Coh}_{\text{支撑固有于 } A}(X, \mathcal{I})
$$
为 $\textit{Coh}(X, \mathcal{I})$ 的全子范畴，
它由满足下述条件的对象 $(\mathcal{F}_n)$ 组成：
$\mathcal{F}_1$ 的支撑在 $\Spec(A)$ 上固有。
由于商模的支撑包含在原模的支撑中，函子
(\ref{spaces-more-morphisms-equation-completion-functor}) 诱导函子
\begin{equation}
\label{spaces-more-morphisms-equation-completion-functor-proper-over-A}
\textit{Coh}_{\text{支撑固有于 }A}(\mathcal{O}_X)
\longrightarrow
\textit{Coh}_{\text{支撑固有于 }A}(X, \mathcal{I})
\end{equation}
第一个结果是此函子全忠实。

\begin{lemma}
\label{spaces-more-morphisms-lemma-fully-faithful}
处于情形 \ref{spaces-more-morphisms-situation-existence}。
设 $\mathcal{F}$、$\mathcal{G}$ 为凝聚 $\mathcal{O}_X$-模。
假设 $\mathcal{F}$ 与 $\mathcal{G}$ 的支撑之交在
$\Spec(A)$ 上固有。则由 (\ref{spaces-more-morphisms-equation-completion-functor}) 得到的映射
$$
\Mor_{\textit{Coh}(\mathcal{O}_X)}(\mathcal{F}, \mathcal{G})
\longrightarrow
\Mor_{\textit{Coh}(X, \mathcal{I})}
(\mathcal{F}^\wedge, \mathcal{G}^\wedge)
$$
是双射。特别地，(\ref{spaces-more-morphisms-equation-completion-functor-proper-over-A})
全忠实。
\end{lemma}

\begin{proof}
令 $\mathcal{H} = \SheafHom_{\mathcal{O}_X}(\mathcal{G}, \mathcal{F})$。
这是凝聚 $\mathcal{O}_X$-模，因为它限制到平展映到 $X$ 的概形后，
由《模》引理
\ref{modules-lemma-internal-hom-locally-kernel-direct-sum} 是凝聚的。
% P09-ERR-0717 P09-ERR-0718
由引理 \ref{spaces-more-morphisms-lemma-completion-internal-hom}，映射
$$
\lim_n H^0(X, \mathcal{H}/\mathcal{I}^n\mathcal{H})
\to
\Mor_{\textit{Coh}(X, \mathcal{I})}
(\mathcal{G}^\wedge, \mathcal{F}^\wedge)
$$
是双射。设 $i : Z \to X$ 为 $\mathcal{H}$ 的概形论支撑。
显然，$Z$ 是闭子空间，且 $|Z|$ 包含在 $\mathcal{F}$ 与
$\mathcal{G}$ 的支撑之交中。因此由假设，$Z \to \Spec(A)$ 固有
（见《空间的导出范畴》第
\ref{spaces-perfect-section-proper-over-base} 节）。
写成 $\mathcal{H} = i_*\mathcal{H}'$，其中
$\mathcal{O}_Z$-模 $\mathcal{H}'$ 是凝聚的。我们有
$i_*(\mathcal{H}'/I^n\mathcal{H}') = \mathcal{H}/I^n\mathcal{H}$。
因此得到
\begin{align*}
\lim_n H^0(X, \mathcal{H}/\mathcal{I}^n\mathcal{H})
& =
\lim_n H^0(Z, \mathcal{H}'/\mathcal{I}^n\mathcal{H}') \\
& =
H^0(Z, \mathcal{H}') \\
& =
H^0(X, \mathcal{H}) \\
& = 
\Mor_{\textit{Coh}(\mathcal{O}_X)}(\mathcal{F}, \mathcal{G})
\end{align*}
其中第二个等式由形式函数定理得到（《空间的上同调》引理
\ref{spaces-cohomology-lemma-spell-out-theorem-formal-functions}）。
这就证明了引理。
\end{proof}

\begin{remark}
\label{spaces-more-morphisms-remark-inverse-systems-kernel-cokernel-annihilated-by}
设 $S$ 为概形。设 $X$ 为 $S$ 上的诺特代数空间，并设
$\mathcal{I}, \mathcal{K} \subset \mathcal{O}_X$ 为拟凝聚理想层。
设 $\alpha : (\mathcal{F}_n) \to (\mathcal{G}_n)$ 为
$\textit{Coh}(X, \mathcal{I})$ 的态射。
给定仿射概形 $U = \Spec(A)$ 以及满平展态射 $U \to X$，
以 $I, K \subset A$ 表示与限制
$\mathcal{I}|_U, \mathcal{K}|_U$ 对应的理想。
以 $\alpha_U : M \to N$ 表示有限 $A^\wedge$-模之间的态射，
它经由《概形的上同调》引理
\ref{coherent-lemma-inverse-systems-affine} 对应于 $\alpha|_U$。
% P09-ERR-0719
我们断言下列条件等价：
\begin{enumerate}
\item 存在整数 $t \geq 1$，使得 $\Ker(\alpha_n)$ 与
$\Coker(\alpha_n)$ 都被 $\mathcal{K}^t$ 零化，对所有 $n \geq 1$；
\item 对任意（或某个）如上的仿射开集
$\Spec(A) = U \subset X$，模 $\Ker(\alpha_U)$ 与
$\Coker(\alpha_U)$ 都被 $K^t$ 零化，其中某个整数 $t \geq 1$。
% P09-ERR-0720
\end{enumerate}
若这些等价条件成立，就称 $\alpha$ 是一个
{\it 其核与余核被 $\mathcal{K}$ 的某次幂零化的映射}。
关于等价性的证明，参见《概形的上同调》评注
\ref{coherent-remark-inverse-systems-kernel-cokernel-annihilated-by}。
\end{remark}

\begin{lemma}
\label{spaces-more-morphisms-lemma-existence-easy}
设 $S$ 为概形。设 $X$ 为 $S$ 上的诺特代数空间，并设
$\mathcal{I} \subset \mathcal{O}_X$ 为拟凝聚理想层。
设 $\mathcal{G}$ 为凝聚 $\mathcal{O}_X$-模，$(\mathcal{F}_n)$ 为
$\textit{Coh}(X, \mathcal{I})$ 的对象，并设
$\alpha : (\mathcal{F}_n) \to \mathcal{G}^\wedge$ 是一个其核与余核
被 $\mathcal{I}$ 的某次幂零化的映射。
则存在唯一的（精确到唯一同构）三元组
$(\mathcal{F}, a, \beta)$，其中
\begin{enumerate}
\item $\mathcal{F}$ 是凝聚 $\mathcal{O}_X$-模；
\item $a : \mathcal{F} \to \mathcal{G}$ 是 $\mathcal{O}_X$-模态射，
其核与余核被 $\mathcal{I}$ 的某次幂零化；
\item $\beta : (\mathcal{F}_n) \to \mathcal{F}^\wedge$ 是同构；
\item $\alpha = a^\wedge \circ \beta$。
\end{enumerate}
\end{lemma}

\begin{proof}
$\textit{Coh}(X, \mathcal{I})$ 和 $\textit{Coh}(\mathcal{O}_X)$
的对象具有唯一性和平展下降，这蕴含只须构造
$(\mathcal{F}, a, \beta)$，而这一构造可在 $X$ 上平展局部进行。
% P09-ERR-0721
因而问题约化到概形情形，即《概形的上同调》引理
\ref{coherent-lemma-existence-easy}。
\end{proof}

\begin{lemma}
\label{spaces-more-morphisms-lemma-existence-tricky}
处于情形 \ref{spaces-more-morphisms-situation-existence}。设
$\mathcal{K} \subset \mathcal{O}_X$ 为拟凝聚理想层。
设 $X_e \subset X$ 为由 $\mathcal{K}^e$ 切出的闭子空间。
令 $\mathcal{I}_e = \mathcal{I}\mathcal{O}_{X_e}$。
设 $(\mathcal{F}_n)$ 为
$\textit{Coh}_{\text{支撑固有于 } A}(X, \mathcal{I})$ 的对象。假设
\begin{enumerate}
\item 函子
$\textit{Coh}_{\text{支撑固有于 } A}(\mathcal{O}_{X_e})
\to \textit{Coh}_{\text{支撑固有于 } A}(X_e, \mathcal{I}_e)$
对所有 $e \geq 1$ 都是范畴等价；
\item 存在对象 $\mathcal{H}$，它属于
$\textit{Coh}_{\text{支撑固有于 } A}(\mathcal{O}_X)$，以及映射
$\alpha : (\mathcal{F}_n) \to \mathcal{H}^\wedge$，其核与余核
被 $\mathcal{K}$ 的某次幂零化。
\end{enumerate}
则 $(\mathcal{F}_n)$ 属于
(\ref{spaces-more-morphisms-equation-completion-functor-proper-over-A}) 的本质像。
\end{lemma}

\begin{proof}
在本证明中，我们将不再另行说明地使用下述事实：对闭浸入
$i : Z \to X$，函子 $i_*$ 给出 $Z$ 上凝聚模的范畴与 $X$ 上
被 $Z$ 的理想层零化的凝聚模范畴之间的等价；见
《空间的上同调》引理
\ref{spaces-cohomology-lemma-i-star-equivalence}。
特别地，我们把
$$
\textit{Coh}_{\text{支撑固有于 } A}(\mathcal{O}_{X_e})
\subset
\textit{Coh}_{\text{支撑固有于 } A}(\mathcal{O}_X)
$$
看成由被 $\mathcal{K}^e$ 零化的模组成的全子范畴，并把
$$
\textit{Coh}_{\text{支撑固有于 } A}(X_e, \mathcal{I}_e)
\subset
\textit{Coh}_{\text{支撑固有于 } A}(X, \mathcal{I})
$$
看成由被 $\mathcal{K}^e$ 零化的对象组成的全子范畴。% P09-ERR-0722
此外，(1) 告诉我们，这两个范畴在完备化函子
(\ref{spaces-more-morphisms-equation-completion-functor-proper-over-A}) 下等价。

\medskip\noindent
应用此等价，得到凝聚 $\mathcal{O}_X$-模 $\mathcal{G}_e$；
它被 $\mathcal{K}^e$ 零化，并对应于系统
$(\mathcal{F}_n/\mathcal{K}^e\mathcal{F}_n)$；该系统属于
$\textit{Coh}_{\text{支撑固有于 } A}(X, \mathcal{I})$。态射
$\mathcal{F}_n/\mathcal{K}^{e + 1}\mathcal{F}_n \to
\mathcal{F}_n/\mathcal{K}^e\mathcal{F}_n$ 对应于典范态射
$\mathcal{G}_{e + 1} \to \mathcal{G}_e$，后者诱导同构
$\mathcal{G}_{e + 1}/\mathcal{K}^e\mathcal{G}_{e + 1} \to \mathcal{G}_e$。
我们得到对象 $(\mathcal{G}_e)$，它属于范畴
$\textit{Coh}_{\text{支撑固有于 } A}(X, \mathcal{K})$。
映射 $\alpha$ 诱导态射系统
$$
\mathcal{F}_n/\mathcal{K}^e\mathcal{F}_n
\longrightarrow
\mathcal{H}/(\mathcal{I}^n + \mathcal{K}^e)\mathcal{H}
$$
从而再次由范畴等价得到态射
$\mathcal{G}_e \to \mathcal{H}/\mathcal{K}^e\mathcal{H}$。
由假设 (2)，存在整数 $t \geq 1$，使得 $\mathcal{K}^t$ 零化所有态射
$\mathcal{F}_n \to \mathcal{H}/\mathcal{I}^n\mathcal{H}$ 的核与余核。
于是 $\mathcal{K}^{2t}$ 零化态射
$\mathcal{F}_n/\mathcal{K}^e\mathcal{F}_n \to
\mathcal{H}/(\mathcal{I}^n + \mathcal{K}^e)\mathcal{H}$
的核与余核（略去细节；见《概形的上同调》评注
\ref{coherent-remark-inverse-systems-kernel-cokernel-annihilated-by}）。
进而可知 $\mathcal{K}^{4t}$ 零化下列态射的核与余核：
$$
\mathcal{G}_e
\longrightarrow
\mathcal{H}/\mathcal{K}^e\mathcal{H},
$$
（略去细节；见《概形的上同调》评注
\ref{coherent-remark-inverse-systems-kernel-cokernel-annihilated-by}）。
% P09-ERR-0723
应用引理 \ref{spaces-more-morphisms-lemma-existence-easy}，得到凝聚
$\mathcal{O}_X$-模 $\mathcal{F}$、态射
$a : \mathcal{F} \to \mathcal{H}$ 以及同构
$\beta : (\mathcal{G}_e) \to (\mathcal{F}/\mathcal{K}^e\mathcal{F})$，
后者位于 $\textit{Coh}(X, \mathcal{K})$ 中。逆向追溯可知，
对给定的 $n$，三元组
$(\mathcal{F}/\mathcal{I}^n\mathcal{F}, a \bmod \mathcal{I}^n, \beta
\bmod \mathcal{I}^n)$ 是上述引理中对应于态射
$\alpha_n \bmod \mathcal{K}^e :
(\mathcal{F}_n/\mathcal{K}^e\mathcal{F}_n) \to
(\mathcal{H}/(\mathcal{I}^n + \mathcal{K}^e)\mathcal{H})$
的三元组，该态射位于 $\textit{Coh}(X, \mathcal{K})$ 中。
因此，引理 \ref{spaces-more-morphisms-lemma-existence-easy} 中的唯一性给出典范同构
$\mathcal{F}/\mathcal{I}^n\mathcal{F} \to \mathcal{F}_n$，
它与这里涉及的所有态射相容。

\medskip\noindent
为完成引理的证明，还须说明 $\mathcal{F}$ 的支撑在 $A$ 上固有。
按构造，$a : \mathcal{F} \to \mathcal{H}$ 的核被
$\mathcal{K}$ 的某次幂零化。因此，此核的支撑包含在
$\mathcal{G}_1$ 的支撑中。由于 $\mathcal{G}_1$ 是
$\textit{Coh}_{\text{支撑固有于 } A}(\mathcal{O}_{X_1})$ 的对象，
可知此支撑在 $A$ 上固有。再结合 $\mathcal{H}$ 的支撑在 $A$
上固有这一事实，由《空间的导出范畴》引理
\ref{spaces-perfect-lemma-union-closed-proper-over-base}，
可知 $\mathcal{F}$ 的支撑在 $A$ 上固有。
\end{proof}

\begin{lemma}
\label{spaces-more-morphisms-lemma-inverse-systems-push-pull}
设 $S$ 为概形。设 $f : X \to Y$ 为 $S$ 上诺特代数空间的
可表固有态射。设
$\mathcal{J}, \mathcal{K} \subset \mathcal{O}_Y$ 为拟凝聚理想层。
假设 $f$ 在 $V = Y \setminus V(\mathcal{K})$ 上是同构。
令 $\mathcal{I} = f^{-1}\mathcal{J} \mathcal{O}_X$。
设 $(\mathcal{G}_n)$ 为 $\textit{Coh}(Y, \mathcal{J})$ 的对象，
设 $\mathcal{F}$ 为凝聚 $\mathcal{O}_X$-模，并设
$\beta : (f^*\mathcal{G}_n)  \to \mathcal{F}^\wedge$ 是
$\textit{Coh}(X, \mathcal{I})$ 中的同构。则存在映射
$$
\alpha :
(\mathcal{G}_n)
\longrightarrow
(f_*\mathcal{F})^\wedge
$$
它位于 $\textit{Coh}(Y, \mathcal{J})$ 中，并且其核与余核
被 $\mathcal{K}$ 的某次幂零化。
\end{lemma}

\begin{proof}
由于 $f$ 是固有态射，$f_*\mathcal{F}$ 是凝聚
$\mathcal{O}_Y$-模（《空间的上同调》引理
\ref{spaces-cohomology-lemma-proper-pushforward-coherent}）。
故引理的陈述有意义。考虑复合
$$
\gamma_n : \mathcal{G}_n \to
f_*f^*\mathcal{G}_n \to
f_*(\mathcal{F}/\mathcal{I}^n\mathcal{F}).
$$
这里第一个映射是伴随映射，第二个是 $f_*\beta_n$。
我们断言，存在唯一的引理所述 $\alpha$，使得复合
$$
\mathcal{G}_n \xrightarrow{\alpha_n}
f_*\mathcal{F}/\mathcal{J}^nf_*\mathcal{F} \to
f_*(\mathcal{F}/\mathcal{I}^n\mathcal{F})
$$
都等于 $\gamma_n$，对所有 $n$。由唯一性以及
$\textit{Coh}(Y, \mathcal{J})$ 的平展下降，只须在 $Y$ 上
平展局部证明。故可假设 $Y$ 是某个诺特环的谱。
由于 $f$ 可表，可知 $X$ 也是概形。这样便约化到概形情形；
见《概形的上同调》引理
\ref{coherent-lemma-inverse-systems-push-pull} 的证明。
\end{proof}

\begin{theorem}[Grothendieck 存在定理]
\label{spaces-more-morphisms-theorem-grothendieck-existence}
在情形 \ref{spaces-more-morphisms-situation-existence} 中，函子
(\ref{spaces-more-morphisms-equation-completion-functor-proper-over-A}) 是范畴等价。
\end{theorem}

\begin{proof}
在定理的证明中，我们将不再另行说明地使用《空间的上同调》引理
\ref{spaces-cohomology-lemma-i-star-equivalence} 的范畴等价。
由引理 \ref{spaces-more-morphisms-lemma-fully-faithful}，该函子全忠实。
因此只须证明该函子本质满。

\medskip\noindent
考虑集合 $\Xi$；其元素为拟凝聚理想层
$\mathcal{K} \subset \mathcal{O}_X$，并满足：断言对每个对象
$(\mathcal{F}_n)$ 都成立，其中该对象属于
$\textit{Coh}_{\text{支撑固有于 }A}(X, \mathcal{I})$，并被
$\mathcal{K}$ 零化。
我们要证明 $(0)$ 属于 $\Xi$。若不然，由于 $X$ 诺特，存在
极大拟凝聚理想层 $\mathcal{K}$，它不属于 $\Xi$；见《空间的上同调》引理
\ref{spaces-cohomology-lemma-acc-coherent}。
把 $X$ 替换为 $X$ 中与 $\mathcal{K}$ 对应的闭子概形后，% P09-ERR-0724
可以假设每个非零 $\mathcal{K}$ 都属于 $\Xi$。
设 $(\mathcal{F}_n)$ 为
$\textit{Coh}_{\text{支撑固有于 }A}(X, \mathcal{I})$ 的对象。
我们将证明此对象属于本质像，从而完成定理的证明。

\medskip\noindent
应用 Chow 引理（引理 \ref{spaces-more-morphisms-lemma-chow-noetherian-separated}），
找到固有满态射 $f : Y \to X$；它在某个稠密开子空间
$U \subset X$ 上是同构，并且 $Y$ 在 $A$ 上 H-拟射影。
注意 $Y$ 是概形，且 $f$ 可表。选取开浸入
$j : Y \to Y'$，其中 $Y'$ 在 $A$ 上射影；见《态射》引理
\ref{morphisms-lemma-H-quasi-projective-open-H-projective}。
设 $T_n$ 为 $\mathcal{F}_n$ 的概形论支撑。
注意 $|T_n| = |T_1|$，故 $T_n$ 在 $A$ 上固有，对所有 $n$
（《空间的态射》引理
\ref{spaces-morphisms-lemma-image-proper-is-proper}）。
于是 $f^*\mathcal{F}_n$ 支撑在闭子概形 $f^{-1}T_n$ 上；
后者在 $A$ 上固有（由《空间的态射》引理
\ref{spaces-morphisms-lemma-composition-proper} 以及 $f$ 的固有性）。
特别地，复合 $f^{-1}T_n \to Y \to Y'$ 是闭态射
（《态射》引理 \ref{morphisms-lemma-image-proper-scheme-closed}）。
设 $T'_n \subset Y'$ 为相应的闭子概形；它包含在开子概形 $Y$ 中，
并且等于 $f^{-1}T_n$，这里把二者视为 $Y$ 的闭子概形。
设 $\mathcal{F}_n'$ 为凝聚 $\mathcal{O}_{Y'}$-模，它对应于
$f^*\mathcal{F}_n$；后者被视为 $Y'$ 上的凝聚模，所用闭浸入为
$f^{-1}T_n = T'_n \subset Y'$。
于是 $(\mathcal{F}_n')$ 是
$\textit{Coh}(Y', I\mathcal{O}_{Y'})$ 的对象。
由射影情形的 Grothendieck 存在定理（《概形的上同调》引理
\ref{coherent-lemma-existence-projective}），存在凝聚
$\mathcal{O}_{Y'}$-模 $\mathcal{F}'$ 以及同构
$(\mathcal{F}')^\wedge \cong (\mathcal{F}'_n)$，
该同构位于 $\textit{Coh}(Y', I\mathcal{O}_{Y'})$ 中。
设 $Z' \subset Y'$ 为 $\mathcal{F}'$ 的概形论支撑。
由于 $\mathcal{F}'/I\mathcal{F}' = \mathcal{F}'_1$，可知
$Z' \cap V(I\mathcal{O}_{Y'}) = T'_1$（作为点集）。
结构态射 $p' : Y' \to \Spec(A)$ 固有，故
$p'(Z' \cap (Y' \setminus Y))$ 在 $\Spec(A)$ 中闭。
若它非空，则它会包含 $V(I)$ 的一点，因为 $I$ 包含在 $A$ 的
Jacobson 根中（《代数》引理
\ref{algebra-lemma-radical-completion}）。
但上面已经看到
$Z' \cap (p')^{-1}V(I) = T'_1 \subset Y$，故可知
$Z' \subset Y$。于是 $\mathcal{F}'|_Y$ 支撑在 $Y$ 的某个
闭子概形上，而该闭子概形在 $A$ 上固有。

\medskip\noindent
设 $\mathcal{K}$ 为切出约化补集 $X \setminus U$ 的拟凝聚理想层。
由《空间的上同调》引理
\ref{spaces-cohomology-lemma-proper-pushforward-coherent}，
$\mathcal{O}_X$-模 $\mathcal{H} = f_*\mathcal{F}'$ 是凝聚的；
并且由引理 \ref{spaces-more-morphisms-lemma-inverse-systems-push-pull}，存在态射
$\alpha : (\mathcal{F}_n) \to \mathcal{H}^\wedge$；它位于范畴
$\textit{Coh}_{\text{支撑固有于 } A}(X, \mathcal{I})$ 中，且其核与余核
被 $\mathcal{K}$ 的某次幂零化。% P09-ERR-0725
设 $Z_0 \subset X$ 为 $\mathcal{H}$ 的概形论支撑。
显然 $|Z_0| \subset f(|Z'|)$。因此
$Z_0 \to \Spec(A)$ 固有（《空间的态射》引理
\ref{spaces-morphisms-lemma-image-proper-is-proper}）。
所以 $\mathcal{H}$ 是
$\textit{Coh}_{\text{支撑固有于 } A}(\mathcal{O}_X)$ 的对象。
由于每个理想层 $\mathcal{K}^e$ 都属于 $\Xi$，
引理 \ref{spaces-more-morphisms-lemma-existence-tricky} 的假设成立，结论随之得到。
\end{proof}

\begin{remark}[展开 Grothendieck 存在定理]
\label{spaces-more-morphisms-remark-reformulate-existence-theorem}
设 $A$ 为关于理想 $I$ 完备的诺特环。
写成 $S = \Spec(A)$ 以及 $S_n = \Spec(A/I^n)$。
设 $X \to S$ 为分离且有限型的代数空间态射。
对 $n \geq 1$，令 $X_n = X \times_S S_n$。图示如下：
$$
\xymatrix{
X_1 \ar[r]_{i_1} \ar[d] & X_2 \ar[r]_{i_2} \ar[d] & X_3 \ar[r] \ar[d] &
\ldots & X \ar[d] \\
S_1 \ar[r] & S_2 \ar[r] & S_3 \ar[r] & \ldots & S
}
$$
在此情形下，考虑系统 $(\mathcal{F}_n, \varphi_n)$，其中
\begin{enumerate}
\item $\mathcal{F}_n$ 是凝聚 $\mathcal{O}_{X_n}$-模；
\item $\varphi_n : i_n^*\mathcal{F}_{n + 1} \to \mathcal{F}_n$
是同构；
\item $\text{Supp}(\mathcal{F}_1)$ 在 $S_1$ 上固有。
\end{enumerate}
定理 \ref{spaces-more-morphisms-theorem-grothendieck-existence} 说明完备化函子
$$
\begin{matrix}
\text{凝聚 }\mathcal{O}_X\text{-模 }\mathcal{F} \\
\text{其支撑固有于 }A
\end{matrix}
\quad
\longrightarrow
\quad
\begin{matrix}
\text{系统 }(\mathcal{F}_n) \\
\text{如上}
\end{matrix}
$$
是范畴等价。在 $X$ 在 $A$ 上固有这一特殊情形中，可以省略
关于支撑的条件。
\end{remark}

\section{Grothendieck 代数化定理}
\label{spaces-more-morphisms-section-algebraization}

\noindent
本节是《概形的上同调》第
\ref{coherent-section-algebraization} 节的类似版本。
不过，本节没有给出带丰沛可逆层的固有代数空间之形变的代数化结果，
因为一个带丰沛可逆层的固有代数空间已经是概形。
我们确实给出了相对维数为 $1$ 的固有代数空间的代数化结果。
第一个结果以闭子概形和有限态射的语言转述 Grothendieck 存在定理。
% P09-ERR-0726

\begin{lemma}
\label{spaces-more-morphisms-lemma-algebraize-formal-closed-subscheme}
设 $A$ 为关于理想 $I$ 完备的诺特环。
写成 $S = \Spec(A)$ 以及 $S_n = \Spec(A/I^n)$。
设 $X \to S$ 为分离且有限型的代数空间态射。
对 $n \geq 1$，令 $X_n = X \times_S S_n$。
设给定代数空间的交换图
$$
\xymatrix{
Z_1 \ar[r] \ar[d] & Z_2 \ar[r] \ar[d] & Z_3 \ar[r] \ar[d] & \ldots \\
X_1 \ar[r]^{i_1} & X_2 \ar[r]^{i_2} & X_3 \ar[r] & \ldots
}
$$
且其中各方块为笛卡尔方块。假设
\begin{enumerate}
\item $Z_1 \to X_1$ 是闭浸入；
\item $Z_1 \to S_1$ 固有。
\end{enumerate}
则存在代数空间的闭浸入 $Z \to X$，使得
$Z_n = Z \times_S S_n$，对所有 $n \geq 1$。
此外，$Z$ 在 $S$ 上固有。
\end{lemma}

\begin{proof}
把竖直态射记为 $j_n : Z_n \to X_n$。
由于陈述中的方块为笛卡尔方块，$j_n$ 到 $X_1$ 的基变换为
$j_1$。因此《空间的极限》引理
\ref{spaces-limits-lemma-check-closed-infinitesimally} 表明
$j_n$ 是闭浸入。
令 $\mathcal{F}_n = j_{n, *}\mathcal{O}_{Z_n}$，于是
$j_n^\sharp$ 是满射 $\mathcal{O}_{X_n} \to \mathcal{F}_n$。
再次使用方块为笛卡尔方块这一事实，可知
$\mathcal{F}_{n + 1}$ 到 $X_n$ 的拉回是 $\mathcal{F}_n$。
因此，由评注 \ref{spaces-more-morphisms-remark-reformulate-existence-theorem} 中转述的
Grothendieck 存在定理，存在映射 $\mathcal{O}_X \to \mathcal{F}$；
它是凝聚 $\mathcal{O}_X$-模之间的映射，并且限制到 $X_n$ 后恢复
$\mathcal{O}_{X_n} \to \mathcal{F}_n$。
此外，$\mathcal{F}$ 的支撑在 $S$ 上固有。
由于完备化函子正合（引理 \ref{spaces-more-morphisms-lemma-exact}），可知
$\mathcal{O}_X \to \mathcal{F}$ 是满射。
因此有 $\mathcal{F} = \mathcal{O}_X/\mathcal{J}$，其中
$\mathcal{J}$ 是某个拟凝聚理想层。
令 $Z = V(\mathcal{J})$，证明即告完成。
\end{proof}

\begin{lemma}
\label{spaces-more-morphisms-lemma-algebraize-formal-algebraic-space-finite-over-proper}
设 $A$ 为关于理想 $I$ 完备的诺特环。
写成 $S = \Spec(A)$ 以及 $S_n = \Spec(A/I^n)$。
设 $X \to S$ 为分离且有限型的代数空间态射。
对 $n \geq 1$，令 $X_n = X \times_S S_n$。
设给定代数空间的交换图
$$
\xymatrix{
Y_1 \ar[r] \ar[d] & Y_2 \ar[r] \ar[d] & Y_3 \ar[r] \ar[d] & \ldots \\
X_1 \ar[r]^{i_1} & X_2 \ar[r]^{i_2} & X_3 \ar[r] & \ldots
}
$$
且其中各方块为笛卡尔方块。假设
\begin{enumerate}
\item $Y_1 \to X_1$ 是有限态射；
\item $Y_1 \to S_1$ 固有。
\end{enumerate}
则存在代数空间的有限态射 $Y \to X$，使得
$Y_n = Y \times_S S_n$，对所有 $n \geq 1$。
此外，$Y$ 在 $S$ 上固有。
\end{lemma}

\begin{proof}
把竖直态射记为 $f_n : Y_n \to X_n$。
由于陈述中的方块为笛卡尔方块，$f_n$ 到 $X_1$ 的基变换为
$f_1$。因此引理 \ref{spaces-more-morphisms-lemma-thicken-property-morphisms-cartesian}
表明 $f_n$ 是有限态射。
令 $\mathcal{F}_n = f_{n, *}\mathcal{O}_{Y_n}$。
使用方块为笛卡尔方块这一事实，可知
$\mathcal{F}_{n + 1}$ 到 $X_n$ 的拉回是 $\mathcal{F}_n$。
因此，由评注 \ref{spaces-more-morphisms-remark-reformulate-existence-theorem} 中转述的
Grothendieck 存在定理，存在凝聚 $\mathcal{O}_X$-模
$\mathcal{F}$，它限制到 $X_n$ 后恢复 $\mathcal{F}_n$。
此外，$\mathcal{F}$ 的支撑在 $S$ 上固有。
由于完备化函子全忠实（定理 \ref{spaces-more-morphisms-theorem-grothendieck-existence}），
可知乘法映射
$\mathcal{F}_n \otimes_{\mathcal{O}_{X_n}} \mathcal{F}_n \to
\mathcal{F}_n$ 相容地给出 $\mathcal{F}$ 上的代数结构。
% P09-ERR-0727
令 $Y = \underline{\Spec}_X(\mathcal{F})$，证明即告完成。
\end{proof}

\begin{lemma}
\label{spaces-more-morphisms-lemma-algebraize-morphism}
设 $A$ 为关于理想 $I$ 完备的诺特环。
写成 $S = \Spec(A)$ 以及 $S_n = \Spec(A/I^n)$。
设 $X$、$Y$ 为 $S$ 上的代数空间。对 $n \geq 1$，令
$X_n = X \times_S S_n$ 以及 $Y_n = Y \times_S S_n$。
设给定相容的交换图系统
$$
\xymatrix{
& & X_{n + 1} \ar[rd] \ar[rr]_{g_{n + 1}} & & Y_{n + 1} \ar[ld] \\
X_n \ar[rru] \ar[rd] \ar[rr]_{g_n} & & Y_n \ar[rru] \ar[ld] & S_{n + 1} \\
& S_n \ar[rru]
}
$$
假设
\begin{enumerate}
\item $X \to S$ 固有；
\item $Y \to S$ 分离且有限型。% P09-ERR-0728
\end{enumerate}
则存在唯一的代数空间态射 $g : X \to Y$（在 $S$ 上），
使得 $g_n$ 是 $g$ 到 $S_n$ 的基变换。
\end{lemma}

\begin{proof}
态射 $(1, g_n) : X_n \to X_n \times_S Y_n$ 是闭浸入，
因为 $Y_n \to S_n$ 分离（《空间的态射》引理
\ref{spaces-morphisms-lemma-section-immersion}）。
因此由引理 \ref{spaces-more-morphisms-lemma-algebraize-formal-closed-subscheme}，
存在闭子空间 $Z \subset X \times_S Y$，它在 $S$ 上固有，
且到 $S_n$ 的基变换恢复
$X_n \subset X_n \times_S Y_n$。第一投影 $p : Z \to X$
是固有态射（因为 $Z$ 在 $S$ 上固有；见《空间的态射》引理
\ref{spaces-morphisms-lemma-universally-closed-permanence}），
并且它到 $S_n$ 的基变换对所有 $n$ 都是同构。
特别地，$p : Z \to X$ 在 $Z$ 的一个开子空间上拟有限，
该开子空间包含 $Z_0$ 的每一点；例如由《空间的态射》引理
\ref{spaces-morphisms-lemma-locally-finite-type-quasi-finite-part} 可得。
% P09-ERR-0729
由于 $Z$ 在 $S$ 上固有，此开邻域就是整个 $Z$。% P09-ERR-0730
因此由 Zariski 主定理可知 $p : Z \to X$ 有限
（例如应用引理 \ref{spaces-more-morphisms-lemma-quasi-finite-separated-pass-through-finite}，
并利用 $Z$ 在 $X$ 上的固有性看出该浸入是闭浸入）。
应用定理 \ref{spaces-more-morphisms-theorem-grothendieck-existence} 的范畴等价，
可知 $p_*\mathcal{O}_Z = \mathcal{O}_X$，因为这在模 $I^n$
之后对所有 $n$ 都成立。因此 $p$ 是同构，并且得到态射
$g$，即复合 $X \cong Z \to Y$。
我们略去唯一性的证明。% P09-ERR-0731
\end{proof}

\begin{remark}
\label{spaces-more-morphisms-remark-weaken-separation-axioms-question}
可以问，在 Grothendieck 代数化定理（采用引理
\ref{spaces-more-morphisms-lemma-algebraize-morphism} 的形式）中，能否对靶采用更弱的分离公理。
更精确地说，设 $A$、$I$、$S$、$S_n$、$X$、$Y$、$X_n$、
$Y_n$ 和 $g_n$ 如引理 \ref{spaces-more-morphisms-lemma-algebraize-morphism} 的陈述，并假设
\begin{enumerate}
\item $X \to S$ 固有；
\item $Y \to S$ 局部有限型。
\end{enumerate}
是否存在代数空间态射 $g : X \to Y$（在 $S$ 上），使得
$g_n$ 是 $g$ 到 $S_n$ 的基变换？一般情形下我们不知道答案；
若你知道，请发送电子邮件至
\href{mailto:stacks.project@gmail.com}{stacks.project@gmail.com}。
若 $Y \to S$ 分离，则由该引理结论成立（可立即约化到
$X$ 在 $S$ 上有限型的情形，具体是选择一个包含 $g_1$ 的像的拟紧开集）。
% P09-ERR-0732
若只假设 $Y \to S$ 拟分离，结论也成立。首先，如前一样，
可假设 $Y$ 既拟紧又拟分离。然后可使用 \cite{Bhatt-Algebraize}
或从 \cite{Hall-Rydh-coherent} 将 $(g_n)$ 代数化。% P09-ERR-0733
具体而言，为应用第一篇参考文献，使用
$$
D_{perf}(X) \to \lim D_{perf}(X_n)
\xrightarrow{\lim Lg_n^*}
\lim D_{perf}(Y_n) = D_{perf}(Y)
$$
其中最后一步使用固有代数空间 $Y$ 在 $R$ 上的导出范畴的
Grothendieck 存在结果（参见《空间上的平坦性》评注
\ref{spaces-flat-remark-correct-generality}）。
所引论文表明，此箭头按所需确定态射 $Y \to X$。% P09-ERR-0734
为应用第二篇参考文献，对凝聚模使用相同论证：
$$
\textit{Coh}(\mathcal{O}_X) \to \lim \textit{Coh}(\mathcal{O}_{X_n})
\xrightarrow{\lim g_n^*}
\lim \textit{Coh}(\mathcal{O}_{Y_n}) = \textit{Coh}(\mathcal{O}_Y)
$$
其中最后一个等式是 Grothendieck 存在定理（定理
\ref{spaces-more-morphisms-theorem-grothendieck-existence}）的推论。
第二篇参考文献说明，此函子对应于态射 $Y \to X$（在 $R$ 上）。
% P09-ERR-0735
若以后需要这一推广，我们将在此处精确陈述并仔细证明结果。
\end{remark}

\begin{lemma}
\label{spaces-more-morphisms-lemma-formal-algebraic-space-proper-reldim-1}
设 $(A, \mathfrak m, \kappa)$ 为完备局部诺特环。
令 $S = \Spec(A)$ 以及 $S_n = \Spec(A/\mathfrak m^n)$。
考虑代数空间的交换图
$$
\xymatrix{
X_1 \ar[r]_{i_1} \ar[d] & X_2 \ar[r]_{i_2} \ar[d] & X_3 \ar[r] \ar[d] &
\ldots \\
S_1 \ar[r] & S_2 \ar[r] & S_3 \ar[r] & \ldots
}
$$
且其中各方块为笛卡尔方块。若 $\dim(X_1) \leq 1$，
则存在概形的射影态射 $X \to S$ 以及同构
$X_n \cong X \times_S S_n$，且这些同构与 $i_n$ 相容。
% P09-ERR-0736
\end{lemma}

\begin{proof}
由《域上的空间》引理
\ref{spaces-over-fields-lemma-codim-1-point-in-schematic-locus}，
代数空间 $X_1$ 是概形。因此，$X_1$ 是维数
$\leq 1$ 的固有概形（在 $\kappa$ 上）。由《簇》引理
\ref{varieties-lemma-dim-1-proper-projective}，可知 $X_1$ 在
$\kappa$ 上 H-射影。设 $\mathcal{L}_1$ 为 $X_1$ 上的丰沛可逆层。

\medskip\noindent
我们将证明 $\mathcal{L}_1$ 可提升为可逆层的相容系统
$\{\mathcal{L}_n\}$；这些层定义在 $\{X_n\}$ 上。注意，由引理 \ref{spaces-more-morphisms-lemma-thickening-scheme}，
$X_n$ 也是概形。回顾 $X_1 \to X_n$ 诱导底拓扑空间之间的同胚。
在证明的其余部分，我们不区分 $X_n$ 上的层和 $X_1$ 上的层。
假设给定 $\mathcal{L}_n$ 到 $X_n$ 的一个提升。% P09-ERR-0737
我们考虑正合序列
$$
1 \to
(1 + \mathfrak m^n\mathcal{O}_{X_{n + 1}})^* \to
\mathcal{O}_{X_{n + 1}}^* \to \mathcal{O}_{X_n}^* \to 1
$$
它由 $X_{n + 1}$ 上的层组成。
$\mathcal{L}_n$ 在 $H^1(X_n, \mathcal{O}_{X_n}^*)$ 中的类（见
《上同调》引理 \ref{cohomology-lemma-h1-invertible}）可提升为
$H^1(X_{n + 1}, \mathcal{O}_{X_{n + 1}}^*)$ 的元素，当且仅当
$H^2(X_{n + 1}, (1 + \mathfrak m^n\mathcal{O}_{X_{n + 1}})^*)$
中的障碍为零。由于 $X_1$ 是维数 $\leq 1$ 的诺特概形，
此上同调群消失（《上同调》命题
\ref{cohomology-proposition-vanishing-Noetherian}）。

\medskip\noindent
由 Grothendieck 代数化定理（《概形的上同调》定理
\ref{coherent-theorem-algebraization}），我们得到概形的射影态射
$X \to S = \Spec(A)$ 以及相容的同构系统
$X_n = S_n \times_S X$。
\end{proof}

\begin{lemma}
\label{spaces-more-morphisms-lemma-projective-over-complete}
设 $(A, \mathfrak m, \kappa)$ 为完备诺特局部环。
设 $X$ 为 $\Spec(A)$ 上的代数空间。若
$X \to \Spec(A)$ 固有且 $\dim(X_\kappa) \leq 1$，则
$X$ 是 $A$ 上的射影概形。
\end{lemma}

\begin{proof}
令 $X_n = X \times_{\Spec(A)} \Spec(A/\mathfrak m^n)$。
由引理 \ref{spaces-more-morphisms-lemma-formal-algebraic-space-proper-reldim-1}，
存在射影态射 $Y \to \Spec(A)$ 以及相容的同构
$Y \times_{\Spec(A)} \Spec(A/\mathfrak m^n) \cong
X \times_{\Spec(A)} \Spec(A/\mathfrak m^n)$。
由引理 \ref{spaces-more-morphisms-lemma-algebraize-morphism}，可知 $X \cong Y$，
证明即告完成。
\end{proof}

\section{正则浸入}
\label{spaces-more-morphisms-section-regular-immersions}

\noindent
本节是《除子》第 \ref{divisors-section-regular-immersions} 节
针对代数空间态射的类似版本。建议读者先在该节了解概形的正则浸入。

\medskip\noindent
在《除子》第 \ref{divisors-section-regular-immersions} 节中，
我们为概形态射定义了四类正则浸入。据我们所知，其中只有三类
关于靶在平展拓扑下局部；通常意义下的正则浸入仍然不是。
正因如此，我们实际上无法为代数空间态射定义正则浸入。
（这类麻烦促使 Grothendieck 及其学派以 Koszul-正则浸入取代
正则浸入的原始概念；见 \cite[Exposee VII, Definition 1.4]{SGA6}。）
% P09-ERR-0738 P09-ERR-0739
不过，我们可以定义 Koszul-正则、$H_1$-正则和拟正则浸入。
还要指出，由于 Koszul-正则浸入不在任意基变换下保持，
我们不能采用《空间的态射》第
\ref{spaces-morphisms-section-representable} 节的策略来定义它们。
类似地，由于 Koszul-正则浸入不关于源平展局部，我们也不能使用
《空间的态射》引理
\ref{spaces-morphisms-lemma-local-source-target} 来定义它们。
我们改以下面的引理代替该引理。

\begin{lemma}
\label{spaces-more-morphisms-lemma-representable-etale-local-target}
设 $\mathcal{P}$ 为概形态射的一项性质，且它关于靶平展局部。
设 $S$ 为概形。设 $f : X \to Y$ 为 $S$ 上代数空间的可表态射。
考虑交换图
$$
\xymatrix{
X \times_Y V \ar[d] \ar[r] & V \ar[d] \\
X \ar[r]^f & Y
}
$$
其中 $V$ 为概形，且 $V \to Y$ 平展。以下条件等价：
\begin{enumerate}
\item 对任意如上图，投影 $X \times_Y V \to V$ 具有性质
$\mathcal{P}$；
\item 对某个如上图，若 $V \to Y$ 满，则投影
$X \times_Y V \to V$ 具有性质 $\mathcal{P}$。
\end{enumerate}
若 $X$ 和 $Y$ 可表，则这些条件还等价于 $f$（作为概形态射）
具有性质 $\mathcal{P}$。
\end{lemma}

\begin{proof}
先证明 (1) 与 (2) 等价。蕴含 (1) $\Rightarrow$ (2) 是立即的。
假设
$$
\xymatrix{
X \times_Y V \ar[d] \ar[r] & V \ar[d] \\
X \ar[r]^f & Y
}
\quad\quad
\xymatrix{
X \times_Y V' \ar[d] \ar[r] & V' \ar[d] \\
X \ar[r]^f & Y
}
$$
是引理中的两个图。假设 $V \to Y$ 满，且
$X \times_Y V \to V$ 具有性质 $\mathcal{P}$。
为证明 (2) 蕴含 (1)，我们必须证明
$X \times_Y V' \to V'$ 具有 $\mathcal{P}$。为此，考虑图
$$
\xymatrix{
X \times_Y V \ar[d] &
(X \times_Y V) \times_X (X \times_Y V') \ar[l] \ar[d] \ar[r] &
X \times_Y V' \ar[d] \\
V &
V \times_Y V' \ar[l] \ar[r] &
V'
}
$$
由我们关于 $\mathcal{P}$ 在源上平展局部的假设，% P09-ERR-0740
可知 $\mathcal{P}$ 在平展基变换下保持；见《下降》引理
\ref{descent-lemma-pullback-property-local-target}。
因此，若左边的竖直箭头具有 $\mathcal{P}$，则中间的竖直箭头也如此。
% P09-ERR-0741
由于 $U \times_X U' \to U'$ 满且平展（因而给出 $U'$ 的平展覆盖），
% P09-ERR-0742
又因为假设 $\mathcal{P}$ 关于靶在平展拓扑下局部，
故这蕴含左边的竖直箭头具有 $\mathcal{P}$。% P09-ERR-0743

\medskip\noindent
若 $X$ 和 $Y$ 可表，则可取 $\text{id}_Y : Y \to Y$ 作为平展覆盖，
从而看出引理的最后一个陈述成立。
\end{proof}

\noindent
注意，“是 Koszul-正则（相应地，$H_1$-正则、拟正则）浸入”
是概形态射的一项性质，并且它关于靶 fpqc 局部；见《下降》引理
\ref{descent-lemma-descending-property-regular-immersion}。
因此，下面的定义是有意义的。

\begin{definition}
\label{spaces-more-morphisms-definition-regular-immersion}
设 $S$ 为概形。设 $i : X \to Y$ 为 $S$ 上代数空间的态射。
\begin{enumerate}
\item 称 $i$ 为{\it Koszul-正则浸入}，如果 $i$ 可表，并且引理
\ref{spaces-more-morphisms-lemma-representable-etale-local-target} 的等价条件对
$\mathcal{P}(f) =$“$f$ 是 Koszul-正则浸入”成立。
\item 称 $i$ 为{\it $H_1$-正则浸入}，如果 $i$ 可表，并且引理
\ref{spaces-more-morphisms-lemma-representable-etale-local-target} 的等价条件对
$\mathcal{P}(f) =$“$f$ 是 $H_1$-正则浸入”成立。
\item 称 $i$ 为{\it 拟正则浸入}，如果 $i$ 可表，并且引理
\ref{spaces-more-morphisms-lemma-representable-etale-local-target} 的等价条件对
$\mathcal{P}(f) =$“$f$ 是拟正则浸入”成立。
\end{enumerate}
\end{definition}

\begin{lemma}
\label{spaces-more-morphisms-lemma-regular-quasi-regular-immersion}
设 $S$ 为概形。设 $i : Z \to X$ 为 $S$ 上代数空间的浸入。
我们有以下蕴含：
$i$ 是 Koszul-正则 $\Rightarrow$
$i$ 是 $H_1$-正则 $\Rightarrow$
$i$ 是拟正则。
\end{lemma}

\begin{proof}
由定义，本引理立即约化为《除子》引理
\ref{divisors-lemma-regular-quasi-regular-immersion}。
\end{proof}

\begin{lemma}
\label{spaces-more-morphisms-lemma-regular-immersion-noetherian}
设 $S$ 为概形。设 $i : Z \to X$ 为 $S$ 上代数空间的浸入。
假设 $X$ 局部诺特。则
$i$ 是 Koszul-正则 $\Leftrightarrow$
$i$ 是 $H_1$-正则 $\Leftrightarrow$
$i$ 是拟正则。
\end{lemma}

\begin{proof}
由定义 \ref{spaces-more-morphisms-definition-regular-immersion}
（以及《空间的性质》第
\ref{spaces-properties-section-types-properties} 节中局部诺特代数空间的定义），
这立即转化为概形情形，即《除子》引理
\ref{divisors-lemma-regular-immersion-noetherian}。
\end{proof}

\begin{lemma}
\label{spaces-more-morphisms-lemma-flat-base-change-regular-immersion}
\begin{slogan}
正则浸入在平坦基变换下稳定。
\end{slogan}
设 $S$ 为概形。设 $i : Z \to X$ 为代数空间的
Koszul-正则、$H_1$-正则或拟正则浸入（在 $S$ 上）。
设 $X' \to X$ 为 $S$ 上代数空间的平坦态射。
则基变换 $i' : Z \times_X X' \to X'$ 是 Koszul-正则、
$H_1$-正则或拟正则浸入。
\end{lemma}

\begin{proof}
由定义 \ref{spaces-more-morphisms-definition-regular-immersion}
（以及《空间的态射》第
\ref{spaces-morphisms-section-flat} 节中代数空间平坦态射的定义），
本引理约化为概形情形；见《除子》引理
\ref{divisors-lemma-flat-base-change-regular-immersion}。
\end{proof}

\begin{lemma}
\label{spaces-more-morphisms-lemma-quasi-regular-immersion}
设 $S$ 为概形。设 $i : Z \to X$ 为 $S$ 上代数空间的浸入。
则 $i$ 是拟正则浸入，当且仅当以下条件成立：
\begin{enumerate}
\item $i$ 局部有限表示；
\item 余法层 $\mathcal{C}_{Z/X}$ 有限局部自由；
\item 映射 (\ref{spaces-more-morphisms-equation-conormal-algebra-quotient}) 是同构。
\end{enumerate}
\end{lemma}

\begin{proof}
通过平展局部化，这由概形情形推出（《除子》引理
\ref{divisors-lemma-quasi-regular-immersion}）；使用定义
\ref{spaces-more-morphisms-definition-regular-immersion} 和引理
\ref{spaces-more-morphisms-lemma-etale-conormal-algebra}。
\end{proof}

\begin{lemma}
\label{spaces-more-morphisms-lemma-transitivity-conormal-quasi-regular}
设 $S$ 为概形。设 $Z \to Y \to X$ 为 $S$ 上代数空间的浸入。
假设 $Z \to Y$ 是 $H_1$-正则的。则引理
\ref{spaces-more-morphisms-lemma-transitivity-conormal} 的典范序列
$$
0 \to i^*\mathcal{C}_{Y/X} \to
\mathcal{C}_{Z/X} \to
\mathcal{C}_{Z/Y} \to 0
$$
正合，并且（平展）局部分裂。
\end{lemma}

\begin{proof}
由于 $\mathcal{C}_{Z/Y}$ 有限局部自由（见引理
\ref{spaces-more-morphisms-lemma-quasi-regular-immersion} 和引理
\ref{spaces-more-morphisms-lemma-regular-quasi-regular-immersion}），只须证明该序列正合。
因为该序列一般已经右正合，只须说明第一个映射是单射。
在 $X$ 上作平展局部化后，这约化为概形情形；见《除子》引理
\ref{divisors-lemma-transitivity-conormal-quasi-regular}。
\end{proof}

\noindent
拟正则浸入的复合未必拟正则；见《代数》评注
\ref{algebra-remark-join-quasi-regular-sequences}。
其他类型的正则浸入在复合下保持。

\begin{lemma}
\label{spaces-more-morphisms-lemma-composition-regular-immersion}
设 $S$ 为概形。设 $i : Z \to Y$ 和 $j : Y \to X$ 为
$S$ 上代数空间的浸入。
\begin{enumerate}
\item 若 $i$ 和 $j$ 都是 Koszul-正则浸入，则 $j \circ i$ 也是。
\item 若 $i$ 和 $j$ 都是 $H_1$-正则浸入，则 $j \circ i$ 也是。
\item 若 $i$ 是 $H_1$-正则浸入且 $j$ 是拟正则浸入，
则 $j \circ i$ 是拟正则浸入。
\end{enumerate}
\end{lemma}

\begin{proof}
这由概形情形立即推出；见《除子》引理
\ref{divisors-lemma-composition-regular-immersion}。
\end{proof}

\begin{lemma}
\label{spaces-more-morphisms-lemma-permanence-regular-immersion}
设 $S$ 为概形。设 $i : Z \to Y$ 和 $j : Y \to X$ 为
$S$ 上代数空间的浸入。假设 $j$ 局部有限表示，并且引理
\ref{spaces-more-morphisms-lemma-transitivity-conormal} 的序列
$$
0 \to i^*\mathcal{C}_{Y/X} \to
\mathcal{C}_{Z/X} \to
\mathcal{C}_{Z/Y} \to 0
$$
正合且局部分裂。
\begin{enumerate}
\item 若 $j \circ i$ 是拟正则浸入，则 $i$ 也是。
\item 若 $j \circ i$ 是 $H_1$-正则浸入，则 $i$ 也是。
\item 若 $j$ 和 $j \circ i$ 都是 Koszul-正则浸入，则 $i$ 也是。
\end{enumerate}
\end{lemma}

\begin{proof}
这由概形情形立即推出；见《除子》引理
\ref{divisors-lemma-permanence-regular-immersion}。
\end{proof}

\begin{lemma}
\label{spaces-more-morphisms-lemma-extra-permanence-regular-immersion-noetherian}
设 $S$ 为概形。设 $i : Z \to Y$ 和 $j : Y \to X$ 为
$S$ 上代数空间的浸入。假设 $X$ 局部诺特。以下条件等价：
\begin{enumerate}
\item $i$ 和 $j$ 是 Koszul 正则浸入；
\item $i$ 和 $j \circ i$ 是 Koszul 正则浸入；
\item $j \circ i$ 是 Koszul 正则浸入，且余法序列
$$
0 \to i^*\mathcal{C}_{Y/X} \to
\mathcal{C}_{Z/X} \to
\mathcal{C}_{Z/Y} \to 0
$$
正合且局部分裂。
\end{enumerate}
% P09-ERR-0744
\end{lemma}

\begin{proof}
这由概形情形立即推出；见《除子》引理
\ref{divisors-lemma-extra-permanence-regular-immersion-noetherian}。
\end{proof}

\section{相对拟凝聚性}
\label{spaces-more-morphisms-section-relative-pseudo-coherence}

\noindent
本节是《态射进阶》第
\ref{more-morphisms-section-relative-pseudo-coherence} 节的类似版本。
不过，在处理代数空间的这些材料时，我们决定只使用导出范畴中
上同调层拟凝聚的对象。理由有二：(1) 这大大简化了叙述；
(2) 我们目前用不到更一般的概念。

\begin{remark}
\label{spaces-more-morphisms-remark-match-relative-pseudo-coherence}
设 $S$ 为概形。设 $f : X \to Y$ 为 $S$ 上可表代数空间的
局部有限型态射。设 $f_0 : X_0 \to Y_0$ 为表示 $f$ 的概形态射
（这一记号略显别扭，但只是暂用）。则 $f_0$ 局部有限型。
若 $E$ 是 $D_\QCoh(\mathcal{O}_X)$ 的对象，则 $E$ 是
唯一对象 $E_0$ 的拉回，该对象属于 $D_\QCoh(\mathcal{O}_{X_0})$；见
《空间的导出范畴》引理
\ref{spaces-perfect-lemma-derived-quasi-coherent-small-etale-site}。
在此情形下，“$E$ 是 $m$-拟凝聚的（相对于 $Y$）”意指
“$E_0$ 是 $m$-拟凝聚的（相对于 $Y_0$）”，其定义见《态射进阶》第
\ref{more-morphisms-section-relative-pseudo-coherence} 节。
\end{remark}

\begin{lemma}
\label{spaces-more-morphisms-lemma-qcoh-relative-pseudo-coherence-characterize}
设 $S$ 为概形。设 $f : X \to Y$ 为 $S$ 上代数空间的
局部有限型态射。设 $m \in \mathbf{Z}$。设
$E \in D_\QCoh(\mathcal{O}_X)$。采用评注
\ref{spaces-more-morphisms-remark-match-relative-pseudo-coherence} 中说明的记号，
以下条件等价：
\begin{enumerate}
\item 对每个交换图
$$
\xymatrix{
U \ar[d] \ar[r] & V \ar[d] \\
X \ar[r] & Y
}
$$
其中 $U$、$V$ 为概形且竖直箭头平展，复形 $E|_U$ 是
$m$-拟凝聚的（相对于 $V$）；
\item 对某个如 (1) 的交换图，若 $U \to X$ 满，则复形
$E|_U$ 是 $m$-拟凝聚的（相对于 $V$）；
\item 对每个如 (1) 的交换图，若 $U$ 和 $V$ 仿射，则复形
$R\Gamma(U, E)$（作为 $\mathcal{O}_X(U)$-模）是
$m$-拟凝聚的（相对于 $\mathcal{O}_Y(V)$）。
\end{enumerate}
\end{lemma}

\begin{proof}
由《态射进阶》引理
\ref{more-morphisms-lemma-qcoh-relative-pseudo-coherence-characterize}，
(1) 蕴含 (3)。

\medskip\noindent
假设 (3)。任取一个如 (1) 的交换图，其中 $U \to X$ 满。
选取仿射开覆盖 $V = \bigcup V_j$，并选取仿射开覆盖
$(U \to V)^{-1}(V_j) = \bigcup U_{ij}$。由 (3) 及《态射进阶》引理
\ref{more-morphisms-lemma-qcoh-relative-pseudo-coherence-characterize}，
可知 $E|_U$ 是 $m$-拟凝聚的（相对于 $V$）。因此 (3) 蕴含 (2)。

\medskip\noindent
假设 (2)。选取交换图
$$
\xymatrix{
U \ar[d] \ar[r] & V \ar[d] \\
X \ar[r] & Y
}
$$
其中 $U$、$V$ 为概形，竖直箭头平展，态射 $U \to X$ 满，
且 $E|_U$ 是 $m$-拟凝聚的（相对于 $V$）。接着，设给定第二个交换图
$$
\xymatrix{
U' \ar[d] \ar[r] & V' \ar[d] \\
X \ar[r] & Y
}
$$
其中竖直箭头平展，且 $U', V'$ 为概形。我们要证明
$E|_{U'}$ 是 $m$-拟凝聚的（相对于 $V'$）。
态射 $U'' = U \times_X U' \to U'$ 满且平展，而
$U'' \to V'$ 分解通过 $V'' = V' \times_Y V$；后者在 $V'$ 上平展。
因此只须证明 $E|_{U''}$ 是 $m$-拟凝聚的（相对于 $V''$）；见
《态射进阶》引理
\ref{more-morphisms-lemma-relative-pseudo-coherent-descends-fppf} 和
\ref{more-morphisms-lemma-relative-pseudo-coherent-post-compose}。
再使用一次第二个引理，只须证明 $E|_{U''}$ 是 $m$-拟凝聚的
（相对于 $V$）。由《态射进阶》引理
\ref{more-morphisms-lemma-pull-relative-pseudo-coherent}，
以及概形的平展态射拟凝聚这一事实，此断言成立；后者见《态射进阶》引理
\ref{more-morphisms-lemma-flat-finite-presentation-pseudo-coherent}。
\end{proof}

\begin{definition}
\label{spaces-more-morphisms-definition-relative-pseudo-coherence}
设 $S$ 为概形。设 $f : X \to Y$ 为 $S$ 上代数空间的
局部有限型态射。设 $E$ 为 $D_\QCoh(\mathcal{O}_X)$ 的对象。
设 $\mathcal{F}$ 为拟凝聚 $\mathcal{O}_X$-模。固定
$m \in \mathbf{Z}$。
\begin{enumerate}
\item 称 $E$ {\it 为 $m$-拟凝聚的（相对于 $Y$）}，如果引理
\ref{spaces-more-morphisms-lemma-qcoh-relative-pseudo-coherence-characterize} 的等价条件成立。
\item 称 $E$ {\it 相对于 $Y$ 拟凝聚}，如果 $E$ 是
$m$-拟凝聚的（相对于 $Y$），对所有 $m \in \mathbf{Z}$。
\item 称 $\mathcal{F}$ {\it 为 $m$-拟凝聚的（相对于 $Y$）}，
如果把 $\mathcal{F}$ 视为 $D_\QCoh(\mathcal{O}_X)$ 的对象时，
它是 $m$-拟凝聚的（相对于 $Y$）。
\item 称 $\mathcal{F}$ {\it 相对于 $Y$ 拟凝聚}，
如果把 $\mathcal{F}$ 视为 $D_\QCoh(\mathcal{O}_X)$ 的对象时，
它相对于 $Y$ 拟凝聚。
\end{enumerate}
\end{definition}

\noindent
复形相对于基拟凝聚的大多数性质将立即由概形情形的相应性质推出。
我们会在需要时在此加入相关引理。

\begin{lemma}
\label{spaces-more-morphisms-lemma-relative-pseudo-coherent-is-moot}
设 $S$ 为概形。设 $f : X \to Y$ 为 $S$ 上代数空间的态射。
设 $E$ 为 $D_\QCoh(\mathcal{O}_X)$ 的对象。% P09-ERR-0745
若 $f$ 平坦且局部有限表示，则以下条件等价：
\begin{enumerate}
\item $E$ 相对于 $Y$ 拟凝聚；
\item $E$ 在 $X$ 上拟凝聚。
\end{enumerate}
\end{lemma}

\begin{proof}
由平展局部化和定义，可以假设 $X$ 和 $Y$ 是概形。
对概形情形，此结论由《态射进阶》引理
\ref{more-morphisms-lemma-check-relative-pseudo-coherence-on-charts} 推出。
\end{proof}

\section{拟凝聚态射}
\label{spaces-more-morphisms-section-pseudo-coherent}

\noindent
本节是《态射进阶》第
\ref{more-morphisms-section-pseudo-coherent} 节针对概形态射的类似版本。
建议读者先在该节了解概形的拟凝聚态射。

\medskip\noindent
概形态射的“拟凝聚”性质关于源和靶平展局部。为此，使用
《态射进阶》引理
\ref{more-morphisms-lemma-descending-property-pseudo-coherent} 和
\ref{more-morphisms-lemma-pseudo-coherent-fppf-local-source}，
以及《下降》引理 \ref{descent-lemma-etale-local-source-target}。
由《空间的态射》引理
\ref{spaces-morphisms-lemma-local-source-target}，
可以如下定义代数空间拟凝聚态射；当所涉及的代数空间可表时，
它与《态射进阶》第
\ref{more-morphisms-section-pseudo-coherent} 节已经定义的概念一致。

\begin{definition}
\label{spaces-more-morphisms-definition-pseudo-coherent}
设 $S$ 为概形。设 $f : X \to Y$ 为 $S$ 上代数空间的态射。
\begin{enumerate}
\item 称 $f$ {\it 拟凝聚}，如果《空间的态射》引理
\ref{spaces-morphisms-lemma-local-source-target} 的等价条件对
$\mathcal{P} =$“拟凝聚”成立。
\item 设 $x \in |X|$。称 $f$ {\it 在 $x$ 处拟凝聚}，如果存在
开邻域 $X' \subset X$（它含有 $x$），使得
$f|_{X'} : X' \to Y$ 拟凝聚。
\end{enumerate}
\end{definition}

\noindent
注意，拟凝聚态射的基变换一般并不拟凝聚。

\begin{lemma}
\label{spaces-more-morphisms-lemma-flat-base-change-pseudo-coherent}
拟凝聚态射的平坦基变换是拟凝聚的。
\end{lemma}

\begin{proof}
从略。提示：使用本引理的概形版本；见《态射进阶》引理
\ref{more-morphisms-lemma-flat-base-change-pseudo-coherent}。
\end{proof}

\begin{lemma}
\label{spaces-more-morphisms-lemma-composition-pseudo-coherent}
拟凝聚态射的复合是拟凝聚的。
\end{lemma}

\begin{proof}
从略。提示：使用本引理的概形版本；见《态射进阶》引理
\ref{more-morphisms-lemma-composition-pseudo-coherent}。
\end{proof}

\begin{lemma}
\label{spaces-more-morphisms-lemma-pseudo-coherent-finite-presentation}
拟凝聚态射局部有限表示。
\end{lemma}

\begin{proof}
这由定义立即得到。
\end{proof}

\begin{lemma}
\label{spaces-more-morphisms-lemma-flat-finite-presentation-pseudo-coherent}
局部有限表示的平坦态射是拟凝聚的。
\end{lemma}

\begin{proof}
从略。提示：使用本引理的概形版本；见《态射进阶》引理
\ref{more-morphisms-lemma-flat-finite-presentation-pseudo-coherent}。
\end{proof}

\begin{lemma}
\label{spaces-more-morphisms-lemma-permanence-pseudo-coherent}
设 $f : X \to Y$ 为代数空间的态射，并且它在某个基代数空间
$B$ 上拟凝聚。则 $f$ 拟凝聚。
\end{lemma}

\begin{proof}
从略。提示：使用本引理的概形版本；见《态射进阶》引理
\ref{more-morphisms-lemma-permanence-pseudo-coherent}。
\end{proof}

\begin{lemma}
\label{spaces-more-morphisms-lemma-Noetherian-pseudo-coherent}
设 $S$ 为概形。设 $f : X \to Y$ 为 $S$ 上代数空间的态射。
若 $Y$ 局部诺特，则 $f$ 拟凝聚，当且仅当 $f$ 局部有限型。
\end{lemma}

\begin{proof}
从略。提示：使用本引理的概形版本；见《态射进阶》引理
\ref{more-morphisms-lemma-Noetherian-pseudo-coherent}。
\end{proof}

\section{完美态射}
\label{spaces-more-morphisms-section-perfect}

\noindent
本节是《态射进阶》第
\ref{more-morphisms-section-perfect} 节关于概形态射的内容
在代数空间情形的对应版本。建议读者先阅读该节关于概形完美态射的内容。

\medskip\noindent
概形态射的“完美”性质在源与靶上平展局部。为看出这一点，使用
《态射进阶》引理
\ref{more-morphisms-lemma-descending-property-perfect} 和
\ref{more-morphisms-lemma-perfect-fppf-local-source}
以及《下降》引理
\ref{descent-lemma-etale-local-source-target}。
由《空间的态射》引理
\ref{spaces-morphisms-lemma-local-source-target}，
可以如下定义代数空间的完美态射；当所涉及的代数空间可表时，
它与《态射进阶》第 \ref{more-morphisms-section-perfect} 节
已经定义的概念一致。

\begin{definition}
\label{spaces-more-morphisms-definition-perfect}
设 $S$ 为概形。设 $f : X \to Y$ 为 $S$ 上代数空间的态射。
\begin{enumerate}
\item 称 $f$ {\it 完美}，如果《空间的态射》引理
\ref{spaces-morphisms-lemma-local-source-target} 的等价条件对
$\mathcal{P} =$“完美”成立。
\item 设 $x \in |X|$。称 $f$ {\it 在 $x$ 处完美}，如果存在
开邻域 $X' \subset X$，它含有 $x$，并且使得
$f|_{X'} : X' \to Y$ 完美。
\end{enumerate}
\end{definition}

\noindent
注意，完美态射是拟凝聚的，因而局部有限表示。要注意，
完美态射的基变换一般并不完美。

\begin{lemma}
\label{spaces-more-morphisms-lemma-flat-base-change-perfect}
完美态射的平坦基变换是完美的。
\end{lemma}

\begin{proof}
从略。提示：使用本引理的概形版本；见《态射进阶》引理
\ref{more-morphisms-lemma-flat-base-change-perfect}。
\end{proof}

\begin{lemma}
\label{spaces-more-morphisms-lemma-composition-perfect}
完美态射的复合是完美的。
\end{lemma}

\begin{proof}
从略。提示：使用本引理的概形版本；见《态射进阶》引理
\ref{more-morphisms-lemma-composition-perfect}。
\end{proof}

\begin{lemma}
\label{spaces-more-morphisms-lemma-flat-finite-presentation-perfect}
设 $S$ 为概形。设 $f : X \to Y$ 为 $S$ 上代数空间的态射。
下列条件等价：
\begin{enumerate}
\item $f$ 平坦且完美；
\item $f$ 平坦且局部有限表示。
\end{enumerate}
\end{lemma}

\begin{proof}
从略。提示：使用本引理的概形版本；见《态射进阶》引理
\ref{more-morphisms-lemma-flat-finite-presentation-perfect}。
\end{proof}

\begin{lemma}
\label{spaces-more-morphisms-lemma-perfect-proper-perfect-direct-image}
设 $S$ 为概形。设 $Y$ 为 $S$ 上的诺特代数空间。
设 $f : X \to Y$ 为代数空间的完美固有态射。
设 $E \in D(\mathcal{O}_X)$ 为完美的。则
$Rf_*E$ 是 $D(\mathcal{O}_Y)$ 的完美对象。
\end{lemma}

\begin{proof}
我们断言《空间的导出范畴》引理
\ref{spaces-perfect-lemma-perfect-direct-image} 适用。
条件 (1) 和 (2) 是立即的。条件 (3) 在 $X$ 上是局部的。
因此可以假设 $X$ 和 $Y$ 仿射，并且 $E$ 由一个
$\mathcal{O}_X$-模的严格完美复形表示。
所以只需证明：在平展位点上，$\mathcal{O}_X$ 作为
$f^{-1}\mathcal{O}_Y$-模层具有有限 Tor 维数。由《空间的导出范畴》引理
\ref{spaces-perfect-lemma-tor-dimension-rel}，只需在 Zariski 位点上检验这一点。
对于概形的有限型态射，由《态射进阶》引理
\ref{more-morphisms-lemma-check-perfect-stalks}，
这等价于完美性。
\end{proof}

\section{局部完全交态射}
\label{spaces-more-morphisms-section-lci}

\noindent
本节是《态射进阶》第 \ref{more-morphisms-section-lci} 节
关于概形态射的内容在代数空间情形的对应版本。
建议读者先阅读该节关于概形局部完全交态射的内容。

\medskip\noindent
概形态射的“为局部完全交态射”这一性质在源与靶上平展局部。
为看出这一点，使用《态射进阶》引理
\ref{more-morphisms-lemma-descending-property-lci} 和
\ref{more-morphisms-lemma-lci-syntomic-local-source}
以及《下降》引理 \ref{descent-lemma-etale-local-source-target}。
由《空间的态射》引理
\ref{spaces-morphisms-lemma-local-source-target}，
可以如下定义代数空间的局部完全交态射；当所涉及的代数空间可表时，
它与《态射进阶》第 \ref{more-morphisms-section-lci} 节
已经定义的概念一致。

\begin{definition}
\label{spaces-more-morphisms-definition-lci}
设 $S$ 为概形。设 $f : X \to Y$ 为 $S$ 上代数空间的态射。
\begin{enumerate}
\item 称 $f$ 是一个 {\it Koszul 态射}，或称 $f$ 是一个
{\it 局部完全交态射}，如果《空间的态射》引理
\ref{spaces-morphisms-lemma-local-source-target} 的等价条件对
$\mathcal{P}(f) =$“$f$ 是一个局部完全交态射”成立。
\item 设 $x \in |X|$。称 $f$ {\it 在 $x$ 处为 Koszul}，如果存在
开邻域 $X' \subset X$，它含有 $x$，并且使得
$f|_{X'} : X' \to Y$ 是局部完全交态射。
\end{enumerate}
\end{definition}

\noindent
在某种意义下，下面的引理给出了局部完全交态射的定义性质。

\begin{lemma}
\label{spaces-more-morphisms-lemma-lci}
设 $S$ 为概形。设 $f : X \to Y$ 为 $S$ 上代数空间的
局部完全交态射。设 $P$ 为在 $Y$ 上光滑的代数空间。
设 $U \to X$ 为代数空间的平展态射，并设
% P09-ERR-0746
$i : U \to P$ 为 $Y$ 上代数空间的浸入。图示如下：
$$
\xymatrix{
X \ar[rd] & U \ar[l] \ar[d] \ar[r]_i & P \ar[ld] \\
& Y
}
$$
则 $i$ 是代数空间的 Koszul-正则浸入。
\end{lemma}

\begin{proof}
取概形 $V$ 和满平展态射 $V \to Y$。
取概形 $W$ 和满平展态射 $W \to P \times_Y V$。
令 $U' = U \times_P W$；它是一个在 $U$ 上平展的概形。
我们必须证明 $U' \to W$ 是概形的 Koszul-正则浸入；见
定义 \ref{spaces-more-morphisms-definition-regular-immersion}。
由上面的定义 \ref{spaces-more-morphisms-definition-lci}，概形态射 $U' \to V$
是局部完全交态射。因此结论由《态射进阶》引理
\ref{more-morphisms-lemma-lci} 得到。
\end{proof}

\noindent
在此汇集所有 Koszul 态射共有的一些性质看来是合适的。

\begin{lemma}
\label{spaces-more-morphisms-lemma-lci-properties}
设 $S$ 为概形。设 $f : X \to Y$ 为 $S$ 上代数空间的
局部完全交态射。则
\begin{enumerate}
\item $f$ 局部有限表示；
\item $f$ 拟凝聚；
\item $f$ 完美。
\end{enumerate}
\end{lemma}

\begin{proof}
从略。提示：使用本引理的概形版本；见《态射进阶》引理
\ref{more-morphisms-lemma-lci-properties}。
\end{proof}

\noindent
要注意，Koszul 态射的基变换一般并非 Koszul。

\begin{lemma}
\label{spaces-more-morphisms-lemma-flat-base-change-lci}
局部完全交态射的平坦基变换是局部完全交态射。
\end{lemma}

\begin{proof}
从略。提示：使用本引理的概形版本；见《态射进阶》引理
\ref{more-morphisms-lemma-flat-base-change-lci}。
\end{proof}

\begin{lemma}
\label{spaces-more-morphisms-lemma-composition-lci}
局部完全交态射的复合是局部完全交态射。
\end{lemma}

\begin{proof}
从略。提示：使用本引理的概形版本；见《态射进阶》引理
\ref{more-morphisms-lemma-composition-lci}。
\end{proof}

\begin{lemma}
\label{spaces-more-morphisms-lemma-flat-lci}
\begin{slogan}
Syntomic 等于平坦加 lci（对代数空间而言）。
\end{slogan}
设 $S$ 为概形。设 $f : X \to Y$ 为 $S$ 上代数空间的态射。
下列条件等价：
\begin{enumerate}
\item $f$ 平坦且为局部完全交态射；
\item $f$ 为 syntomic。
\end{enumerate}
\end{lemma}

\begin{proof}
从略。提示：使用本引理的概形版本；见《态射进阶》引理
\ref{more-morphisms-lemma-flat-lci}。
\end{proof}

\begin{lemma}
\label{spaces-more-morphisms-lemma-regular-immersion-lci}
设 $S$ 为概形。$S$ 上代数空间的 Koszul-正则浸入是局部完全交态射。
\end{lemma}

\begin{proof}
设 $i : X \to Y$ 为 $S$ 上代数空间的 Koszul-正则浸入。
按定义，存在满平展态射 $V \to Y$，其中 $V$ 为概形，
使得 $X \times_Y V$ 是概形，并且基变换
$X \times_Y V \to V$ 是概形的 Koszul-正则浸入。
由《态射进阶》引理
\ref{more-morphisms-lemma-regular-immersion-lci} 可知，
$X \times_Y V \to V$ 是局部完全交态射。
由定义 \ref{spaces-more-morphisms-definition-lci} 得到 $i$ 是代数空间的局部完全交态射。
\end{proof}

\begin{lemma}
\label{spaces-more-morphisms-lemma-lci-permanence}
设 $S$ 为概形。设
$$
\xymatrix{
X \ar[rr]_f \ar[rd] & & Y \ar[ld] \\
& Z
}
$$
为 $S$ 上代数空间态射的交换图。假设 $Y \to Z$ 光滑，
并且 $X \to Z$ 是局部完全交态射。则
$f : X \to Y$ 是局部完全交态射。
\end{lemma}

\begin{proof}
取概形 $W$ 和满平展态射 $W \to Z$。
取概形 $V$ 和满平展态射 $V \to W \times_Z Y$。
取概形 $U$ 和满平展态射 $U \to V \times_Y X$。
于是 $U \to W$ 是概形的局部完全交态射，且
$V \to W$ 是概形的光滑态射。由概形情形的结论
（《态射进阶》引理 \ref{more-morphisms-lemma-lci-permanence}），
得到 $U \to V$ 是局部完全交态射。按定义，这意味着
$f$ 是局部完全交态射。
\end{proof}

\begin{lemma}
\label{spaces-more-morphisms-lemma-descending-property-lci}
性质 $\mathcal{P}(f) =$“$f$ 是局部完全交态射”在基底上 fpqc 局部。
\end{lemma}

\begin{proof}
设 $S$ 为概形。设 $f : X \to Y$ 为 $S$ 上代数空间的态射。
设 $\{Y_i \to Y\}$ 为 fpqc 覆盖（《空间上的拓扑》定义
\ref{spaces-topologies-definition-fpqc-covering}）。
设 $f_i : X_i \to Y_i$ 为 $f$ 沿 $Y_i \to Y$ 的基变换。
如果 $f$ 是局部完全交态射，则由引理
\ref{spaces-more-morphisms-lemma-flat-base-change-lci}，每个 $f_i$ 都是局部完全交态射。

\medskip\noindent
反过来，假设每个 $f_i$ 都是局部完全交态射。
可以用一个加细替换该覆盖（这仍是因为平坦基变换保持
为局部完全交态射这一性质）。因此可以假设 $Y_i$ 对每个 $i$ 都是概形；见
《空间上的拓扑》引理 \ref{spaces-topologies-lemma-refine-fpqc-schemes}。
取概形 $V$ 和满平展态射 $V \to Y$。
取概形 $U$ 和满平展态射 $U \to V \times_Y X$。
我们必须证明 $U \to V$ 是概形的局部完全交态射。
由《空间上的拓扑》引理
\ref{spaces-topologies-lemma-recognize-fpqc-covering}，
$\{Y_i \times_Y V \to V\}$ 是概形的 fpqc 覆盖。
由概形情形（《态射进阶》引理
\ref{more-morphisms-lemma-descending-property-lci}），只需证明
% P09-ERR-0747
$$
U \times_Y Y_i = U \times_V (V \times_Y Y_i) \longrightarrow V
$$
这一 $U \to V$ 沿 $V \times_Y Y_i \to V$ 的基变换
是局部完全交态射。可以把它写成如下复合：
$$
U \times_Y Y_i \longrightarrow
(V \times_Y X) \times_Y Y_i = V \times_Y X_i \longrightarrow
V \times_Y Y_i
$$
第一支箭头是概形的平展态射（因为它是
$U \to V \times_Y X$ 的基变换），第二支箭头是概形的
局部完全交态射，因为它是 $f_i$ 的平坦基变换。
结论成立，因为为局部完全交态射这一性质在源上 syntomic 局部，
并且平展态射是 syntomic（《态射进阶》引理
\ref{more-morphisms-lemma-lci-syntomic-local-source}
及《态射》引理 \ref{morphisms-lemma-etale-syntomic}）。
\end{proof}

\begin{lemma}
\label{spaces-more-morphisms-lemma-lci-syntomic-local-source}
性质 $\mathcal{P}(f) =$“$f$ 是局部完全交态射”在源上 syntomic 局部。
\end{lemma}

\begin{proof}
这由《空间上的下降》引理
\ref{spaces-descent-lemma-transfer-from-schemes}
和《态射进阶》引理
\ref{more-morphisms-lemma-lci-syntomic-local-source} 得到。
\end{proof}

\begin{lemma}
\label{spaces-more-morphisms-lemma-base-change-lci-fibres}
设 $S$ 为概形。考虑代数空间的交换图
$$
\xymatrix{
X \ar[rr]_f \ar[rd]_p & & Y \ar[ld]^q \\
& Z
}
$$
其中各态射在 $S$ 上。
假设 $p$ 和 $q$ 都平坦且局部有限表示。
则存在开子空间 $U(f) \subset X$，使得 $|U(f)| \subset |X|$
是 $f$ 为 Koszul 之点的集合。此外，对代数空间的任意态射
$Z' \to Z$，如果 $f' : X' \to Y'$ 是 $f$ 沿 $Z' \to Z$
的基变换，则 $U(f')$ 是 $U(f)$ 在投影 $X' \to X$ 下的逆像。
\end{lemma}

\begin{proof}
本引理是《态射进阶》引理
\ref{more-morphisms-lemma-base-change-lci-fibres} 的对应版本，
事实上我们将由该引理推出本引理。按定义 \ref{spaces-more-morphisms-definition-lci}，
集合 $\{x \in |X| : f \text{ 为 Koszul 于点 }x\}$
在 $|X|$ 中开，因而由《空间的性质》引理
\ref{spaces-properties-lemma-open-subspaces}，
它对应于开子空间 $U(f)$；这是 $X$ 的一个开子空间。
所以只需证明最后一个断言。

\medskip\noindent
取概形 $W$ 和满平展态射 $W \to Z$。
取概形 $V$ 和满平展态射 $V \to W \times_Z Y$。
取概形 $U$ 和满平展态射 $U \to V \times_Y X$。
最后，取概形 $W'$ 和满平展态射 $W' \to W \times_Z Z'$。
令 $V' = W' \times_W V$ 且 $U' = W' \times_W U$，
从而得到满平展态射 $V' \to Y'$ 和 $U' \to X'$。
以下将不再说明地使用这一事实：代数空间的平展态射在相伴拓扑空间上
诱导开映射（见《空间的性质》引理
\ref{spaces-properties-lemma-etale-open}）。
注意，按定义，$U(f)$ 在 $|X|$ 中表示集合 $T$ 的像；
该集合由 $U$ 中使概形态射 $U \to V$ 为 Koszul 的那些点组成。
类似地，$U(f')$ 在 $|X'|$ 中表示集合 $T'$ 的像；
该集合由 $U'$ 中使概形态射 $U' \to V'$ 为 Koszul 的那些点组成。
现在，按构造，图
$$
\xymatrix{
U' \ar[r] \ar[d] & U \ar[d] \\
V' \ar[r] & V
}
$$
是笛卡尔的（在概形范畴中）。因此，上述《态射进阶》引理
\ref{more-morphisms-lemma-base-change-lci-fibres} 适用，并说明
$T'$ 是 $T$ 的逆像。由于 $|U'| \to |X'|$ 满，这就推出了本引理。
\end{proof}

\begin{lemma}
\label{spaces-more-morphisms-lemma-unramified-lci}
设 $S$ 为概形。设 $f : X \to Y$ 为 $S$ 上代数空间的
局部完全交态射。则 $f$ 非分歧，当且仅当 $f$ 形式非分歧；
在这种情况下，余法层 $\mathcal{C}_{X/Y}$ 在 $X$ 上有限局部自由。
\end{lemma}

\begin{proof}
这由概形态射的相应结论经平展局部化得到；见《态射进阶》引理
\ref{more-morphisms-lemma-unramified-lci} 和引理
\ref{spaces-more-morphisms-lemma-universal-thickening-localize}。
（注意，在该引理的情形中，态射 $V \to U$
按定义非分歧且为局部完全交态射。）
\end{proof}

\begin{lemma}
\label{spaces-more-morphisms-lemma-transitivity-conormal-lci}
设 $S$ 为概形。设 $Z \to Y \to X$ 为 $S$ 上代数空间的
形式非分歧态射。假设 $Z \to Y$ 是局部完全交态射。
% P09-ERR-0748
引理 \ref{spaces-more-morphisms-lemma-transitivity-conormal} 的正合序列
$$
0 \to i^*\mathcal{C}_{Y/X} \to
\mathcal{C}_{Z/X} \to
\mathcal{C}_{Z/Y} \to 0
$$
是短正合的。
\end{lemma}

\begin{proof}
取概形 $U$ 和满平展态射 $U \to X$。
取概形 $V$ 和满平展态射 $V \to U \times_X Y$。
取概形 $W$ 和满平展态射 $W \to V \times_Y Z$。
由引理 \ref{spaces-more-morphisms-lemma-universal-thickening-localize}，
态射 $W \to V$ 和 $V \to U$ 都形式非分歧。此外，序列
$i^*\mathcal{C}_{Y/X} \to \mathcal{C}_{Z/X} \to \mathcal{C}_{Z/Y} \to 0$
限制为相应序列
$i^*\mathcal{C}_{V/U} \to \mathcal{C}_{W/U} \to \mathcal{C}_{W/V} \to 0$；
这是针对 $W \to V \to U$ 而言的。
所以，由概形情形的结论（《态射进阶》引理
\ref{more-morphisms-lemma-transitivity-conormal-lci}）可得所求，
因为按定义，态射 $W \to V$ 是局部完全交态射。
\end{proof}

\section{一个态射何时是同构？}
\label{spaces-more-morphisms-section-when-isomorphism}

\noindent
更一般地，可以问：“一个态射何时具有性质 $\mathcal{P}$？”
更精确的问题如下。假设给定代数空间的交换图
$$
\xymatrix{
X \ar[rr]_f \ar[rd]_p & & Y \ar[ld]^q \\
& Z
}
$$
是否存在代数空间的单态射 $W \to Z$，使得下列两个性质成立：
\begin{enumerate}
\item 基变换 $f_W : X_W \to Y_W$ 具有性质 $\mathcal{P}$；
\item 代数空间的任意态射 $Z' \to Z$ 经由 $W$ 分解，
当且仅当基变换 $f_{Z'} : X_{Z'} \to Y_{Z'}$ 具有性质
$\mathcal{P}$。
\end{enumerate}
在许多情形中，如果 $W \to Z$ 存在，那么它是浸入、开浸入或闭浸入。

\medskip\noindent
这一问题的答案可能依赖于态射 $f$、$p$ 和 $q$ 的辅助性质。
例如 $\mathcal{P}(f) =$“$f$ 平坦”；在 $Y = S$ 的情形中，
我们已经在“平坦性进阶”一章对概形态射详细讨论过这一例子，
起点是《平坦性进阶》第 \ref{flat-section-flattening-functors} 节。

\begin{lemma}
\label{spaces-more-morphisms-lemma-where-unramified}
考虑代数空间的交换图
$$
\xymatrix{
X \ar[rr]_f \ar[rd]_p & & Y \ar[ld]^q \\
& Z
}
$$
假设 $p$ 局部有限型且为闭映射。则存在开子空间 $W \subset Z$，
使得态射 $Z' \to Z$ 经由 $W$ 分解，当且仅当基变换
$f_{Z'} : X_{Z'} \to Y_{Z'}$ 非分歧。
\end{lemma}

\begin{proof}
由《空间的态射》引理
\ref{spaces-morphisms-lemma-where-unramified}，
存在开子空间 $U(f) \subset X$，它是使 $f$ 非分歧的点的集合。
此外，$U(f)$ 的构造与任意基变换交换。令 $W \subset Z$
为如下开子空间（见《空间的性质》引理
\ref{spaces-properties-lemma-open-subspaces}），其底层点集为
$$
|W| = |Z| \setminus |p|\left(|X| \setminus |U(f)|\right)
$$
也就是说，$z \in |Z|$ 是 $W$ 的点，当且仅当 $f$
在 $X$ 中位于 $z$ 上方的每一点处都非分歧。
由于假设 $p$ 为闭映射，该集合为开。因为 $U(f)$ 的构造与任意
基变换交换，立即可见（使用《空间的性质》引理
\ref{spaces-properties-lemma-factor-through-open-subspace}）
$W$ 具有所需的泛性质。
\end{proof}

\begin{lemma}
\label{spaces-more-morphisms-lemma-where-unramified-universally-injective}
考虑代数空间的交换图
$$
\xymatrix{
X \ar[rr]_f \ar[rd]_p & & Y \ar[ld]^q \\
& Z
}
$$
假设
\begin{enumerate}
\item $p$ 局部有限型；
\item $p$ 为闭映射；
\item $p_2 : X \times_Y X \to Z$ 为闭映射。
\end{enumerate}
则存在开子空间 $W \subset Z$，使得态射 $Z' \to Z$ 经由 $W$
分解，当且仅当基变换 $f_{Z'} : X_{Z'} \to Y_{Z'}$
非分歧且普遍单射。
\end{lemma}

\begin{proof}
用引理 \ref{spaces-more-morphisms-lemma-where-unramified} 中得到的开子空间替换 $Z$ 后，
可以假设 $f$ 已经非分歧；注意，这不会破坏假设 (2) 或 (3)。
由《空间的态射》引理
\ref{spaces-morphisms-lemma-diagonal-unramified-morphism}，
可知 $\Delta_{X/Y} : X \to X \times_Y X$ 是开浸入。
任意基变换后这一点仍成立。因此，由《空间的态射》引理
\ref{spaces-morphisms-lemma-universally-injective}，
$f_{Z'}$ 普遍单射，当且仅当对角线的基变换
$X_{Z'} \to (X \times_Y X)_{Z'}$ 是同构。令 $W \subset Z$
为如下开子空间（见《空间的性质》引理
\ref{spaces-properties-lemma-open-subspaces}），其底层点集为
$$
|W| = |Z| \setminus
|p_2|\left(|X \times_Y X| \setminus \Im(|\Delta_{X/Y}|)\right)
$$
也就是说，$z \in |Z|$ 是 $W$ 的点，当且仅当
$|X \times_Y X| \to |Z|$ 在 $z$ 上的纤维包含于
$|X| \to |X \times_Y X|$ 的像中。由上面的讨论立即可知，
限制 $p^{-1}(W) \to q^{-1}(W)$（即 $f$ 的限制）非分歧且普遍单射。

\medskip\noindent
反过来，假设 $f_{Z'}$ 非分歧且普遍单射。为了证明
$Z' \to Z$ 经由 $W$ 分解，只需证明 $|Z'| \to |Z|$
的像包含于 $|W|$；见《空间的性质》引理
\ref{spaces-properties-lemma-factor-through-open-subspace}。
因此，只需证明 $Z'$ 是一个域的谱时的结论。
% P09-ERR-0749
记 $z \in |Z|$ 为 $|Z'| \to |Z|$ 的像。上面的讨论说明
$$
|X_{Z'}| \longrightarrow |(X \times_Y X)_{Z'}|
$$
是满的。由《空间的性质》引理
\ref{spaces-properties-lemma-points-cartesian}，在交换图
$$
\xymatrix{
|X_{Z'}| \ar[d] \ar[r] &
|(X \times_Y X)_{Z'}| \ar[d] \\
|p|^{-1}(\{z\}) \ar[r] &
|p_2|^{-1}(\{z\})
}
$$
中，竖直箭头均为满的。于是得到所需的 $z \in |W|$。
\end{proof}

\begin{lemma}
\label{spaces-more-morphisms-lemma-where-closed-immersion}
考虑代数空间的交换图
$$
\xymatrix{
X \ar[rr]_f \ar[rd]_p & & Y \ar[ld]^q \\
& Z
}
$$
假设
\begin{enumerate}
\item $p$ 局部有限型；
\item $p$ 普遍闭；
\item $q : Y \to Z$ 可分。
\end{enumerate}
则存在开子空间 $W \subset Z$，使得态射 $Z' \to Z$ 经由 $W$
分解，当且仅当基变换 $f_{Z'} : X_{Z'} \to Y_{Z'}$ 是闭浸入。
\end{lemma}

\begin{proof}
我们将使用闭浸入作为普遍闭、非分歧且普遍单射态射的刻画；见引理
\ref{spaces-more-morphisms-lemma-characterize-closed-immersion}。
首先，注意到 $p$ 普遍闭且 $q$ 可分，所以 $f$ 普遍闭；见
《空间的态射》引理
\ref{spaces-morphisms-lemma-universally-closed-permanence}。
于是 $f$ 的任意基变换都普遍闭；见《空间的态射》引理
\ref{spaces-morphisms-lemma-base-change-universally-closed}。
因此，为完成本引理的证明，只需证明引理
\ref{spaces-more-morphisms-lemma-where-unramified-universally-injective} 的假设成立。
投影 $\text{pr}_0 : X \times_Y X \to X$ 作为 $f$ 的基变换而普遍闭；
见《空间的态射》引理
\ref{spaces-morphisms-lemma-base-change-universally-closed}。
所以 $X \times_Y X \to Z$ 作为普遍闭态射的复合而普遍闭
（见《空间的态射》引理
\ref{spaces-morphisms-lemma-composition-universally-closed}）。
这就完成了本引理的证明。
\end{proof}

\begin{lemma}
\label{spaces-more-morphisms-lemma-where-flat}
考虑代数空间的交换图
$$
\xymatrix{
X \ar[rr]_f \ar[rd]_p & & Y \ar[ld]^q \\
& Z
}
$$
假设
\begin{enumerate}
\item $p$ 局部有限表示；
\item $p$ 平坦；
\item $p$ 为闭映射；
\item $q$ 局部有限型。
\end{enumerate}
则存在开子空间 $W \subset Z$，使得态射 $Z' \to Z$ 经由 $W$
分解，当且仅当基变换 $f_{Z'} : X_{Z'} \to Y_{Z'}$ 平坦。
\end{lemma}

\begin{proof}
由引理 \ref{spaces-more-morphisms-lemma-base-change-flatness-fibres}，集合
$$
A = \{x \in |X| : X\text{ 在点 }x \text{ 处平坦于 }Y\}.
$$
在 $|X|$ 中开，并且它的构造与任意基变换交换。令
$W \subset Z$ 为如下开子空间（见《空间的性质》引理
\ref{spaces-properties-lemma-open-subspaces}），其底层点集为
$$
|W| = |Z| \setminus |p|\left(|X| \setminus A\right)
$$
也就是说，$z \in |Z|$ 是 $W$ 的点，当且仅当
$|X| \to |Z|$ 在 $z$ 上的整个纤维包含于 $A$。
由于 $p$ 为闭映射，该集合为开。因为 $A$ 的构造与任意
基变换交换，所以 $W$ 符合要求。
\end{proof}

\begin{lemma}
\label{spaces-more-morphisms-lemma-where-surjective-flat}
考虑代数空间的交换图
$$
\xymatrix{
X \ar[rr]_f \ar[rd]_p & & Y \ar[ld]^q \\
& Z
}
$$
假设
\begin{enumerate}
\item $p$ 局部有限表示；
\item $p$ 平坦；
\item $p$ 为闭映射；
\item $q$ 局部有限型；
\item $q$ 为闭映射。
\end{enumerate}
则存在开子空间 $W \subset Z$，使得态射 $Z' \to Z$ 经由 $W$
分解，当且仅当基变换 $f_{Z'} : X_{Z'} \to Y_{Z'}$ 满且平坦。
\end{lemma}

\begin{proof}
由引理 \ref{spaces-more-morphisms-lemma-where-flat}，可以假设 $f$ 平坦。
由《空间的态射》引理
\ref{spaces-morphisms-lemma-finite-presentation-permanence}，
注意到 $f$ 局部有限表示。所以 $f$ 为开映射；见《空间的态射》引理
\ref{spaces-morphisms-lemma-fppf-open}。令 $W \subset Z$
为如下开子空间（见《空间的性质》引理
\ref{spaces-properties-lemma-open-subspaces}），其底层点集为
$$
|W| = |Z| \setminus |q|\left(|Y| \setminus |f|(|X|)\right).
$$
换言之，对 $z \in |Z|$，有 $z \in |W|$，当且仅当
$|Y| \to |Z|$ 在 $z$ 上的整个纤维包含于
$|X| \to |Y|$ 的像。由于 $q$ 为闭映射，该集合在 $|Z|$ 中开。
按构造，态射 $X_W \to Y_W$ 是满的。最后，假设
$X_{Z'} \to Y_{Z'}$ 是满的。为了证明 $Z' \to Z$ 经由 $W$ 分解，
只需证明 $|Z'| \to |Z|$ 的像包含于 $|W|$；见《空间的性质》引理
\ref{spaces-properties-lemma-factor-through-open-subspace}。
因此，只需证明 $Z'$ 是一个域的谱时的结论。
% P09-ERR-0750
记 $z \in |Z|$ 为 $|Z'| \to |Z|$ 的像。由《空间的性质》引理
\ref{spaces-properties-lemma-points-cartesian}，在交换图
$$
\xymatrix{
|X_{Z'}| \ar[d] \ar[r] &
|Y_{Z'}| \ar[d] \\
|p|^{-1}(\{z\}) \ar[r] &
|q|^{-1}(\{z\})
}
$$
中，竖直箭头均为满的。于是得到所需的 $z \in |W|$。
\end{proof}

\begin{lemma}
\label{spaces-more-morphisms-lemma-where-isomorphism}
考虑代数空间的交换图
$$
\xymatrix{
X \ar[rr]_f \ar[rd]_p & & Y \ar[ld]^q \\
& Z
}
$$
假设
\begin{enumerate}
\item $p$ 局部有限表示；
\item $p$ 平坦；
\item $p$ 普遍闭；
\item $q$ 局部有限型；
\item $q$ 为闭映射；
\item $q$ 可分。
\end{enumerate}
则存在开子空间 $W \subset Z$，使得态射 $Z' \to Z$ 经由 $W$
分解，当且仅当基变换 $f_{Z'} : X_{Z'} \to Y_{Z'}$ 是同构。
\end{lemma}

\begin{proof}
由引理 \ref{spaces-more-morphisms-lemma-where-surjective-flat}，存在开子空间
$W_1 \subset Z$，使得 $f_{Z'}$ 满且平坦，当且仅当
$Z' \to Z$ 经由 $W_1$ 分解。由引理
\ref{spaces-more-morphisms-lemma-where-closed-immersion}，存在开子空间
$W_2 \subset Z$，使得 $f_{Z'}$ 是闭浸入，当且仅当
$Z' \to Z$ 经由 $W_2$ 分解。我们断言 $W = W_1 \cap W_2$ 符合要求。
当然，如果 $f_{Z'}$ 是同构，那么 $Z' \to Z$ 经由 $W$ 分解。
所以只需证明 $f_W$ 是同构。按构造，$f_W$ 是满、平坦的闭浸入。
特别地，$f_W$ 可表。由于概形的满平坦闭浸入是同构
（见《态射》引理
\ref{morphisms-lemma-characterize-flat-closed-immersions}），结论成立。
（注意，事实上 $f_W$ 局部有限表示，从而为开映射；因此若愿意，
可以避免使用这个引理。）
\end{proof}

\begin{lemma}
\label{spaces-more-morphisms-lemma-where-lci}
考虑代数空间的交换图
$$
\xymatrix{
X \ar[rr]_f \ar[rd]_p & & Y \ar[ld]^q \\
& Z
}
$$
假设
\begin{enumerate}
\item $p$ 平坦且局部有限表示；
\item $p$ 为闭映射；
\item $q$ 平坦且局部有限表示。
\end{enumerate}
则存在开子空间 $W \subset Z$，使得态射 $Z' \to Z$ 经由 $W$
分解，当且仅当基变换 $f_{Z'} : X_{Z'} \to Y_{Z'}$
是局部完全交态射。
\end{lemma}

\begin{proof}
由引理 \ref{spaces-more-morphisms-lemma-base-change-lci-fibres}，存在开子空间
$U(f) \subset X$，它是使 $f$ 为 Koszul 的点的集合。
此外，$U(f)$ 的构造与任意基变换交换。令 $W \subset Z$
为如下开子空间（见《空间的性质》引理
\ref{spaces-properties-lemma-open-subspaces}），其底层点集为
$$
|W| = |Z| \setminus |p|\left(|X| \setminus |U(f)|\right)
$$
也就是说，$z \in |Z|$ 是 $W$ 的点，当且仅当 $f$
在 $X$ 中位于 $z$ 上方的每一点处都为 Koszul。
由于假设 $p$ 为闭映射，该集合为开。因为 $U(f)$ 的构造与任意
基变换交换，立即可见（使用《空间的性质》引理
\ref{spaces-properties-lemma-factor-through-open-subspace}）
$W$ 具有所需的泛性质。
\end{proof}

\section{微分层与余法层的正合列}
\label{spaces-more-morphisms-section-exact}

\noindent
本节汇集关于余法层和微分层的正合列的一些结果。
在某种意义下，这些都是与代数空间的可复合态射相伴的
余切复形三角形的实现。

\medskip\noindent
% P09-ERR-0751
下面各个序列中的每个映射都是在引理
\ref{spaces-more-morphisms-lemma-functoriality-differentials} 或引理
\ref{spaces-more-morphisms-lemma-universal-thickening-functorial} 中构造的。
设 $S$ 为概形。设 $g : Z \to Y$ 和 $f : Y \to X$
为 $S$ 上代数空间的态射。
\begin{enumerate}
\item 有典范正合列
$$
g^*\Omega_{Y/X} \to \Omega_{Z/X} \to \Omega_{Z/Y} \to 0,
$$
见引理 \ref{spaces-more-morphisms-lemma-triangle-differentials}。
如果 $g : Z \to Y$ 形式光滑，那么这个序列是短正合列；见引理
\ref{spaces-more-morphisms-lemma-triangle-differentials-formally-smooth}。
\item 如果 $g$ 形式非分歧，那么有典范正合列
$$
\mathcal{C}_{Z/Y} \to g^*\Omega_{Y/X} \to \Omega_{Z/X} \to 0,
$$
见引理 \ref{spaces-more-morphisms-lemma-universally-unramified-differentials-sequence}。
如果 $f \circ g : Z \to X$ 形式光滑，那么这个序列是短正合列；见引理
\ref{spaces-more-morphisms-lemma-differentials-formally-unramified-formally-smooth}。
\item 如果 $g$ 和 $f \circ g$ 都形式非分歧，那么有典范正合列
$$
\mathcal{C}_{Z/X} \to \mathcal{C}_{Z/Y} \to g^*\Omega_{Y/X} \to 0,
$$
见引理 \ref{spaces-more-morphisms-lemma-two-unramified-morphisms}。
如果 $f : Y \to X$ 形式光滑，那么这个序列是短正合列；见引理
\ref{spaces-more-morphisms-lemma-two-unramified-morphisms-formally-smooth}。
\item 如果 $g$ 和 $f$ 都形式非分歧，那么有典范正合列
$$
g^*\mathcal{C}_{Y/X} \to \mathcal{C}_{Z/X} \to \mathcal{C}_{Z/Y} \to 0.
$$
见引理 \ref{spaces-more-morphisms-lemma-transitivity-conormal-unramified}。
如果 $g : Z \to Y$ 是局部完全交态射，那么这个序列是短正合列；
见引理 \ref{spaces-more-morphisms-lemma-transitivity-conormal-lci}。
\end{enumerate}

\section{拟凝聚复形的刻画，II}
\label{spaces-more-morphisms-section-characterize-pseudo-coherent}

\noindent
本节用上同调讨论拟凝聚复形的一种刻画。关于代数空间上拟凝聚复形的
早先材料可见《空间的导出范畴》第
\ref{spaces-perfect-section-spell-out} 节和《空间的导出范畴》第
\ref{spaces-perfect-section-pseudo-coherent-hocolim} 节。
本节在概形情形的对应版本是《态射进阶》第
\ref{more-morphisms-section-characterize-pseudo-coherent} 节。
一个基本工具是利用 Chow 引理的导出版本约化到射影空间的情形；
见引理 \ref{spaces-more-morphisms-lemma-derived-chow}。

\begin{lemma}
\label{spaces-more-morphisms-lemma-case-of-tor-independence}
设 $S$ 为概形。考虑代数空间的交换图
$$
\xymatrix{
Z' \ar[d] \ar[r] & Y' \ar[d] \\
X' \ar[r] & B'
}
$$
它在 $S$ 上。设 $B \to B'$ 为态射。记 $X$ 和 $Y$ 分别为
$X'$ 和 $Y'$ 到 $B$ 的基变换。假设 $Y' \to B'$ 和
$Z' \to X'$ 平坦。则 $X \times_B Y$ 和 $Z'$
在 $X' \times_{B'} Y'$ 上 Tor 独立。
\end{lemma}

\begin{proof}
由《空间的导出范畴》引理
\ref{spaces-perfect-lemma-tor-independent}，可以在
$X \times_B Y$ 和 $Z'$ 上平展局部地检验 Tor 独立性。
这\footnote{更详细的论证如下。取满平展态射 $W' \to B'$，
其中 $W'$ 为概形。取满平展态射 $W \to B \times_{B'} W'$，
其中 $W$ 为概形。取满平展态射
$U' \to X' \times_{B'} W'$，其中 $U'$ 为概形。取满平展态射
$V' \to Y' \times_{B'} W'$，其中 $V'$ 为概形。
注意，$U' \times_{W'} V' \to X' \times_{B'} Y'$ 满平展。
取满平展态射
$T' \to Z' \times_{X' \times_{B'} Y'} U' \times_{W'} V'$，
其中 $T'$ 为概形。
% P09-ERR-0752
记 $U$ 和 $V$ 为 $U'$ 和 $V'$ 到 $W$ 的基变换。
于是本引理断言：作为代数空间，$X \times_B Y$ 和 $Z'$
在 $X' \times_{B'} Y'$ 上 Tor 独立，当且仅当作为概形，
$U \times_W V$ 和 $T'$ 在 $U' \times_{W'} V'$ 上 Tor 独立。
因此，只需对以 $T', U', V', W'$ 为四个角的方块及基变换
$W \to W'$ 证明本引理。$Y' \to B'$ 和 $Z' \to X'$ 的平坦性
蕴含 $V' \to W'$ 和 $T' \to U'$ 的平坦性。}
把本引理约化到概形的情形，而该情形就是《态射进阶》引理
\ref{more-morphisms-lemma-case-of-tor-independence}。
\end{proof}

\begin{lemma}[导出 Chow 引理]
\label{spaces-more-morphisms-lemma-derived-chow}
设 $A$ 为环。设 $X$ 为 $A$ 上可分且有限表示的代数空间。
设 $x \in |X|$。则存在一个 $n \geq 0$、一个闭子空间
$Z \subset X \times_A \mathbf{P}^n_A$、一个点 $z \in |Z|$、
一个开集 $V \subset \mathbf{P}^n_A$，以及一个对象
$E$（它属于 $D(\mathcal{O}_{X \times_A \mathbf{P}^n_A})$），使得
\begin{enumerate}
\item $Z \to X \times_A \mathbf{P}^n_A$ 有限表示；
\item $c : Z \to \mathbf{P}^n_A$ 在 $V$ 上是闭浸入；
令 $W = c^{-1}(V)$；
\item $b : Z \to X$ 到 $W$ 的限制平展，
$z \in W$，并且 $b(z) = x$；
\item $E|_{X \times_A V} \cong
(b, c)_*\mathcal{O}_Z|_{X \times_A V}$；
\item $E$ 拟凝聚且支撑在 $Z$ 上。
\end{enumerate}
\end{lemma}

\begin{proof}
可以找到有限型 $\mathbf{Z}$-子代数 $A' \subset A$
以及代数空间 $X'$，它在 $A'$ 上可分且有限表示，并使它到 $A$
的基变换为 $X$。见《空间的极限》引理
\ref{spaces-limits-lemma-descend-finite-presentation} 和
\ref{spaces-limits-lemma-descend-separated-morphism}。
设 $x' \in |X'|$ 为 $x$ 的像。如果能对 $(X'/A', x')$
证明本引理，那么对 $(X/A, x)$ 的结论随即成立。
确实，如果 $n', Z', z', V', E'$ 给出 $(X'/A', x')$ 的解，
那么可以令 $n = n'$；令 $Z \subset X \times \mathbf{P}^n$
为 $Z'$ 的逆像；
% P09-ERR-0753
令 $z \in Z$ 为映到 $x$ 的唯一点；令
$V \subset \mathbf{P}^n_A$ 为 $V'$ 的逆像；并令 $E$
为 $E'$ 的导出拉回。注意，由《位点上的上同调》引理
\ref{sites-cohomology-lemma-pseudo-coherent-pullback}，$E$ 拟凝聚。
% P09-ERR-0754
只需检验 (5)。为此，令 $W = c^{-1}(V)$ 和
$W' = (c')^{-1}(V')$，并考虑笛卡尔方块
$$
\xymatrix{
W \ar[d]_{(b, c)} \ar[r] & W' \ar[d]^{(b', c')} \\
X \times_A V \ar[r] & X' \times_{A'} V'
}
$$
由引理 \ref{spaces-more-morphisms-lemma-case-of-tor-independence}，$X \times_A V$ 和 $W'$
在 $X' \times_{A'} V'$ 上 Tor 独立。因此，由《空间的导出范畴》引理
\ref{spaces-perfect-lemma-compare-base-change}，
$(b', c')_*\mathcal{O}_{W'}$ 到 $X \times_A V$ 的导出拉回是
$(b, c)_*\mathcal{O}_W$。这里还使用了
$R(b', c')_*\mathcal{O}_{Z'} = (b', c')_*\mathcal{O}_{Z'}$，
因为 $(b', c')$ 是闭浸入；对 $(b, c)_*\mathcal{O}_Z$ 也同样如此。
% P09-ERR-0755
由于 $E'|_{U' \times_{A'} V'} =
(b', c')_*\mathcal{O}_{W'}$，得到
$E|_{U \times_A V} = (b, c)_*\mathcal{O}_W$，从而 (5) 成立。
这把我们约化到下一段所述的情形。

\medskip\noindent
假设 $A$ 在 $\mathbf{Z}$ 上有限型。取平展态射 $U \to X$，
其中 $U$ 为仿射概形，并取一点 $u \in U$，它映到 $x$。
于是 $U$ 在 $A$ 上有限型。取闭浸入 $U \to \mathbf{A}^n_A$，并记
% P09-ERR-0756
$j : U \to \mathbf{P}^n_A$ 为将其与开浸入
$\mathbf{A}^n_A \to \mathbf{P}^n_A$ 复合而得到的浸入。
令 $Z$ 为
$$
(\text{id}_U, j) : U \longrightarrow X \times_A \mathbf{P}^n_A
$$
的概形论闭包。设 $z \in Z$ 为 $u$ 的像。
设 $Y \subset \mathbf{P}^n_A$ 为 $j$ 的概形论闭包。
于是显然，$Z \subset X \times_A Y$ 是
$(\text{id}_U, j) : U \to X \times_A Y$ 的概形论闭包。
由于 $X$ 可分，态射 $X \times_A Y \to Y$ 也可分。
因此，由《空间的态射》引理
\ref{spaces-morphisms-lemma-scheme-theoretic-image-of-partial-section}，
可见 $Z \to Y$ 在开子概形 $j(U) \subset Y$ 上是同构。
取开集 $V \subset \mathbf{P}^n_A$，使 $V \cap Y = j(U)$。
于是 (2) 成立，且 $W = (\text{id}_U, j)(U)$，从而 (3) 成立。
(1) 成立是因为 $A$ 诺特。

\medskip\noindent
因为 $A$ 诺特，所以 $X$ 和 $X \times_A \mathbf{P}^n_A$
都是诺特代数空间。因此在这种情形可以取
$E = (b, c)_*\mathcal{O}_Z$：(4) 是显然的；关于 (5)，见
《空间的导出范畴》引理
\ref{spaces-perfect-lemma-identify-pseudo-coherent-noetherian}。
这就完成了证明。
\end{proof}

\begin{lemma}
\label{spaces-more-morphisms-lemma-compute-Fourier-Mukai-for-derived-chow}
设 $X/A$、$x \in |X|$ 以及 $n, Z, z, V, E$
如引理 \ref{spaces-more-morphisms-lemma-derived-chow} 中所述。对任意
$K \in D_\QCoh(\mathcal{O}_X)$，有
$$
Rq_*(Lp^*K \otimes^\mathbf{L} E)|_V = R(W \to V)_*K|_W
$$
其中 $p : X \times_A \mathbf{P}^n_A \to X$ 和
$q : X \times_A \mathbf{P}^n_A \to \mathbf{P}^n_A$ 是投影，
而态射 $W \to V$ 是有限表示闭浸入 $c|_W : W \to V$。
\end{lemma}

\begin{proof}
由于 $W = c^{-1}(V)$，并且 $c$ 在 $V$ 上是闭浸入，
可知 $c|_W$ 是闭浸入。它有限表示，因为 $W$ 和 $V$
都在 $A$ 上有限表示；见《空间的态射》引理
\ref{spaces-morphisms-lemma-finite-presentation-permanence}。
首先有
$$
Rq_*(Lp^*K \otimes^\mathbf{L} E)|_V =
Rq'_*\left((Lp^*K \otimes^\mathbf{L} E)|_{X \times_A V}\right)
$$
其中 $q' : X \times_A V \to V$ 是投影，因为总正像的构造与局部化交换。
% P09-ERR-0757
记 $i = (b, c)|_W : W \to X \times_A V$ 为给定的闭浸入。则
$$
Rq'_*\left((Lp^*K \otimes^\mathbf{L} E)|_{X \times_A V}\right) =
Rq'_*(Lp^*K|_{X \times_A V} \otimes^\mathbf{L} i_*\mathcal{O}_W)
$$
这由性质 (5) 得到。由于 $i$ 是闭浸入，有
$i_*\mathcal{O}_W = Ri_*\mathcal{O}_W$。使用《空间的导出范畴》引理
\ref{spaces-perfect-lemma-cohomology-base-change}，可以把它改写为
$$
Rq'_* Ri_* Li^* Lp^*K|_{X \times_A V} =
R(q' \circ i)_* Lb^*K|_W =
R(W \to V)_* K|_W
$$
这正是所要证明的。（注意，限制到 $W$ 与经由 $W \to X$
导出拉回是同一件事，因为 $W$ 在 $X$ 上平展。）
\end{proof}

\begin{lemma}
\label{spaces-more-morphisms-lemma-characterize-pseudo-coherent}
设 $A$ 为环。设 $X$ 为 $A$ 上可分且有限表示的代数空间。
设 $K \in D_\QCoh(\mathcal{O}_X)$。如果
$R\Gamma(X, E \otimes^\mathbf{L} K)$ 是 $D(A)$ 中的拟凝聚对象，
并且这对每个拟凝聚的 $E$（它属于 $D(\mathcal{O}_X)$）都成立，
那么 $K$ 相对于 $A$ 拟凝聚
（定义 \ref{spaces-more-morphisms-definition-relative-pseudo-coherence}）。
\end{lemma}

\begin{proof}
假设 $K \in D_\QCoh(\mathcal{O}_X)$，并且
$R\Gamma(X, E \otimes^\mathbf{L} K)$ 是 $D(A)$ 中的拟凝聚对象，
其中 $E$ 是 $D(\mathcal{O}_X)$ 中任意拟凝聚对象。
设 $x \in |X|$。我们将证明 $K$ 相对于 $A$ 在 $x$
的某个平展邻域中拟凝聚。由相对拟凝聚性的定义，这将证明本引理。

\medskip\noindent
取 $n, Z, z, V, E$ 如引理 \ref{spaces-more-morphisms-lemma-derived-chow} 中所述。
记 $p : X \times \mathbf{P}^n \to X$ 和
$q : X \times \mathbf{P}^n \to \mathbf{P}^n_A$ 为投影。
于是对任意 $i \in \mathbf{Z}$，有
\begin{align*}
& R\Gamma(\mathbf{P}^n_A,
Rq_*(Lp^*K \otimes^\mathbf{L} E)
\otimes^\mathbf{L}
\mathcal{O}_{\mathbf{P}^n_A}(i)) \\
& =
R\Gamma(X \times \mathbf{P}^n,
Lp^*K \otimes^\mathbf{L} E \otimes^\mathbf{L}
Lq^*\mathcal{O}_{\mathbf{P}^n_A}(i)) \\
& =
R\Gamma(X,
K \otimes^\mathbf{L} Rp_*(E \otimes^\mathbf{L}
Lq^*\mathcal{O}_{\mathbf{P}^n_A}(i)))
\end{align*}
这由《空间的导出范畴》引理
\ref{spaces-perfect-lemma-cohomology-base-change} 得到。
由《空间的导出范畴》引理
\ref{spaces-perfect-lemma-flat-proper-pseudo-coherent-direct-image-general}，
复形 $Rp_*(E \otimes^\mathbf{L} Lq^*\mathcal{O}_{\mathbf{P}^n_A}(i))$
在 $X$ 上拟凝聚。因此，假设告诉我们，上面显示公式中的表达式
是 $D(A)$ 中的拟凝聚对象。由《概形的导出范畴》引理
\ref{perfect-lemma-pseudo-coherent-on-projective-space}，
得到 $Rq_*(Lp^*K \otimes^\mathbf{L} E)$ 在
$\mathbf{P}^n_A$ 上拟凝聚。由引理
\ref{spaces-more-morphisms-lemma-compute-Fourier-Mukai-for-derived-chow}，有
% P09-ERR-0758
$$
Rq_*(Lp^*K \otimes^\mathbf{L} E)|_{X \times_A V} =
R(W \to V)_*K|_W
$$
由于 $W \to V$ 是到 $\mathbf{P}^n_A$ 的开子概形中的闭浸入，
这意味着 $K|_W$ 相对于 $A$ 拟凝聚；例如可由《态射进阶》引理
\ref{more-morphisms-lemma-check-relative-pseudo-coherence-on-charts} 看出。
\end{proof}

\begin{lemma}
\label{spaces-more-morphisms-lemma-characterize-pseudo-coh-improved}
设 $A$ 为环。设 $X$ 为 $A$ 上可分且有限表示的代数空间。
设 $K \in D_\QCoh(\mathcal{O}_X)$。如果
$R \Gamma (X, E \otimes ^{\mathbf{L}} K)$ 是 $D(A)$ 中的拟凝聚对象，
并且这对每个完美的 $E \in D(\mathcal{O}_X)$ 都成立，
那么 $K$ 相对于 $A$ 拟凝聚。
\end{lemma}

\begin{proof}
鉴于引理 \ref{spaces-more-morphisms-lemma-characterize-pseudo-coherent}，只需证明
$R \Gamma (X, E \otimes ^{\mathbf{L}} K)$ 是 $D(A)$ 中的拟凝聚对象，
并且这对每个拟凝聚的 $E \in D(\mathcal{O}_X)$ 都成立。
由《空间的导出范畴》命题
\ref{spaces-perfect-proposition-detecting-bounded-above}，
得到 $K \in D^-_\QCoh (\mathcal{O}_X)$。现在结论由
《空间的导出范畴》引理
\ref{spaces-perfect-lemma-perfect-enough} 得到。
\end{proof}

\section{相对完美对象}
\label{spaces-more-morphisms-section-relatively-perfect}

\noindent
本节引入 \cite{lieblich-complexes} 中的一个概念。
《概形的导出范畴》第 \ref{perfect-section-relatively-perfect} 节
已经对概形态射讨论过这个概念。

\begin{definition}
\label{spaces-more-morphisms-definition-relatively-perfect}
设 $S$ 为概形。设 $f : X \to Y$ 为 $S$ 上代数空间的态射，
并且它平坦且局部有限表示。称对象 $E$（它属于 $D(\mathcal{O}_X)$）
{\it 相对于 $Y$ 完美}，或称它 {\it $Y$-完美}，如果 $E$ 拟凝聚
（《位点上的上同调》定义
\ref{sites-cohomology-definition-pseudo-coherent}），并且 $E$
作为 $D(f^{-1}\mathcal{O}_Y)$ 的对象局部具有有限 Tor 维数
（《位点上的上同调》定义
\ref{sites-cohomology-definition-tor-amplitude}）。
\end{definition}

\noindent
有关讨论请见《概形的导出范畴》评注
\ref{perfect-remark-discuss-rel-perfect}；这里仅指出：由引理
\ref{spaces-more-morphisms-lemma-relative-pseudo-coherent-is-moot}，$E$ 拟凝聚等价于
$E$ 相对于 $Y$ 拟凝聚。此外，$E$ 的拟凝聚性蕴含
$E \in D_\QCoh(\mathcal{O}_X)$；见《空间的导出范畴》引理
\ref{spaces-perfect-lemma-pseudo-coherent}。

\begin{example}
\label{spaces-more-morphisms-example-relatively-perfect-field}
设 $k$ 为域。设 $X$ 为 $k$ 上有限表示的代数空间
（特别地，$X$ 拟紧）。则对象 $E$（它属于 $D(\mathcal{O}_X)$）
为 $k$-完美，当且仅当它有界且拟凝聚（按定义），也就是说，
当且仅当它属于 $D^b_{\textit{Coh}}(X)$（由《空间的导出范畴》引理
\ref{spaces-perfect-lemma-identify-pseudo-coherent-noetherian}）。
所以相对完美{\bf 并不}意味着“在纤维上完美”。
\end{example}

\noindent
相应的代数概念在《代数进阶》第
\ref{more-algebra-section-relatively-perfect} 节中研究。
可以如下把代数空间的概念与代数概念联系起来。

\begin{lemma}
\label{spaces-more-morphisms-lemma-affine-locally-rel-perfect}
设 $S$ 为概形。设 $f : X \to Y$ 为 $S$ 上代数空间的态射，
并且它平坦且局部有限表示。设
$E \in D_\QCoh(\mathcal{O}_X)$。下列条件等价：
\begin{enumerate}
\item $E$ 是 $Y$-完美的；
\item 对每个交换图
$$
\xymatrix{
U \ar[d] \ar[r]_g & V \ar[d] \\
X \ar[r]^f & Y
}
$$
其中 $U$、$V$ 是概形且竖直箭头平展，复形 $E|_U$
在《概形的导出范畴》定义
\ref{perfect-definition-relatively-perfect} 的意义下是 $V$-完美的；
\item 对某个如 (2) 的交换图，其中 $U \to X$ 为满态射，
复形 $E|_U$ 在《概形的导出范畴》定义
\ref{perfect-definition-relatively-perfect} 的意义下是 $V$-完美的；
\item 对每个如 (2) 的交换图，其中 $U$ 和 $V$ 仿射，
复形 $R\Gamma(U, E)$ 是 $\mathcal{O}_Y(V)$-完美的。
\end{enumerate}
\end{lemma}

\begin{proof}
为使本引理的 (2)、(3)、(4) 有意义，注意对象
$E|_U$（它属于 $D_\QCoh(\mathcal{O}_U)$）对应于对象
$E_0$（它属于 $D_\QCoh(\mathcal{O}_{U_0})$），其中 $U_0$
表示 $U$ 的底层概形；见《空间的导出范畴》引理
\ref{spaces-perfect-lemma-derived-quasi-coherent-small-etale-site}。
此外，在这种情形中，$E_0$ 拟凝聚，当且仅当 $E|_U$ 拟凝聚；
见《空间的导出范畴》引理
\ref{spaces-perfect-lemma-descend-pseudo-coherent}。
同样，由《空间的导出范畴》引理
\ref{spaces-perfect-lemma-tor-dimension-rel}，$E|_U$ 在
$f^{-1}\mathcal{O}_Y|_U = g^{-1}\mathcal{O}_V$ 上局部具有有限 Tor 维数，
当且仅当 $E_0$ 在 $g_0^{-1}\mathcal{O}_{V_0}$ 上局部具有有限 Tor 维数。
这里 $g_0 : U_0 \to V_0$ 是表示 $g : U \to V$ 的概形态射
（记号如《空间的导出范畴》评注
\ref{spaces-perfect-remark-match-total-direct-images}）。
最后，注意“拟凝聚”是平展局部性质，当然，“局部具有有限 Tor 维数”
也是平展局部性质。因此，只需平展局部地检验 $Y$-完美性；
由上面的讨论，(1) 蕴含 (2)，而 (3) 蕴含 (1)。由于
《概形的导出范畴》引理
\ref{perfect-lemma-affine-locally-rel-perfect} 说明 (4) 与 (2)、(3) 等价，
证明完毕。
\end{proof}

\begin{lemma}
\label{spaces-more-morphisms-lemma-triangulated}
设 $S$ 为概形。设 $f : X \to Y$ 为 $S$ 上代数空间的态射，
并且它平坦且局部有限表示。$D(\mathcal{O}_X)$ 中由
$Y$-完美对象组成的全子范畴是饱和的\footnote{《导出范畴》定义
\ref{derived-definition-saturated}。}三角子范畴。
\end{lemma}

\begin{proof}
这由《位点上的上同调》引理
\ref{sites-cohomology-lemma-cone-pseudo-coherent}、
\ref{sites-cohomology-lemma-summands-pseudo-coherent}、
\ref{sites-cohomology-lemma-cone-tor-amplitude} 和
\ref{sites-cohomology-lemma-summands-tor-amplitude} 得到。
\end{proof}

\begin{lemma}
\label{spaces-more-morphisms-lemma-perfect-relatively-perfect}
设 $S$ 为概形。设 $f : X \to Y$ 为 $S$ 上代数空间的态射，
并且它平坦且局部有限表示。$D(\mathcal{O}_X)$ 的完美对象是
$Y$-完美的。如果 $K, M \in D(\mathcal{O}_X)$，则
$K \otimes_{\mathcal{O}_X}^\mathbf{L} M$ 为 $Y$-完美，只要
$K$ 完美且 $M$ 为 $Y$-完美。
\end{lemma}

\begin{proof}
使用引理 \ref{spaces-more-morphisms-lemma-affine-locally-rel-perfect} 约化到概形的情形，
然后应用《概形的导出范畴》引理
\ref{perfect-lemma-perfect-relatively-perfect}。
\end{proof}

\begin{lemma}
\label{spaces-more-morphisms-lemma-base-change-relatively-perfect}
设 $S$ 为概形。设 $f : X \to Y$ 为 $S$ 上代数空间的态射，
并且它平坦且局部有限表示。设 $g : Y' \to Y$
为 $S$ 上代数空间的态射。令 $X' = Y' \times_Y X$，并记
% P09-ERR-0759
$g' : X' \to X$ 为投影。如果 $K \in D(\mathcal{O}_X)$
为 $Y$-完美的，那么 $L(g')^*K$ 为 $Y'$-完美的。
\end{lemma}

\begin{proof}
使用引理 \ref{spaces-more-morphisms-lemma-affine-locally-rel-perfect} 约化到概形的情形，
然后应用《概形的导出范畴》引理
\ref{perfect-lemma-base-change-relatively-perfect}。
\end{proof}

\begin{situation}
\label{spaces-more-morphisms-situation-relative-descent}
设 $S$ 为概形。设 $Y = \lim_{i \in I} Y_i$ 为 $S$ 上
代数空间有向系的极限，其仿射转移态射为
$g_{i'i} : Y_{i'} \to Y_i$。假设 $Y_i$ 对每个 $i \in I$
都拟紧且拟分离。记 $g_i : Y \to Y_i$ 为投影。固定元素
$0 \in I$ 和有限表示平坦态射 $X_0 \to Y_0$。
令 $X_i = Y_i \times_{Y_0} X_0$ 和 $X = Y \times_{Y_0} X_0$，
并记转移态射 $f_{i'i} : X_{i'} \to X_i$ 以及投影
$f_i : X \to X_i$。
\end{situation}

\begin{lemma}
\label{spaces-more-morphisms-lemma-relative-descend-homomorphisms}
在情形 \ref{spaces-more-morphisms-situation-relative-descent} 中，设 $K_0$ 和 $L_0$
为 $D(\mathcal{O}_{X_0})$ 的对象。令
$K_i = Lf_{i0}^*K_0$ 和 $L_i = Lf_{i0}^*L_0$（其中 $i \geq 0$），并令
$K = Lf_0^*K_0$ 和 $L = Lf_0^*L_0$。则映射
$$
\colim_{i \geq 0} \Hom_{D(\mathcal{O}_{X_i})}(K_i, L_i)
\longrightarrow
\Hom_{D(\mathcal{O}_X)}(K, L)
$$
是同构，条件是 $K_0$ 拟凝聚，并且
$L_0 \in D_\QCoh(\mathcal{O}_{X_0})$ 作为
$D((X_0 \to Y_0)^{-1}\mathcal{O}_{Y_0})$ 的对象
（局部）具有有限 Tor 维数。
% P09-ERR-0760
\end{lemma}

\begin{proof}
对每个拟紧拟分离对象 $U_0$（它属于 $(X_0)_{spaces, \etale}$），
考虑条件 $P$：
$$
\colim_{i \geq 0} \Hom_{D(\mathcal{O}_{U_i})}(K_i|_{U_i}, L_i|_{U_i})
\longrightarrow
\Hom_{D(\mathcal{O}_U)}(K|_U, L|_U)
$$
是同构，其中 $U = X \times_{X_0} U_0$ 且
$U_i = X_i \times_{X_0} U_0$。我们将证明 $P$ 对每个 $U_0$ 成立。

\medskip\noindent
假设 $(U_0 \subset W_0, V_0 \to W_0)$ 是
$(X_0)_{spaces, \etale}$ 中的初等特异方形，并且 $P$ 对
$U_0, V_0, U_0 \times_{W_0} V_0$ 成立。
则由导出范畴中 Hom 的 Mayer--Vietoris 性质，$P$ 对 $W_0$ 成立；
见《空间的导出范畴》引理
\ref{spaces-perfect-lemma-mayer-vietoris-hom}。

\medskip\noindent
首先考虑 $U_0 = W_0 \times_{Y_0} X_0$，其中 $W_0$ 是
$(Y_0)_{spaces, \etale}$ 的拟紧拟分离对象。
把《空间的导出范畴》引理
\ref{spaces-perfect-lemma-induction-principle} 的归纳原理
应用于这些 $W_0$ 以及上一段，便发现只需证明 $P$，
而只考虑 $U_0 = W_0 \times_{Y_0} X_0$ 且 $W_0$ 仿射的情形。
换言之，已经约化到 $Y_0$ 仿射的情形。接着再次应用归纳原理，
这次应用于 $X_0$ 的所有拟紧拟分离开集，从而进一步约化到
$X_0$ 也仿射的情形。

\medskip\noindent
如果 $X_0$ 和 $Y_0$ 仿射，就回到了概形情形；该情形已在
《概形的导出范畴》引理
\ref{perfect-lemma-relative-descend-homomorphisms} 中证明。
读者可以使用《空间的导出范畴》引理
\ref{spaces-perfect-lemma-pseudo-coherent}、
\ref{spaces-perfect-lemma-derived-quasi-coherent-small-etale-site}、
\ref{spaces-perfect-lemma-descend-pseudo-coherent} 和
\ref{spaces-perfect-lemma-tor-dimension-rel}，把陈述转换为仅涉及
概形及概形上模的导出范畴的陈述。
\end{proof}

\begin{lemma}
\label{spaces-more-morphisms-lemma-descend-relatively-perfect}
在情形 \ref{spaces-more-morphisms-situation-relative-descent} 中，
由 $Y$-完美对象组成的范畴（这些对象属于 $D(\mathcal{O}_X)$）
是由 $Y_i$-完美对象组成的范畴（这些对象属于
$D(\mathcal{O}_{X_i})$）的余极限。
\end{lemma}

\begin{proof}
对每个拟紧拟分离对象 $U_0$（它属于 $(X_0)_{spaces, \etale}$），
考虑条件 $P$：函子
$$
\colim_{i \geq 0} D_{Y_i\text{-完美}}(\mathcal{O}_{U_i})
\longrightarrow
D_{Y\text{-完美}}(\mathcal{O}_U)
$$
是等价，其中 $U = X \times_{X_0} U_0$ 且
$U_i = X_i \times_{X_0} U_0$。由引理
\ref{spaces-more-morphisms-lemma-relative-descend-homomorphisms}，已知这个函子全忠实。
因此只需证明本质满。

\medskip\noindent
假设 $(U_0 \subset W_0, V_0 \to W_0)$ 是
$(X_0)_{spaces, \etale}$ 中的初等特异方形，并且 $P$ 对
$U_0, V_0, U_0 \times_{W_0} V_0$ 成立。我们断言 $P$ 对 $W_0$ 成立。
将使用记号 $U_i = X_i \times_{X_0} U_0$、
$U = X \times_{X_0} U_0$，并对 $V_0$ 和 $W_0$ 使用类似记号。
我们将滥用符号 $f_i$ 来表示所有态射 $U \to U_i$、$V \to V_i$、
$U \times_W V \to U_i \times_{W_i} V_i$ 和 $W \to W_i$。
% P09-ERR-0761
设 $E$ 是 $Y$-完美对象，它属于 $D(\mathcal{O}_W)$。
目标：证明 $E$ 在该函子的本质像中。由假设，可以找到
$i \geq 0$、$Y_i$-完美对象 $E_{U, i}$（在 $U_i$ 上）、
$Y_i$-完美对象 $E_{V, i}$（在 $V_i$ 上），以及同构
$Lf_i^*E_{U, i} \to E|_U$ 和 $Lf_i^*E_{V, i} \to E|_V$。
令
% P09-ERR-0762
$$
a : E_{U, i} \to (Rf_{i, *}E)|_{U_i}
\quad\text{且}\quad
b : E_{V, i} \to (Rf_{i, *}E)|_{V_i}
$$
为与同构 $Lf_i^*E_{U, i} \to E|_U$ 和
$Lf_i^*E_{V, i} \to E|_V$ 伴随的映射。
% P09-ERR-0763
由全忠实性，增大 $i$ 后，可以找到同构
$c : E_{U, i}|_{U_i \times_{W_i} V_i} \to E_{V, i}|_{U_i \times_{W_i} V_i}$，
其拉回为如下复合识别：
$$
Lf_i^*E_{U, i}|_{U \times_W V} \to E|_{U \times_W V} \to
Lf_i^*E_{V, i}|_{U \times_W V}.
$$
应用《空间的导出范畴》引理
\ref{spaces-perfect-lemma-glue}，得到对象 $E_i$（在 $W_i$ 上）
以及映射 $d : E_i \to Rf_{i, *}E$，它限制为映射 $a$ 和 $b$，
分别在 $U_i$ 和 $V_i$ 上。于是显然，$E_i$ 是 $Y_i$-完美的
（因为相对完美是平展局部性质），而 $d$ 与同构
$Lf_i^*E_i \to E$ 伴随。

\medskip\noindent
完全照搬引理 \ref{spaces-more-morphisms-lemma-relative-descend-homomorphisms} 的证明中
使用归纳原理（《空间的导出范畴》引理
\ref{spaces-perfect-lemma-induction-principle}）的论证，约化到
$X_0$ 和 $Y_0$ 都仿射的情形：先处理
$(Y_0)_{spaces, \etale}$ 中的拟紧拟分离对象，把问题约化到
$Y_0$ 仿射；再处理 $(X_0)_{spaces, \etale}$ 中的拟紧拟分离对象，
% P09-ERR-0764
把问题约化到 $X_0$ 仿射。在仿射情形，结论由概形情形得到，
即《概形的导出范畴》引理
\ref{perfect-lemma-descend-relatively-perfect}。
到概形情形的转换由引理
\ref{spaces-more-morphisms-lemma-affine-locally-rel-perfect} 完成。
\end{proof}

\begin{lemma}
\label{spaces-more-morphisms-lemma-derived-pushforward-rel-perfect}
设 $S$ 为概形。设 $f : X \to Y$ 为 $S$ 上代数空间的态射，
并且它平坦、固有且有限表示。设 $E \in D(\mathcal{O}_X)$
为 $Y$-完美的。则 $Rf_*E$ 是 $D(\mathcal{O}_Y)$ 的完美对象，
并且它的构造与任意基变换交换。
\end{lemma}

\begin{proof}
关于基变换的陈述是《空间的导出范畴》引理
\ref{spaces-perfect-lemma-base-change-tensor}（其中
$\mathcal{G}^\bullet$ 等于 $\mathcal{O}_X$，并集中在次数 $0$）。
所以只需证明 $Rf_*E$ 是完美对象。我们将通过极限论证约化到
$Y$ 为诺特仿射的情形。

\medskip\noindent
这个问题在 $Y$ 上平展局部，故可假设 $Y$ 仿射。
记 $Y = \Spec(R)$。把 $R = \colim R_i$ 写成诺特环 $R_i$
的滤过余极限。由《空间的极限》引理
\ref{spaces-limits-lemma-descend-finite-presentation}，存在 $i$
和代数空间 $X_i$，它在 $R_i$ 上有限表示，并使它到 $R$
的基变换为 $X$。由《空间的极限》引理
\ref{spaces-limits-lemma-eventually-proper} 和
\ref{spaces-limits-lemma-descend-flat}，可以假设 $X_i$
在 $R_i$ 上固有且平坦。由引理
\ref{spaces-more-morphisms-lemma-descend-relatively-perfect}，可以假设存在
% P09-ERR-0765
$R_i$-完美对象 $E_i$（它属于 $D(\mathcal{O}_{X_i})$），
其到 $X$ 的拉回为 $E$。把《空间的导出范畴》引理
\ref{spaces-perfect-lemma-perfect-direct-image} 应用于
$X_i \to \Spec(R_i)$ 和 $E_i$，再使用已经证明的基变换性质，
便得到结论。
\end{proof}

\begin{lemma}
\label{spaces-more-morphisms-lemma-compute-ext-rel-perfect}
设 $S$ 为概形。设 $f : X \to Y$ 为 $S$ 上代数空间的态射。
设 $E, K \in D(\mathcal{O}_X)$。假设
\begin{enumerate}
\item $Y$ 拟紧且拟分离；
\item $f$ 固有、平坦且有限表示；
\item $E$ 是 $Y$-完美的；
\item $K$ 拟凝聚。
\end{enumerate}
则存在拟凝聚的 $L \in D(\mathcal{O}_Y)$，使得
$$
Rf_*R\SheafHom(K, E) = R\SheafHom(L, \mathcal{O}_Y)
$$
而且任意基变换后同样成立：给定
$$
\vcenter{
\xymatrix{
X' \ar[r]_{g'} \ar[d]_{f'} &
X \ar[d]^f \\
Y' \ar[r]^g &
Y
}
}
\quad\quad
\begin{matrix}
\text{若该图为笛卡尔图，则有} \\
Rf'_*R\SheafHom(L(g')^*K, L(g')^*E) \\
= R\SheafHom(Lg^*L, \mathcal{O}_{Y'})
\end{matrix}
$$
\end{lemma}

\begin{proof}
由于 $Y$ 拟紧且拟分离，$X$ 也如此。
由《空间的导出范畴》引理
\ref{spaces-perfect-lemma-pseudo-coherent-hocolim}，可以写成
$K = \text{hocolim} K_n$，其中 $K_n$ 完美，并且 $K_n \to K$
在截断 $\tau_{\geq -n}$ 上诱导同构。设 $K_n^\vee$ 为对偶完美复形
（《位点上的上同调》引理
\ref{sites-cohomology-lemma-dual-perfect-complex}）。得到完美对象的逆系
$\ldots \to K_3^\vee \to K_2^\vee \to K_1^\vee$。
由引理 \ref{spaces-more-morphisms-lemma-perfect-relatively-perfect}，
$K_n^\vee \otimes_{\mathcal{O}_X} E$ 是 $Y$-完美的。
因此，可以把引理 \ref{spaces-more-morphisms-lemma-derived-pushforward-rel-perfect}
应用于 $K_n^\vee \otimes_{\mathcal{O}_X} E$，得到逆系
$$
\ldots \to M_3 \to M_2 \to M_1
$$
它由 $Y$ 上的完美复形组成，其中
$$
M_n = Rf_*(K_n^\vee \otimes_{\mathcal{O}_X}^\mathbf{L} E) =
Rf_*R\SheafHom(K_n, E)
$$
此外，这些复形的构造与任意基变换交换；具体地，
$Lg^*M_n =
Rf'_*((L(g')^*K_n)^\vee \otimes_{\mathcal{O}_{X'}}^\mathbf{L} L(g')^*E) =
Rf'_*R\SheafHom(L(g')^*K_n, L(g')^*E)$。

\medskip\noindent
由于 $K_n \to K$ 在 $\tau_{\geq -n}$ 上诱导同构，可知
$K_n \to K_{n + 1}$ 在 $\tau_{\geq -n}$ 上诱导同构。
因此，$K_{n + 1}^\vee \to K_n^\vee$ 在 $\tau_{\leq n}$ 上诱导同构，
因为 $K_n^\vee = R\SheafHom(K_n, \mathcal{O}_X)$。
假设 $E$ 的 Tor 振幅在 $[a, b]$ 中，这里把它看作
$f^{-1}\mathcal{O}_Y$-模复形。则任意基变换后也是如此；见《空间的导出范畴》引理
\ref{spaces-perfect-lemma-tor-independence-and-tor-amplitude}。
可以看出映射
$K_{n + 1}^\vee \otimes_{\mathcal{O}_X} E \to
K_n^\vee \otimes_{\mathcal{O}_X} E$
在 $\tau_{\leq n + a}$ 上诱导同构，任意基变换后也同样成立。
应用右导出函子 $Rf_*$，得到映射 $M_{n + 1} \to M_n$
在 $\tau_{\leq n + a}$ 上诱导同构，任意基变换后也同样成立。
取特异三角
$$
M_{n + 1} \to M_n \to C_n \to M_{n + 1}[1]
$$
取 $y \in |Y|$。选择初等平展邻域
$(U, u) \to (Y, y)$；这是可行的，见《良好空间》引理
\ref{decent-spaces-lemma-decent-space-elementary-etale-neighbourhood}。
取 $Y'$ 为 $u$ 处剩余域的谱。拉回后可见
$C_n|_U \otimes_{\mathcal{O}_U}^\mathbf{L} \kappa(u)$
仅在次数 $\geq n + a$ 具有非零上同调。由《代数进阶》引理
\ref{more-algebra-lemma-lift-perfect-from-residue-field}，
可见完美复形 $C_n|_U$ 的 Tor 振幅在 $[n + a, m_n]$ 中，
其中 $m_n$ 为某个整数；必要时再缩小 $U$。
所以 $C_n$ 的 Tor 振幅在 $[n + a, m_n]$ 中，其中 $m_n$
为某个整数（因为 $Y$ 拟紧）。特别地，对偶完美复形
$C_n^\vee$ 的 Tor 振幅在 $[-m_n, -n - a]$ 中。

\medskip\noindent
令 $L_n = M_n^\vee$ 为对偶完美复形。上一段讨论的结论是，
$L_n \to L_{n + 1}$ 在 $\tau_{\geq -n - a}$ 上诱导同构。
所以 $L = \text{hocolim} L_n$ 拟凝聚；见《空间的导出范畴》引理
\ref{spaces-perfect-lemma-pseudo-coherent-hocolim}。由于
$$
R\SheafHom(K, E) = R\SheafHom(\text{hocolim} K_n, E) =
R\lim R\SheafHom(K_n, E) = R\lim K_n^\vee \otimes_{\mathcal{O}_X} E
$$
（《位点上的上同调》引理
\ref{sites-cohomology-lemma-colim-and-lim-of-duals}），并且
$R\lim$ 与 $Rf_*$ 交换，得到
$$
Rf_*R\SheafHom(K, E) = R\lim M_n = R\lim R\SheafHom(L_n, \mathcal{O}_Y) =
R\SheafHom(L, \mathcal{O}_Y)
$$
这证明了 $Y$ 上的公式。由于 $M_n$ 的构造与
% P09-ERR-0766
基变换相容，任意基变换后公式仍成立。
\end{proof}

\begin{remark}
\label{spaces-more-morphisms-remark-compare-L}
读者可能已经注意到引理 \ref{spaces-more-morphisms-lemma-compute-ext-rel-perfect}
与《空间的导出范畴》引理
\ref{spaces-perfect-lemma-compute-ext} 之间的相似性。
事实上，引理 \ref{spaces-more-morphisms-lemma-compute-ext-rel-perfect} 的拟凝聚复形 $L$
可刻画为 $Y$ 上唯一的拟凝聚复形，使得存在函子性同构
$$
\Ext^i_{\mathcal{O}_Y}(L, \mathcal{F}) \longrightarrow
\Ext^i_{\mathcal{O}_X}(K,
E \otimes_{\mathcal{O}_X}^\mathbf{L} Lf^*\mathcal{F})
$$
当 $\mathcal{F}$ 遍历 $\QCoh(\mathcal{O}_Y)$ 时，这些同构与边界映射相容。
如果以后需要这一点，我们将在此陈述精确结果并给出详细证明。
\end{remark}

\begin{lemma}
\label{spaces-more-morphisms-lemma-bounded-on-fibres}
设 $S$ 为概形。设 $X$ 为 $S$ 上代数空间，使得结构态射
$f : X \to S$ 平坦且局部有限表示。设 $E$ 为
$D(\mathcal{O}_X)$ 的拟凝聚对象。下列条件等价：
\begin{enumerate}
\item $E$ 是 $S$-完美的；
\item $E$ 局部下有界，并且对每个点 $s \in S$，
对象 $L(X_s \to X)^*E$（它属于 $D(\mathcal{O}_{X_s})$）局部下有界。
\end{enumerate}
\end{lemma}

\begin{proof}
由于一切都是局部的，立即约化到 $X$ 和 $S$ 仿射的情形；
见引理 \ref{spaces-more-morphisms-lemma-affine-locally-rel-perfect}。这一情形由
《概形的导出范畴》引理 \ref{perfect-lemma-bounded-on-fibres} 处理。
\end{proof}

\begin{lemma}
\label{spaces-more-morphisms-lemma-characterize-relatively-perfect}
设 $A$ 为环。设 $X$ 为在 $A$ 上可分、有限表示且平坦的代数空间。
设 $K \in D_\QCoh(\mathcal{O}_X)$。如果
$R \Gamma (X, E \otimes^\mathbf{L} K)$ 是 $D(A)$ 中的完美对象，
并且这对每个完美的 $E \in D(\mathcal{O}_X)$ 都成立，
那么 $K$ 是 $\Spec(A)$-完美的。
\end{lemma}

\begin{proof}
由引理 \ref{spaces-more-morphisms-lemma-characterize-pseudo-coh-improved}，$K$
相对于 $A$ 拟凝聚。由引理
\ref{spaces-more-morphisms-lemma-relative-pseudo-coherent-is-moot}，$K$
在 $D(\mathcal{O}_X)$ 中拟凝聚。由《空间的导出范畴》命题
\ref{spaces-perfect-proposition-detecting-bounded-below}，
可见 $K$ 属于 $D^-(\mathcal{O}_X)$。设 $\mathfrak{p}$ 为 $A$ 的素理想，并记
% P09-ERR-0767
$i : Y \to X$ 为 $\mathfrak{p}$ 上概形论纤维的包含；也就是说，
% P09-ERR-0768
$Y$ 是 $\kappa(\mathfrak p)$ 上的概形。由引理
\ref{spaces-more-morphisms-lemma-bounded-on-fibres}，只要能证明 $Li^*(K)$ 下有界，
结论就成立。设 $G \in D_{perf} (\mathcal{O}_X)$ 为生成
$D_\QCoh (\mathcal{O}_X)$ 的完美复形；见《空间的导出范畴》定理
\ref{spaces-perfect-theorem-bondal-van-den-Bergh}。有
\begin{align*}
R\Hom _{\mathcal{O}_Y}(Li^*(G), Li^*(K))
& =
R\Gamma(Y, Li^*(G ^\vee \otimes ^\mathbf{L} K)) \\
& =
R\Gamma(X, G^\vee \otimes ^{\mathbf{L}} K)
\otimes^\mathbf{L}_A \kappa(\mathfrak{p})
\end{align*}
第一个等式使用了 $Li^*$ 保持完美对象和对偶这一事实以及
《位点上的上同调》引理
\ref{sites-cohomology-lemma-dual-perfect-complex}；略去一些细节。
第二个等式由《空间的导出范畴》引理
\ref{spaces-perfect-lemma-compare-base-change} 得到，因为 $X$
在 $A$ 上平坦。由假设，这给出 $D(\kappa(\mathfrak{p}))$
中的完美对象。对象 $Li^*(G) \in D_{perf}(\mathcal{O}_Y)$
由《空间的导出范畴》评注
\ref{spaces-perfect-remark-pullback-generator} 生成
$D_\QCoh(\mathcal{O}_Y)$。于是《空间的导出范畴》命题
\ref{spaces-perfect-proposition-detecting-bounded-below}
蕴含 $Li^*(K)$ 下有界，从而结论成立。
\end{proof}

\section{立方定理}
\label{spaces-more-morphisms-section-theorem-cube}

\noindent
本节是《态射进阶》第
\ref{more-morphisms-section-theorem-cube} 节的对应版本。
下面的引理说明，在引理所作的假设下，Picard 函子的对角线
可由局部闭浸入表示。

\begin{lemma}
\label{spaces-more-morphisms-lemma-diagonal-picard-flat-proper}
设 $S$ 为概形。设 $f : X \to Y$ 为 $S$ 上代数空间的
平坦有限表示固有态射。设 $\mathcal{E}$ 为有限局部自由
$\mathcal{O}_X$-模。对态射 $g : Y' \to Y$，考虑基变换图
$$
\xymatrix{
X' \ar[d]_{f'} \ar[r]_{g'} & X \ar[d]^f \\
Y' \ar[r]^g & Y
}
$$
假设 $\mathcal{O}_{Y'} \to f'_*\mathcal{O}_{X'}$
对所有 $g : Y' \to Y$ 都是同构。
则存在有限表示浸入 $j : Z \to Y$，使得态射
$g : Y' \to Y$ 经由 $Z$ 分解，当且仅当存在有限局部自由
$\mathcal{O}_{Y'}$-模 $\mathcal{N}$，满足
% P09-ERR-0769
$(f')^*\mathcal{N} \cong (g')^*\mathcal{L}$。
\end{lemma}

\begin{proof}
设 $y : \Spec(k) \to Y$ 为域值点。纤维 $X_y$
（即 $f$ 在 $y$ 处的纤维）连通，这是因为我们的假设
$H^0(X_y, \mathcal{O}_{X_y}) = k$。所以 $\mathcal{E}$ 的秩在各纤维上恒定。
由于 $f$ 为开映射（《空间的态射》引理
\ref{spaces-morphisms-lemma-fppf-open}）且为闭映射，得到分解
$Y = \coprod Y_r$，把 $Y$ 分解为开闭子空间，使得
$\mathcal{E}$ 恒秩 $r$，这是在 $Y_r$ 的逆像上而言的。
所以可以假设 $\mathcal{E}$ 恒秩 $r$。记
% P09-ERR-0770
$\mathcal{E}^\vee = \SheafHom(\mathcal{E}, \mathcal{O}_X)$
为秩 $r$ 的对偶模。

\medskip\noindent
由上同调与基变换（更精确地，由《空间的导出范畴》引理
\ref{spaces-perfect-lemma-flat-proper-perfect-direct-image-general}），
可见 $E = Rf_*\mathcal{E}$ 是 $Y$ 的导出范畴中的完美对象，
并且它的构造与任意基变换交换。对
$E' = Rf_*\mathcal{E}^\vee$ 也同样如此。
由于在次数 $< 0$ 从来没有上同调，$E$ 和 $E'$ 的 Tor 振幅
（局部）包含在 $[0, b]$ 中，其中 $b$ 是某个数。
注意，对任意 $g : Y' \to Y$，有
$f'_*((g')^*\mathcal{E}) = H^0(Lg^*E)$ 和
$f'_*((g')^*\mathcal{E}^\vee) = H^0(Lg^*E')$。
设 $j : Z \to Y$ 和 $j' : Z' \to Y$ 为《空间的导出范畴》引理
\ref{spaces-perfect-lemma-locally-closed-where-H0-locally-free}
对 $E$ 和 $E'$（取 $a = 0$ 和 $r = r$）构造的局部闭浸入；
它们由如下性质刻画：
% P09-ERR-0771
$H^0(Lj^*E)$ 和 $H^0((j')^*E')$ 是秩 $r$ 的局部自由模，
并且与拉回相容。

\medskip\noindent
设 $g : Y' \to Y$ 为态射。如果存在引理中所述的
$\mathcal{N}$，那么使用投影公式《位点上的上同调》引理
\ref{sites-cohomology-lemma-projection-formula}，可见模
$$
f'_*((g')^*\mathcal{E}) \cong
f'_*((f')^*\mathcal{N}) \cong
\mathcal{N} \otimes_{\mathcal{O}_{Y'}} f'_*\mathcal{O}_{X'} \cong
\mathcal{N}\quad\text{且类似地}\quad
f'_*((g')^*\mathcal{E}^\vee) \cong \mathcal{N}^\vee
$$
局部自由且秩为 $r$，并且仍局部自由且秩为 $r$，
即使经过任意进一步的基变换 $Y'' \to Y'$ 也是如此。
所以在这种情形，$g : Y' \to Y$ 经由 $j$ 和 $j'$ 分解。
因此，可以替换 $Y$，所用空间为 $Z \times_Y Z'$，并假设
$f_*\mathcal{E}$ 和 $f_*\mathcal{E}^\vee$ 是局部自由
$\mathcal{O}_Y$-模且秩为 $r$，其构造与任意基变换交换。

\medskip\noindent
在这种情形下，如果 $g : Y' \to Y$ 为态射，并且存在引理中所述的
$\mathcal{N}$，那么映射（次数 $0$ 的杯积）
$$
f'_*((g')^*\mathcal{E})
\otimes_{\mathcal{O}_{Y'}}
f'_*((g')^*\mathcal{E}^\vee)
\longrightarrow \mathcal{O}_{Y'}
$$
是完美配对。反过来，如果这个杯积映射是完美配对，
那么在 $Y'$ 上局部地可找到截面的基
% P09-ERR-0769
$\sigma_1, \ldots, \sigma_r$（属于 $f'_*((g')^*\mathcal{L})$）以及
$\tau_1, \ldots, \tau_r$（属于
$f'_*((g')^*\mathcal{E}^\vee)$），其乘积满足
$\sigma_i \tau_j = \delta_{ij}$。把 $\sigma_i$ 看作
$(g')^*\mathcal{L}$ 的截面（在 $X'$ 上），并把 $\tau_j$ 看作
$(g')^*\mathcal{L}^\vee$ 的截面（在 $X'$ 上），得到
$$
\sigma_1, \ldots, \sigma_r :
\mathcal{O}_{X'}^{\oplus r}
\longrightarrow
(g')^*\mathcal{E}
$$
是同构，其逆由
$$
\tau_1, \ldots, \tau_r :
(g')^*\mathcal{E}
\longrightarrow
\mathcal{O}_{X'}^{\oplus r}
$$
给出。换言之，
$(f')^*f'_*(g')^*\mathcal{E} \cong (g')^*\mathcal{E}$。
而杯积非退化这一条件切出一个逆紧开子概形
（即某个适当行列式非零的轨迹），证明完毕。
\end{proof}

\section{复形有限性性质的下降}
\label{spaces-more-morphisms-section-descent-finiteness}

\noindent
本节是《态射进阶》第
\ref{more-morphisms-section-descent-finiteness} 节和
《概形的导出范畴》第
\ref{perfect-section-descent-finiteness} 节的对应版本。

\begin{lemma}
\label{spaces-more-morphisms-lemma-pseudo-coherent-descends-fpqc}
设 $S$ 为概形。设 $\{f_i : X_i \to X\}$ 为 $S$ 上代数空间的
fpqc 覆盖。设 $E \in D_\QCoh(\mathcal{O}_X)$，并设
$m \in \mathbf{Z}$。则 $E$ 是 $m$-拟凝聚的，当且仅当每个
$Lf_i^*E$ 都是 $m$-拟凝聚的。
\end{lemma}

\begin{proof}
拉回总保持 $m$-拟凝聚性，见《位点上的上同调》引理
\ref{sites-cohomology-lemma-pseudo-coherent-pullback}。
所以只需假设 $Lf_i^*E$ 是 $m$-拟凝聚的，并证明 $E$
是 $m$-拟凝聚的。首先，可以假设 $X_i$ 对所有 $i$ 都是概形，见
《空间上的拓扑》引理
\ref{spaces-topologies-lemma-refine-fpqc-schemes}。
% P09-ERR-0772
接着，选取满 étale 态射 $U \to X$，其中 $U$ 是概形。
则 $U_i = U \times_X X_i$ 是概形，并得到概形的 fpqc 覆盖
$\{U_i \to U\}$，见《空间上的拓扑》引理
\ref{spaces-topologies-lemma-recognize-fpqc-covering}。
由概形情形可知结论对
$\{U_i \to U\}_{i \in I}$ 成立，见《概形的导出范畴》引理
\ref{perfect-lemma-pseudo-coherent-descends-fpqc}。
另一方面，限制 $E|_U$ 来自
$D_\QCoh(\mathcal{O}_U)$ 中的一个对象（这里使用 Zariski
拓扑和 $U$ 的“通常”结构层来定义），见《空间的导出范畴》引理
\ref{spaces-perfect-lemma-derived-quasi-coherent-small-etale-site}。
由于两种拟凝聚概念（étale 与 Zariski）由《空间的导出范畴》引理
\ref{spaces-perfect-lemma-descend-pseudo-coherent}
可知相同，故引理得证。
\end{proof}

\begin{lemma}
\label{spaces-more-morphisms-lemma-tor-amplitude-descends-fppf}
设 $S$ 为概形。设 $\{g_i : Y_i \to Y\}$ 为 $S$ 上代数空间的
fpqc 覆盖。设 $f : X \to Y$ 为代数空间的态射，并置
$X_i = Y_i \times_Y X$，投影记为 $f_i : X_i \to Y_i$
和 $g'_i : X_i \to X$。设 $E \in D_\QCoh(\mathcal{O}_X)$，
并设 $a, b \in \mathbf{Z}$。则下列条件等价：
\begin{enumerate}
\item $E$ 具有包含在 $[a, b]$ 中的 Tor 振幅，这是作为
$D(f^{-1}\mathcal{O}_Y)$ 的对象而言的；以及
% P09-ERR-0773
\item $L(g'_i)^*E$ 具有包含在 $[a, b]$ 中的 Tor 振幅，
这是作为 $D(f_i^{-1}\mathcal{O}_{Y_i})$ 的对象而言的，
并且对所有 $i$ 都如此。
\end{enumerate}
把“具有包含在 $[a, b]$ 中的 Tor 振幅”替换为
“局部具有有限 Tor 维数”，结论也成立。
\end{lemma}

\begin{proof}
% P09-ERR-0774
由《空间的导出范畴》引理
\ref{spaces-perfect-lemma-tor-independence-and-tor-amplitude}
可知，拉回保持“具有包含在 $[a, b]$ 中的 Tor 振幅”。
注意，$Y_i$ 和 $X$ 在 $Y$ 上 Tor 无关，因为
$Y_i \to Y$ 平坦。假设 (2)，来证明 (1)。Tor 维数可在茎处计算，见
《位点上的上同调》引理
\ref{sites-cohomology-lemma-tor-amplitude-stalk}
和《空间的性质》定理
\ref{spaces-properties-theorem-exactness-stalks}。
设 $\overline{x}$ 为 $X$ 的几何点。选取 $i$ 和几何点
$\overline{x}_i$（它位于 $X_i$ 中），使其像为
$\overline{x}$（在 $X$ 中）。则
% P09-ERR-0775
$$
(L(g_i')^*E)_{\overline{x}_i} =
E_{\overline{x}}
\otimes_{\mathcal{O}_{X, \overline{x}}}^\mathbf{L}
\mathcal{O}_{X_, \overline{x}_i}
$$
设 $\overline{y}_i$（在 $Y_i$ 中）和 $\overline{y}$（在 $Y$ 中）
分别为 $\overline{x}_i$ 和 $\overline{x}$ 的像。
由于 $X$ 和 $Y_i$ 在 $Y$ 上 Tor 无关，可应用《代数进阶》引理
\ref{more-algebra-lemma-base-change-comparison}，从而上式右端等于
$E_{\overline{x}}
\otimes_{\mathcal{O}_{Y, \overline{y}}}^\mathbf{L}
\mathcal{O}_{Y_i, \overline{y}_i}$，这是在
$D(\mathcal{O}_{Y_i, \overline{y}_i})$ 中而言的。
% P09-ERR-0776
由于已假设它的 Tor 振幅包含在 $[a, b]$ 中，由《代数进阶》引理
\ref{more-algebra-lemma-flat-descent-tor-amplitude}
可知 $E_{\overline{x}}$ 在
$D(\mathcal{O}_{Y, \overline{y}})$ 中的 Tor 振幅包含在
$[a, b]$ 中。因此 (1) 成立。

\medskip\noindent
再使用一些初等拓扑，“局部有限 Tor 维数”的情形也随之成立。
\end{proof}

\noindent
下面的引理其实并不属于本节。

\begin{lemma}
\label{spaces-more-morphisms-lemma-thickening-pseudo-coherent}
设 $S$ 为概形。设 $i : X \to X'$ 为代数空间的有限阶加厚。
设 $K' \in D(\mathcal{O}_{X'})$ 为对象，并且
$K = Li^*K'$ 是拟凝聚的。则 $K'$ 是拟凝聚的。
\end{lemma}

\begin{proof}
先证明 $K'$ 具有拟凝聚上同调层；建议读者跳过这一部分。
为此，由第 \ref{spaces-more-morphisms-section-thickenings} 节可约化到一阶加厚的情形。
设 $\mathcal{I} \subset \mathcal{O}_{X'}$ 为切出 $X$
的拟凝聚理想层。把短正合列
$$
0 \to \mathcal{I} \to \mathcal{O}_{X'} \to i_*\mathcal{O}_X \to 0
$$
与 $K'$ 作张量，得到特异三角
$$
K' \otimes_{\mathcal{O}_{X'}}^\mathbf{L} \mathcal{I}
\to K' \to
K' \otimes_{\mathcal{O}_{X'}}^\mathbf{L} i_*\mathcal{O}_X
\to
(K' \otimes_{\mathcal{O}_{X'}}^\mathbf{L} \mathcal{I})[1]
$$
由于 $i_* = Ri_*$，并且可以把 $\mathcal{I}$ 看作拟凝聚
$\mathcal{O}_X$-模（因为这里是一阶加厚），可将此改写为
$$
i_*(K \otimes_{\mathcal{O}_X}^\mathbf{L} \mathcal{I})
\to K' \to
i_*K \to
i_*(K \otimes_{\mathcal{O}_X}^\mathbf{L} \mathcal{I})[1]
$$
这些项的识别请使用《空间的上同调》引理
\ref{spaces-cohomology-lemma-projection-formula-finite}。
由于 $K$ 属于 $D_\QCoh(\mathcal{O}_X)$，可知
$K'$ 属于 $D_\QCoh(\mathcal{O}_{X'})$；这里使用《空间的导出范畴》引理
\ref{spaces-perfect-lemma-pseudo-coherent}、
\ref{spaces-perfect-lemma-quasi-coherence-tensor-product} 和
\ref{spaces-perfect-lemma-quasi-coherence-direct-image}。

\medskip\noindent
假设 $K'$ 属于 $D_\QCoh(\mathcal{O}_{X'})$。这个问题在 $X'$
上是 étale 局部的，所以可以假设 $X'$ 仿射。
在这种情形，结论由概形情形得出（《态射进阶》引理
\ref{more-morphisms-lemma-thickening-pseudo-coherent}）。
将其转写为概形语言时，使用《空间的导出范畴》引理
\ref{spaces-perfect-lemma-derived-quasi-coherent-small-etale-site} 和
\ref{spaces-perfect-lemma-descend-pseudo-coherent}，以及注记
\ref{spaces-perfect-remark-match-total-direct-images}。
\end{proof}

\begin{lemma}
\label{spaces-more-morphisms-lemma-thickening-relatively-perfect}
设 $S$ 为概形。考虑笛卡儿图
$$
\xymatrix{
X \ar[r]_i \ar[d]_f & X' \ar[d]^{f'} \\
Y \ar[r]^j & Y'
}
$$
这是 $S$ 上代数空间的图。假设 $X' \to Y'$ 平坦且局部有限表示，而 $Y \to Y'$
为有限阶加厚。设 $E' \in D(\mathcal{O}_{X'})$。如果
$E = Li^*(E')$ 是 $Y$-完美的，则 $E'$ 是 $Y'$-完美的。
\end{lemma}

\begin{proof}
回忆一下，$Y$-完美性对于 $E$ 意味着 $E$ 拟凝聚，并且作为
$f^{-1}\mathcal{O}_Y$-模的复形在局部具有有限 Tor 维数
（定义 \ref{spaces-more-morphisms-definition-relatively-perfect}）。
由引理 \ref{spaces-more-morphisms-lemma-thickening-pseudo-coherent}，$E'$ 拟凝聚。
特别地，$E'$ 属于 $D_\QCoh(\mathcal{O}_{X'})$，见
《空间的导出范畴》引理
\ref{spaces-perfect-lemma-pseudo-coherent}。
由引理 \ref{spaces-more-morphisms-lemma-affine-locally-rel-perfect}，这把问题约化到概形情形。
概形情形是《态射进阶》引理
\ref{more-morphisms-lemma-thickening-relatively-perfect}。
\end{proof}

\begin{lemma}
\label{spaces-more-morphisms-lemma-henselian-relatively-perfect}
设 $(R, I)$ 为由环和包含在 Jacobson 根中的理想 $I$ 所组成的偶。
置 $S = \Spec(R)$ 和 $S_0 = \Spec(R/I)$。设 $X$ 为 $R$
上的代数空间，其结构态射 $f : X \to S$ 固有、平坦且有限表示。
% P09-ERR-0777
记 $X_0 = S_0 \times_S X$。设
$E \in D(\mathcal{O}_X)$ 拟凝聚。如果导出限制 $E_0$
（即 $E$ 到 $X_0$ 的导出限制）是 $S_0$-完美的，
则 $E$ 是 $S$-完美的。
\end{lemma}

\begin{proof}
选取满 étale 态射 $U \to X$，其中 $U$ 仿射。
选取闭浸入 $U \to \mathbf{A}^d_S$。置
$U_0 = S_0 \times_S U$。复形 $E_0|_{U_0}$
的 Tor 振幅包含在某个区间 $[a, b]$ 中，其中
$a, b \in \mathbf{Z}$。设 $\overline{x}$ 为 $X$ 的几何点。
我们将证明 $E_{\overline{x}}$ 在 $R$ 上的 Tor 振幅包含在
$[a - d, b]$ 中。这样即可完成证明，因为由《位点上的上同调》引理
\ref{sites-cohomology-lemma-tor-amplitude-stalk}
和《空间的性质》定理
\ref{spaces-properties-theorem-exactness-stalks}，
Tor 振幅可从茎读出。

\medskip\noindent
设 $x \in |X|$ 为由 $\overline{x}$ 确定的点。
回忆 $|X| \to |S|$ 为闭映射（根据固有态射的定义）。
由于 $I$ 包含在 Jacobson 根中，$S$ 的每个非空闭子集
都含有闭子概形 $S_0$ 的一个点。因此，可找到特化
$x \leadsto x_0$（在 $|X|$ 中），其中 $x_0 \in |X_0|$。
选取 $u_0 \in U_0$，它映到 $x_0$。由《适度空间》引理
\ref{decent-spaces-lemma-generalizations-lift-flat}
（或由直接适用于 étale 态射的《适度空间》引理
\ref{decent-spaces-lemma-specialization}），可找到特化
$u \leadsto u_0$（在 $U$ 中），使得 $u$ 映到 $x$。可以把
$\overline{x}$ 提升为几何点 $\overline{u}$，它是 $U$ 的点且位于
$u$ 上方。
于是 $E_{\overline{x}} = (E|_U)_{\overline{u}}$。

\medskip\noindent
% P09-ERR-0778
写成 $U = \Spec(A)$。则 $A$ 是平坦有限表示的 $R$-代数，
并且是多项式 $R$-代数的商，此多项式代数含 $d$ 个变量。
限制 $E|_U$（由《空间的导出范畴》引理
\ref{spaces-perfect-lemma-pseudo-coherent}、
\ref{spaces-perfect-lemma-derived-quasi-coherent-small-etale-site} 和
\ref{spaces-perfect-lemma-descend-pseudo-coherent}，
以及《概形的导出范畴》引理
\ref{perfect-lemma-affine-compare-bounded} 和
\ref{perfect-lemma-pseudo-coherent-affine}）对应于拟凝聚对象 $K$，
它属于 $D(A)$。
注意，$E_0$ 对应于 $K \otimes_A^\mathbf{L} A/IA$。
设 $\mathfrak q \subset \mathfrak q_0 \subset A$
为对应于 $u \leadsto u_0$ 的素理想。则
$$
E_{\overline{x}} =
(E|_U)_{\overline{u}} =
E_u \otimes_{\mathcal{O}_{U, u}}^\mathbf{L} \mathcal{O}_{U, \overline{u}} =
K_{\mathfrak q} \otimes_{A_\mathfrak q}^\mathbf{L} A_{\mathfrak q}^{sh}
$$
（略去了一些细节）。由于
$A_\mathfrak q \to A_\mathfrak q^{sh}$ 平坦，它作为
$R$-模的 Tor 振幅等于 $K_\mathfrak q$ 作为 $R$-模的
Tor 振幅（《代数进阶》引理
\ref{more-algebra-lemma-no-change-tor-amplitude}）。
此外，$K_{\mathfrak q}$ 是 $K_{\mathfrak q_0}$ 的局部化。
所以只需证明 $K_{\mathfrak q_0}$ 具有包含在
$[a - d, b]$ 中的 Tor 振幅，这是作为 $R$-模的复形而言的。

\medskip\noindent
设 $I \subset \mathfrak p_0 \subset R$ 为对应于
$f(x_0)$ 的素理想。则有
\begin{align*}
K \otimes_R^\mathbf{L} \kappa(\mathfrak p_0)
& =
(K \otimes_R^\mathbf{L} R/I) \otimes_{R/I}^\mathbf{L}
\kappa(\mathfrak p_0) \\
& =
(K \otimes_A^\mathbf{L} A/IA) \otimes_{R/I}^\mathbf{L} \kappa(\mathfrak p_0)
\end{align*}
% P09-ERR-0779
第二个等式成立是因为 $R \to A$ 平坦。
根据对 $a,b$ 的选择，这个复形仅在区间 $[a,b]$
内的次数上具有上同调。最后，将《代数进阶》引理
\ref{more-algebra-lemma-lift-from-fibre-relatively-perfect}
应用于 $R \to A$、$\mathfrak q_0$、$\mathfrak p_0$ 和 $K$，
即可得到结论。
\end{proof}

\section{结点曲线的范族}
\label{spaces-more-morphisms-section-families-nodal}

\noindent
本节接续《代数曲线》第
\ref{curves-section-families-nodal} 节。
关于我们的术语选择，也请参阅该节。

\medskip\noindent
概形态射的“相对维数为 $1$ 且至多带结点奇性”这一性质，
在源和目标上都是 étale 局部的，见《下降》引理
\ref{descent-lemma-etale-local-source-target}，以及
《代数曲线》引理
\ref{curves-lemma-nodal-family-etale-local-source}、
\ref{curves-lemma-nodal-family-fpqc-local-target} 和
\ref{curves-lemma-nodal-family-postcompose-etale}。
它还在基变换下稳定，并且在目标上是 fpqc 局部的，见
《代数曲线》引理 \ref{curves-lemma-base-change-nodal-family} 和
\ref{curves-lemma-nodal-family-fpqc-local-target}。因此，根据
《空间的态射》引理 \ref{spaces-morphisms-lemma-local-source-target}，
可如下为代数空间定义相对维数为 $1$ 且至多带结点奇性的态射；
当该态射可表时，这与《空间的态射》第
\ref{spaces-morphisms-section-representable} 节中已经定义的概念一致。

\begin{definition}
\label{spaces-more-morphisms-definition-nodal-family}
设 $S$ 是概形。
设 $f : X \to Y$ 是 $S$ 上代数空间的态射。
称 $f$ 是{\it 相对维数为 $1$ 且至多带结点奇性的}，
如果《空间的态射》引理
\ref{spaces-morphisms-lemma-local-source-target} 的等价条件对
$\mathcal{P} =$“相对维数为 $1$ 且至多带结点奇性”成立。
\end{definition}

\begin{lemma}
\label{spaces-more-morphisms-lemma-base-change-nodal}
相对维数为 $1$ 且至多带结点奇性这一性质在基变换下保持。
\end{lemma}

\begin{proof}
见《空间的态射》评注 \ref{spaces-morphisms-remark-base-change-P}
和《代数曲线》引理 \ref{curves-lemma-base-change-nodal-family}。
\end{proof}

\begin{lemma}
\label{spaces-more-morphisms-lemma-nodal-local}
设 $S$ 是概形。
设 $f : X \to Y$ 是 $S$ 上代数空间的态射。
下列条件等价：
\begin{enumerate}
\item $f$ 相对维数为 $1$ 且至多带结点奇性；
\item 对每个概形 $Z$ 和任意态射 $Z \to Y$，态射
$Z \times_Y X \to Z$ 都相对维数为 $1$ 且至多带结点奇性；
\item 对每个仿射概形 $Z$ 和任意态射
$Z \to Y$，态射 $Z \times_Y X \to Z$ 都
相对维数为 $1$ 且至多带结点奇性；
\item 存在概形 $V$ 和满 étale 态射
$V \to Y$，使得 $V \times_Y X \to V$
相对维数为 $1$ 且至多带结点奇性；
\item 存在概形 $U$ 和满 étale 态射
$\varphi : U \to X$，使得复合态射 $f \circ \varphi$
相对维数为 $1$ 且至多带结点奇性；
\item 对每个交换图
$$
\xymatrix{
U \ar[d] \ar[r] & V \ar[d] \\
X \ar[r] & Y
}
$$
% P09-ERR-0780
其中 $U$、$V$ 是概形，竖直箭头均为 étale 态射，则上方的
水平箭头相对维数为 $1$ 且至多带结点奇性；
\item 存在交换图
$$
\xymatrix{
U \ar[d] \ar[r] & V \ar[d] \\
X \ar[r] & Y
}
$$
% P09-ERR-0781
其中 $U$、$V$ 是概形，竖直箭头均为 étale 态射，
$U \to X$ 为满射，并且上方的水平箭头
相对维数为 $1$ 且至多带结点奇性；以及
\item 存在 Zariski 覆盖 $Y = \bigcup_{i \in I} Y_i$，
% P09-ERR-0782
以及 $f^{-1}(Y_i) = \bigcup X_{ij}$，使得每个态射
$X_{ij} \to Y_i$ 都相对维数为 $1$ 且至多带结点奇性。
\end{enumerate}
\end{lemma}

\begin{proof}
% P09-ERR-0783
略去。
\end{proof}

\noindent
下面的引理说明，是否相对维数为 $1$ 且至多带结点奇性，
可以在纤维上检验。

\begin{lemma}
\label{spaces-more-morphisms-lemma-locus-where-nodal}
设 $S$ 是概形。
设 $f : X \to Y$ 是 $S$ 上代数空间的态射，并且它平坦且局部有限表示。
则存在最大的开子空间 $X' \subset X$，使得
$f|_{X'} : X' \to Y$ 相对维数为 $1$ 且至多带结点奇性。
此外，$X'$ 的构造与任意基变换交换。
\end{lemma}

\begin{proof}
选取交换图
$$
\xymatrix{
U \ar[d] \ar[r]_h & V \ar[d] \\
X \ar[r]^f & Y
}
$$
其中 $U$、$V$ 是概形，竖直箭头均为 étale 态射，并且
$U \to X$ 是满射。由概形情形的引理
（《代数曲线》引理 \ref{curves-lemma-locus-where-nodal}），
可找到最大的开子概形 $U' \subset U$，使得
$h|_{U'} : U' \to V$ 相对维数为 $1$ 且至多带结点奇性，
并且 $U'$ 的构造与基变换交换。
令 $X' \subset X$ 为其点对应于开子集
$\Im(|U'| \to |X|)$ 的开子空间。
由引理 \ref{spaces-more-morphisms-lemma-nodal-local}，$X' \to Y$
相对维数为 $1$ 且至多带结点奇性，而且 $X'$ 是具有此性质的
最大开子空间（这也推出 $U'$ 是 $X'$ 在 $U$ 中的逆像，但我们
不需要这一点）。因为基变换后同样如此，证明完成。
\end{proof}




\section{分解性质}
\label{spaces-more-morphisms-section-resolution-property}

\noindent
我们继续《空间的导出范畴》第
\ref{spaces-perfect-section-resolution-property} 节中的讨论。

\begin{situation}
\label{spaces-more-morphisms-situation-descend-resolution-property}
设 $S$ 是概形。设 $X$ 是 $S$ 上拟紧且拟分离的代数空间。
设 $V \to X$ 是满 étale 态射，其中 $V$ 是仿射概形
（这样的态射由《空间的性质》引理
\ref{spaces-properties-lemma-quasi-compact-affine-cover} 保证存在）。
选取交换图
$$
\xymatrix{
V \ar[rd]_\varphi \ar[rr]_j & & Y \ar[ld]^\pi \\
& X
}
$$
其中 $j$ 是开浸入，而 $\pi$ 是代数空间的有限态射
（这样的图由引理
\ref{spaces-more-morphisms-lemma-quasi-finite-separated-pass-through-finite} 保证存在）。
设 $\mathcal{I} \subset \mathcal{O}_Y$ 是 $Y$ 上有限型拟凝聚理想层，满足
$V(\mathcal{I}) = Y \setminus j(V)$（这样的理想层由《空间的极限》引理
\ref{spaces-limits-lemma-quasi-coherent-finite-type-ideals} 保证存在）。
\end{situation}

\begin{lemma}
\label{spaces-more-morphisms-lemma-surjection-in-Noetherian-case}
在情形 \ref{spaces-more-morphisms-situation-descend-resolution-property} 中，假设
$X$ 是诺特的。那么对任意凝聚 $\mathcal{O}_X$-模
$\mathcal{F}$，存在 $r \geq 0$、整数
$n_1, \ldots, n_r \geq 0$，以及满射
$$
\bigoplus\nolimits_{i = 1, \ldots, r} \pi_*(\mathcal{I}^{n_i})
\longrightarrow \mathcal{F}
$$
这是 $\mathcal{O}_X$-模意义下的满射。
\end{lemma}

\begin{proof}
% P09-ERR-0784
以 $\omega_{Y/X}$ 表示如下凝聚 $\mathcal{O}_Y$-模：存在
同构
$$
\pi_*\omega_{Y/X} \cong
\SheafHom_{\mathcal{O}_X}(\pi_*\mathcal{O}_Y, \mathcal{O}_X)
$$
它是 $\pi_*\mathcal{O}_Y$-模的同构，见《空间的态射》引理
\ref{spaces-morphisms-lemma-affine-equivalence-modules} 和
《空间上的下降》引理
\ref{spaces-descent-lemma-finite-over-finite-module}。
典范映射 $\mathcal{O}_X \to \pi_*\mathcal{O}_Y$ 给出典范映射
$$
\text{Tr}_\pi : \pi_*\omega_{Y/X} \longrightarrow \mathcal{O}_X
$$
% P09-ERR-0785
因为 $V$ 是诺特仿射概形，我们可以选取截面
$$
s_1, \ldots, s_r \in
\Gamma(V, \pi^*\mathcal{F} \otimes_{\mathcal{O}_Y} \omega_{Y/X})
$$
使它们生成凝聚模
% P09-ERR-0786
$\pi^*\mathcal{F} \otimes_{\mathcal{O}_X} \omega_{Y/X}$
在 $V$ 上的限制。
由《空间的上同调》引理
\ref{spaces-cohomology-lemma-homs-over-open}，可以选取整数
$n_i \geq 0$，使得 $s_i$ 延拓为映射
$s_i' : \mathcal{I}^{n_i} \to
\pi^*\mathcal{F} \otimes_{\mathcal{O}_Y} \omega_{Y/X}$。
将其推前到 $X$，得到映射
$$
\sigma_i :
\pi_*\mathcal{I}^{n_i}
\xrightarrow{\pi_*s'_i}
\pi_*(\pi^*\mathcal{F} \otimes_{\mathcal{O}_Y} \omega_{Y/X}) =
\mathcal{F} \otimes_{\mathcal{O}_X} \pi_*\omega_{Y/X}
\xrightarrow{\text{Tr}_\pi} \mathcal{F}
$$
其中等号来自《空间的上同调》引理
\ref{spaces-cohomology-lemma-finite-rings}。
为完成证明，我们将说明这些映射之和是满射。

\medskip\noindent
设 $x \in |X|$ 是 $X$ 的一点。设 $v \in |V|$ 是映到 $x$ 的一点。
可以选取 étale 邻域 $(U, u) \to (X, x)$，使得
$$
U \times_X Y = W \coprod W'
$$
（代数空间的不交并），其中 $W \to U$ 是同构，并且
唯一一点 $w \in W$ 位于 $u$ 上方，并映到一点 $v$；
后者属于 $V \subset Y$。
为说明这一点，可使用引理
\ref{spaces-more-morphisms-lemma-etale-splits-off-just-one-quasi-finite-part} 和
《Étale 态射》引理 \ref{etale-lemma-etale-etale-local}。
必要时进一步缩小 $U$，可假设 $W$ 通过投影映入
$V \subset Y$。由于 $\omega_{Y/X}$ 的构造与 étale 局部化交换，
可见
$$
\pi_*\omega_{Y/X}|_U = (\pi|_W)_*\omega_{W/U} \oplus (\pi|_{W'})_*\omega_{W'/U}
$$
我们有 $(\pi|_W)_*\omega_{W/U} = \mathcal{O}_U$，而这个同构由
迹映射 $\text{Tr}_\pi|_U$ 在上述分解的第一个直和项上的限制给出。
因为 $W$ 映入 $V$，可见
$\mathcal{I}^{n_i}|_W = \mathcal{O}_W$。因此
$$
\pi_*(\mathcal{I}^{n_i})|_U = \mathcal{O}_U \oplus
(W' \to U)_*(\mathcal{I}^{n_i}|_{W'})
$$
追踪交换图可知（略去细节），$\sigma_i|_U$ 在直和项
$\mathcal{O}_U$ 上就是映射 $\mathcal{O}_U \to \mathcal{F}$，
它对应于截面
$$
s_i|_W  \in
\Gamma(W,\pi^*\mathcal{F} \otimes_{\mathcal{O}_Y} \omega_{Y/X})
=
\Gamma(W, \mathcal{F}|_W \otimes_{\mathcal{O}_W} \omega_{W/U})
=
\Gamma(U, \mathcal{F})
$$
因为截面 $s_i$ 生成模
$\pi^*\mathcal{F} \otimes_{\mathcal{O}_Y} \omega_{Y/X}$
在 $V$ 上的限制，
并且 $W$ 映入 $V$，所以
$\bigoplus \sigma_i$ 在 $U$ 上的限制是满射。
由于 $x$ 是任意一点，证明完成。
\end{proof}

\begin{lemma}
\label{spaces-more-morphisms-lemma-resolution-property-in-Noetherian-case}
在情形 \ref{spaces-more-morphisms-situation-descend-resolution-property} 中，假设
$X$ 是诺特的。那么 $X$ 具有分解性质，当且仅当
$\pi_*\mathcal{I}$ 是某个有限秩局部自由
$\mathcal{O}_X$-模的商。
\end{lemma}

\begin{proof}
由《空间的上同调》引理
\ref{spaces-cohomology-lemma-finite-pushforward-coherent}，
模 $\pi_*\mathcal{I}$ 是凝聚的。因此，如果 $X$ 具有分解性质，
则 $\pi_*\mathcal{I}$ 是某个有限秩局部自由
$\mathcal{O}_X$-模的商。
反过来，假设给定满射 $\mathcal{E} \to \pi_*\mathcal{I}$，
其源是一个有限秩局部自由 $\mathcal{O}_X$-模 $\mathcal{E}$。
注意，对所有 $n \geq 1$，都有满射
$$
\pi_*\mathcal{I} \otimes_{\mathcal{O}_X} \pi_*\mathcal{I}^n
\longrightarrow
\pi_*\mathcal{I}^{n + 1}
$$
因此，$\mathcal{E}^{\otimes n}$ 满射到
$\pi_*\mathcal{I}^n$，而且这对所有 $n \geq 1$ 成立。将此与引理
\ref{spaces-more-morphisms-lemma-surjection-in-Noetherian-case} 的结论结合，
便知 $X$ 具有分解性质。
\end{proof}

\begin{lemma}
\label{spaces-more-morphisms-lemma-resolution-property-one-sheaf}
在情形 \ref{spaces-more-morphisms-situation-descend-resolution-property} 中，
代数空间 $X$ 具有分解性质，当且仅当
$\pi_*\mathcal{I}$ 是某个有限秩局部自由
$\mathcal{O}_X$-模的商。
\end{lemma}

\begin{proof}
对于推前 $\pi_*\mathcal{G}$，若这里的源是有限型拟凝聚
$\mathcal{O}_Y$-模 $\mathcal{G}$，则由《空间上的下降》引理
\ref{spaces-descent-lemma-finite-over-finite-module}，
可知它是有限型拟凝聚 $\mathcal{O}_X$-模。
特别地，如果 $X$ 具有分解性质，则由《空间的导出范畴》定义
\ref{spaces-perfect-definition-resolution-property}，
$\pi_*\mathcal{I}$ 是某个有限秩局部自由
$\mathcal{O}_X$-模的商。

\medskip\noindent
假设有满射 $\mathcal{E} \to \pi_*\mathcal{I}$，
其源是有限秩局部自由 $\mathcal{O}_X$-模 $\mathcal{E}$。
在证明的其余部分中，我们把问题约化到引理
\ref{spaces-more-morphisms-lemma-resolution-property-in-Noetherian-case}
已经处理的诺特情形，以说明 $X$ 具有分解性质。
我们建议读者跳过证明的其余部分。

\medskip\noindent
第一个约化是，可以把 $X$ 视为
$\Spec(\mathbf{Z})$ 上的代数空间，见
《空间》定义 \ref{spaces-definition-base-change}。
（这不影响引理的条件或结论。）

\medskip\noindent
由《空间的极限》引理
\ref{spaces-limits-lemma-finite-in-finite-and-finite-presentation}，
可以写成 $Y = \lim Y_i$，其中 $Y_i$ 在 $X$ 上有限且有限表示，
而转移映射是闭浸入。考虑闭子空间
$Z = V(\mathcal{I})$，它位于 $Y$ 中。由于 $\mathcal{I}$ 是有限型的，
态射 $Z \to Y$ 是有限表示的。因此，由《空间的极限》引理
\ref{spaces-limits-lemma-descend-finite-presentation}，
可以找到 $i$ 和有限表示态射 $Z_i \to Y_i$，其到 $Y$ 的
基变换是 $Z \to Y$。
对 $i' \geq i$，记
$Z_{i'} = Z_i \times_{Y_i} Y_{i'}$。
增大 $i$ 后，由《空间的极限》引理
\ref{spaces-limits-lemma-descend-closed-immersion}，
可假设 $Z_i \to Y_i$ 是（有限表示的）闭浸入。
% P09-ERR-0787
以 $\mathcal{I}_i \subset \mathcal{O}_{Y_i}$ 表示
$Z_i$ 的理想层，并以 $\pi_i : Y_i \to X$ 表示结构态射。
对 $i' \geq i$ 也同样记号。由于
$Z = \lim_{i' \geq i} Z_{i'}$，我们有
$$
\pi_*\mathcal{I} = \colim \pi_{i', *}\mathcal{I}_{i'}
$$
% P09-ERR-0788
这个系统中的转移映射全都是满射；这由映射
$\pi_{i, *}\mathcal{O}_{Y_i} \to \pi_{i', *}\mathcal{O}_{Y_{i'}}$
的满射性以及事实
$Z_{i'} = Z_i \times_{Y_i} Y_{i'}$ 推出。
由《空间的上同调》引理
\ref{spaces-cohomology-lemma-finite-presentation-quasi-compact-colimit}，
对某个 $i' \geq i$，映射
$\mathcal{E} \to \pi_*\mathcal{I}$ 可提升为映射
$\mathcal{E} \to \pi_{i', *}\mathcal{I}_{i'}$。
增大 $i'$ 后，该映射
$\mathcal{E} \to \pi_{i', *}\mathcal{I}_{i'}$
成为满射（否则，具有满转移映射的余核系统之余极限非零）。
这把我们约化到下一段讨论的情形。

\medskip\noindent
假设 $X$ 是 $\mathbf{Z}$ 上的代数空间，并且
$Y \to X$ 是有限表示的。由绝对诺特逼近，可以把
$X = \lim X_i$ 写成有向极限，其中每个 $X_i$ 都是
$\mathbf{Z}$ 上有限型的拟分离代数空间，而转移态射是仿射的，
见《空间的极限》命题 \ref{spaces-limits-proposition-approximate}。
由于 $\pi : Y \to X$ 是有限表示的，由《空间的极限》引理
\ref{spaces-limits-lemma-descend-finite-presentation}，
可以找到 $i$ 和有限表示态射
$\pi_i : Y_i \to X_i$，其到 $X$ 的基变换是 $\pi$。
增大 $i$ 后，由《空间的极限》引理
\ref{spaces-limits-lemma-descend-finite}，
可假设 $\pi_i$ 是有限的。
接着，由《空间的极限》引理
\ref{spaces-limits-lemma-descend-invertible-modules}，
可假设存在有限秩局部自由
$\mathcal{O}_{X_i}$-模 $\mathcal{E}_i$，
其到 $X$ 的拉回是 $\mathcal{E}$。
还可假设存在映射
$\mathcal{E}_i \to \pi_{i, *}\mathcal{O}_{Y_i}$，
其到 $X$ 的拉回是复合映射
$\mathcal{E} \to \pi_*\mathcal{I} \to \pi_*\mathcal{O}_Y$，
见《空间的极限》引理
\ref{spaces-limits-lemma-descend-modules-finite-presentation}。
余核
$$
\mathcal{E}_i \to \pi_{i, *}\mathcal{O}_{Y_i} \to \mathcal{Q}_i \to 0
$$
% P09-ERR-0789
是凝聚的 $\mathcal{O}_{Y_i}$-模，其到 $X$ 的拉回是（有限表示）余核
$\mathcal{Q}$；这个余核来自映射
$\mathcal{E} \to \pi_*\mathcal{O}_Y$。换言之，我们有
$\mathcal{Q} = \pi_*(\mathcal{O}_Y/\mathcal{I})$。考虑映射
$$
\mathcal{E}_i
\otimes_{\mathcal{O}_{X_i}}
\pi_{i, *} \mathcal{O}_{Y_i}
\longrightarrow
\pi_{i, *}\mathcal{O}_{Y_i}
\otimes_{\mathcal{O}_{X_i}}
\pi_{i, *} \mathcal{O}_{Y_i}
\to
\pi_{i, *} \mathcal{O}_{Y_i}
\to
\mathcal{Q}_i
$$
其中第二个箭头由
$\pi_{i, *}\mathcal{O}_{Y_i}$ 的代数结构给出。
% P09-ERR-0790
这个映射到 $Y$ 的拉回为零，因为
$\mathcal{E} \to \pi_*\mathcal{O}_Y$ 的像是理想
$\pi_*\mathcal{I}$。因此，由《空间的极限》引理
\ref{spaces-limits-lemma-descend-modules-finite-presentation}，
增大 $i$ 后，可假设所显示的复合映射为零。
% P09-ERR-0791
这恰好意味着 $\mathcal{E}_i \to \pi_{i, *}\mathcal{O}_{Y_i}$
的像具有形式 $\pi_{i, *}\mathcal{I}_i$，
其中 $\mathcal{I}_i \subset \mathcal{O}_{Y_i}$ 是某个凝聚理想层。
由于 $\mathcal{E}_i \to \pi_{i, *}\mathcal{O}_{Y_i}$
拉回为 $\mathcal{E} \to \pi_*\mathcal{O}_Y$，可见
$\mathcal{I}_i$ 到 $Y$ 的拉回生成 $\mathcal{I}$。
% P09-ERR-0792
以 $V_i \subset Y_i$ 表示其补集为
$V(\mathcal{I}_i) \subset Y_i$ 的开子空间。
由上述评注，$V$ 是 $V_i$ 的逆像。
增大 $i$ 后，可假设 $V_i$ 是仿射的，并且
$\pi_i|_{V_i} : V_i \to X_i$ 是 étale 的，见
《空间的极限》引理
\ref{spaces-limits-lemma-limit-is-affine} 和
\ref{spaces-limits-lemma-descend-etale}。
至此，可应用引理
\ref{spaces-more-morphisms-lemma-resolution-property-in-Noetherian-case}
断定 $X_i$ 具有分解性质。由于 $X \to X_i$ 是仿射的，
由《空间的导出范畴》引理
\ref{spaces-perfect-lemma-resolution-property-goes-up-affine}，
可断定 $X$ 也具有分解性质。
\end{proof}

\begin{lemma}
\label{spaces-more-morphisms-lemma-resolution-property-descends}
设 $S$ 是概形。
设 $X = \lim X_i$ 是 $S$ 上拟紧且拟分离代数空间的一个
直系统的极限，其转移态射都是仿射的。
% P09-ERR-0793
那么 $X$ 具有分解性质，当且仅当 $X_i$ 对某个 $i$ 具有分解性质。
\end{lemma}

\begin{proof}
如果 $X_i$ 具有分解性质，那么由《空间的导出范畴》引理
\ref{spaces-perfect-lemma-resolution-property-goes-up-affine}，
$X$ 也具有分解性质。
假设 $X$ 具有分解性质。
选取 $i \in I$。可以选取仿射概形 $V_i$ 和满 étale 态射
$V_i \to X_i$
（《空间的性质》引理
\ref{spaces-properties-lemma-quasi-compact-affine-cover}）。
可以选取嵌入 $j : V_i \to Y_i$，其中 $Y_i$ 在 $X_i$ 上
有限且有限表示
（引理 \ref{spaces-more-morphisms-lemma-quasi-finite-separated-pass-through-finite-addendum}）。
可以选取有限型拟凝聚理想
$\mathcal{I}_i \subset \mathcal{O}_{Y_i}$，使得
$V_i = Y_i \setminus V(\mathcal{I}_i)$
（《空间的极限》引理
\ref{spaces-limits-lemma-quasi-coherent-finite-type-ideals}）。
% P09-ERR-0794
以 $V \to Y \to X$ 表示 $V_i \to Y_i \to X_i$ 到 $X$ 的基变换。
以 $\mathcal{I} \subset \mathcal{O}_Y$ 表示理想
$\mathcal{I}_i$ 的拉回。由引理
\ref{spaces-more-morphisms-lemma-resolution-property-one-sheaf} 的容易方向，
存在有限秩局部自由 $\mathcal{O}_X$-模
$\mathcal{E}$ 和满射
$\mathcal{E} \to \pi_*\mathcal{I}$。
注意，由于 $\pi_i : Y_i \to X_i$ 是有限且有限表示的，
$\pi : Y \to X$ 也有限且有限表示，并且
$\mathcal{O}_{X_i}$-模
$\pi_{i, *}\mathcal{O}_{Y_i}$ 和
$\pi_{i, *}(\mathcal{O}_{Y_i}/\mathcal{I}_i)$
% P09-ERR-0795
都是有限表示的；它们拉回到 $X$ 后分别给出
$\pi_*\mathcal{O}_Y$ 和
$\pi_*(\mathcal{O}_Y/\mathcal{I})$。因此，由《空间的极限》引理
\ref{spaces-limits-lemma-descend-modules-finite-presentation}，
增大 $i$ 后，可以找到有限秩局部自由
$\mathcal{O}_{X_i}$-模 $\mathcal{E}_i$ 和映射
$\mathcal{E}_i \to \pi_{i, *}\mathcal{O}_{Y_i}$，
其到 $X$ 的基变换恢复复合映射
$\mathcal{E} \to \pi_*\mathcal{I} \to \pi_*\mathcal{O}_Y$。
有限表示 $\mathcal{O}_{X_i}$-模
$\Coker(\mathcal{E}_i \to \pi_{i, *}\mathcal{O}_{Y_i})$ 和
$\pi_{i, *}(\mathcal{O}_{Y_i}/\mathcal{I}_i)$
到 $X$ 的拉回作为 $\pi_*\mathcal{O}_Y$ 的商彼此相同。
因此，由《空间的极限》引理
\ref{spaces-limits-lemma-descend-modules-finite-presentation}，
可以假设二者相同；换言之，
% P09-ERR-0796
$\mathcal{E}_i \to \pi_{i, *}\mathcal{O}_{X_i}$ 的像等于
$\pi_{i, *}\mathcal{I}_i$。
于是，由引理 \ref{spaces-more-morphisms-lemma-resolution-property-one-sheaf}，
可断定 $X_i$ 具有分解性质。
\end{proof}

\begin{lemma}
\label{spaces-more-morphisms-lemma-resolution-property-affine-diagonal}
设 $S$ 是概形。
设 $X$ 是具有分解性质的拟紧且拟分离代数空间。
那么 $X$ 在 $\mathbf{Z}$ 上具有仿射对角线
（定义见《空间的性质》定义
\ref{spaces-properties-definition-separated}）。
\end{lemma}

\begin{proof}
我们可以像概形情形那样证明，但这里改为从概形情形推出本引理。
首先，可以并且确实假设
$S = \Spec(\mathbf{Z})$。
接着，选取概形 $Y$ 和满整态射
$f : Y \to X$，见《合宜空间》引理
\ref{decent-spaces-lemma-there-is-a-scheme-integral-over}。
于是 $f$ 是仿射的，故由《空间的导出范畴》引理
\ref{spaces-perfect-lemma-resolution-property-goes-up-affine}，
$Y$ 具有分解性质。因此，由概形情形，
概形 $Y$ 具有仿射对角线，见《概形的导出范畴》引理
\ref{perfect-lemma-resolution-property-affine-diagonal}。
接着，考虑交换图
$$
\xymatrix{
Y \ar[d] \ar[rr]_{\Delta_Y} & & Y \times_{\mathbf{Z}} Y \ar[d] \\
X \ar[rr]^{\Delta_X} & & X \times_{\mathbf{Z}} X
}
$$
注意，右侧竖直箭头是整的，因而是仿射的。
设 $W \to X \times_{\mathbf{Z}} X$ 是态射，其中 $W$ 是仿射的。
于是可见
$$
Y \times_{X \times_{\mathbf{Z}} X} W =
Y \times_{\Delta_Y, Y \times_{\mathbf{Z}} Y}
(Y \times_{\mathbf{Z}} Y) \times_{X \times_{\mathbf{Z}} X} W
$$
是仿射的。另一方面，$Y \to X$ 是整且满的，因此
$$
Y \times_{X \times_{\mathbf{Z}} X} W
\longrightarrow
X \times_{X \times_{\mathbf{Z}} X} W
$$
% P09-ERR-0797
是整满态射，因为它是 $Y \to X$ 到 $W$ 的基变换。
由《空间的极限》命题
\ref{spaces-limits-proposition-affine}，
可断定这个箭头的目标是仿射的。
于是 $\Delta_X$ 如所需是仿射的。
\end{proof}






\section{爆破与分解性质}
\label{spaces-more-morphisms-section-resolution-property-by-blowing-up}

\noindent
我们证明，作一次爆破后便满足分解性质。

\begin{lemma}
\label{spaces-more-morphisms-lemma-blow-up-resolution-property}
设 $S$ 为概形。设 $X$ 为 $S$ 上拟紧且拟分离的
代数空间。假设 $|X|$ 只有有限多个不可约分支。
则存在稠密拟紧开子空间 $U \subset X$ 和一个
$U$-容许爆破 $X' \to X$，使得代数空间
$X'$ 具有分解性质。
\end{lemma}

\begin{proof}
由《空间的极限》引理
\ref{spaces-limits-lemma-there-is-a-scheme-finite-over-refined}，
存在满、有限且有限表示的态射
$f : Y \to X$，其中 $Y$ 是概形；
% P09-ERR-0798
还存在拟紧稠密开子空间 $U \subset X$，使得
$f^{-1}(U) \to U$ 有限平展。
由《态射详论》引理 \ref{more-morphisms-lemma-blow-up-ample-family}，
存在拟紧稠密开子空间 $V \subset Y$ 和一个
$V$-容许爆破 $Y' \to Y$，使得 $Y'$ 具有丰沛可逆模族。
缩小 $U$ 后，可以假设 $f^{-1}(U) \subset V$
（细节略）。因此 $f' : Y' \to X$ 在 $U$ 上有限平展；
特别地，态射 $(f')^{-1}(U) \to U$ 有限局部自由。
由引理 \ref{spaces-more-morphisms-lemma-finite-after-blowing-up}，存在一个
$U$-容许爆破 $X' \to X$，使得严格变换 $Y''$
（即 $Y'$ 的严格变换）在 $X'$ 上有限局部自由。
% P09-ERR-0799
图示如下：
$$
\xymatrix{
Y'' \ar[d] \ar[r]_g &
Y' \ar[r] &
Y \ar[d] \\
X' \ar[rr] & & X
}
$$
由于 $g : Y'' \to Y'$ 是一次爆破
（《空间上的除子》引理 \ref{spaces-divisors-lemma-strict-transform}），
其中心是 $X' \to X$ 的爆破中心之逆像，
可知 $g : Y'' \to Y'$ 是射影态射，并且存在某个
$g$-丰沛可逆模于 $Y''$ 上。因此，由《态射详论》引理
\ref{more-morphisms-lemma-ample-family-ample-relative}，
可知 $Y''$ 具有丰沛可逆模族。
所以 $Y''$ 具有分解性质；见《概形的导出范畴》引理
\ref{perfect-lemma-resolution-property-ample-family}。
最后，由《空间的导出范畴》引理
\ref{spaces-perfect-lemma-resolution-property-goes-down-finite-flat}，
可断定 $X'$ 具有分解性质。
\end{proof}

\begin{lemma}
\label{spaces-more-morphisms-lemma-blow-up-resolution-property-filtered}
设 $S$ 为概形。设 $X$ 为 $S$ 上拟紧且拟分离的
代数空间。则存在 $t \geq 0$ 及闭子空间
$$
X \supset Z_0 \supset Z_1 \supset \ldots \supset Z_t = \emptyset
$$
使得 $Z_i \to X$ 有限表示，
$Z_0 \subset X$ 是加厚，并且对于每个
% P09-ERR-0800
$i = 0, \ldots t - 1$，存在一个
% P09-ERR-0801
$(Z_i \setminus Z_{i - 1})$-容许爆破
$Z'_i \to Z_i$，使得 $Z'_i$ 具有分解性质。
\end{lemma}

\begin{proof}
本段使用绝对诺特逼近，把问题约化到在
$\Spec(\mathbf{Z})$ 上有限表示的代数空间情形。
我们可以把 $X$ 看成 $\Spec(\mathbf{Z})$ 上的代数空间；见
《空间》定义 \ref{spaces-definition-base-change} 和
《空间的性质》定义 \ref{spaces-properties-definition-separated}。
因此可以应用《空间的极限》命题
\ref{spaces-limits-proposition-approximate}。
由此可找到仿射态射 $X \to X_0$，其中 $X_0$
在 $\mathbf{Z}$ 上有限表示。
如果能对 $X_0$ 证明本引理，那么把分层和各爆破中心拉回到 $X$，
就得到 $X$ 的结论；这里用到分解性质沿仿射态射上升
（《空间的导出范畴》引理
\ref{spaces-perfect-lemma-resolution-property-goes-up-affine}），
以及仿射态射的严格变换仍为仿射态射——细节略。
这就把问题约化到下一段讨论的情形。

\medskip\noindent
假设 $X$ 在 $\mathbf{Z}$ 上有限表示。
那么 $X$ 是诺特的，并且 $|X|$ 是有限维诺特拓扑空间
（具有有限多个不可约分支）。
因此可以对 $\dim(|X|)$ 作归纳。
由引理 \ref{spaces-more-morphisms-lemma-blow-up-resolution-property}，
存在稠密开子空间 $U \subset X$ 和一个 $U$-容许爆破
$X' \to X$，使得 $X'$ 具有分解性质。
令 $Z_0 = X$，并令 $Z_1 \subset X$ 为满足
$|Z_1| = |X| \setminus |U|$ 的约化闭子空间。
由归纳假设，可找到整数 $t \geq 0$ 及由闭子空间组成的滤过
$$
Z_1 \supset Z_{1, 0} \supset Z_{1, 1} \supset \ldots
\supset Z_{1, t} = \emptyset
$$
其中 $Z_{1, 0} \to Z_1$ 是加厚，并且存在
$(Z_{1, i} \setminus Z_{1, i + 1})$-容许爆破
$Z'_{1, i} \to Z_{1, i}$，使得 $Z'_{1, i}$
具有分解性质。由于 $Z_1$ 是约化的，故 $Z_1 = Z_{1, 0}$。
因此，可以令 $Z_i = Z_{1, i - 1}$ 和
$Z'_i = Z'_{1, i - 1}$（$i \geq 1$）；引理得证。
\end{proof}
